
\medterm{P Arm Of A Chromosome} The shorter arm of a chromosome, labeled p, opposite the longer q arm. A genetic report may use this label to show where a gene change or missing or extra DNA segment is located.

\medterm{P Wave} The first small wave on an ECG, showing the electrical activation of the heart’s upper chambers. Its shape and timing can help identify abnormal heart rhythms, enlargement of a chamber, or problems with electrical conduction.

\medterm{P-Glycoprotein} A transport protein in cell membranes that pumps many medicines and other substances out of cells. It affects how medicines are absorbed and removed and can cause interactions when its activity is blocked or increased.


\medterm{P53 Tumor Suppressor} A protein that helps prevent damaged cells from dividing or surviving. Changes in the TP53 gene can allow abnormal cells to multiply and contribute to many cancers, but a gene change alone does not determine a person's outcome.

\medterm{PA X-Ray} A posteroanterior X-ray is taken with the beam traveling from the back toward the front, commonly for a chest image. This position reduces enlargement of the heart’s outline compared with some other views.

\medterm{PA/IVS} Alternate index label; see \textit{Pulmonary atresia with intact ventricular septum}.

\medterm{Pacemaker} A small implanted device that sends electrical signals to the heart when its natural rhythm is too slow or unreliable. It may treat fainting or symptoms from slow heartbeats and certain conduction blocks; regular checks monitor the battery and wires.

\synonyms
Pacemaker, artificial

\medterm{Pacemaker Complications} Problems after pacemaker placement can include bleeding, infection, a collapsed lung, heart or blood-vessel injury, movement or breakage of a lead, blood clots, battery failure, or an abnormal heartbeat. New chest pain, fainting, fever, or worsening breathlessness needs urgent assessment.

\medterm{Pachyderma} Abnormal thickening and coarsening of the skin, which may occur with chronic inflammation, lymphedema, endocrine disease, or pachydermoperiostosis.

\medterm{Pachygyria} A birth-related brain-development abnormality in which the brain’s outer surface has fewer, broader folds than usual. It can cause developmental delay, learning difficulties, seizures, movement problems, or other neurologic effects of varying severity.

\medterm{Pachymeningitis} Inflammation and thickening of the tough outer covering of the brain and spinal cord. It may cause headache, nerve problems, or pressure on the brain and can result from infection, immune disease, or other inflammation.

\medterm{Pachyonychia Congenita} An inherited skin disorder causing very thick, abnormal nails and painful thick skin on the palms and soles. Some people also develop blisters, cysts, excess sweating, or white patches in the mouth.


\medterm{Pacing} Planning activity, rest, and effort so symptoms do not flare while endurance gradually improves. It may help with fatigue, chronic pain, heart or lung disease, and recovery after infection.

\medterm{Pacinian Corpuscle} A specialized nerve ending deep in the skin and connective tissue that detects deep pressure and vibration. It helps the nervous system sense contact, movement, and the handling of objects.


\medterm{Paclitaxel} Paclitaxel is a taxane chemotherapy drug that stabilizes microtubules and prevents cell division; it treats several cancers and can cause myelosuppression, neuropathy, alopecia, myalgia, and hypersensitivity reactions.

\medterm{Paclitaxel Liposome} A liposomal formulation of paclitaxel designed to alter drug distribution and delivery; paclitaxel disrupts microtubules and is used as chemotherapy in selected cancers.


\medterm{PACS} A computer system used by healthcare services to store, retrieve, view, and share medical images such as X-rays, CT scans, MRI scans, and ultrasound studies.





\medterm{PAES} Alternate index label; see \textit{Popliteal artery entrapment}.

\medterm{Page Kidney} High blood pressure caused when blood or another mass presses on a kidney from the outside. The pressure activates the kidney’s blood-pressure system and may follow injury, bleeding, or treatment with blood thinners.

\medterm{Paget Disease Of Bone} A long-term disorder in which bone is broken down and rebuilt too quickly and unevenly. Affected bones may become enlarged, weak, painful, or bent and can fracture or press on nearby nerves; many people have no symptoms.

\synonyms
Paget's disease; Paget disease of the bone; Osteitis deformans

\medterm{Paget Disease Of The Anus} A rare cancer in the skin around the anus that may look like persistent eczema, redness, itching, soreness, or an erosion. A biopsy is needed, and doctors usually check for cancer in nearby or internal organs.

\medterm{Paget Disease Of The Nipple} A rare breast cancer involving the nipple and areola, often causing persistent eczema-like scaling, redness, itching, burning, or discharge. It may be associated with underlying ductal carcinoma in situ or invasive breast cancer.

\medterm{Paget Disease Of The Vulva} A rare cancer in the skin of the vulva that may cause a persistent itchy, red, scaly, sore, or moist patch. A biopsy is needed, and assessment may look for another cancer beneath or elsewhere.

\medterm{Pagetoid Reticulosis} A rare localized variant of cutaneous T-cell lymphoma with solitary slowly enlarging scaly plaques containing atypical pagetoid lymphocytes.

\medterm{Pain} An unpleasant physical and emotional experience linked to actual or possible tissue injury. Pain can be short-lived or persistent and is influenced by nerves, inflammation, health, emotions, past experiences, and social circumstances.

\medterm{Pain (Geriatric)} Pain in an older adult, which may be underreported or shown as reduced activity, sleep change, or confusion. Assessment considers communication, memory, function, medicines, and causes; treatment is individualized and monitored for side effects.

\medterm{Pain (Vascular)} Pain caused by impaired blood flow, often appearing in a limb with exertion as claudication or, in more severe disease, at rest.


\medterm{Pain Assessment Scales} Tools that help a person describe pain intensity or its effect on daily life. They may use numbers, words, faces, or observation of behavior when self-report is not possible. A score tracks change but does not identify the cause or fully measure suffering.

\medterm{Pain Cancer (Cancer Pain)} Alternate index label for Cancer Pain.

\medterm{Pain Elbow (Elbow Pain)} Alternate index label for Elbow Pain.

\medterm{Pain Gallbladder (Gallbladder Pain (Gall Bladder Pain))} Alternate index label for Gallbladder Pain (Gall Bladder Pain).

\medterm{Pain Heel (Heel Spurs)} Alternate index label for Heel Spurs.

\medterm{Pain Kidney (Kidney Pain)} Alternate index label for Kidney Pain.


\medterm{Pain Lower Right Abdomen (Appendicitis)} Alternate index label for Appendicitis.

\medterm{Pain Management} Care that identifies the cause of pain and reduces pain while improving function and quality of life. It may combine medicines, physical treatment, psychological support, procedures, and treatment of the underlying condition, with risks reviewed regularly.

\medterm{Pain Medications - Opioid Analgesics} Opioid analgesics can relieve selected moderate-to-severe pain but may cause constipation, sedation, tolerance, dependence, overdose, and slowed breathing. They require individualized prescribing and monitoring.

\synonyms
Pain Medications - Narcotics

\medterm{Pain Medicine} The medical specialty concerned with assessing and treating acute, chronic, and complex pain through medication, procedures, rehabilitation, and multidisciplinary care.

\synonyms
Pain management; pain medicine specialty
\medterm{Pain Nerve (Neuropathic Pain (Nerve Pain))} Alternate index label for Neuropathic Pain (Nerve Pain).

\medterm{Pain Perception} The process by which the nervous system notices possible or actual tissue injury and the brain interprets it as pain. Attention, emotions, past experiences, inflammation, and medicines can change how pain feels.

\medterm{Pain Relief (Analgesia)} Reduction of pain using medicines, procedures, physical treatment, psychological support, or a combination. The best approach depends on the cause, duration, severity, function, other health conditions, and risks of treatment.
\synonyms
Analgesia

\medterm{Pain Scrotum (Testicular Disorders)} Alternate index label for Testicular Disorders.

\medterm{Pain Tailbone (Coccydynia)} Alternate index label for Coccydynia.

\medterm{Pain Threshold} The minimum intensity of a stimulus that a person perceives as painful; it varies with tissue, context, attention, inflammation, medications, and neurologic state.

\medterm{Pain Unit} A clinical service specializing in assessment and treatment of acute, chronic, cancer, neuropathic, or procedural pain using medicines, procedures, rehabilitation, and psychological care.

\medterm{Pain Vaginal (Vaginal Pain (Vulvodynia))} Alternate index label for Vaginal Pain (Vulvodynia).

\medterm{Pain Whiplash (Whiplash)} Alternate index label for Whiplash.

\medterm{Pain, Back} Alternate index label; see \textit{Back pain}.

\medterm{Pain, Breast} Alternate index label; see \textit{Breast pain}.

\medterm{Painful Bladder Syndrome} Alternate index label; see \textit{Interstitial cystitis}.

\medterm{Painful Intercourse (Dyspareunia)} Persistent or recurrent genital pain before, during, or after sexual intercourse, with physical, hormonal, emotional, or relationship-related causes.

\medterm{Painful Menstrual Periods} Painful menstrual periods, or dysmenorrhea, are cramps during menstruation caused by uterine contractions or conditions such as endometriosis or fibroids.

\medterm{Painful Swallowing} Painful swallowing, or odynophagia, is discomfort when swallowing and may result from infection, inflammation, injury, reflux, or a blockage.

\medterm{Paint, Lacquer, And Varnish Remover Poisoning} Poisoning caused by ingestion, inhalation, or skin or eye exposure to solvents and other chemicals in paint, lacquer, or varnish removers. Effects may include irritation, central nervous system depression, respiratory injury, metabolic abnormalities, or organ damage; severity depends on the substance, dose, route, and timing.

\medterm{Palatal Muscle} Muscle of the soft palate that elevates or tenses the palate and helps separate the nasal and oral cavities during swallowing and speech.

\medterm{Palatal Myoclonus} Repeated, involuntary twitching of the soft palate, sometimes causing a clicking sound in the ear or difficulty speaking and swallowing. It may follow a brainstem injury or occur without an identified cause.

\medterm{Palate} The roof of the mouth, consisting of a bony hard palate and muscular soft palate.
\medterm{Palate Neoplasm} An abnormal growth in the roof of the mouth. It may be noncancerous or cancerous and can cause a lump, pain, bleeding, loose teeth, difficulty swallowing, or speech changes.

\medterm{Palatine} Relating to the palate, the roof of the mouth, or the palatine bones. Palatine structures include the tonsils, nerves, and blood vessels of the mouth; problems may include tonsillitis, cleft palate, or infection near a tooth.
\medterm{Palatine Tonsil} One of two masses of lymphoid tissue at the sides of the throat that help detect inhaled or swallowed pathogens.


\medterm{Palatorrhaphy} Palatorrhaphy is surgical repair of a cleft or defect of the palate to restore separation between the oral and nasal cavities and improve feeding and speech.

\medterm{Palatoschisis} A congenital cleft or failure of fusion of the palate, commonly called cleft palate, which may affect feeding, speech, hearing, and dental development.

\medterm{Paleness} Unusual lightening of the skin or mucous membranes, which may result from reduced blood flow, anemia, shock, cold, or normal variation.



\medterm{Palindromic Rheumatism} Repeated attacks of joint pain and swelling that begin suddenly and resolve completely within hours or days, leaving no lasting joint damage between attacks. Some people later develop rheumatoid arthritis, especially when specific antibodies are present.

\medterm{Paliperidone} An antipsychotic medicine used to treat schizophrenia and schizoaffective disorder. It changes the action of certain brain chemicals and is available as tablets or long-acting injections; possible effects include sleepiness, movement problems, weight gain, hormone changes, and heart-rhythm changes.
\medterm{Palivizumab} A monoclonal antibody against respiratory syncytial virus fusion protein, used prophylactically in selected high-risk infants and young children during RSV season.


\medterm{Pallesthesia} The ability to feel vibration, usually tested with a vibrating tuning fork on the fingers or toes. Reduced vibration sense may result from nerve disease, spinal-cord problems, poor circulation, or other neurologic conditions.
\medterm{Palliative Care} Medical care for serious illness that relieves pain and other symptoms and supports daily life, emotional health, and family needs. It can begin at any stage and can be provided alongside treatment intended to control or cure disease.

\medterm{Palliative Rehabilitation} Rehabilitation for a person with serious or life-limiting illness that focuses on comfort, independence, meaningful activities, and quality of life. It may address weakness, breathlessness, pain, fatigue, movement, communication, swallowing, or daily tasks.

\medterm{Palliative Radiation} Radiation therapy given to relieve symptoms caused by cancer, such as pain, bleeding, obstruction, or pressure from a tumor. It is usually intended to improve comfort or function rather than cure the cancer and may be delivered in a short treatment course.

\medterm{Palliative Surgery} Surgery intended to relieve symptoms, prevent complications, or improve function when the underlying disease may not be curable. It can ease blockage, bleeding, pain, breathing difficulty, or problems caused by advanced cancer.
\medterm{Palliative Treatment} Treatment intended to relieve symptoms and improve comfort or quality of life rather than cure the underlying disease.

\medterm{Pallidotomy} Brain surgery that makes a carefully placed small injury in an area involved in movement control. It may reduce severe Parkinson disease symptoms such as involuntary movements when medicines no longer provide adequate control.
\medterm{Pallister-Hall Syndrome} A rare inherited condition that may cause an extra finger or toe, a noncancerous growth near the brain’s hormone-control center, and abnormalities of the anus, kidneys, airway, or pituitary hormones.

\medterm{Pallor} Unusual paleness of skin or mucous membranes, often due to reduced blood flow, anemia, or other illness.
\medterm{Palm} The inner surface of the hand between the wrist and fingers. It contains thick skin, creases, muscles, nerves, and blood vessels that help with grip, touch, movement, and protection.

\medterm{PALM-COEIN Classification} A system clinicians use to organize the causes of abnormal bleeding from the uterus. PALM covers structural causes such as polyps, adenomyosis, fibroids, and cancer; COEIN covers clotting, ovulation, lining, medicine-related, and other causes. It helps guide evaluation and treatment.

\synonyms
FIGO AUB Classification; PALM-COEIN System

\medterm{Palm Sweating Excessive (Hyperhidrosis)} Alternate index label for Hyperhidrosis.

\medterm{Palmar} Relating to the palm or front surface of the hand. The palm has thick skin, sweat glands, muscles, nerves, and blood vessels; changes such as redness, pain, numbness, or a tightening band may help identify an underlying condition.
\medterm{Palmar Erythema Of Pregnancy} Redness of both palms that can occur during pregnancy, often because of normal blood-vessel and hormone changes. It is usually harmless and fades after birth, but new itching, yellowing, pain, or other symptoms need assessment.

\medterm{Palmar Grasp Reflex} A primitive reflex in which an infant flexes the fingers around an object placed in the palm.
It is normally present at birth and gradually diminishes as voluntary grasp develops;
asymmetry or persistence may warrant neurologic assessment.

\medterm{Palmar Reflex} The palmar reflex is an infant response in which stroking the palm causes finger flexion; persistence or asymmetry beyond the expected age may indicate neurologic dysfunction.

\medterm{Palmitate} Palmitate is the salt or ester form of palmitic acid and participates in fatty-acid metabolism, lipid storage, and protein modification.

\medterm{Palmitic Acid} A 16-carbon saturated fatty acid common in dietary and stored fats; excess saturated-fat intake can contribute to adverse lipid profiles, although the body also synthesizes it normally.

\medterm{Palmoplantar Hyperkeratosis} Excessive thickening of the skin on the palms and soles, inherited or acquired, causing calluses, fissures, pain, impaired walking, and sometimes infection.

\medterm{Palonosetron Hydrochloride} A long-acting 5-HT3 serotonin-receptor antagonist used to prevent nausea and vomiting caused by chemotherapy, surgery, or selected other treatments.

\medterm{Palpable} Palpable means detectable by touch during physical examination, such as a mass, enlarged organ, lymph node, pulse, tenderness, or abnormal surface.

\medterm{Palpate} To examine part of the body by feeling it carefully with the hands. A healthcare professional may assess size, shape, tenderness, warmth, movement, swelling, firmness, pulses, or unusual lumps.

\medterm{Palpation} Examination by touching the body to assess tenderness, temperature, texture, masses, pulses, organ size, and movement.

\medterm{Palpebra} The anatomical term for an eyelid. Eyelids protect the eye, spread tears over the cornea, and limit light entry; disorders include drooping, inward or outward turning, inflammation, styes, and blocked oil glands.
\medterm{Palpebral Conjunctiva} The thin, moist membrane lining the inner surfaces of the eyelids. It joins the membrane covering the front of the eye and helps protect, lubricate, and move tears across the eye’s surface.

\medterm{Palpebral Slant - Eye} The direction of the opening between the upper and lower eyelids. Its shape varies normally, but an unusual slant may be associated with a genetic or developmental condition.

\medterm{Palpitation} Awareness of the heart beating fast, forcefully, irregularly, or as if it skipped a beat. Causes include stress, stimulants, anemia, thyroid disease, medicines, and abnormal rhythms.


\medterm{PALS} Alternate name; see \textit{Pediatric Advanced Life Support}.

\medterm{Palsy} Weakness or loss of movement in a body part, often affecting muscles supplied by a particular nerve. It may be temporary or lasting and can result from nerve injury, infection, stroke, inflammation, or a problem present from birth.
\medterm{Palsy Cerebral (Cerebral Palsy)} Alternate index label for Cerebral Palsy.

\medterm{Palsy Of The Long Thoracic Nerve} Weakness of the nerve that activates a muscle holding the shoulder blade against the ribs. It can cause a winged shoulder blade, difficulty raising the arm, and pain after injury, infection, surgery, or pressure on the nerve.

\medterm{Palsy, Bell's} Alternate index label; see \textit{Bell's palsy}.

\medterm{Palsy, Cerebral} Alternate index label; see \textit{Cerebral palsy}.

\medterm{Palsy, Facial} Alternate index label; see \textit{Bell's palsy}.

\medterm{Pamidronic Acid} An intravenous bisphosphonate that inhibits osteoclast-mediated bone resorption and is used for selected hypercalcemia, bone complications of cancer, and Paget disease.


\medterm{Panacinar Emphysema} A form of emphysema in which air sacs and small airways are damaged throughout affected areas of the lung. It can cause breathlessness and reduced lung function and is strongly associated with severe alpha-1 antitrypsin deficiency.

\synonyms
Panlobular Emphysema; Lower Lobe Emphysema


\medterm{Pancoast Syndrome} Symptoms caused by a cancer at the top of the lung growing into nearby nerves or other tissues. It may cause severe shoulder or arm pain, hand weakness, numbness, and drooping of one eyelid with a smaller pupil.

\medterm{Pancoast Tumor} A lung cancer at the top of the lung that can grow into nearby nerves, blood vessels, or the spine. It may cause severe shoulder or arm pain, hand weakness, or drooping of one eyelid.

\medterm{Pancolitis} Inflammation affecting the entire large intestine. It may result from ulcerative colitis, infection, reduced blood flow, or another inflammatory condition and can cause diarrhea, abdominal pain, fever, or bleeding.
\medterm{Pancreas} An organ behind the stomach that makes digestive juices and hormones. Its digestive enzymes enter the small intestine, while insulin and other hormones help control blood sugar and digestion.

\medterm{Pancreas Divisum} A birth-related difference in which two drainage tubes of the pancreas do not join normally. Most people have no symptoms, but some develop repeated pancreatitis because pancreatic fluid does not drain freely.

\synonyms
Divisum Pancreas (Pancreas Divisum)

\medterm{Pancreas Enzyme} An enzyme produced by the pancreas, such as amylase, lipase, or proteases, that helps digest carbohydrate, fat, or protein in the small intestine.

\medterm{Pancreas Infection} Infection involving the pancreas or a pancreatic collection, usually occurring with pancreatitis, surgery, trauma, or spread from nearby infection.

\medterm{Pancreas Inflammation} Alternate index label; see \textit{Pancreatitis}.

\medterm{Pancreas Inflammation (Pancreatitis)} Alternate index label for Pancreatitis.

\medterm{Pancreas Inflammatory Cysts (Pancreatic Cysts)} Alternate index label for Pancreatic Cysts.

\medterm{Pancreas Intraductal Papillary Mucinous Neoplasm (Pancreatic Cysts)} Alternate index label for Pancreatic Cysts.

\medterm{Pancreas IPMN (Pancreatic Cysts)} Alternate index label for Pancreatic Cysts.

\medterm{Pancreas Mucinous Cyst Adenomas (Pancreatic Cysts)} Alternate index label for Pancreatic Cysts.

\medterm{Pancreas Neoplasm} An abnormal growth in the pancreas that may be noncancerous, precancerous, or cancerous. Its effects depend on its location and may include pain, jaundice, digestive problems, or diabetes.

\medterm{Pancreas Non-Inflammatory Cysts (Pancreatic Cysts)} Alternate index label for Pancreatic Cysts.

\medterm{Pancreas Pseudocysts (Pancreatic Cysts)} Alternate index label for Pancreatic Cysts.

\medterm{Pancreas Serous Cyst Adenomas (Pancreatic Cysts)} Alternate index label for Pancreatic Cysts.

\medterm{Pancreas Solid Pseudopapillary Tumor (Pancreatic Cysts)} Alternate index label for Pancreatic Cysts.

\medterm{Pancreas Transplant} A procedure that replaces a diseased pancreas with a healthy donor pancreas, usually to restore insulin production in selected people with severe diabetes.

\medterm{Pancreas Transplantation} Alternate index label; see \textit{Pancreas Transplant}.

\medterm{Pancreas, Diseases Of} Disorders affecting the pancreas, including inflammation, cysts, tumors, diabetes, and poor digestive-enzyme production. They may cause abdominal pain, vomiting, weight loss, jaundice, abnormal blood sugar, or digestive problems.
\medterm{Pancreatectomy} Surgery to remove part or all of the pancreas. It may treat a pancreatic tumor, severe injury, or another serious disease; removal of pancreatic tissue can affect digestion and blood-sugar control.

\medterm{Pancreatic Abscess} A pocket of pus in or near the pancreas, usually after severe or infected pancreatitis. It may cause continuing fever, worsening abdominal pain, vomiting, or sepsis and often needs drainage and antibiotics.

\medterm{Pancreatic Adenocarcinoma} The most common cancer of the pancreas, arising from cells in its digestive-juice ducts. It may cause jaundice, abdominal or back pain, weight loss, or no early symptoms; smoking, chronic pancreatitis, diabetes, obesity, and inherited risks can contribute.

\medterm{Pancreatic Cancer} Cancer that begins in the pancreas, most often in cells lining its digestive-juice ducts. It may cause jaundice, abdominal or back pain, weight loss, or digestive problems, but early cancer may cause no symptoms; treatment depends on type and spread.

\synonyms
Cancer, Pancreatic

\medterm{Pancreatic Carcinoma} Pancreatic carcinoma is cancer arising in the pancreas, most often from duct cells, and may cause jaundice, pain, weight loss, or digestive problems.

\medterm{Pancreatic Cysts} Fluid-filled sacs in or around the pancreas. Some follow inflammation or injury, while others are growths that can become cancerous; imaging and sometimes fluid or tissue testing guide follow-up or treatment.

\synonyms
Cyst, Pancreatic
Cysts Pancreatic Inflammatory (Pancreatic Cysts)
Cysts True Pancreas (Pancreatic Cysts)
Intraductal Papillary Mucinous Neoplasm Pancreas (Pancreatic Cysts)
IPMN Pancreas (Pancreatic Cysts)
Mucinous Cyst Adenomas Pancreas (Pancreatic Cysts)
Pancreas Inflammatory Cysts (Pancreatic Cysts)
Pancreas Intraductal Papillary Mucinous Neoplasm (Pancreatic Cysts)
Pancreas IPMN (Pancreatic Cysts)
Pancreas Mucinous Cyst Adenomas (Pancreatic Cysts)
Pancreas Non-Inflammatory Cysts (Pancreatic Cysts)
Pancreas Pseudocysts (Pancreatic Cysts)
Pancreas Serous Cyst Adenomas (Pancreatic Cysts)
Pancreas Solid Pseudopapillary Tumor (Pancreatic Cysts)
Serous Cyst Adenomas Pancreas (Pancreatic Cysts)

\medterm{Pancreatic Diseases} Disorders affecting the pancreas’s digestive or hormone-producing functions, including inflammation, cysts, tumors, duct blockage, and diabetes-related damage. They may cause abdominal pain, vomiting, weight loss, jaundice, abnormal blood sugar, or poor digestion.

\medterm{Pancreatic Duct} The main tube carrying digestive fluid from the pancreas into the first part of the small intestine. It often joins the bile duct before opening into the intestine; blockage can cause pain, jaundice, or pancreatitis.

\medterm{Pancreatic Ductal Adenocarcinoma} The most common exocrine pancreatic malignancy, arising from ductal epithelium and often presenting late with painless jaundice, weight loss, pain, or pancreatitis. Diagnosis and staging rely on imaging, tissue sampling when needed, and multidisciplinary assessment.

\medterm{Pancreatic Elastase} A stool test that estimates how well the pancreas makes digestive enzymes. A low result supports pancreatic insufficiency, which can cause greasy stools, weight loss, and vitamin deficiency, but watery stool and other factors can affect interpretation.

\medterm{Pancreatic Enzyme} A digestive substance made by the pancreas and released into the small intestine. Amylase breaks down starch, lipase breaks down fat, and proteases break down protein; bicarbonate helps neutralize stomach acid.

\medterm{Pancreatic Exocrine Insufficiency} Alternate index label for Exocrine Pancreatic Insufficiency.

\medterm{Pancreatic Fistula} A pancreatic fistula is an abnormal passage through which pancreatic juice leaks from the pancreas, often after trauma or surgery, causing fluid collections, ascites, or skin drainage.

\medterm{Pancreatic Hormones} Hormones produced by pancreatic islet cells, including insulin, glucagon, somatostatin, and pancreatic polypeptide, which regulate glucose, digestion, and energy metabolism.

\medterm{Pancreatic Insufficiency} Failure of the pancreas to produce or deliver enough digestive enzymes or other secretions for normal digestion and absorption. It may cause greasy stools, diarrhea, bloating, weight loss, or vitamin deficiency; exocrine insufficiency specifically refers to inadequate digestive enzymes.

\medterm{Pancreatic Islet} A pancreatic islet is a cluster of endocrine cells that releases hormones such as insulin, glucagon, and somatostatin to regulate blood glucose and digestion.


\medterm{Pancreatic Juice} Pancreatic juice is an alkaline digestive secretion containing bicarbonate and enzymes such as amylase, lipase, and proteases, delivered into the duodenum.

\medterm{Pancreatic Mass} A lump or abnormal area seen in the pancreas on an imaging test. It may be a cyst, inflammation, scar, or tumor; further scans, blood tests, or a tissue sample may be needed to determine its cause.
\medterm{Pancreatic Neuroendocrine Tumor} A tumor arising from hormone-related cells in the pancreas. It may release excess hormones or cause no hormone symptoms and can grow slowly or aggressively; its behavior depends on size, spread, cell features, and growth rate.

\medterm{Pancreatic Neuroendocrine Tumors} Tumors arising from hormone-making cells in the pancreas. Some release excess hormones, such as insulin or gastrin, while others do not. Their behavior depends on size, spread, cell appearance, growth rate, and hormone production.

\synonyms
Islet Cell Cancer

\medterm{Pancreatic Necrosis} Death of pancreatic or nearby fatty tissue, usually during severe acute pancreatitis. The dead tissue may remain sterile or become infected, causing fever, worsening pain, sepsis, and organ failure.

\synonyms
Necrotizing pancreatitis; Acute necrotizing pancreatitis

\medterm{Pancreatic Pseudocyst} A walled-off collection of pancreatic fluid that usually develops after acute or chronic pancreatitis. It has no true cell lining and may cause abdominal pain, vomiting, early fullness, infection, bleeding, or blockage of nearby organs.

\medterm{Pancreatic Surgery} An operation on the pancreas to remove diseased tissue, drain a collection, repair an injury, or treat cancer. The type of surgery depends on the problem and its location.
\medterm{Pancreatic Trauma} Injury to the pancreas caused by blunt or penetrating trauma. Assessment may require CT, laboratory testing, and observation or intervention according to the injury.

\medterm{Pancreaticobiliary Cytology} Cytologic evaluation of pancreatic and biliary lesions, usually from fine-needle aspiration, brushings, or bile-duct sampling. It is used to assess neoplasia, inflammation, and obstructive lesions alongside imaging and histology.

\medterm{Pancreaticoduodenectomy} Pancreaticoduodenectomy, or the Whipple procedure, removes the head of the pancreas and nearby parts of the digestive tract to treat selected tumors and other serious disease.

\medterm{Pancreatitis} Inflammation of the pancreas that may be acute or chronic. Causes include gallstones, alcohol, medicines, metabolic disorders, and other conditions. Acute disease may cause severe upper-abdominal pain, vomiting, fluid loss, infection, tissue death, or organ failure; chronic disease can permanently scar the pancreas and cause poor digestion, weight loss, or diabetes.

\synonyms
Inflamed Pancreas
Pancreas Inflammation
Pancreas Inflammation (Pancreatitis)

\medterm{Pancreatitis - Children} Pancreatitis in children is inflammation of the pancreas causing abdominal pain, nausea, vomiting, and elevated pancreatic enzymes from structural, genetic, medication, infectious, or other causes.


\medterm{Pancreatoblastoma} A rare cancer of the pancreas that occurs mainly in children. It may cause an abdominal lump, pain, vomiting, weight loss, or hormone-related symptoms and requires specialist imaging, biopsy, staging, and treatment.

\medterm{Pancreatography} Imaging of the tubes that carry digestive juices from the pancreas, usually after contrast dye is introduced. It can show narrowing, blockage, leaks, stones, or other changes and may be done during ERCP or with MRI; ERCP also carries a risk of pancreatitis.

\medterm{Pancreazymin} An older name for cholecystokinin, a hormone that stimulates pancreatic enzyme secretion and gallbladder contraction.
\medterm{Pancrelipase} Pancrelipase is a mixture of pancreatic digestive enzymes used as replacement therapy for exocrine pancreatic insufficiency, including that caused by cystic fibrosis or pancreatic disease.

\medterm{Pancuronium} A long-acting medicine that temporarily prevents skeletal muscles from contracting during anesthesia or assisted breathing. It causes paralysis but does not relieve pain or cause unconsciousness and must be given with reliable airway and breathing support.

\medterm{Pancytopenia} A reduction in red blood cells, white blood cells, and platelets at the same time. It can cause tiredness, infections, and bleeding and may result from medicines, infection, immune disease, or a problem in the bone marrow.
\medterm{PANDAS} A proposed condition in which obsessive-compulsive symptoms or movement tics begin suddenly or worsen after a streptococcal infection in a child. The diagnosis remains debated, so clinicians must also assess common neurologic, psychiatric, and infectious causes.

\medterm{Pandemic} An epidemic that spreads across multiple countries or continents and affects many people. It may occur when an infectious agent spreads widely among populations with limited previous immunity.

\medterm{Panendoscopy} Examination of several connected areas of the throat and upper swallowing and breathing passages with a viewing instrument, often including the throat, voice box, and swallowing tube. Tissue samples may be taken during the procedure.
\medterm{Paneth Cell} A cell at the base of a small-intestinal gland that releases substances helping control bacteria and other microbes in the gut. It supports the intestinal lining’s normal defense system.
\medterm{Panhypopituitarism} Failure of the pituitary gland to make most or all of its important hormones. It can cause tiredness, low blood pressure, infertility, growth or thyroid problems, and inability to respond normally to illness; lifelong hormone replacement may be needed.

\medterm{Panic} Panic is sudden intense fear or distress with autonomic symptoms such as racing heart, sweating, trembling, breathlessness, or a sense of impending danger.

\medterm{Panic Attack} A sudden episode of intense fear or discomfort, often with racing heart, shortness of breath, trembling, or dizziness.

\medterm{Panic Attacks And Panic Disorder} Panic attacks are sudden episodes of intense fear; panic disorder involves recurrent attacks and persistent concern about having more.

\medterm{Panic Disorder} Repeated unexpected panic attacks followed by ongoing worry about more attacks or major changes in behavior to avoid them. Attacks peak within minutes and may cause a racing heart, breathlessness, trembling, dizziness, or fear of dying.

\medterm{Panlobular Emphysema} A type of permanent lung damage affecting the whole small air-sac unit, often more severe in the lower lungs. It is strongly associated with alpha-1 antitrypsin deficiency and can cause breathlessness, cough, and reduced exercise ability.

\medterm{Panmyelopathy} A term for widespread disease affecting the bone marrow and sometimes the spinal cord. The exact meaning depends on context and may involve reduced production of several blood-cell types together with neurologic problems.
\medterm{Panniculectomy} Surgical removal of a hanging apron of excess lower-abdominal skin and fat, usually after substantial weight loss; it is distinct from a full abdominoplasty.

\medterm{Panniculitis} Inflammation of the layer of fat beneath the skin, usually causing tender red or purple lumps, often on the legs. Causes include infection, immune disease, medicines, injury, and pancreatic disease; treatment depends on the cause.

\medterm{Panniculitis Idiopathic Nodular (Weber-Christian Disease)} Alternate index label for Weber-Christian Disease.

\medterm{Panniculus} A panniculus is a fold or layer of subcutaneous fatty tissue, especially the lower abdominal apron that may overhang the pubic region.

\medterm{Pannus} An abnormal layer of inflamed tissue that grows over a joint or the clear front surface of the eye. It can damage cartilage or reduce vision and may occur with rheumatoid arthritis, trachoma, or other chronic inflammation.

\medterm{Panophthalmitis} Severe infection and inflammation involving the entire eye and its surrounding tissues. It can cause intense pain, redness, swelling, discharge, and major vision loss and requires emergency treatment to control infection and protect life and sight.
\medterm{Panoramic Radiograph} An extraoral radiograph displaying both jaws, teeth, temporomandibular regions, and
surrounding structures in one image. It is useful for surveying dentition, development,
impaction, and gross osseous pathology but has less fine detail than intraoral
radiographs.

\medterm{Panretinal Photocoagulation} A laser treatment that places many small burns in the outer retina to reduce the abnormal blood-vessel growth caused by severe retinal disease, especially diabetic retinopathy. It can preserve central vision but may reduce side or night vision.

\medterm{Pantothenic Acid} Vitamin B5, needed to help cells make energy and produce fatty acids and certain hormones.
\medterm{Pantothenic Acid And Biotin} Two B vitamins involved in energy metabolism, fatty-acid synthesis, and other enzyme reactions; deficiency is uncommon but can occur with severe malnutrition, malabsorption, or specific inherited disorders.

\medterm{Panuveitis} Inflammation affecting all the main internal layers around the eye, including the iris, the gel inside the eye, and the tissues near the retina. It can cause pain, redness, light sensitivity, floaters, and vision loss and needs urgent eye care.

\medterm{Pap Smear} A screening test in which cells are collected from the cervix and examined for early changes that could become cervical cancer or for cancer itself. It may be done alone or with a test for high-risk human papillomavirus (HPV).

\medterm{Pap Test} A screening test for cervical cancer and precancerous changes in which cells are collected from the cervix and examined under a microscope. It may also reveal inflammation or some infections, but an abnormal result does not by itself diagnose cancer.

\medterm{Papain} A protein-breaking enzyme found in papaya latex. It is used in some foods and products, but evidence for many medicinal claims is limited; it can irritate tissues or cause allergic reactions.
\medterm{Papanicolaou (Pap) Test} A screening test that collects cells from the cervix and examines them for abnormal changes. It can detect precancerous cells and cancer early, allowing treatment before serious disease develops.
\medterm{Papillary Fibroelastoma} A usually noncancerous, small growth attached to a heart valve. It often causes no symptoms but can release a clot and lead to a stroke or other blockage, especially when it is mobile; treatment depends on its size, location, and risk.

\medterm{Papillary Muscle} Muscles inside the lower chambers of the heart that connect by thin cords to the heart valves. They tighten when the heart contracts to keep the valves from turning inside out and prevent blood from flowing backward.

\medterm{Papillary Thyroid Carcinoma} The most common type of thyroid cancer. It often grows slowly and may spread to nearby neck lymph nodes; treatment and outlook depend on the tumor's size, spread, and other features.

\medterm{Papillectomy} Surgical removal of a papilla or papillary growth, such as a tumor near the bile and pancreatic ducts.
\medterm{Papilledema} Swelling of the area where the optic nerve enters the eye because pressure inside the skull is raised. It can threaten vision and may result from bleeding, a mass, infection, or other serious illness.

\medterm{Papilloedema} Swelling of the optic discs caused by raised intracranial pressure.
\medterm{Papilloma} A usually noncancerous growth from a surface lining, often appearing as a wart, skin tag, or small finger-like projection. Some are linked to human papillomavirus; the site, appearance, and tissue examination determine whether further treatment is needed.

\medterm{Papillomatosis Cutis Lymphostatica} Thick, rough, wart-like skin caused by long-standing swelling from poor lymph or vein drainage, usually in the legs. The skin may become cobblestoned and prone to infection; controlling the swelling and protecting the skin are important.

\medterm{Papillomavirus} A papillomavirus, especially human papillomavirus, which can cause warts and some anogenital, oropharyngeal, and other cancers.
\medterm{Papular} A papule: a small, solid, elevated skin lesion <1 cm in diameter (palpable), of various
colors and etiologies.

\medterm{Papular Urticaria} A hypersensitivity reaction to insect bites, especially from fleas, bedbugs, or
mosquitoes, causing recurrent intensely itchy papules or small vesicles, often clustered
on exposed areas in children. Diagnosis is clinical.

\medterm{Papule} A small, solid, raised skin lesion, usually less than 1 centimetre across and without visible fluid. Papules can result from irritation, inflammation, infection, an overgrowth of skin cells, or another skin disorder.
\medterm{PAPVR} Alternate index label; see \textit{Partial anomalous pulmonary venous return}.


\medterm{Paracentesis} A procedure in which a clinician uses a needle to remove fluid from the abdomen. The fluid may be tested to find its cause or removed to relieve pressure and discomfort.

\medterm{Paracetamol Overdose} Taking too much paracetamol, also called acetaminophen, can seriously damage the liver even before symptoms appear. The risk depends on the amount, timing, repeated doses, and health of the person; urgent assessment is needed because an antidote can prevent liver failure.

\medterm{Paracetamol Toxicity} Liver injury caused by excessive acetaminophen (paracetamol) exposure. Severe poisoning may
cause nausea, abdominal pain, liver failure, or death and requires urgent assessment.

\medterm{Paracusia} A hearing disturbance in which sounds are perceived unusually, inaccurately, or differently from their actual qualities. It may occur with hearing loss, inner-ear disease, neurological conditions, medicines, or mental health disorders.
\medterm{Paradichlorobenzene Poisoning} Harm from mothballs or other products containing paradichlorobenzene. It may cause headache, nausea, dizziness, liver injury, or destruction of red blood cells; significant exposure needs poison-control advice.

\medterm{Paraesophageal Hernia} A hiatal hernia in which part of the stomach moves through the diaphragm beside the esophagus.
It may cause obstruction or strangulation even when reflux is absent.

\medterm{Paraesthesia} An unusual skin sensation such as tingling, pins and needles, burning, prickling, or numbness. It can result from pressure on a nerve, nerve disease, poor circulation, medicines, vitamin deficiency, or other causes.
\medterm{Paraffin} A mixture of hydrocarbons used in laboratory tissue processing, topical preparations, and other applications; some forms are flammable.
\medterm{Paraffin Poisoning} Harm from swallowing or breathing paraffin wax or liquid paraffin. Small swallowed amounts may cause stomach upset, but liquid entering the lungs can cause chemical pneumonia; do not induce vomiting.

\medterm{Paraganglioma} A rare tumor arising from specialized nerve-related cells outside the adrenal glands. Some release hormones that cause high blood pressure, sweating, headaches, or a fast heartbeat; inherited risk may require genetic assessment.

\medterm{Parageusia (Dysgeusia)} An altered sense of taste in which food may taste metallic, bitter, salty, unpleasant, or different from usual. Causes include infections, medicines, nutritional deficiencies, head injury, and mouth, nerve, or salivary-gland problems.
\medterm{Paragonimiasis} A parasitic infection caused by lung flukes, usually acquired by eating raw or undercooked freshwater crab or crayfish. It commonly affects the lungs and may cause a long-lasting cough, chest pain, fever, or blood-stained sputum.

\medterm{Paragraphia} A language problem in which a person writes an unintended word, letter, or symbol instead of the one intended. It may occur with damage to language areas of the brain, such as after a stroke, and is assessed with other speaking, reading, and writing abilities.
\medterm{Parainfectious} Describing a problem that occurs during or soon after an infection, often because the immune system reacts to the infection rather than because germs directly invade the affected tissue.

\medterm{Parainfluenza} Parainfluenza viruses are respiratory viruses that commonly cause colds, croup, bronchitis, or pneumonia, especially in children; illness ranges from mild to severe airway disease.

\medterm{Paraldehyde} An older sedative and anticonvulsant medicine with limited current use.
\medterm{Paralysis} Loss of voluntary muscle movement in part or all of the body due to nervous-system or muscle disease.
\medterm{Paralysis Agitans} An older name for Parkinson disease, characterized by tremor, rigidity, and bradykinesia.
\medterm{Paralysis, Stomach} Stomach paralysis, or gastroparesis, is delayed gastric emptying without a mechanical blockage, causing early fullness, nausea, vomiting, bloating, and variable glucose control.

\medterm{Paralytic Ileus} Failure of intestinal peristalsis without a mechanical blockage, causing abdominal distension, pain, vomiting, and absent or reduced bowel sounds.
\medterm{Paramedical} Relating to healthcare professionals or services that support medical practice, such as physiotherapy, radiography, or emergency medical care.
\medterm{Parametritis} Inflammation or infection of the tissue beside the uterus, often after pelvic infection, childbirth, miscarriage, or surgery. It may cause pelvic pain, fever, tenderness, abnormal discharge, and requires treatment based on its cause.
\medterm{Paramnesia} A memory disturbance involving false, distorted, or inaccurately familiar recollections.
\medterm{Paramyotonia Congenita} A rare inherited muscle disorder in which muscles become stiff and slow to relax, especially after repeated movement or exposure to cold. The face, eyelids, neck, and hands are often affected; temporary weakness may follow an attack. Avoiding cold and strenuous triggers helps, and medicines may reduce symptoms.

\synonyms
Eulenburg Disease; Paramyotonia Congenita of Eulenburg

\medterm{Paranasal Sinuses} Air-filled spaces in the bones around the nose that connect with the nasal passages. Their lining produces mucus, which drains into the nose; inflammation or blockage can cause facial pressure, congestion, and pain.

\medterm{Paraneoplastic Disorder} A health problem caused indirectly by cancer rather than by the tumor invading or spreading to the affected organ. Tumors may release hormones or trigger immune reactions that affect the nerves, skin, blood, joints, or other distant organs.

\medterm{Paraneoplastic Neurologic Syndrome} A nervous-system problem caused by the immune response to a cancer rather than by the tumor pressing on the nerves. It may cause weakness, loss of balance, memory problems, seizures, or abnormal sensation and can occur before the cancer is found.

\medterm{Paraneoplastic Neurological Disorder} Alternate index label; see \textit{Paraneoplastic syndromes of the nervous system}.

\medterm{Paraneoplastic Neurological Syndrome} Alternate index label; see \textit{Paraneoplastic syndromes of the nervous system}.

\medterm{Paraneoplastic Syndrome} A group of symptoms caused by substances or immune reactions from a cancer, rather than by the tumor pressing on or spreading to the affected area. It may cause hormone, nerve, skin, blood, or joint problems; treating the cancer may improve them.

\medterm{Paraneoplastic Syndromes Of The Nervous System} Nervous-system problems caused by an immune reaction to a cancer rather than by the cancer growing into the nerves or brain. They may cause weakness, poor balance, seizures, memory change, or abnormal sensation.

\synonyms
Cancer, Paraneoplastic Syndromes
Paraneoplastic Neurological Disorder
Paraneoplastic Neurological Syndrome

\medterm{Paranoia} A thought process characterized by excessive or irrational suspiciousness and distrust
of others, often manifesting as delusions of persecution, reference, or grandeur.

\medterm{Paranoid} Characterized by persistent or excessive suspicion and distrust, often involving beliefs that others intend harm. Severe or fixed beliefs may occur in a psychotic disorder and need professional assessment.
\medterm{Paranoid Personality Disorder} A longstanding pattern of distrust and suspiciousness in which other people’s motives are interpreted as threatening or deceitful.

\medterm{Paranoid Psychosis} An older term for psychosis dominated by delusions of persecution, reference, or grandeur.
\medterm{Paraparesis} Weakness in both legs caused by disease affecting the spinal cord, brain, nerves, or muscles.
\medterm{Paraphasia} An unintended substitution, omission, or distortion of a sound or word during speech. It commonly occurs with aphasia after brain injury or stroke and may affect naming, conversation, reading, or writing.
\medterm{Paraphilia} A persistent and intense sexual interest involving atypical objects, situations, or activities. A paraphilia by itself is not a mental disorder. It becomes a paraphilic disorder when it causes significant personal distress or problems in daily life, apart from distress caused only by social disapproval, or when acting on it involves harm, risk of harm, or people who do not or cannot consent.

\medterm{Paraphilic Disorder} Alternate index label for Paraphilic Disorders.

\medterm{Paraphilic Disorders} The group of mental disorders in which a paraphilia causes significant personal distress or problems in daily life, apart from distress caused only by social disapproval, or involves harm, risk of harm, or people who do not or cannot consent. The group includes voyeuristic, exhibitionistic, frotteuristic, fetishistic, sexual masochism, sexual sadism, pedophilic, and transvestic disorders, as well as other specified and unspecified paraphilic disorders. A paraphilia without these consequences is not, by itself, a paraphilic disorder.

\medterm{Paraphimosis} Paraphimosis occurs when a retracted foreskin cannot be returned over the glans, causing painful swelling and possible loss of blood flow; it requires urgent reduction.

\medterm{Paraplegia} Paralysis of both legs and usually part of the lower trunk. Paralysis of both legs and often part of the trunk, usually from spinal-cord or neurologic disease.
\medterm{Parapneumonic Effusion} A collection of fluid between the lung and chest wall that develops with pneumonia. Small uncomplicated collections may improve with antibiotics, while infected or loculated fluid can cause persistent fever, chest pain, or breathlessness and may require drainage.

\medterm{Parapneumonic Pleural Effusion} Pleural fluid accumulation associated with bacterial pneumonia. It may be uncomplicated
or progress to a complicated effusion or empyema, requiring imaging, pleural-fluid
analysis, antibiotics, and sometimes drainage.

\medterm{Paraquat} A highly toxic herbicide. Ingestion can cause injury to the mouth and gastrointestinal tract, kidney failure, and delayed lung injury; severe poisoning has a high risk of death.

\medterm{Paraquat Poisoning} A medical emergency caused by swallowing or significant exposure to the herbicide paraquat. It can burn the mouth and digestive tract and later cause lung, kidney, or liver failure; immediate emergency care is essential.

\medterm{Paraseptal Emphysema} A form of emphysema affecting air sacs near the outer surface of the lung. It may cause few symptoms but can form fragile air pockets that rupture and produce a collapsed lung, especially in tall young adults or people who smoke.

\medterm{Parasite} An organism that lives on or inside another organism, called the host, and obtains food or shelter from it. Parasites may cause disease in the skin, intestine, blood, or organs.

\medterm{Parasitemia} The presence and amount of parasites in the bloodstream. It is especially used for malaria and some other blood infections; a higher amount may indicate more severe disease, but interpretation depends on the parasite and the person’s symptoms.

\medterm{Parasitic} Parasitic means relating to an organism that lives in or on a host and obtains nutrients at the host’s expense, sometimes causing disease.

\medterm{Parasitic Infection} An illness caused by a parasite---an organism that lives on or inside another living creature (the host) and survives by taking nutrients at the host's expense. Human parasites belong to three main groups: single-celled protozoa (such as \textit{Plasmodium} causing malaria, \textit{Entamoeba} causing amebic dysentery, and \textit{Giardia} causing giardiasis, which can multiply rapidly inside the body), helminths or parasitic worms (such as roundworms, tapeworms, pinworms, hookworms, and flukes, which are multicellular and can live in tissues for years), and ectoparasites (such as scabies mites, lice, fleas, and ticks that feed on or burrow into the skin). Parasites enter the body through contaminated water or unwashed food, undercooked meat, insect bites, walking barefoot on contaminated soil, or direct physical contact. Symptoms depend on the parasite and can range from watery or bloody diarrhea, stomach cramps, and nausea, to prolonged fevers, extreme fatigue, severe anemia, skin rashes, and chronic organ enlargement or damage. Parasitic infections are diagnosed using stool tests, blood smears, or tissue exams, and are treated with targeted antiparasitic medicines (such as antiprotozoals or worm-clearing anthelmintics). Access to clean drinking water, proper sanitation, thorough cooking, and insect protection are key to prevention.

\synonyms
Parasitic Infections; Parasitic Disease; Parasitosis; Parasitic Infestation

\medterm{Parasitic Worm} A worm that lives in or on a host and uses the host for nourishment. Human infections may affect the intestine, blood, lymph system, lungs, liver, skin, or other tissues and can cause anemia, pain, or organ damage.

\medterm{Parasitology} The branch of microbiology that studies parasites, including protozoa, helminths, and
ectoparasites. Parasitic infections are a major global health problem, particularly in
tropical and subtropical regions. Protozoan parasites include Plasmodium (malaria),
Trypanosoma (sleeping sickness, Chagas disease), and Leishmania (leishmaniasis).

\medterm{Parasitosis} Infestation or infection by a parasite; the term may describe a confirmed parasitic disease or, in delusional parasitosis, a fixed false belief of infestation without evidence.

\medterm{Parasomnia} A parasomnia is an undesirable movement, behavior, emotion, or experience occurring during sleep or transitions between sleep and wakefulness.

\medterm{Parasomnias} Unwanted movements, behaviors, emotions, or experiences that occur during sleep or at the transition between sleep and wakefulness.

\medterm{Parasuicide} Parasuicide is an older term for a suicide attempt or intentional self-harm without fatal outcome; current assessment focuses on self-harm, intent, safety, and urgent support.

\medterm{Parasympathetic Nervous System} The part of the automatic nervous system that supports resting functions such as digestion, saliva, urination, and slowing the heartbeat. It also helps the body conserve energy and recover after activity.

\medterm{Parathyroid} One of the usually four small glands behind the thyroid gland. They release parathyroid hormone, which helps keep calcium and phosphate levels balanced in the blood and bones.

\medterm{Parathyroid Adenoma} A parathyroid adenoma is a usually benign tumor of a parathyroid gland that can secrete excess parathyroid hormone and cause primary hyperparathyroidism with high calcium.

\medterm{Parathyroid Cancer} Parathyroid cancer is a rare malignant tumor of the parathyroid gland that often causes excessive parathyroid hormone and high blood calcium.

\medterm{Parathyroid Carcinoma} A rare malignant tumor of the parathyroid gland that commonly causes excessive parathyroid-hormone secretion, hypercalcemia, kidney stones, bone disease, or related symptoms.

\medterm{Parathyroid Crisis} A life-threatening rise in blood calcium caused by severe overactivity of the parathyroid glands. It may cause extreme thirst, vomiting, constipation, confusion, weakness, abnormal heart rhythm, kidney injury, or coma and requires emergency treatment.

\medterm{Parathyroid Gland} A small gland behind the thyroid that releases parathyroid hormone to regulate calcium and phosphate.

\medterm{Parathyroid Gland Removal} Surgery to remove one or more overactive parathyroid glands, usually for persistent high blood calcium. It can improve symptoms and protect bones or kidneys; low calcium, bleeding, voice-nerve injury, or persistent disease may occur.

\medterm{Parathyroid Glands} Usually four small endocrine glands behind the thyroid that secrete parathyroid hormone to regulate blood calcium and phosphate levels.

\medterm{Parathyroid Hormone} A hormone that regulates calcium and phosphate levels in the blood and supports bone remodeling.
\medterm{Parathyroid Hormone Blood Test} A blood test measuring parathyroid hormone, interpreted with calcium, phosphate, vitamin D, and kidney function to evaluate parathyroid and calcium disorders.

\medterm{Parathyroid Hormone Physiology} When blood calcium falls, the parathyroid glands release more parathyroid hormone. The hormone acts on bone and kidneys and increases active vitamin D, helping raise calcium and maintain the balance needed for nerves, muscles, and bones.

\medterm{Parathyroid Hormone-Related Protein Blood Test} A blood test measuring parathyroid hormone-related protein, used mainly to evaluate humoral hypercalcemia of malignancy and selected disorders of calcium regulation.

\medterm{Parathyroid Hyperplasia} Enlargement of one or more small glands behind the thyroid, causing them to release too much parathyroid hormone. This can raise blood calcium and may occur with kidney disease or inherited multiple-gland disorders.

\medterm{Parathyroid Surgery} Operations on the parathyroid glands, usually performed for selected cases of primary or
secondary hyperparathyroidism. The decision depends on symptoms, calcium levels, kidney involvement, bone disease,
overall health, and current clinical guidance.

\medterm{Parathyroidectomy} Surgery to remove one or more overactive parathyroid glands. It may treat primary or kidney-related hyperparathyroidism when hormone excess causes high calcium, kidney stones, bone disease, or other complications; calcium levels must be monitored afterward.

\medterm{Paratubal Cyst} A usually benign fluid-filled cyst adjacent to the fallopian tube, arising from paratubal or embryologic remnants; most are asymptomatic but large cysts may cause pain or torsion.

\medterm{Paratyphoid Fever} An infection caused by Salmonella Paratyphi bacteria, usually spread through contaminated food or water. It can cause prolonged fever, headache, weakness, abdominal symptoms, and sometimes serious complications without antibiotic treatment.
\medterm{Parenchyma} The working tissue of an organ, distinguished from its supporting framework, blood vessels, and connective tissue. For example, lung parenchyma performs gas exchange, while liver parenchyma processes nutrients and harmful substances.
\medterm{Parenchymal} Parenchymal means relating to the functional cells of an organ, distinguished from its supporting connective tissue, vessels, and framework.

\medterm{Parenchymal Diffuse Lung Disease (Interstitial Lung Disease (Interstitial Pneumonia))} Alternate index label for Interstitial Lung Disease (Interstitial Pneumonia).


\medterm{Parenteral} Parenteral means delivered by a route other than the gastrointestinal tract, commonly by injection or infusion such as intravenous, intramuscular, or subcutaneous administration.

\medterm{Parenteral Nutrition} Nutrition given through a vein when the digestive tract cannot safely absorb enough nutrients. The solution can contain water, glucose, amino acids, fats, vitamins, and minerals. It can cause bloodstream infection, liver problems, or abnormal blood sugar and salts, so it requires monitoring.

\medterm{Paresis} Partial weakness or incomplete paralysis of voluntary muscles. It may affect one limb, one side of the body, both legs, or all four limbs, depending on the nerves, brain, or spinal cord involved.
\medterm{Paresthesia} An unusual sensation such as tingling, prickling, burning, pins and needles, or numbness. Brief symptoms often follow pressure on a nerve, while persistent or spreading symptoms need assessment for nerve or circulation problems.

\medterm{Paricalcitol} Paricalcitol is a vitamin-D receptor activator used to reduce secondary hyperparathyroidism in selected chronic kidney disease, with monitoring for calcium and phosphate changes.

\medterm{Parietal} Relating to a wall of a body cavity or to the parietal bone or lobe of the skull and brain.
\medterm{Parietal Bone} One of two broad bones forming much of the skull's sides and roof. They protect the brain, meet along the top midline, and connect with several neighboring skull bones.
\medterm{Parietal Lobe} A part of the brain that processes touch, pain, temperature, body position, and spatial awareness. Damage may cause numbness, difficulty recognizing objects by touch, poor coordination, or trouble understanding where parts of the body are.

\medterm{Parietal Pericardium} The outer serous layer of the pericardium that lines the inner surface of the fibrous sac surrounding the heart.

\medterm{Parietal Pleura} The pleural layer lining the chest wall, diaphragm, and mediastinum; it is sensitive to pain when inflamed.


\medterm{Parinaud Oculoglandular Syndrome} Conjunctivitis with nearby lymph-node swelling, usually caused by infections such as Bartonella, tularemia, or certain viruses.

\medterm{Parinaud Syndrome} A disorder caused by dysfunction of the dorsal midbrain, characterized by impaired upward gaze, abnormal pupil
responses, and convergence-retraction eye movements. Causes include tumors, stroke,
multiple sclerosis, and hydrocephalus.

\medterm{Paris System For Reporting Urinary Cytology} A standardized framework for urinary cytology that emphasizes detection of high-grade urothelial carcinoma and uses defined categories to improve diagnostic consistency and clinical communication.

\medterm{Parity} Parity is the number of pregnancies that reached a viable gestational age, regardless of whether the infants were born alive or stillborn; it differs from gravidity.

\medterm{Park-Bench Position} A modified lateral position in which the upper body is partly rotated toward prone; it is used for selected posterior brain and neck procedures.

\medterm{Parkinson Disease} Alternate spelling; see \textit{Parkinson's Disease}.



\medterm{Parkinson's Disease} A progressive disorder of the nervous system that mainly affects movement. It develops when nerve cells in a part of the brain called the substantia nigra become damaged or die. These cells normally make dopamine, a chemical messenger that helps produce smooth, purposeful movement. Reduced dopamine can cause tremor, slowness of movement, muscle stiffness, balance or walking problems, a softer voice, and smaller handwriting. Non-movement symptoms may include constipation, loss of smell, tiredness, sleep problems, depression, or changes in thinking. Symptoms and their rate of progression vary from person to person.

\synonyms
Parkinson Disease
Parkinsons Disease (Parkinson's Disease)

\medterm{Parkinsonism} A group of movement problems involving slowness, stiffness, shaking, or balance difficulty. Parkinson disease is one cause, but medicines, other brain disorders, and several less common conditions can produce similar symptoms.

\medterm{Parkinsons Disease (Parkinson's Disease)} Alternate index label for Parkinson's Disease.


\medterm{Paromomycin} Paromomycin is an aminoglycoside antibiotic used mainly in the intestine for selected parasitic infections and sometimes for other infections under specialist care.

\medterm{Paronychia} Infection or inflammation of the skin around a fingernail or toenail, causing pain, redness, swelling, and sometimes pus. It may follow nail biting, injury, a hangnail, or repeated water exposure.

\medterm{Parosteal Osteosarcoma} Parosteal osteosarcoma is a usually lower-grade bone cancer that grows on the outer surface of a bone, often near the back of the thigh.

\medterm{Parotid Gland} The largest major saliva-producing gland, located in front of and below each ear. It releases saliva into the mouth near the upper back teeth and can become swollen from infection, blockage, inflammation, or a tumor.

\medterm{Parotid Neoplasm} An abnormal growth in a parotid salivary gland in front of or below the ear. Most are noncancerous, but some are cancerous; a painless lump, facial weakness, pain, or rapid growth needs assessment.

\medterm{Parotid Tumor} An abnormal growth in a parotid salivary gland. Most parotid tumors are noncancerous, but some are malignant and may cause a lump, pain, or facial-nerve weakness.

\medterm{Parotid Tumors} Abnormal growths arising in the parotid salivary glands, located in front of and below the ears.

\medterm{Parotidectomy} Surgery to remove part or all of a parotid salivary gland, usually because of a tumor. Possible complications include facial-nerve weakness, bleeding, infection, saliva leakage, numbness, or sweating while eating.

\medterm{Parotitis} Inflammation or infection of one or both parotid salivary glands, which lie in front of and below the ears. It can cause painful swelling, fever, dry mouth, or pus near the upper back teeth. Causes include viruses such as mumps, bacteria, blocked ducts, dehydration, and autoimmune disease.
\medterm{Parotitis (Mumps)} Alternate index label for Mumps.

\medterm{Paroxetine Hydrochloride} Paroxetine hydrochloride is an SSRI antidepressant used for depression and several anxiety-related disorders; it can cause sexual side effects and withdrawal symptoms if stopped abruptly.

\medterm{Paroxysm} A sudden, often brief attack or burst of symptoms or body activity that may recur. Examples include an episode of pain, coughing, rapid heartbeat, seizure, or involuntary movement; the cause depends on the symptom.
\medterm{Paroxysmal Atrial Tachycardia} Sudden episodes of rapid, regular atrial rhythm that start and stop abruptly, often producing palpitations, lightheadedness, chest discomfort, or shortness of breath.

\medterm{Paroxysmal Cold Hemoglobinuria} A rare immune disorder in which cold temperatures trigger destruction of red blood cells when the blood warms again. It can cause sudden anemia, fever, back pain, jaundice, or dark urine, often after an infection, and needs medical assessment.

\medterm{Paroxysmal Hemicrania} A rare headache disorder causing brief, repeated attacks of severe one-sided head or facial pain.
\medterm{Paroxysmal Nocturnal Dyspnea} Sudden breathlessness that wakes a person after they have been asleep for a while, often improving after sitting or standing. It commonly occurs with heart failure but may have other heart or lung causes and needs assessment.

\medterm{Paroxysmal Nocturnal Haemoglobinuria} A rare acquired blood disorder in which abnormal blood-forming cells make red blood cells vulnerable to immune destruction. It can cause anemia, dark urine, blood clots, low blood counts, kidney injury, or bone-marrow failure.
\medterm{Paroxysmal Nocturnal Hemoglobinuria} A rare acquired blood disorder in which red cells are destroyed by the immune system. It can cause dark urine, anemia, blood clots in unusual locations, low blood counts, and kidney or other organ problems.

\medterm{Paroxysmal Nocturnal Hemoglobinuria (Paroxysmal Nocturnal Hemoglobinuria PNH)} Alternate index label for Paroxysmal Nocturnal Hemoglobinuria PNH.

\medterm{Paroxysmal Supraventricular Tachycardia} Repeated episodes of a very fast, usually regular heartbeat that start and stop suddenly and begin in the upper chambers or connecting electrical tissues of the heart. It may cause pounding heartbeat, dizziness, chest discomfort, breathlessness, or fainting.
\synonyms
PSVT (Paroxysmal Supraventricular Tachycardia (PSVT))
Supraventricular Tachycardia Paroxysmal (Paroxysmal Supraventricular Tachycardia (PSVT))
Tachycardia Paroxysmal Atrial (Paroxysmal Supraventricular Tachycardia (PSVT))
Tachycardia Paroxysmal Supraventricular (Paroxysmal Supraventricular Tachycardia (PSVT))

\medterm{Pars} A part or portion of an anatomical structure, used in names such as pars interarticularis.
\medterm{Pars Compacta} A dopamine-producing part of the substantia nigra, a movement-control area deep in the brain. Loss of these nerve cells contributes to the slowness, stiffness, and tremor of Parkinson disease.

\medterm{Pars Planitis} A form of inflammation inside the eye, near the edge of the retina. It may cause floaters, blurred vision, and swelling or other retinal complications and is treated according to its severity and cause.


\medterm{Partial} Partial means incomplete or involving only a portion, such as partial loss, partial response, partial thickness, or partial seizure; the specific meaning depends on context.

\medterm{Partial Agonist} A medicine that attaches to a cell receptor and activates it, but produces a smaller effect than a full activator even when all available receptors are occupied. It can reduce the effect of a stronger activator at the same receptor.

\medterm{Partial Anomalous Pulmonary Venous Connection} Alternate index label; see \textit{Partial anomalous pulmonary venous return}.

\medterm{Partial Breast Brachytherapy} Radiation treatment placed inside or next to the area where a breast cancer was removed, usually after breast-conserving surgery. It treats selected early cancers over a shorter course but may cause skin changes, breast firmness, infection, or delayed healing.

\medterm{Partial Breast Radiation Therapy - External Beam} Radiation delivered from outside the body to the part of the breast where a cancer was removed. It may be suitable for selected early cancers and can cause temporary skin irritation, swelling, fatigue, or later breast firmness.

\medterm{Partial Cystectomy} Surgery to remove part of the urinary bladder, usually for a tumor or localized disease while leaving the rest of the bladder in place. The remaining bladder may need monitoring, and urinary or surgical complications can occur.

\medterm{Partial Denture} A removable replacement for some missing teeth, supported by remaining teeth, gums, or both. It may use metal or plastic and can improve chewing, speech, appearance, and tooth stability.

\medterm{Partial Hysterectomy} A common but imprecise term usually referring to supracervical or subtotal hysterectomy, in which the uterine body is removed while the cervix remains. The fallopian tubes and ovaries may or may not be removed separately.

\medterm{Partial Knee Replacement} Surgery replacing only the damaged compartment of the knee with an artificial surface while preserving healthy bone, ligaments, and cartilage. It may relieve pain from localized arthritis; infection, clots, stiffness, loosening, or later arthritis elsewhere can occur.

\medterm{Partial Nephrectomy} Partial nephrectomy removes a diseased part of a kidney while preserving the remaining healthy tissue.
\medterm{Partial Remission} A reduction in signs, symptoms, or measurable disease burden that does not meet criteria for complete remission; the definition depends on the disease and response criteria used.

\medterm{Partial Response} A measurable improvement after treatment in which a tumor or other disease becomes smaller or less active but does not disappear completely. The exact amount of improvement required is defined by the cancer type, test method, or study criteria.

\medterm{Partial Seizure} A seizure beginning in a localized region of one cerebral hemisphere, now called a focal seizure, with symptoms determined by the involved network and awareness either preserved or impaired.

\medterm{Partial Thromboplastin Time} A blood test measuring how long a clot takes to form through several clotting pathways. It helps investigate unusual bleeding and monitor unfractionated heparin; abnormal results may reflect a clotting-factor problem, medicine effect, or another illness.

\medterm{Partial Volume Artifact} A scan distortion that occurs when one small image area contains different tissues. Their signals are blended, which can blur borders or hide a small finding; thinner slices or higher-resolution images may reduce the effect.

\medterm{Partial-Volume Effect} A scan limitation in which a small structure and the surrounding tissue are measured together, making the structure appear blurred or its activity seem too high or low. Thinner slices or correction methods may reduce the effect.

\medterm{Participation} A person’s involvement in life situations and activities, a central concept in rehabilitation and the International Classification of Functioning that extends beyond impairment alone.

\medterm{Particle (Vascular)} A small material placed into a blood vessel to block it during a procedure. Particles may be used to control bleeding or reduce blood flow to a tumor, fibroid, or other abnormal area and must be directed carefully to avoid harming healthy tissue.

\synonyms
Vascular embolization particle

\medterm{Particle Size} Particle size is the diameter or size distribution of particles in a powder, suspension, aerosol, or formulation, influencing dissolution, absorption, stability, and delivery.

\medterm{Particulate Matter} Small solid or liquid particles suspended in air. Fine particles can penetrate the lungs and increase
cardiovascular and respiratory risk.

\medterm{Partograph (Partogram)} A graphic record of cervical change, fetal descent, contractions, and maternal or fetal observations during
labor. It can support communication and recognition of abnormal progress, but findings
must be interpreted with current labor guidance and the overall clinical situation.

\synonyms
Partogram

\medterm{Parturition} The process of childbirth, including cervical ripening, labor contractions, delivery of the fetus, and expulsion of the placenta.

\medterm{Parturition (Childbirth)} The process of giving birth, encompassing three stages: (1) First stage: from the onset
of regular uterine contractions (cervical effacement and dilation) to full dilation (10
cm). Latent phase—slow dilation (0--4 cm) over hours, variable duration. Active
phase—more rapid dilation (4--10 cm), typically 1--1.5 cm/hour in nulliparous women.

\synonyms
Childbirth

\medterm{Parvimonas Micra} An anaerobic gram-positive bacterium found in the mouth and gastrointestinal tract that can cause dental, bone, joint, or bloodstream infections.

\medterm{Parvoviridae} A family of very small DNA viruses; human parvovirus B19 can cause fifth disease, anemia, or fetal complications.

\medterm{Parvovirus} A small DNA virus; parvovirus B19 is the main human pathogen and can cause rash, joint symptoms, or red-cell aplasia.

\medterm{Parvovirus B19} A human virus spread mainly through respiratory droplets. It commonly causes a mild rash illness in children, but can cause joint pain or severe anemia and may harm a fetus during pregnancy.

\medterm{Parvovirus B19 Infection} An infection that may cause fever, a slapped-cheek rash, joint pain, or temporary reduction in red-blood-cell production. It can be serious in people with certain blood disorders or weakened immunity and requires pregnancy assessment after significant exposure.

\medterm{Parvovirus Infection} An infection caused by parvovirus B19 that commonly causes mild fever and rash but may be serious in some people.

\synonyms
Slapped Cheek Disease

\medterm{Passive Exercise} Movement of a person's joint or limb by another person, a machine, or gravity without active muscle effort from the patient. It can maintain flexibility and reduce stiffness when illness or weakness limits movement, but forceful movement can injure tissues.

\medterm{Passive Immunity} Protection produced by transferred antibodies rather than by the person's own immune response, such as maternal antibodies or immunoglobulin treatment.
\medterm{Passive Immunization} Protection produced by giving preformed antibodies, such as immunoglobulin or antivenom, which acts quickly but is temporary because the recipient does not generate durable immune memory.


\medterm{Pasteurella} A genus of bacteria carried by animals that can cause wound infections after bites and other human disease.
\medterm{Pasteurella Aerogenes} Pasteurella aerogenes is a rare Pasteurella species; its recovery from a clinical sample should be interpreted with the source and evidence of infection.


\medterm{Pasteurella Canis} Pasteurella canis is a bacterium found in dogs and cats that can cause wound or soft-tissue infection after an animal bite.

\medterm{Pasteurella Multocida} A bacterium commonly carried by animals that can cause wound infections after bites or scratches and,
less often, respiratory or bloodstream infection.

\medterm{Pasteurellaceae} Pasteurellaceae is a family of bacteria that includes Pasteurella and related animal-associated organisms capable of causing bite-wound and respiratory infections.

\medterm{Pasteurellosis} Infection caused by Pasteurella bacteria, commonly introduced through animal bites or scratches and causing cellulitis, abscesses, or occasionally invasive disease.

\medterm{Pasteurization} A heat treatment reducing microbial load in liquids by heating below boiling point to
inactivate vegetative pathogens and spoilage organisms while preserving quality.

\medterm{Patau Syndrome} A genetic condition caused by an extra copy of chromosome 13. It can cause severe problems involving brain development, growth, the eyes, heart, kidneys, hands or feet, and the lip or palate. Many affected pregnancies end before birth, and surviving infants need complex specialist care.

\medterm{Patch} A flat, circumscribed area of altered skin color larger than approximately 1 cm. It is
the larger counterpart of a macule and may be caused by pigmentary, vascular,
inflammatory, or neoplastic disease.

\medterm{Patch (Vascular)} A piece of biological or synthetic material sewn onto a blood vessel to widen or repair it.

\synonyms
Vascular patch; angioplasty patch

\medterm{Patch Graft} A piece of natural or synthetic material sewn onto a blood vessel to widen or repair it during patch angioplasty. Materials may include a patient's tissue, bovine pericardium, or synthetic PTFE.

\medterm{Patch Test} A skin test used to identify delayed allergic contact reactions to substances applied under adhesive patches.

\medterm{Patch Testing} Patch testing identifies allergens causing allergic contact dermatitis. The standard
series includes 35 allergens. Additional custom trays include occupational and personal
allergens. Patches are applied to the upper back for 48 hours and removed. Reading at 48
hours (initial) and 72 to 96 hours (delayed) identifies positive reactions.

\medterm{Patches} Patches are flat, distinct areas of altered skin color larger than a macule; their appearance, distribution, and associated symptoms help identify the cause.

\medterm{Patchy Hair Loss (Alopecia Areata)} Alternate index label for Alopecia Areata.

\medterm{Patchy Skin Color} Areas of skin that are lighter, darker, redder, or otherwise different in color from the surrounding skin. Causes include vitiligo, birthmarks, inflammation, infection, sun exposure, injury, and changes in blood flow; new or changing patches should be assessed.

\medterm{Patella} The kneecap, a triangular bone at the front of the knee. It protects the joint and helps the thigh muscles straighten the leg by guiding and improving the force of the knee tendon.

\medterm{Patella (Kneecap)} The largest sesamoid bone in the human body, embedded within the quadriceps tendon
at the front of the knee. It improves the mechanical efficiency of knee extension, protects
the front of the joint, and distributes forces across the femoral trochlea.

\synonyms
Kneecap

\medterm{Patellar Dislocation} Complete displacement of the kneecap from the femoral groove, usually laterally, causing acute pain, deformity, swelling, and difficulty extending the knee.

\medterm{Patellofemoral Instability} A tendency of the patella to move abnormally or dislocate from the femoral trochlea.
Risk factors include abnormal trochlear shape, a high-riding patella, ligament insufficiency,
altered limb alignment, and previous dislocation.

\medterm{Patellofemoral Pain Syndrome} Pain around or behind the kneecap, often worsened by running, stairs, squatting, or sitting with the knee bent. It commonly results from repeated loading, muscle weakness, movement changes, or kneecap tracking problems.

\synonyms
Runner's Knee

\medterm{Patellofemoral Syndrome} Pain around or behind the kneecap related to movement and loading of the patellofemoral joint, often without a major structural injury.

\medterm{Patent (Adjective)} Patent means open, unobstructed, or functioning, as in a patent airway, vessel, duct, or congenital fetal channel.

\medterm{Patent Ductus Arteriosus} A blood-vessel connection between the heart’s main lung artery and the aorta that remains open after birth instead of closing. It can send extra blood to the lungs and strain the heart; effects range from none to heart failure.

\medterm{Patent Foramen Ovale} A persistent flap-like opening between the atria that normally permits fetal blood flow
and usually closes after birth. It is often an incidental finding but may allow a clot to pass from the venous to the
arterial circulation in selected circumstances.

\medterm{Patent Urachus} A patent urachus is persistence of the fetal channel between the bladder and umbilicus after birth, allowing urine to drain from the umbilicus and predisposing to infection.

\medterm{Patent Urachus Repair} Surgery to close an abnormal open connection between the bladder and the navel that remains from fetal development. It may stop urine leaking from the umbilicus and prevent infection; treatment depends on the child's anatomy and health.

\medterm{Paternal Uniparental Disomy 14 (UDP 14 Pat) Thoracic Dysplasia} A rare genetic condition resulting from inheritance of both copies of chromosome 14 from
the father, leading to dysregulated expression of imprinted genes on chromosome 14. It
presents with severe thoracic hypoplasia, pulmonary hypoplasia, and characteristic
facial features. It is a lethal condition in most cases.

\synonyms
UDP 14 pat

\medterm{Paternity Testing} A DNA test that assesses whether a person is the biological parent of a child. Samples are usually collected from the child and possible parent, and the result reports either exclusion or a calculated probability of biological relationship.

\medterm{Pathogen} An organism or infectious agent that can enter the body and cause disease. Examples include viruses, bacteria, fungi, and parasites.

\medterm{Pathogenesis} The way a disease begins and develops, including its causes, changes in cells or tissues, and progression over time. Understanding pathogenesis helps explain symptoms, complications, and possible treatment targets.
\medterm{Pathogenic} Pathogenic means capable of causing disease in a susceptible host; pathogenicity depends on the organism or factor, dose, route, and host defenses.

\medterm{Pathogenic Variant} A change in DNA with strong evidence that it causes or contributes to a disease. Its importance depends on the gene, the person’s symptoms and family history, and whether one or both copies of the gene are affected.

\medterm{Pathogenicity} The capacity of a microorganism to cause disease in a susceptible host. Pathogenicity is
a qualitative property, whereas virulence describes the relative degree of
disease-causing capacity among strains or organisms.

\medterm{Pathognomonic} Describes a finding considered highly specific for a particular disease. Truly pathognomonic findings are uncommon; most findings must still be interpreted with the history, examination, and other test results.
\medterm{Pathologic} Pathologic means related to disease or an abnormal structural or functional change, rather than normal anatomy or physiology.

\medterm{Pathologic Dilatation} Abnormal widening of a hollow organ, duct, blood vessel, or other structure beyond its expected size because of obstruction, weakness, pressure, or disease.

\medterm{Pathologic Fracture} A break in a bone weakened by disease, such as cancer spread, myeloma, infection, a benign tumor, or a metabolic bone disorder. It may happen after little or no injury and requires treatment of both the break and its cause.

\medterm{Pathologic Hypertrophy} Enlargement of an organ because its cells become larger, usually after long-term extra workload or hormone stimulation. It can help at first but may eventually impair function, as in an overworked heart.

\medterm{Pathologic Ossification} Abnormal formation of bone in soft tissue or in a location where bone should not develop, as in heterotopic ossification after trauma, neurologic injury, or surgery.

\medterm{Pathologic Process} A pathologic process is an abnormal biological change that produces disease, such as infection, inflammation, ischemia, degeneration, abnormal growth, or immune injury.

\medterm{Pathologic Response} A pathologic response is an abnormal reaction of cells, tissues, or organs to injury, infection, immune activation, toxins, or other stressors that contributes to disease.

\medterm{Pathological Fracture} A fracture occurring in bone weakened by disease, such as osteoporosis, tumor, infection, or metabolic bone disorder, rather than by normal force alone.

\medterm{Pathological Gambling} Legacy, nonpreferred label; see \textit{Gambling Disorder}.

\medterm{Pathological Laughter And Crying} Alternate index label; see \textit{Pseudobulbar affect}.

\medterm{Pathological Stealing} Alternate index label; see \textit{Kleptomania}.

\medterm{Pathologist} A pathologist is a physician who diagnoses disease by examining tissues, cells, blood, or body fluids and may guide laboratory testing, cancer classification, and treatment decisions.

\medterm{Pathology} Study of disease. Includes general pathology (cell injury, inflammation, neoplasia) and systemic pathology (organ-specific diseases).
\medterm{Pathology Artifact} A change introduced while collecting, preserving, cutting, staining, or imaging a sample rather than by disease. Recognizing these false changes prevents an incorrect diagnosis.

\medterm{Pathophysiology} The study of the functional changes produced by disease or injury and the mechanisms through which they cause signs and symptoms.

\medterm{Patient} A person receiving medical care, advice, examination, or treatment. Good care includes listening to the patient's concerns, explaining choices clearly, respecting informed decisions, protecting privacy, and reducing avoidable harm.

\medterm{Patient Advocacy} Acting on behalf of a patient to ensure that the patient's needs, rights, and
preferences are addressed in healthcare.

\medterm{Patient Autonomy} A patient’s right and capacity to make informed voluntary decisions about healthcare after receiving understandable information about options, benefits, risks, and alternatives.


\medterm{Patient Communication} The exchange of information and understanding between a patient, caregivers, and healthcare professionals during care.
\medterm{Patient Education} The process of explaining a health condition, treatment, medicines, exercises, equipment, warning signs, and prevention strategies in a way the patient can understand and use.

\medterm{Patient Observation} Systematic monitoring of a patient’s symptoms, vital signs, examination findings, behavior, and response to treatment to detect deterioration or guide management.

\medterm{Patient Positioning} Placing and supporting a patient in a position that permits safe care, accurate examination or imaging, and adequate comfort.

\medterm{Patient Preparation} The instructions and assessments completed before a test or procedure, such as fasting, medicine review, hydration, consent, or skin preparation when required.

\medterm{Patient Rights} Ethical and legal protections for people receiving care, including informed consent, privacy, access to information, respectful treatment, participation in decisions, and refusal of treatment where lawful.

\medterm{Patient Safety} Preventing avoidable harm during healthcare. It includes accurate identification, safe prescribing and medicine use, infection prevention, clear communication, fall prevention, and checking procedures and results.

\medterm{Patient-Centered Care} Healthcare planned around a person's needs, values, preferences, and goals. It includes clear communication, shared decisions, respect for cultural and personal choices, coordination between professionals, and support for the person and family.

\medterm{Patient-Centered Outcome} An outcome that reflects what matters to the patient, including symptoms, function,
quality of life, treatment burden, and ability to achieve personal goals, rather than
relying only on laboratory or imaging measures.

\medterm{Patient-Controlled Analgesia} A method in which a patient self-administers prescribed doses of analgesic medication
using a programmable pump with preset safety limits. The medication, dose, and lockout interval
are individualized according to the patient's condition and monitoring needs.

\medterm{Patient-Reported Outcome Measure} A validated questionnaire completed directly by a patient to assess symptoms, function,
health status, or quality of life without clinician interpretation of the response.

\medterm{Patlak Plot (Patlak-Rutland Analysis)} A graph used with some nuclear-medicine scans to estimate how quickly a tracked substance moves from the blood into tissue and stays there. It can help measure tissue activity, including activity in some tumors.

\medterm{Pattern Recognition Receptors} Proteins on or inside immune cells that detect common features of germs or damaged tissue. Their activation helps start the body’s early immune response.
\medterm{Pattern-Recognition Receptors} Sensors on or inside immune cells that detect common features of germs or damaged cells. Their activation starts early immune responses, including inflammation and signals that help coordinate protection and tissue repair.

\medterm{Pavor Nocturnus (Night Terrors)} Episodes of sudden intense fear, screaming, and autonomic arousal during sleep, usually arising from deep non-REM sleep and often followed by little or no memory.

\synonyms
Night Terrors

\medterm{Paxillin} Paxillin is a focal-adhesion adaptor protein that links integrin signaling to the actin cytoskeleton and helps regulate cell adhesion, migration, and mechanical sensing.

\medterm{Pazopanib} An oral tyrosine-kinase inhibitor targeting vascular endothelial growth factor and other receptors, used for selected advanced renal-cell carcinoma and soft-tissue sarcoma.

\medterm{PBA} Alternate index label; see \textit{Pseudobulbar affect}.

\medterm{PBG Urine Test} A urine test measuring porphobilinogen, used urgently to evaluate suspected acute hepatic porphyria during compatible neurovisceral or abdominal attacks.

\medterm{PCBs} PCBs, or polychlorinated biphenyls, are persistent industrial chemicals that accumulate in tissues and can affect the nervous, endocrine, immune, and reproductive systems.

\medterm{PCL Injuries} Sprains or tears of the posterior cruciate ligament, a strong ligament inside the knee. They may follow a blow to the shin or forceful bending and can cause pain, swelling, stiffness, or a feeling that the knee gives way.

\medterm{PCO} PCO may mean polycystic ovaries or posterior capsular opacification; the intended condition depends on the clinical context.

\medterm{PCO Disease} Polycystic ovary syndrome, a hormonal and metabolic disorder involving irregular ovulation, androgen excess, and sometimes polycystic-appearing ovaries.

\medterm{PCOS} A common hormone disorder that can cause irregular ovulation or periods, higher levels of androgen hormones, acne, or excess facial hair. Some people have ovaries with many small follicles, but that finding is not required for diagnosis.

\medterm{PCP} Pneumocystis jirovecii pneumonia (formerly Pneumocystis carinii pneumonia); opportunistic infection in immunocompromised patients, especially transplant recipients.
\medterm{PDA} Alternate index label; see \textit{Patent ductus arteriosus}.

\medterm{PDE5 Inhibitors} Medicines that improve blood flow into the penis by relaxing blood vessels and are mainly used for erectile dysfunction. They must not be combined with nitrate medicines because blood pressure can fall dangerously low.

\medterm{PE Anticoagulation} Anticoagulant treatment for pulmonary embolism to prevent clot extension and recurrence. Options
include direct oral anticoagulants, heparins, and vitamin K antagonists; the choice and
duration depend on bleeding risk, pregnancy, kidney function, cancer, provoking factors,
and other clinical features.

\medterm{Peak Expiratory Flow Rate} The fastest speed at which a person can blow air out after taking a full breath. A handheld meter measures it in litres per minute; lower or changing values can help monitor airway narrowing, especially in asthma.

\medterm{Peak Flow Meter} A handheld device that measures peak expiratory flow rate, the fastest speed of air expelled after a full breath, usually reported in liters per minute. It provides a simple measure of large-airway flow but does not measure the full range of airflow and lung volumes assessed by spirometry.
\medterm{Peanut Allergy} An immune reaction to peanut proteins that can cause symptoms ranging from hives to life-threatening anaphylaxis.

\synonyms
Allergy, Peanut

\medterm{Peanut Oil} Peanut oil is oil extracted from peanuts; refined oil usually contains little allergenic protein, whereas unrefined products may trigger reactions in susceptible people.

\medterm{Pearly Penile Papules} Small, benign, dome-shaped papules arranged circumferentially around the glans penis,
representing normal anatomical variants requiring no treatment.

\medterm{Peau D'Orange} Peau d’orange is dimpled skin resembling an orange peel, caused by edema tethered around hair follicles; it can occur with breast cancer, infection, or lymphatic obstruction.


\medterm{Pectinate Line} A boundary inside the anal canal where the lining, nerve supply, blood vessels, and lymph drainage change. Conditions above and below this line may cause different symptoms and may be treated differently.

\medterm{Pectoral} Relating to the chest muscles or the front of the chest. The pectoral muscles move the arm and help stabilize the shoulder. Pain in this area may come from muscle strain, the ribs, breast tissue, or organs inside the chest.
\medterm{Pectoral Girdle} The pectoral girdle consists of the clavicles and scapulae, attaching the upper limbs to the trunk and providing broad surfaces for shoulder and arm movement.

\medterm{Pectoralis} Relating to the chest muscles, especially pectoralis major and pectoralis minor. These muscles help move the arm, support the shoulder, and connect the chest wall with the upper limb.
\medterm{Pectoralis Major Muscle} The large chest muscle that moves the arm across the body, turns it inward, and helps with pushing. It attaches to the collarbone, breastbone, ribs, and upper arm.

\medterm{Pectoralis Minor Muscle} A small chest muscle beneath the larger pectoralis major. It helps draw the shoulder blade forward and downward, hold it against the ribs, and assist forceful breathing; tightness can contribute to rounded shoulders.

\medterm{Pectoralis Muscles} The pectoralis major and minor muscles of the chest, which move and stabilize the shoulder and assist with movements such as pushing, adduction, and breathing.

\medterm{Pectoriloquy} Abnormally clear transmission of spoken sounds through the chest wall, often suggesting lung consolidation or another change in lung tissue.

\medterm{Pectus Carinatum} A chest-wall shape present from childhood in which the breastbone pushes outward, sometimes called pigeon chest. It may cause no health problems, but some people have chest discomfort, reduced exercise tolerance, or emotional distress.

\medterm{Pectus Excavatum} A chest-wall shape present from childhood in which the breastbone is pushed inward, sometimes called funnel chest. It may cause no symptoms, but severe cases can affect exercise, breathing, posture, or emotional wellbeing.

\synonyms
Sunken Chest


\medterm{Pectus Excavatum Repair} Surgical correction of a sunken sternum and adjacent costal cartilage, performed when chest-wall deformity causes significant symptoms, physiologic impairment, or psychosocial impact.


\medterm{Pediatric Advanced Life Support} Emergency training and guidance for recognizing and treating life-threatening breathing, circulation, or heart problems in infants and children. It covers resuscitation, airway support, shock, abnormal rhythms, and care after recovery.

\synonyms
PALS (Pediatric Advanced Life Support); Pediatric Resuscitation Protocol

\medterm{Pediatric AIDS} Pediatric AIDS is advanced HIV infection in a child, with severe immune suppression or AIDS-defining illness; antiretroviral therapy and prevention of opportunistic infection are essential.

\medterm{Pediatric Anesthesia} Anesthesia for infants and children uses age- and weight-appropriate medicines, equipment, and monitoring. Children’s airways and breathing reserves differ from adults, so they can develop low oxygen levels or changes in blood pressure quickly and need specialized care.

\medterm{Pediatric Appendicitis} Acute inflammation of the vermiform appendix, usually caused by luminal obstruction.
Children may have evolving periumbilical-to-right-lower-quadrant pain, anorexia,
vomiting, fever, and guarding, but presentation can be atypical in young children;
timely surgical and antimicrobial assessment reduces perforation risk.

\medterm{Pediatric Assessment Triangle} A rapid first look at a sick child that does not require touching the child. Clinicians observe appearance, effort of breathing, and skin color or circulation. The findings show how urgently help is needed and guide the next examination and treatment.

\synonyms
PAT (Pediatric Assessment Triangle); Pediatric Visual Triage Triangle

\medterm{Pediatric Brain Tumors} Tumors that develop in a child's brain or nearby tissues. They may be noncancerous or cancerous and can cause headache, vomiting, seizures, vision changes, weakness, or changes in behavior; treatment depends on the type, location, size, and age of the child.

\medterm{Pediatric Cancer} Cancer occurring in infants, children, or adolescents. Common types include leukemia, brain tumors, lymphoma, neuroblastoma, and Wilms tumor; prognosis varies by type and stage.

\medterm{Pediatric Dentistry} Dental care for infants, children, and adolescents, including prevention and treatment of tooth decay, care of developing teeth, growth-related problems, behavior support, and adaptations for children with special healthcare needs.

\medterm{Pediatric Dermatology} The medical specialty concerned with skin, hair, and nail disorders in infants, children,
and adolescents. It includes congenital lesions, inflammatory diseases, infections, vascular lesions, and tumors.

\medterm{Pediatric Dermatology Conditions} Skin conditions in children may include birthmarks, eczema, infections, vascular growths, pigment changes, and inherited disorders. Most are manageable, but a rapidly changing, painful, bleeding, or infected lesion needs medical assessment.

\medterm{Pediatric Diabetic Ketoacidosis} A life-threatening complication of diabetes in which very high blood sugar and a lack of insulin make the blood acidic. A child may urinate and drink excessively, vomit, have abdominal pain, deep breathing, fruity-smelling breath, sleepiness, or coma and needs emergency care.

\medterm{Pediatric Early Warning Score} A hospital scoring system that combines a child’s behavior, breathing, heart rate, circulation, and need for oxygen to detect worsening illness early. A rising score prompts closer observation or urgent review; it supports, but does not replace, clinical judgment.

\synonyms
PEWS; Pediatric Early Warning System

\medterm{Pediatric Eczema} Atopic dermatitis in children commonly begins in infancy with facial and extensor
involvement; flexural disease predominates in older children. Pruritus, xerosis, sleep
disturbance, and scratching are characteristic.

\medterm{Pediatric Fever} An elevated body temperature in a child, usually caused by infection but sometimes by inflammation,
medicines, or environmental heat. Evaluation depends on age, appearance, duration, and associated findings.

\medterm{Pediatric Health} The physical, mental, emotional, and social wellbeing of infants, children, and adolescents. Care includes prevention, vaccinations, growth and development checks, treatment of illness, and support for the child and family.

\medterm{Pediatric Heart Surgery} An operation to repair or improve a child's heart or major blood vessels, often for a congenital heart defect. Procedures may close abnormal openings, improve blood flow, or replace a valve; recovery and risks depend on the defect and operation.


\medterm{Pediatric Meningitis} Inflammation or infection of the coverings of a child’s brain and spinal cord. Fever, headache, stiff neck, vomiting, light sensitivity, confusion, seizures, or a rash may occur; babies may be irritable, feed poorly, or have a bulging soft spot. Suspected meningitis is an emergency.

\medterm{Pediatric Nursing} Pediatric nursing provides age-appropriate care to infants, children, and adolescents, accounting for growth, development, family involvement, communication, safety, and medication differences.

\medterm{Pediatric Obesity} Excess body fat in a child or adolescent, assessed by comparing weight and height with other children of the same age and sex. It can affect breathing, joints, blood sugar, sleep, and emotional wellbeing; assessment should be supportive and non-stigmatizing.

\medterm{Pediatric Obstructive Sleep Apnea} Repeated narrowing or blockage of a child’s throat during sleep, causing pauses in breathing, snoring, restless sleep, daytime behavior or attention problems, and sometimes poor growth. Enlarged tonsils or adenoids are common causes.

\medterm{Pediatric Orthopedics} The branch of orthopedics focused on bone, joint, muscle, and movement problems in infants, children, and adolescents. It includes developmental conditions, limb differences, fractures, sports injuries, spine deformities, and growth-related disorders.

\medterm{Pediatric Pneumonia} Infection or inflammation of the lung tissue in an infant, child, or adolescent. It may cause cough, fever, fast or difficult breathing, chest pain, or low oxygen and can be caused by bacteria, viruses, fungi, or aspiration.

\medterm{Pediatric Sepsis} Sepsis occurring in an infant, child, or adolescent, in which the body’s response to infection causes organ dysfunction. Signs vary with age and may include abnormal temperature, breathing difficulty, unusual sleepiness, poor feeding, or reduced urine.

\medterm{Pediatric Poisoning} Harm caused when a child is exposed to a medicine, chemical, plant, gas, or another toxic substance. Effects depend on the substance, amount, route, and time of exposure and may include vomiting, sleepiness, breathing problems, seizures, or organ injury.
\medterm{Pediatric Ophthalmology} The subspecialty caring for eye and vision disorders in infants, children, and adolescents, including strabismus, amblyopia, congenital abnormalities, refractive error, and childhood eye disease.

\medterm{Pediatric Cataract} Clouding of the natural lens in an infant or child. It may be present at birth or develop after birth and can interfere with visual development, causing amblyopia or reduced vision.

\medterm{Pediatric Glaucoma} Glaucoma occurring in an infant or child, often involving abnormal development or blockage of the eye’s drainage system. Increased eye pressure can enlarge or damage the eye and impair vision.

\medterm{Persistent Fetal Vasculature} Persistence after birth of fetal blood vessels that normally regress before or soon after birth. It usually affects one eye and may cause a small eye, cataract, abnormal pupil reflection, retinal changes, or reduced vision.

\medterm{Congenital Nasolacrimal Duct Obstruction} Blockage of the tear-drainage channel present at birth, commonly causing persistent watering or discharge from one or both eyes. Most cases occur because a membrane at the lower end of the duct has not opened normally.

\medterm{Acquired Nasolacrimal Duct Obstruction} Blockage of the tear-drainage system that develops after birth. It may cause excessive tearing, discharge, or repeated inflammation of the lacrimal sac and can result from aging, infection, injury, inflammation, scarring, or a mass.

\medterm{Pediatric Rehabilitation} Therapy that helps children with conditions present from birth, injuries, illness, or developmental disabilities gain skills and independence. Care considers growth, family goals, school participation, communication, movement, learning, and daily activities.

\medterm{Pediatric Shock} A dangerous condition in which a child’s organs do not receive enough blood flow or oxygen. Causes include severe fluid loss, infection, heart failure, or a blockage. Early signs may include a fast heartbeat, pale or cold skin, unusual sleepiness, confusion, or little urine.

\medterm{Pediatric Skin Tumors} Skin growths in children include harmless birth-related or blood-vessel growths, cysts, and, rarely, cancer. A lump that grows quickly, bleeds, ulcerates, changes color, or causes pain should be assessed by a clinician.

\medterm{Pediatric Sleep Apnea} Pediatric sleep apnea is repeated interruption or reduction of breathing during sleep, commonly from enlarged tonsils or adenoids, causing snoring, restless sleep, behavior changes, or impaired growth.

\medterm{Pediatric Toxicology} The medical study and treatment of poisoning or harmful exposure in children. Common exposures include medicines, iron, household chemicals, and plants. Symptoms vary with the substance and may include vomiting, sleepiness, breathing problems, seizures, or burns; urgent poison-center or emergency advice may be needed.
\medterm{Pediatrician} A medical doctor who specializes in the health and medical care of children from birth
through adolescence (typically up to age 16-18).

\synonyms
Paediatrician

\medterm{Pediatrics} The branch of medicine concerned with the health and medical care of infants, children,
and adolescents from birth to age 18 (or up to age 16 in the UK).

\synonyms
Paediatrics

\medterm{Pedicle} A pedicle is a stalk or attachment containing vessels, nerves, or other supporting structures, such as the base of a polyp, flap, or vertebral arch.

\medterm{Pediculosis} Infestation with lice, small insects that live on the scalp, body, or pubic hair. It causes itching and spreads mainly through close contact or shared personal items.
\medterm{Pediculosis Capitis} Infestation of the scalp by head lice, causing itching and attached eggs or lice in the hair.

\medterm{Pediculus Humanus Capitis} A blood-sucking parasite commonly known as the louse (plural, lice) that can cause an
itchy scalp; infestations are highly contagious and especially common in school-age
children.

\medterm{Pedigree} A family diagram that shows relationships and who has a particular trait or disease. Doctors and genetic counselors use it to look for inheritance patterns and estimate the chance that relatives may be affected.

\medterm{Pedodontics} The branch of dentistry concerned with infants, children, and adolescents.
\medterm{Pedophilia} A continuing sexual interest in children who have not reached puberty. The term describes an interest and is not synonymous with pedophilic disorder; sexual contact with a child is abuse and cannot be consensual.

\medterm{Pedophilic Disorder} A mental-health disorder involving recurrent, intense sexual fantasies, urges, or behaviors toward children who have not reached puberty. Diagnostic criteria include clinically significant distress or impairment, or acting on the urges, together with required age and duration criteria.

\medterm{Pedunculate} Pedunculate describes a structure, lesion, or polyp attached by a narrow stalk rather than a broad base.

\medterm{Pedunculated} Attached to the skin or another surface by a narrow stalk, rather than having a broad base.

\medterm{Peeling Skin} Loss or shedding of the outer layer of skin, which may result from dryness, sunburn, irritation, allergy, infection, inflammation, or a systemic illness.

\medterm{PEEP} Pressure kept in the airways at the end of breathing out, usually by a ventilator or breathing machine. It helps keep the lung air sacs open and improve oxygen levels, but too much pressure can injure the lungs or lower blood flow.


\medterm{PEF} PEF, or peak expiratory flow, is the fastest airflow produced during a forceful exhalation and is used chiefly to monitor variable airway obstruction such as asthma.

\medterm{PEG Tube} A flexible feeding tube placed through the abdominal wall into the stomach with camera guidance. It provides nutrition, fluids, or medicines when swallowing is unsafe or inadequate, and requires regular skin care.




\medterm{Peliosis Hepatis} A usually uncommon condition in which multiple blood-filled spaces form within the liver. It may be linked to medicines, hormones, toxins, infections, cancer, or other diseases. It often causes no symptoms but can rarely cause liver failure or life-threatening bleeding if a space ruptures.


\medterm{Pelizaeus-Merzbacher Disease} A rare inherited disorder in which the brain does not make myelin, the protective covering around nerve fibers, normally. It usually begins in infancy with poor muscle tone, unusual eye movements, and delayed development, followed by stiffness, balance problems, and increasing movement difficulty. There is no cure; care supports movement, feeding, breathing, and communication.

\synonyms
PMD; PLP1-Related Hypomyelinating Leukodystrophy

\medterm{Pellagra} A disease caused by too little niacin, leading to skin changes, diarrhea, and problems with thinking.
\medterm{Pellet} A pellet is a small compacted mass of medication, material, or tissue; drug pellets may be formulated for implantation or controlled release.

\medterm{Pelvic} Relating to the pelvis, the bony ring and associated organs, muscles, vessels, and nerves in the lower trunk.

\medterm{Pelvic Binder} A circumferential device applied around the pelvis to reduce motion and hemorrhage from
suspected unstable pelvic-ring injury. It should be centered over the greater
trochanters and used with definitive imaging and hemorrhage-control management.

\medterm{Pelvic Bone} The ring of bones formed by the paired hip bones, sacrum, and coccyx that supports the trunk, protects pelvic organs, and transmits weight to the lower limbs.

\synonyms
Hip Bone
Hip Bones

\medterm{Pelvic Cavity} The space inside the lower part of the body enclosed by the pelvic bones. It contains the bladder, rectum, reproductive organs, blood vessels, nerves, and supporting muscles.

\medterm{Pelvic Congestion} Enlargement or pooling of blood in the veins of the pelvis, which may cause pelvic pressure, discomfort, or pain.

\medterm{Pelvic Congestion Syndrome} Chronic pelvic pain associated with enlarged or inefficient pelvic veins, often worsening with prolonged standing,
before menstruation, or after intercourse. Diagnosis requires evaluation for other causes of pelvic pain and may
involve pelvic imaging; treatment is individualized by gynecology or vascular specialists.

\medterm{Pelvic CT Scan} Computed tomography of the pelvis using x-rays and computer reconstruction to evaluate bones, organs, soft tissues, masses, trauma, infection, and bleeding.

\medterm{Pelvic Diaphragm} The group of muscles and supporting tissue forming the floor of the pelvis. It supports the bladder, bowel, and reproductive organs and helps control urine and stool.

\medterm{Pelvic Examination} Clinical inspection and palpation of the external genitalia, vagina, cervix, uterus, and adnexa to evaluate bleeding, pain, discharge, lesions, pregnancy-related findings, or infection.

\medterm{Pelvic Exenteration} Very extensive surgery removing some pelvic organs and nearby tissues, usually for selected cancers that have returned or grown locally when less extensive treatment cannot control them. Recovery is prolonged and reconstruction may be needed.

\medterm{Pelvic Floor} A group of muscles and connective tissues supporting the pelvic organs and controlling urinary and fecal continence.

\medterm{Pelvic Floor Muscles} Muscles at the bottom of the pelvis that support the bladder, bowel, and reproductive organs and help control urine, stool, and sexual function.

\medterm{Pelvic Floor Dysfunction} Weakness, overactivity, poor coordination, or injury of the muscles and connective tissues supporting the bladder,
bowel, and reproductive organs. Symptoms may include urinary or fecal problems, pelvic pain, constipation,
pain with sex, or pelvic organ prolapse; assessment may involve pelvic-health clinicians.

\medterm{Pelvic Floor Health And Disorders} The pelvic floor is a group of muscles and supporting tissues that hold the bladder, bowel, and reproductive organs in place and help control urine, stool, and sexual function. Weakness, tightness, injury, or poor coordination can cause pain, leakage, constipation, or prolapse.

\medterm{Pelvic Floor Muscle Training Exercises} Pelvic floor muscle exercises strengthen the muscles supporting the bladder, bowel, and reproductive organs and may improve selected forms of urinary or fecal leakage.

\medterm{Pelvic Girdle} The ring of bones connecting the lower limbs to the trunk, consisting of the hip bones, sacrum, and associated joints.

\medterm{Pelvic Girdle Pain In Pregnancy} Pain arising from the pelvic joints and surrounding tissues during pregnancy, often affecting the pubic symphysis, sacroiliac joints, hips, or lower back. It may impair walking, turning in bed, or climbing stairs.

\medterm{Pelvic Infection} An infection involving pelvic organs or tissues, including pelvic inflammatory disease, urinary infection, abscess, or postoperative infection.

\medterm{Pelvic Inflammatory Disease} Infection and inflammation of the upper female genital tract (uterus, fallopian tubes,
ovaries, and peritoneum), most commonly caused by sexually transmitted infections
(Chlamydia trachomatis, Neisseria gonorrhoeae). Presents with lower abdominal pain,
cervical motion tenderness, adnexal tenderness, and fever.

\medterm{Pelvic Inlet} The upper opening of the bony birth canal, formed by the lower spine, hip bones, and front pelvic joint. Its shape and size help determine how a baby can pass through the pelvis during childbirth.

\medterm{Pelvic Kidney} A kidney that remains in the pelvis rather than ascending to its usual lumbar position during development; it may be asymptomatic or associated with abnormal drainage, stones, infection, or vascular anatomy.

\medterm{Pelvic Laparoscopy} A minimally invasive procedure using a camera inserted through small abdominal incisions to inspect and treat pelvic organs and structures.

\medterm{Pelvic Lymphadenectomy} Pelvic lymphadenectomy is surgical removal of lymph nodes in the pelvis to check for or treat spread of cancer.

\medterm{Pelvic Mass} A lump or abnormal enlargement in the pelvis arising from reproductive, urinary, gastrointestinal, lymphatic, or other tissues.

\medterm{Pelvic Neoplasm} A pelvic neoplasm is an abnormal growth in the pelvis, arising from organs, bones, soft tissues, or lymph nodes; it may be benign or malignant.

\medterm{Pelvic Organ Prolapse} Dropping or bulging of the bladder, uterus, rectum, or nearby tissue into or toward the vagina because supporting muscles and tissues have weakened. It may cause pressure, a vaginal bulge, urinary or bowel problems, or discomfort during sex.

\medterm{Pelvic Organ Prolapse Quantification System} A standard examination used to measure how far the vagina, uterus, bladder, or rectum has moved downward. Measurements are taken while bearing down and describe the severity of prolapse to help plan treatment.

\synonyms
POP-Q System; POP-Q Classification

\medterm{Pelvic Outlet} The lower opening of the bony pelvis and birth canal, bordered by the pubic bone, sitting bones, and tailbone. Its shape and size affect passage through the pelvis during childbirth.

\medterm{Pelvic Pain} Pain in the lower abdomen, pelvis, or area between the genitals and anus. Causes include reproductive, urinary, bowel, muscle, nerve, and other conditions. Sudden severe pain, pregnancy, fever, fainting, vomiting, or heavy bleeding needs urgent assessment.

\medterm{Pelvic Pain During Early Pregnancy} Pain in the pelvis or lower abdomen during early pregnancy, caused by normal uterine changes, miscarriage, ectopic pregnancy, ovarian disease, infection, or other conditions. Severe or one-sided pain, shoulder pain, fainting, fever, or vaginal bleeding requires emergency assessment.

\medterm{Pelvic Pain, Chronic} Alternate index label; see \textit{Chronic pelvic pain}.


\medterm{Pelvic Spleen} A spleen located unusually low in the abdomen or pelvis because its supporting tissues are loose or absent. It may feel like a pelvic lump and can twist, lose its blood supply, or rupture, causing sudden pain and requiring urgent care.

\medterm{Pelvic Support Problems, Uterine Prolapse} Alternate index label; see \textit{Uterine prolapse}.

\medterm{Pelvic Ultrasound - Abdominal} Ultrasound imaging of pelvic organs through the lower abdominal wall, usually with a full bladder, to evaluate the uterus, ovaries, bladder, and surrounding structures.

\medterm{Pelvic Viscus} An organ located in the pelvic cavity, such as the bladder, rectum, uterus, ovaries, prostate, or related structures.

\medterm{Pelvimeter} An instrument used to estimate the size and shape of the pelvis, especially during obstetric assessment. Pelvic measurements are only one part of judging whether vaginal birth is likely and are interpreted with the clinical situation.
\medterm{Pelvimetry} Measurement of the size and shape of the pregnant person's pelvis, usually for selected childbirth concerns. Results are considered with the baby's position, size, labor progress, and overall clinical situation.

\medterm{Pelvis} The bony ring and associated soft tissues at the base of the trunk, connecting the spine with the lower limbs.
\medterm{Pelvis MRI Scan} Magnetic resonance imaging of the pelvis used to evaluate organs, muscles, bones, joints, nerves, blood vessels, and soft-tissue disease without ionizing radiation.

\medterm{Pelvis X-Ray} Plain radiographic imaging of the pelvic bones and hips used to evaluate fractures, alignment, arthritis, deformity, and selected bone lesions.

\medterm{Pemberton Sign} Facial redness or bluish discoloration, neck-vein swelling, or breathing difficulty that appears when both arms are raised. It may suggest that an enlarged thyroid or another chest mass is pressing on major veins or the airway.

\medterm{Pemberton’S Sign} Facial redness or bluish discoloration, neck-vein swelling, or breathing difficulty that appears when both arms are raised. It may suggest that an enlarged thyroid or another chest mass is pressing on major veins or the airway.

\medterm{Pemetrexed} Pemetrexed is a folate-antimetabolite chemotherapy drug that inhibits several enzymes involved in nucleotide synthesis, used especially for mesothelioma and nonsquamous non-small-cell lung cancer; folic acid and vitamin B12 reduce toxicity.

\medterm{Pemetrexed Disodium} The disodium salt of pemetrexed, an antifolate chemotherapy medicine used with other treatment for selected non-small-cell lung cancer and mesothelioma.

\medterm{Pemoline} Pemoline is a stimulant formerly used for attention-deficit disorders; serious liver toxicity led to withdrawal or severe restriction in many countries.

\medterm{Pemphigoid} A group of immune-system diseases that cause firm, often itchy blisters when antibodies attack the deeper layers of skin or mucous membranes. The type, location, age, and associated illness determine treatment.
\medterm{Pemphigoid Bullous (Bullous Pemphigoid)} Alternate index label for Bullous Pemphigoid.

\medterm{Pemphigoid Gestationis} An autoimmune blistering disease of pregnancy characterised by intensely pruritic
urticarial plaques and tense vesicles or bullae, typically arising in the second or
third trimester around the umbilicus and spreading to the trunk and extremities.

\medterm{Pemphigus} A group of autoimmune diseases causing fragile blisters when antibodies attack connections between skin cells.
\medterm{Pemphigus Vulgaris} A rare immune-system disease that causes fragile blisters and painful raw areas on the skin and mouth or other moist linings. Blisters break easily; untreated disease can cause fluid loss, infection, and difficulty eating or drinking.

\medterm{Penbutolol} Penbutolol is a beta-blocker formerly used for high blood pressure; it can slow the heart rate and worsen asthma in susceptible people.

\medterm{Pencil Eraser Swallowing} Swallowing a small pencil eraser usually causes little poisoning risk, but it may lodge in the airway or digestive tract. Choking, breathing difficulty, persistent vomiting, or severe pain requires emergency care.

\medterm{Pencil Swallowing} Swallowing a pencil or part of one can injure the mouth, throat, or digestive tract and may block the airway. Do not force food or drink; breathing trouble, bleeding, or severe pain needs emergency care.


\medterm{Pendular} Moving or swinging freely like a pendulum; pendular nystagmus has a similar speed in both directions.
\medterm{Pendulous} Pendulous means hanging loosely from a supporting structure, as in a pedunculated lesion, drooping tissue, or an organ suspended by a stalk.

\medterm{Penetrance} The proportion of people with a particular gene change who develop the related trait or condition. Penetrance may be complete, incomplete, or influenced by other factors.
\medterm{Penetrating Abdominal Trauma} Injury to the abdomen caused by a projectile, sharp object, or other penetrating
force. It may damage abdominal organs or blood vessels and requires urgent assessment for internal bleeding or organ
injury; management depends on the mechanism, examination, stability, and imaging findings.

\medterm{Penetrating Wounds} Injuries caused when an object enters the skin and damages underlying tissue. They may injure blood vessels, nerves, organs, or bones and require assessment for bleeding, contamination, retained objects, and infection.
\medterm{Penicillin} A family of antibiotics that kill bacteria by stopping them from making cell walls. Different types treat different bacterial infections; they do not work against viruses. Allergy, resistance, and kidney problems can affect the choice of medicine.
\medterm{Penicillin Allergy} An immune reaction to penicillin antibiotics that may cause rash, swelling, breathing difficulty, or anaphylaxis.

\synonyms
Allergy, Penicillin
\medterm{Penicillin G} Penicillin G is an injectable penicillin antibiotic used for infections such as syphilis and selected streptococcal or other susceptible bacterial infections.

\medterm{Penicillin V} Penicillin V is an oral penicillin antibiotic used for selected susceptible infections, including some throat and dental infections.

\medterm{Penicillium} A genus of molds that includes species used to make medicines and foods and species that can cause allergy or opportunistic infection.

\medterm{Penile} Penile means relating to the penis, the external male organ involved in urination and sexual reproduction.

\medterm{Penile Adhesions} Abnormal attachment of penile skin to the glans, often after circumcision or during childhood. They may cause irritation or conceal part of the glans and are assessed and treated according to severity.

\medterm{Penile Cancer} An uncommon cancer of the penis, usually arising from surface skin cells. It may appear as a persistent sore, lump, bleeding area, or skin change; HPV, smoking, foreskin problems, and chronic inflammation can increase risk.

\medterm{Penile Augmentation} A cosmetic or reconstructive procedure intended to increase penile length, girth, or both. Techniques and results vary, and complications may include scarring, deformity, infection, altered sensation, or erectile problems.


\medterm{Penile Disease} A disorder affecting the penis, including infection, inflammation, erectile or structural disease, trauma, skin disease, or malignancy.

\medterm{Penile Edema} Swelling of the penis caused by inflammation, infection, trauma, allergic reaction, impaired lymphatic drainage, or a constricting ring or retracted foreskin. Sudden painful swelling with an irreducible foreskin is an emergency.

\medterm{Penile Erection} Stiffening and enlargement of the penis when blood fills erectile tissue and is temporarily held there. It can occur with sexual or other stimulation; persistent difficulty getting or keeping an erection may be erectile dysfunction.


\medterm{Penile Fracture} A tear in the tough covering of an erect penis, usually after forceful bending. A popping sound, immediate loss of erection, swelling, and bruising may occur; urgent surgery can prevent lasting problems.

\medterm{Penile Hygiene} Routine cleaning and drying of the penis and foreskin, with gentle retraction only when it moves freely and replacement afterward. Good hygiene helps reduce irritation and infection without harsh soaps or forceful retraction.

\medterm{Penile Implant} A surgically implanted device used to produce penile rigidity in selected people with erectile dysfunction that has not responded to other
treatments.

\medterm{Penile Infection} Infection of the penis or foreskin caused by bacteria, fungi, viruses, or parasites, producing pain, redness, discharge, sores, or swelling.

\medterm{Penile Intraepithelial Neoplasia} A precancerous change in the surface cells of the penis. It may appear as a persistent red, white, scaly, or wart-like patch on the glans or foreskin and may be linked to HPV or chronic inflammation. Examination, biopsy, treatment, and follow-up reduce the risk of invasive cancer.

\synonyms
PeIN; Penile Precancerous Lesion

\medterm{Penile Leukoplakia} A persistent white patch on penile skin or mucosa that may reflect irritation, inflammation, or a precancerous disorder and requires assessment.

\medterm{Penile Prosthesis} An implanted device used to produce an erection when erectile dysfunction has not improved with other treatments. It may be inflatable or bendable. Surgery can cause infection, bleeding, pain, device failure, or erosion into nearby tissue.

\medterm{Penile Torsion} Rotation of the penis around its longitudinal axis, usually a congenital condition in which the glans or urethral opening is rotated. Symptomatic or marked torsion may require urologic evaluation and surgical correction.

\medterm{Penile Trauma} Injury to the penis caused by blunt force, bending, penetration, crushing, burns, bites, or medical instruments. It may damage the skin, erectile tissue, blood vessels, or urethra. A penile fracture is a tear in the covering of the erectile tissue, usually after forceful bending of an erect penis. It may cause pain, swelling, bruising, bleeding, or difficulty passing urine.

\medterm{Penis} An external genital organ that contains the outlet for urine and can deliver sperm during sexual activity. It becomes firm when filled with blood during an erection and contains sensitive tissue, skin, blood vessels, and nerves.
\medterm{Penis Carcinoma} Penis carcinoma is cancer of the penis, most often arising from skin or squamous cells and sometimes linked to HPV infection.


\medterm{Penis Disorder (Priapism (Penis Disorder))} Alternate index label for Priapism (Penis Disorder).

\medterm{Penis Disorder Balantis (Balanitis (Penis Disorder))} Alternate index label for Balanitis (Penis Disorder).

\medterm{Penis Pain} Penis pain may result from infection, injury, inflammation, urinary problems, or sexual conditions; severe pain, swelling, discharge, or inability to urinate requires prompt assessment.

\medterm{Penis, Erection Of The} Penile rigidity produced by increased arterial inflow and reduced venous outflow after sexual or reflex stimulation; persistent difficulty is erectile dysfunction.

\medterm{Penis, Small} A penis that is substantially below expected stretched length for age, usually assessed as micropenis in a newborn or child after excluding measurement and developmental variation.

\medterm{Penoscrotal Transposition} A congenital difference in which the penis and scrotum are positioned abnormally relative to each other, often associated with hypospadias or other genital and urinary anomalies.

\medterm{Pentad Of TTP} The five findings once considered typical of thrombotic thrombocytopenic purpura: low platelets, destruction of red blood cells, neurologic symptoms, kidney problems, and fever. Most patients do not have all five, so urgent blood testing is essential when TTP is suspected.


\medterm{Pentamidine} An antimicrobial medicine used for certain parasitic infections and Pneumocystis pneumonia, with potentially serious side effects.
\medterm{Pentamidine Isethionate} An antiprotozoal and antimicrobial drug used in selected cases of Pneumocystis pneumonia, leishmaniasis, trypanosomiasis, and other infections; toxicity includes hypoglycemia, kidney injury, and arrhythmia.

\medterm{Pentatrichomonas Hominis} Pentatrichomonas hominis is a flagellated intestinal protozoan generally regarded as nonpathogenic; finding it in stool usually indicates exposure to fecal contamination.

\medterm{Pentazocine} An opioid pain medicine with mixed agonist and antagonist effects, used for some moderate or severe pain. It can cause drowsiness, dependence, and slowed breathing and may trigger withdrawal in a person dependent on other opioids.
\medterm{Pentazocine Overdose} Taking too much pentazocine, an opioid pain medicine, can cause extreme sleepiness, confusion, slow breathing, pinpoint pupils, seizures, or coma. Call emergency services immediately and give naloxone when available and appropriate.

\medterm{Pentetic Acid} A chelating compound that binds metal ions and is used in selected radiopharmaceuticals, contrast agents, or laboratory applications.

\medterm{Pentobarbital Overdose} Taking too much pentobarbital, a sedative, can severely slow brain activity and breathing, causing low blood pressure, coma, or death. It requires immediate emergency treatment, especially with alcohol or other sedatives.

\medterm{Pentosan Polysulfate} A semi-synthetic sulfated polysaccharide used for selected bladder pain syndrome and historically as an anticoagulant; prolonged use can cause pigmentary maculopathy and bleeding.

\medterm{Pentose} A pentose is a five-carbon sugar, including ribose in RNA and deoxyribose in DNA, and participates in nucleotide and carbohydrate metabolism.

\medterm{Pentose Phosphate Pathway} A cell pathway that uses glucose to make NADPH, which protects cells from chemical damage, and ribose, which helps build DNA and RNA. It is especially important in the liver and red blood cells.

\medterm{Pentoxifylline} A methylxanthine derivative that improves red-cell flexibility and is used in some countries for peripheral vascular disease; benefit varies by condition.


\medterm{Penumbra} Brain tissue surrounding an area of infarction that is functionally impaired but potentially salvageable if blood flow is restored promptly.

\medterm{Peplomycin} Peplomycin is a bleomycin-related chemotherapy medicine used in some countries for selected cancers; lung toxicity is an important risk.

\medterm{Peppermint} Mentha piperita, used as a flavoring and in some preparations for dyspeptic symptoms; concentrated oil can be hazardous.
\medterm{Peppermint Oil Overdose} Swallowing a large amount of peppermint oil can cause nausea, vomiting, abdominal pain, dizziness, allergic reactions, seizures, or lung injury if aspirated. Contact poison control promptly and do not induce vomiting.

\medterm{Pepsin} A stomach enzyme that begins breaking proteins into smaller fragments. It is released as inactive pepsinogen and activated by stomach acid, helping digestion before food enters the small intestine.
\medterm{Pepsin A} Pepsin A is a stomach protease activated by acid that begins digestion of dietary proteins into smaller peptides.

\medterm{Pepsinogen} An inactive precursor released by gastric chief cells and converted to pepsin in acidic conditions.
\medterm{Peptic Esophagitis} Inflammation and injury of the esophageal lining caused mainly by repeated gastroesophageal reflux of acid or bile, producing heartburn, chest discomfort, swallowing pain, or strictures.

\medterm{Peptic Ulcer} An open sore in the lining of the stomach or duodenum, most commonly caused by \textit{Helicobacter pylori} infection or NSAID use.

\synonyms
Stomach Ulcer (Peptic Ulcer)
Ulcer, Peptic

\medterm{Peptic Ulcer Bleeding} Haemorrhage caused by erosion of an artery or other blood vessel beneath a gastric or
duodenal peptic ulcer. It is a common cause of upper gastrointestinal bleeding and may
present with haematemesis, coffee-ground vomiting, melaena, or haemodynamic instability.

\medterm{Peptic Ulcer Disease} One or more open sores in the lining of the stomach or first part of the small intestine. Common causes include Helicobacter pylori infection and anti-inflammatory medicines; ulcers may cause burning pain, bleeding, or a dangerous perforation.



\medterm{Peptidase} An enzyme that breaks peptide bonds in proteins or peptides.
\medterm{Peptide} A short chain of amino acids linked by peptide bonds. Peptides can form parts of proteins or act independently as hormones, neurotransmitters, immune signals, or medicines.
\medterm{Peptide Biosynthesis} The process by which cells make peptides, short chains of amino acids, by linking amino acids according to genetic instructions. Some peptides act as hormones, neurotransmitters, or signals between cells.

\medterm{Peptide Bond} The chemical link joining two amino acids. Repeated peptide bonds form short peptides and long proteins, which cells use for structure, chemical reactions, communication, transport, and many other functions.

\medterm{Peptide Hormone} A hormone made of amino acids that binds a cell-surface receptor and regulates processes such as metabolism, growth, reproduction, stress, and fluid balance.


\medterm{Peptide Metabolism} The body's production, breakdown, and use of peptides. These processes regulate signaling molecules, hormones, digestion, and protein turnover.

\medterm{Peptide Transport} The movement of peptides into, out of, or within cells and tissues through transport proteins, vesicles, or other cellular mechanisms. It helps control peptide signaling, digestion, and removal from the body.

\medterm{Peptide Yy} Peptide YY (PYY) is a hormone released mainly by cells in the intestine after eating. It helps reduce appetite and slow movement of food through the digestive tract and is being studied in relation to body weight and metabolism.

\medterm{Peptidoglycan} A strong mesh-like material forming much of the bacterial cell wall. Some antibiotics kill bacteria by preventing them from making or repairing this protective layer.

\medterm{Peptococcaceae} A family of anaerobic gram-positive cocci found in the human microbiota and sometimes associated with infection.

\medterm{Peptococcus} A genus of anaerobic gram-positive cocci that can contribute to mixed infections, especially in damaged or poorly oxygenated tissue.

\medterm{Peptococcus Niger} An anaerobic gram-positive coccus found in the human microbiota and rarely associated with mixed infection.
\medterm{Peptoniphilus} A genus of anaerobic gram-positive cocci found on skin and mucosal surfaces that can cause opportunistic infection.

\medterm{Peptoniphilus Harei} An anaerobic gram-positive coccus occasionally recovered from wounds, abscesses, or other clinical specimens. An anaerobic Gram-positive bacterium sometimes found in wounds and abscesses, where it may contribute to mixed infections.
\medterm{Peptoniphilus Indolicus} An anaerobic gram-positive coccus found in the genital and other microbiota that can occasionally contribute to mixed infection.

\medterm{Peptostreptococcus} A group of anaerobic gram-positive cocci associated with dental, skin, abdominal, pelvic, and other mixed infections.

\medterm{Peracetic Acid} A strong oxidizing disinfectant and sterilant used for equipment and surfaces; concentrated exposure irritates or burns skin, eyes, and airways.

\medterm{Peracute} Peracute describes disease or a clinical event that develops extremely rapidly, often over minutes to hours, with little time for adaptation or diagnosis.

\medterm{PERC Rule} A checklist used by emergency clinicians for people already judged very unlikely to have a blood clot in the lungs. If all criteria are reassuring, testing may be avoided; it is not suitable for home use or for people with a strong suspicion of clot.

\synonyms
Pulmonary Embolism Rule-out Criteria; PERC Score

\medterm{Percentile} A percentile indicates the percentage of a reference population with a value below a measurement, such as height, weight, or test score.

\medterm{Perception} The brain’s process of organizing and interpreting sensory information into an experience of objects, events, body sensations, or the environment.

\medterm{Perceptual} Relating to the process of detecting and interpreting sensory information.
\medterm{Perceptual Disorder} A clinically significant disturbance in the way sensory information is interpreted, including altered visual, auditory, tactile, or body perception.

\medterm{Perceptual Distortion} An inaccurate or altered interpretation of a real sensory stimulus, which may occur with neurologic disease, psychiatric illness, drugs, migraine, seizures, or sensory impairment.

\medterm{Perceptual Masking} Reduction in the ability to detect or identify one sensory stimulus because another stimulus overlaps it in time, space, or processing.

\medterm{Percussion} Physical examination in which the body surface is tapped and the resulting sound or vibration is interpreted to estimate underlying air, fluid, or solid tissue.
\medterm{Percussion Note} The sound heard when the chest is gently tapped during examination. A normal lung sounds hollow; a dull sound may suggest fluid or infection, while an unusually hollow sound may suggest excess air or a collapsed lung.

\medterm{Percutaneous} Performed through the skin, usually using a needle, catheter, wire, or other instrument. This approach can diagnose or treat disease while avoiding a larger open operation, but bleeding and infection remain possible.
\medterm{Percutaneous Biopsy} Removal of tissue through the skin using a needle or other instrument, often guided by ultrasound, CT, or another imaging method.

\medterm{Percutaneous Coronary Intervention} A procedure that guides a thin tube through a blood vessel to a narrowed heart artery, then widens it with a balloon and often supports it with a stent. It can restore blood flow during a heart attack or treat selected severe coronary disease.

\medterm{Percutaneous Discectomy} Percutaneous discectomy removes or decompresses part of a herniated spinal disc through a small needle or incision.

\medterm{Percutaneous Diskectomy} A minimally invasive procedure removing part of a herniated spinal disk through a small skin opening to relieve nerve pressure.

\medterm{Percutaneous Drainage} Placement of a needle or catheter through the skin to drain a fluid collection or obstructed body space, often with imaging
guidance.

\medterm{Percutaneous Electrical Nerve Stimulation} Electrical stimulation delivered through the skin to reduce pain signals.

\synonyms
PENS

\medterm{Percutaneous Endoscopic Gastrostomy} Placement of a feeding tube through the abdominal wall into the stomach using a camera passed through the mouth. It may provide long-term nutrition, fluids, or medicines when swallowing is unsafe or inadequate.

\medterm{Percutaneous Kidney Procedures} Procedures performed through a small needle or tube passed through the skin into or near the kidney. They may obtain a biopsy, drain infection or urine, or remove stones; bleeding, infection, urine leakage, or injury to nearby organs are possible.

\medterm{Percutaneous Nephrolithotomy} A procedure that removes a large or complex kidney stone through a small opening in the back. A narrow instrument reaches the kidney, breaks or removes the stone, and may leave a temporary drainage tube; bleeding, infection, and injury are possible.

\medterm{Percutaneous Nephrolithotripsy} Percutaneous nephrolithotripsy removes or fragments kidney stones through a tract made through the skin into the kidney, usually for larger or complex stones.

\medterm{Percutaneous Nephrostomy} An image-guided interventional procedure in which a flexible drainage tube (nephrostomy catheter) is placed through the skin of the flank directly into the renal pelvis of the kidney. Using ultrasound and fluoroscopy under local anesthesia, an interventional radiologist or urologist punctures the kidney's urine-collecting system, threads a guide wire, and advances the catheter, which is secured to an external urine-collection bag. Percutaneous nephrostomy is primarily performed to decompress an obstructed, infected, or damaged kidney (hydronephrosis) when a ureteral stone, pelvic tumor, pregnancy, or stricture prevents normal urine outflow to the bladder. Decompressing the blocked kidney provides immediate relief from severe flank pain, treats or prevents life-threatening urosepsis, and preserves kidney function until the underlying obstruction can be surgically resolved.

\synonyms
PCN; Nephrostomy Tube Placement; Percutaneous Renal Drainage

\medterm{Percutaneous Transhepatic Cholangiogram} A percutaneous transhepatic cholangiogram uses a needle through the liver to inject contrast into bile ducts for imaging and may guide drainage or stent placement.

\medterm{Percutaneous Transhepatic Cholangiography} PTC is imaging of the bile ducts after contrast is injected through a needle passed through the skin and liver; drainage or stenting may follow.
\medterm{Percutaneous Transluminal Angioplasty} A catheter-based procedure that inflates a balloon inside a narrowed artery to improve blood flow. A stent may also be placed; risks include bleeding, clotting, vessel injury, and re-narrowing.

\medterm{Percutaneous Transluminal Coronary Angioplasty} Balloon-based treatment that widens a narrowed coronary artery, often followed by stent placement. A catheter procedure using balloon inflation to widen a narrowed coronary artery, sometimes followed by stent placement.
\synonyms
PTCA; balloon angioplasty

\medterm{Percutaneous Umbilical Blood Sampling} A procedure that uses ultrasound to guide a needle into the umbilical cord and collect a fetal blood sample during pregnancy. It may investigate anemia, infection, or genetic problems and carries risks such as bleeding, infection, and pregnancy loss.


\medterm{Percutaneously Inserted Central Catheter - Infants} A thin catheter inserted through a small vein in an infant’s arm or leg and advanced toward a large vein near the heart. It provides longer-term access for fluids, nutrition, or medicines, but requires careful monitoring for infection, blockage, bleeding, and displacement.

\medterm{Perforant Pathway} A group of nerve fibers carrying information from a brain area involved in memory to the hippocampus. It helps the brain form and retrieve memories and can be affected in some dementias and seizure disorders.

\medterm{Perforated Bowel} A perforated bowel is a hole through the intestinal wall allowing contents to leak into the abdomen, causing peritonitis and potentially life-threatening sepsis.

\medterm{Perforated Eardrum} Alternate index label; see \textit{Ruptured eardrum (perforated eardrum)}.

\medterm{Perforation} A hole or complete tear through the wall of an organ or tissue. It can release contents into surrounding spaces, cause severe infection or bleeding, and often requires urgent medical treatment.
\medterm{Perforation Of A Viscus} A hole through the wall of a hollow abdominal organ, such as the stomach, small intestine, or colon. Contents can leak into the abdomen and cause severe infection; ulcers, inflammation, poor blood flow, and injury are possible causes.

\medterm{Performance} How effectively a person, organ, treatment, test, or system carries out a defined task. Examples include exercise performance, heart performance, test performance, and response to treatment.

\medterm{Performance Anxiety} Anxiety about sexual performance that can interfere with sexual function.

\medterm{Performance Status} A rating of how well a person can move, work, and perform daily activities while ill. Clinicians use it to estimate prognosis, decide whether some treatments are safe, and track change over time, especially during cancer care.

\medterm{Perfusion} The flow of blood through the small vessels of a tissue or organ. Good perfusion supplies oxygen and nutrients and removes waste; poor perfusion can cause pain, weakness, tissue injury, or organ failure.

\medterm{Perfusion Defect} An area receiving unusually little blood flow, identified by an imaging or function test. It may reflect a blocked artery, narrowed vessel, poor heart output, inflammation, or damaged tissue.

\medterm{Perfusion Imaging} Imaging that shows how blood flows through an organ or tissue. It can help assess problems such as stroke, reduced heart blood flow, or tumor blood supply.
\medterm{Perfusion PET} Positron-emission tomography that measures tissue blood flow using a suitable radioactive tracer. Cardiac perfusion PET can show reduced blood flow, viability, and sometimes the effects of treatment.

\medterm{Perianal} Located around the anus; the term describes the skin, glands, pain, inflammation, abscesses, fistulas, and other findings in that region.

\medterm{Perianal Abscess} A localized collection of pus near the anus caused by infected anal glands or surrounding tissue, producing severe throbbing pain, swelling, fever, and sometimes a fistula.

\medterm{Perianal Glands} Small glands near the anus that contribute secretions to the local tissue; obstruction or infection can contribute to perianal abscesses and fistulas.

\medterm{Perianal Streptococcal Cellulitis} A bacterial skin infection around the anus, usually caused by group A Streptococcus. It may cause bright-red skin, pain, itching, discharge, or pain during bowel movements.

\medterm{Periapical Abscess} An acute or chronic localized collection of pus at the apex of a tooth's root, arising
from pulpal infection and necrosis (usually due to untreated dental caries or trauma).

\medterm{Periapical Granuloma} A chronic inflammatory lesion at the apex of a tooth, usually resulting from pulpal
necrosis and persistent apical infection. It may be asymptomatic and is diagnosed
through clinical, radiographic, and sometimes histopathologic assessment.

\medterm{Peri-Implantitis} Inflammation around a dental implant with progressive loss of the supporting bone. It may cause bleeding, swelling, pus, pain, or implant looseness and requires dental assessment.

\medterm{Periappendiceal} Periappendiceal means located around the appendix, as in inflammation, fluid, an abscess, or fat stranding near the appendiceal tissue.

\medterm{Periarteritis Nodosa} An older name for polyarteritis nodosa, a necrotizing vasculitis of medium-sized arteries that can cause skin, nerve, muscle, kidney, and abdominal ischemic symptoms without typical glomerulonephritis.

\medterm{Periarticular} Periarticular means located around a joint, including nearby tendons, bursae, ligaments, capsule, and soft tissues rather than the joint surface itself.

\medterm{Pericardial} Relating to the pericardium, the thin two-layered sac surrounding the heart. A small amount of fluid between its layers allows the heart to beat smoothly, while excess fluid or inflammation can interfere with heart function.

\medterm{Pericardial Cavity} The thin fluid-containing space between the layers of the sac surrounding the heart. A small amount of fluid lets the heart move smoothly; excess fluid or blood can press on the heart and impair its pumping.
\medterm{Pericardial Disease (Pericarditis)} Alternate index label for Pericarditis.

\medterm{Pericardial Effusion} Extra fluid collecting in the sac around the heart. A small amount may cause no symptoms, but rapid or large accumulation can compress the heart, lower blood pressure, and become life-threatening.
\medterm{Pericardial Effusion Etiology} Fluid around the heart may result from infection, cancer, kidney failure, low thyroid activity, autoimmune disease, heart injury, surgery, radiation, trauma, or medicines. Tests and the person’s history help identify the cause.

\medterm{Pericardial Fluid Culture} A microbiologic culture of fluid from the pericardial space to identify bacteria, fungi, or other organisms in suspected infectious pericarditis.

\medterm{Pericardial Fluid Gram Stain} A microscopic stain of pericardial fluid that can rapidly suggest bacterial organisms, although a negative result does not exclude infection.

\medterm{Pericardial Inflammation} Alternate index label; see \textit{Pericarditis}.

\medterm{Pericardial Knock} An extra early heart sound caused when a stiff, scarred sac around the heart suddenly stops the heart from filling. It may suggest constrictive pericarditis and must be interpreted with other examination and test findings.

\medterm{Pericardial Mesothelioma} A rare malignant tumor arising from the mesothelial lining around the heart, which can cause effusion, constriction, or impaired cardiac filling.

\medterm{Pericardial Metastases} Cancer that has spread to the sac around the heart from another part of the body. It may cause inflammation, fluid or blood around the heart, breathlessness, chest discomfort, or impaired heart filling.

\medterm{Pericardial Stripping} Surgery to remove part or most of the thickened or scarred sac around the heart. It may relieve constrictive pericarditis, in which the stiff sac prevents the heart from filling normally; bleeding, rhythm problems, and other major surgical risks are possible.

\medterm{Pericardial Tamponade} Dangerous pressure on the heart caused by rapid or large fluid or blood buildup in the sac around it. The heart cannot fill normally, causing breathlessness, low blood pressure, fainting, or shock and requiring emergency treatment.

\medterm{Pericardial Wall} The layered tissue forming the protective sac around the heart. Its inner surfaces produce a small amount of lubricating fluid that lets the heart beat with little friction.
\medterm{Pericardial Window} An operation that creates an opening in the sac around the heart so extra fluid can drain into a nearby space and be absorbed. It may prevent repeated fluid buildup and allow a tissue sample to be taken.

\synonyms
Subxiphoid Pericardiostomy; Pleuropericardial Window; Pericardiostomy

\medterm{Pericardiectomy} Surgery to remove part or most of the sac around the heart when thickening or scarring prevents the heart from filling normally. It is major surgery and may improve constrictive pericarditis but carries bleeding and other serious risks.

\medterm{Pericardiocentesis} A needle or catheter is used to remove fluid or blood from the sac around the heart, usually to relieve dangerous pressure or investigate an effusion. Imaging guidance improves safety, but bleeding, abnormal rhythm, organ injury, or infection can occur.

\medterm{Pericardiostomy} A surgical opening made in the pericardium to drain fluid or prevent recurrent accumulation around the heart.

\medterm{Pericarditis} Inflammation of the sac around the heart. It often causes sharp chest pain that worsens with deep breathing or lying down and improves when sitting up and leaning forward. It may be caused by infections, autoimmune diseases, heart injury, kidney disease, cancer, or some medicines. Severe cases can cause fluid buildup around the heart.

\synonyms
Inflamed Pericardium
Pericardial Inflammation
Pericardial Disease (Pericarditis)

\medterm{Post-Myocardial Infarction Pericarditis} Inflammation of the sac around the heart that occurs after a heart attack. It may begin soon after the heart attack or develop later as part of a delayed immune reaction.

\medterm{Pericarditis - Constrictive} A chronic condition in which a thickened, scarred, or calcified pericardium restricts ventricular filling and causes systemic venous congestion and heart-failure symptoms.

\medterm{Pericardium} A two-layered protective sac surrounding the heart and beginnings of major blood vessels, reducing friction during each heartbeat.

\medterm{Perichondrial Hematoma (Cauliflower Ear)} Alternate index label for Cauliflower Ear.

\medterm{Perichondritis} Perichondritis is inflammation or infection of tissue surrounding cartilage, often affecting the outer ear after trauma or piercing and potentially damaging cartilage.

\medterm{Perichondritis Of The Ear} Infection or inflammation of the auricular perichondrium, often after trauma, piercing,
or surgery. Pseudomonas aeruginosa is an important cause; delayed treatment can produce
cartilage necrosis and permanent deformity.

\medterm{Perichondrium} A tough layer of tissue surrounding most cartilage. It supplies nutrients, contains cells that help cartilage grow and repair, and can become inflamed or infected, especially around the outer ear.
\medterm{Pericoronitis} Inflammation of the soft tissue around the crown of a partially erupted tooth, most
often a mandibular third molar. Food retention and bacterial accumulation may cause
pain, swelling, trismus, fever, or spread into deeper spaces.

\medterm{Pericytes} Pericytes are contractile support cells wrapped around small blood vessels that help regulate capillary stability, blood flow, repair, and the blood-brain barrier.

\medterm{Periductal Mastitis} Inflammation around the major subareolar ducts, often associated with smoking, causing a
painful periareolar mass, discharge, or recurrent breast abscess.


\medterm{Perilymph} Fluid in the bony labyrinth of the inner ear that surrounds the membranous labyrinth and helps transmit sound and balance signals.

\medterm{Perilymphatic Fistula} An abnormal leak between the fluid spaces of the inner ear and the middle ear, often after injury, pressure changes, or surgery. It may cause dizziness, hearing loss, ringing, or symptoms triggered by sound or pressure.

\medterm{Perimetry} A test that measures the full field of vision, including central and side vision. It helps detect and monitor problems such as glaucoma, optic-nerve disease, retinal disease, and disorders affecting the visual pathways in the brain.

\medterm{Perimenopausal} Perimenopausal describes the years of hormonal and menstrual transition before and around menopause, when cycles, bleeding, hot flashes, sleep, and fertility may change.

\medterm{Perimenopause} The transition time in a woman’s life that begins when ovaries produce less estrogen and
menstruation becomes less frequent, and ends when the ovaries no longer produce eggs and
menstruation stops.

\medterm{Perinatal} Perinatal refers to the period around birth, conventionally including late pregnancy and the first days or weeks after delivery depending on the definition used.

\medterm{Perinatal Asphyxia} Too little oxygen or inadequate breathing before, during, or soon after birth. Severe cases can injure the brain and other organs; rapid newborn assessment, breathing support, and treatment of the cause are essential.

\medterm{Perinatal Death} Death of a fetus or newborn during the late pregnancy, birth, or early newborn period; the exact time range varies by definition.

\medterm{Perinatal Psychiatry} The subspecialty focused on psychiatric disorders during pregnancy and the postpartum
period. Common conditions include perinatal depression, perinatal anxiety, perinatal
OCD, and postpartum psychosis. Medication selection must balance maternal mental health
with fetal or neonatal exposure risks.

\medterm{Perinatal Transmission} Transmission of an infection or condition from mother to infant during pregnancy, labor, delivery, or breastfeeding, as with HIV, hepatitis B, syphilis, or certain genetic and congenital disorders.

\medterm{Perinatologist} A perinatologist is an obstetrician specializing in high-risk pregnancy, fetal medicine, prenatal diagnosis, and complications affecting the pregnant person or fetus.

\medterm{Perindopril} Perindopril is an ACE inhibitor used for high blood pressure and cardiovascular risk reduction; it can cause cough, high potassium, or rarely angioedema.

\medterm{Perineal} Relating to the area between the external genitals and anus, which can be affected by injury or infection.

\medterm{Perineal Laceration} Tearing of the perineal tissues during vaginal delivery. The injury is graded by depth:
it may involve skin, perineal muscle, the anal sphincter, or the rectal mucosa.

\medterm{Perineal Laceration Classification} A system grading tears during childbirth by depth: first degree affects skin, second degree also affects muscle, third degree involves the anal-control muscle, and fourth degree extends into the rectum. Accurate examination and repair reduce later pain and bowel-control problems.

\medterm{Perineal Protection} Measures used during the pushing stage of childbirth to reduce severe tearing of the tissue between the vagina and anus. They may include controlled delivery of the baby’s head, support of the perineum, a warm compress, and a selective approach to episiotomy.

\medterm{Perineal Tear} A split in the tissue between the vagina and anus during vaginal birth. First-degree tears affect skin, second-degree tears reach muscle, third-degree tears involve the anal-control muscle, and fourth-degree tears extend into the rectum; severe tears need careful repair and follow-up.

\medterm{Perineum} The area between the external genitalia and the anus. It contains skin, muscles, nerves, and connective tissue and may be affected by injury, infection, childbirth, or surgery.
\medterm{Perineural} Perineural means around a nerve; the term may describe normal connective tissue, inflammation, tumor spread, or imaging findings surrounding a nerve.

\medterm{Perineuronal Satellitosis} A microscopic pattern in which supporting brain cells cluster around a nerve cell. It can occur in some brain tumors, inflammation, or reactive changes and has meaning only with the surrounding tissue findings.

\medterm{Perinuclear Space} The narrow space between the two membranes surrounding a cell nucleus. It connects with part of the cell’s internal membrane network and helps form the nuclear envelope.


\medterm{Periodic Abstinence} Periodic abstinence is avoiding intercourse during estimated fertile days to reduce pregnancy likelihood; cycle timing is imperfect and does not prevent infections.

\medterm{Periodic Acid} A chemical used to stain tissue samples. In the periodic-acid–Schiff test, it helps highlight sugars, mucus, certain fungi, and supporting membranes so they can be seen under a microscope.

\medterm{Periodic Breathing} Repeated cycles of shallow or absent breaths followed by deeper breaths. It can occur normally for short periods in some infants or during sleep, but persistent episodes may reflect a brain, heart, or lung disorder.

\medterm{Periodic Fever, Aphthous Stomatitis, Pharyngitis, Adenitis Syndrome} A recurrent fever syndrome, usually beginning in childhood, characterized by periodic fever with some combination of mouth ulcers, sore throat, and swollen neck lymph nodes; it is commonly abbreviated PFAPA.

\medterm{PFAPA Syndrome} Abbreviation for periodic fever, aphthous stomatitis, pharyngitis, and adenitis syndrome, a recurrent fever disorder that usually begins in childhood.

\medterm{Periodic Limb Movement Disorder} Repeated involuntary limb movements during sleep that may disrupt rest.

\synonyms
PLMD
Sleep Disorder Periodic Limb Movement (Periodic Limb Movement Disorder)

\medterm{Periodic Paralysis Syndrome} Periodic paralysis syndromes are rare inherited or acquired disorders causing episodic muscle weakness associated with abnormal potassium handling or thyroid disease. Attacks may follow rest after exercise, meals, or stress; ECG, potassium testing, and genetic or endocrine evaluation guide care.

\medterm{Periodontal Abscess} A localized accumulation of pus within periodontal tissues, commonly associated with a
periodontal pocket, foreign material, or impaired drainage. It may cause pain, swelling,
suppuration, tooth mobility, and systemic symptoms and requires urgent dental
assessment.

\medterm{Periodontal Cyst} A periodontal cyst is a fluid-filled epithelial-lined lesion near the tooth-supporting tissues, with effects ranging from incidental discovery to local bone destruction.

\medterm{Periodontal Disease} Infection or inflammation of the tissues supporting teeth, including gums and jawbone. It can cause bleeding, swelling, bad breath, loose teeth, bone loss, and tooth loss, but treatment can slow progression.
\medterm{Periodontal Disease, Gingivitis} Alternate index label; see \textit{Gingivitis}.

\medterm{Periodontal Disease, Periodontitis} Alternate index label; see \textit{Periodontitis}.

\medterm{Periodontal Dressing} A protective material placed over the gums after selected dental or periodontal procedures. It can protect the surgical area while it heals, reduce irritation, and support comfort; it should be kept clean and reviewed if it becomes loose, painful, foul-smelling, or causes swelling.

\medterm{Periodontal Flap Surgery} Dental surgery in which the gum is gently lifted so a clinician can clean infection and deposits below the gumline and, when needed, reshape accessible bone before replacing the gum.
\medterm{Periodontal Ligament} Strong tissue fibers that attach a tooth’s root to the surrounding jawbone. They cushion chewing forces, help the body sense pressure on the tooth, and can become inflamed or damaged in gum disease or grinding.

\medterm{Periodontal Pocket} An abnormally deep space between a tooth and the surrounding gum caused by loss of attachment. Plaque can collect there, leading to bleeding, infection, bone loss, loose teeth, and eventual tooth loss.

\medterm{Periodontal Probing} A dentist gently measures the space between the gum and tooth with a marked instrument. Deeper spaces, bleeding, and loss of support can indicate gum disease; measurements help assess severity and monitor treatment.

\medterm{Periodontics} The dental specialty focused on preventing, diagnosing, and treating gum disease and tissues supporting the teeth.
\medterm{Periodontist} A periodontist is a dentist specializing in prevention, diagnosis, and treatment of gum disease and supporting tissues, including periodontal surgery and implants.

\medterm{Periodontitis} A chronic inflammatory disease caused by dysbiotic dental plaque that destroys the periodontal ligament and alveolar bone supporting the teeth. It can cause bleeding gums, periodontal pockets, tooth mobility, recession, bad breath, and tooth loss; severity and progression vary and require dental evaluation.

\synonyms
Periodontal Disease, Periodontitis

\medterm{Periodontitis (Gum Disease)} Alternate index label for Gum Disease.

\medterm{Periodontium} The tissues that hold and support a tooth, including gums, ligament, root covering, and jawbone. Disease of these tissues can cause bleeding, gum recession, loose teeth, bone loss, and tooth loss.
\medterm{Periodontium Disease} Disease of the tissues supporting the teeth, including gingiva, periodontal ligament, cementum, and alveolar bone; periodontitis can cause attachment loss and tooth loss.

\medterm{Perionychium} The skin and soft tissue surrounding a fingernail or toenail, including the nail folds and nearby cuticle; infection or inflammation causes paronychia.

\medterm{Perioperative} Perioperative means relating to the period before, during, and after an operation, including preparation, anesthesia, recovery, and complication prevention.

\medterm{Perioral Dermatitis} A papulopustular eruption around the mouth, sparing the vermilion border. It is more
common in young women and associated with topical corticosteroid use. Other triggers
include fluoridated toothpaste, cosmetics, and inhaler steroids.

\medterm{Periorbital} Located around the eye socket or eyelids. Periorbital swelling may result from injury, allergy, infection, or fluid retention.

\medterm{Periorbital Cellulitis} A bacterial infection of the eyelid and tissues around the eye that causes redness, swelling, and tenderness. It must be distinguished from orbital cellulitis, which can threaten vision and spread more deeply.

\medterm{Periorbital Edema} Swelling of the tissues around the eye, caused by allergy, infection, trauma, venous or lymphatic congestion, kidney disease, thyroid eye disease, or other systemic conditions.

\medterm{Periorbital Infection} Infection of the tissues around the eye; it can cause swelling and redness and must be distinguished from orbital cellulitis, which threatens vision.

\medterm{Periosteal Osteosarcoma} A cancerous bone tumor that develops beneath the thin tissue covering a bone, usually along the shaft of a long bone. It often makes cartilage and causes pain or swelling, requiring specialist cancer treatment.

\medterm{Periosteum} A blood-vessel-rich, pain-sensitive tissue layer covering the outer surface of bone, except where joint cartilage covers it. It supplies nutrients, supports repair, and provides attachment for tendons and ligaments.
\medterm{Periostitis} Inflammation of the periosteum, the pain-sensitive tissue covering the outer surface of bone. It may cause localized bone pain, tenderness, or swelling and can result from repeated stress, infection, injury, or another disease.
\medterm{Peripancreatic Malignancy} Cancer involving the tissues around the pancreas, either by spreading from the pancreas or from another organ. It may cause abdominal or back pain, weight loss, jaundice, or newly developed diabetes and usually needs detailed imaging and specialist care.

\medterm{Peripartum} Relating to the period around childbirth, including late pregnancy, labor, delivery, and the early weeks after birth. Clinicians may use a more specific time range for a particular condition or study.

\medterm{Peripartum Cardiomyopathy} A weakening and enlargement of the heart muscle that develops near the end of pregnancy or during the months after birth, without another clear cause. It can cause breathlessness, swelling, fatigue, or rapid heartbeat and needs prompt cardiac care.

\medterm{Peripartum Cardiomyopathy Prognosis} Recovery of ventricular function varies widely. Prognosis is influenced by the initial ejection
fraction, ventricular size, symptoms, treatment, complications, and follow-up; persistent
dysfunction requires ongoing cardiac care.

\medterm{Peripheral} Peripheral means away from the center or central nervous system; in clinical use it may describe nerves, vessels, lesions, circulation, or symptoms outside the brain, spinal cord, or central structures.

\medterm{Peripheral Arterial Line - Infants} A small catheter placed in an infant's artery to continuously measure blood pressure and obtain blood samples. It requires careful monitoring for bleeding, infection, or reduced blood flow to the limb.

\medterm{Peripheral Artery Bypass - Leg} Surgery that creates a new route around a blocked leg artery to improve blood flow. A vein or artificial graft may be used as the bypass.


\medterm{Peripheral Artery Disease} Narrowing of arteries supplying the limbs, usually caused by atherosclerosis. It may
cause exertional leg pain relieved by rest, reduced pulses, nonhealing wounds, or chronic limb-threatening ischemia, a severe reduction in blood flow associated with persistent pain or tissue loss. Also called peripheral arterial disease.

\synonyms
Peripheral Arterial Disease; PAD

\medterm{Peripheral Artery Disease - Legs} Peripheral artery disease of the legs is atherosclerotic reduction in arterial blood flow causing exertional claudication, diminished pulses, nonhealing wounds, or critical limb ischemia.


\medterm{Peripheral Blood} Peripheral blood is blood obtained from a vein or capillary rather than directly from an organ or marrow, used for counts, smears, chemistry, and other testing.

\medterm{Peripheral Blood Smear} Microscopic examination of stained blood spread on a slide. It shows the size, shape, and appearance of red cells, white cells, and platelets and can help investigate anemia, infection, and blood disorders.

\medterm{Peripheral Blood Stem Cells} Immature blood-forming cells collected from the bloodstream after medicines encourage them to leave the bone marrow. They can be returned to the same person or given to another person to rebuild blood production after intensive treatment.

\medterm{Peripheral Edema} Swelling due to accumulation of fluid in the interstitial space of the limbs, most
commonly the lower extremities.
\medterm{Peripheral Intravenous Line - Infants} A small tube placed in an infant's vein, usually in the hand, foot, or scalp, to give medicines or fluids. The site should be checked for swelling, leakage, or infection.

\medterm{Peripheral Ischemia} Inadequate arterial blood flow to an arm, leg, digit, or other peripheral tissue, causing pain, pallor, coolness, weak pulses, numbness, ulcers, or tissue loss.

\medterm{Peripheral Nerve} A nerve outside the brain and spinal cord that carries motor, sensory, and autonomic signals between the central nervous system and the rest of the body.

\medterm{Peripheral Nerve Blocks} Injections of local anesthetic near a nerve or group of nerves to numb a specific area and reduce pain during or after a procedure. Ultrasound guidance may improve accuracy and safety.

\medterm{Peripheral Nerve Injuries} Damage to nerves outside the brain and spinal cord caused by trauma, compression, stretching, disease, or toxins.
It may cause pain, numbness, weakness, or loss of function; recovery depends on the
severity and location of the injury.

\medterm{Peripheral Nerve Tumors} Growths arising from or around nerves outside the brain and spinal cord; they may be benign or cancerous.

\medterm{Peripheral Nervous System} All nerves outside the brain and spinal cord. It carries messages between the central nervous system and the skin, muscles, organs, and glands, controlling movement, sensation, and automatic functions such as heartbeat and digestion.

\medterm{Peripheral Neuritis (Peripheral Neuropathy)} Alternate index label for Peripheral Neuropathy.

\medterm{Peripheral Neurofibromatosis} Alternate index label; see \textit{Neurofibromatosis type 1}.

\medterm{Peripheral Neuropathy} Damage or dysfunction of nerves outside the brain and spinal cord. It may cause burning pain, tingling, numbness, weakness, or problems with sweating and other automatic body functions, often beginning in the feet or hands.

\synonyms
Neuropathy, Peripheral
Peripheral Neuritis (Peripheral Neuropathy)

\medterm{Peripheral Neuropathy Diabetic (Diabetic Neuropathy)} Alternate index label for Diabetic Neuropathy.

\medterm{Peripheral Oedema} Alternate spelling of Peripheral Edema; swelling caused by excess interstitial fluid, usually in the limbs.

\medterm{Peripheral Vascular Disease} Disease affecting narrowed or blocked blood vessels outside the heart and brain. It commonly reduces blood flow to the legs, causing pain with walking, wounds that heal slowly, coldness, or tissue loss.

\synonyms
Peripheral arterial disease
PVD (Peripheral Vascular Disease)

\medterm{Peripheral Vision} The parts of the visual field seen away from the direct line of sight. Loss can result from eye disease, nerve damage, brain injury, or glaucoma and may increase falls or collisions.

\synonyms
Side vision

\medterm{Peripherally Inserted Central Catheter} A long, flexible catheter inserted into the basilic or cephalic vein and advanced so the
tip sits in the superior vena cava. Inserted at the bedside, often by a specialized PICC
nurse, with ultrasound and ECG or X-ray confirmation. Indications: long-term IV
antibiotics (>2 weeks), TPN, chemotherapy, and frequent blood draws.

\medterm{Peripherally Inserted Central Catheter - Dressing Change} Replacing the sterile covering over a peripherally inserted central catheter while checking the skin and catheter position. A clean technique reduces infection risk; redness, drainage, pain, fever, leakage, or a changed catheter length requires prompt review.

\medterm{Peripherally Inserted Central Catheter - Flushing} Passing a prescribed sterile solution through a peripherally inserted central catheter to keep it open and clear blood or medicine residue. The catheter should be flushed only with the specified solution and technique; resistance, pain, swelling, leakage, or fever needs assessment.

\medterm{Peripherally Inserted Central Catheter - Insertion} Placement of a long flexible tube through a vein in the arm until its tip lies near the heart. It allows repeated intravenous medicines, nutrition, or blood tests, but can cause infection, blockage, bleeding, or a blood clot.

\medterm{Periplasm} The periplasm is the space between the inner and outer membranes of Gram-negative bacteria, containing enzymes, transport proteins, and structural components.

\medterm{Periprosthetic Joint Infection} Infection involving an artificial joint and nearby tissues. It may cause persistent pain, swelling, drainage, fever, or loosening. Diagnosis uses examination, blood tests, joint-fluid testing, cultures, and imaging; treatment may require antibiotics and surgery.

\medterm{Perirenal Abscess} A perirenal abscess is a collection of pus around the kidney, usually from spread of a urinary or kidney infection, causing fever, flank pain, and possible sepsis.

\medterm{Perirenal Extramedullary Hematopoiesis} Production of blood cells in tissue around the kidney instead of mainly in the bone marrow. It is an unusual response to severe long-term anemia or bone-marrow disease and may appear as masses on imaging.

\medterm{Perirenal Hematoma} A collection of blood around the kidney, usually caused by trauma, a bleeding lesion, a procedure, or anticoagulation.

\medterm{Peristalsis} Peristalsis is coordinated, wave-like contraction of smooth muscle that propels food, fluid, or other contents through the gastrointestinal tract and some ducts.

\medterm{Peritoneal} Peritoneal means relating to the peritoneum, the thin membrane lining the abdominal cavity and covering many abdominal organs.

\medterm{Peritoneal Carcinomatosis} Cancer that has spread widely across the lining of the abdomen, often from the bowel, stomach, ovary, or another organ. It may cause abdominal swelling, pain, fluid buildup, weight loss, or bowel blockage; scans and fluid or tissue tests guide treatment.

\medterm{Peritoneal Cavity} The narrow space inside the abdominal lining around the abdominal and pelvic organs. A small amount of fluid normally lubricates it; blood, pus, excess fluid, or cancer cells can collect there during disease.

\medterm{Peritoneal Dialysis} A treatment for kidney failure that uses the lining of the abdomen to remove waste and extra water. Clean dialysis fluid enters through a permanent tube and is later drained; it can be done at home but may cause infection, fluid imbalance, or catheter problems.

\medterm{Peritoneal Dialysis Catheter} A soft tube placed through the abdominal wall into the abdominal cavity for dialysis. Dialysis fluid enters and leaves through it to remove waste and extra water; infection, blockage, leakage, or displacement can occur.

\medterm{Peritoneal Diseases} Disorders affecting the thin lining of the abdomen, including infection, inflammation, cancer, bleeding, scar tissue, or abnormal fluid buildup. Symptoms depend on the cause and may include abdominal pain, swelling, fever, or vomiting.
\medterm{Peritoneal Fluid} The small amount of slippery fluid normally present inside the abdominal cavity. Excess fluid, blood, pus, or cancer cells can collect there and may cause swelling, pain, infection, or other signs of abdominal disease.

\medterm{Peritoneal Fluid Analysis} Laboratory testing of fluid collected from the abdominal cavity. Cell counts, protein, albumin, cultures, and other tests help determine whether the fluid is caused by liver disease, infection, cancer, bleeding, or another problem.

\medterm{Peritoneal Fluid Culture} Culture of ascitic fluid to detect microorganisms, particularly in suspected spontaneous bacterial peritonitis or secondary peritoneal infection.

\medterm{Peritoneal Infection} Infection of the lining of the abdominal cavity, often called peritonitis. Causes include a ruptured
appendix or bowel, abdominal surgery, pelvic infection, or contamination during dialysis.
Severe abdominal pain, fever, vomiting, or a rigid abdomen requires urgent assessment.

\medterm{Peritoneal Lavage} Washing or sampling the abdominal cavity with sterile fluid. It was once used more often to detect internal bleeding or intestinal contents after injury; imaging and other tests now commonly provide this information.

\medterm{Peritoneal Necrosis} Death of peritoneal tissue caused by ischemia, severe infection, inflammation, trauma, or tumor, potentially leading to peritonitis, adhesions, sepsis, or organ failure.

\medterm{Peritoneal Sac} The space formed by the abdominal lining around the abdominal and pelvic organs. It normally contains only a small amount of lubricating fluid, but blood, pus, excess fluid, or cancer can collect there.

\medterm{Peritoneum} A thin, slippery lining that covers the inside of the abdomen and many abdominal organs. It allows organs to move smoothly; inflammation, infection, bleeding, fluid, or cancer in this lining can cause severe abdominal symptoms.

\medterm{Peritoneum Neoplasm} An abnormal growth in the lining of the abdominal cavity. It may begin there or spread from cancer in another organ and can cause abdominal pain, swelling, fluid buildup, or bowel blockage.

\medterm{Peritonitis} Inflammation or infection of the abdominal lining, often caused by a perforated organ, surgery, appendicitis, abdominal injury, or infected abdominal fluid. It can cause severe pain, a swollen or rigid abdomen, fever, vomiting, sepsis, or shock and needs urgent care.

\medterm{Peritonitis - Secondary} Infection of the abdominal lining caused by leakage from a perforated or injured digestive or other abdominal organ. It can rapidly cause severe pain, a rigid abdomen, sepsis, or shock and requires urgent treatment.

\medterm{Peritonitis - Spontaneous Bacterial} Infection of fluid in the abdominal cavity without an obvious hole or other surgical source. It occurs most often in people with advanced liver disease and can cause pain, fever, confusion, or sepsis.

\medterm{Peritonitis In PD} Infection of the abdominal lining in a person receiving peritoneal dialysis. Cloudy drainage, abdominal pain, fever, nausea, or feeling unwell needs prompt testing of the drainage and antibiotics; severe or repeated infection may require catheter removal.

\medterm{Peritonsillar Abscess} A collection of pus beside a tonsil, usually after a throat infection. It can cause severe one-sided throat pain, fever, muffled speech, drooling, difficulty opening the mouth, or trouble breathing and requires urgent treatment and drainage.

\medterm{Periungual Fibroma} A noncancerous, firm growth arising beside or beneath a fingernail or toenail. It may distort the nail and can occur alone or as one of several skin findings in tuberous sclerosis complex.

\medterm{Periurethral} Located around the urethra, the tube that carries urine out of the body. The term may describe glands, infection, an abscess, a urethral diverticulum, or tissue and treatments used for urinary leakage.
\medterm{Periventricular Leukomalacia} Injury to the brain’s developing nerve fibers near its fluid-filled spaces, most often in premature infants. It may later affect movement, muscle tone, vision, learning, or development and is associated with cerebral palsy.

\medterm{Perivitelline Space} The perivitelline space is the narrow region between an egg’s plasma membrane and surrounding vitelline or zona material, important in fertilization and early development.

\medterm{Perleche} Inflammation, soreness, and cracking at one or both corners of the mouth. Saliva, ill-fitting dentures, yeast or bacterial infection, skin irritation, and iron or vitamin shortages can contribute; treatment depends on the cause and keeping the area dry.
\medterm{Permanent Teeth} The adult set of teeth that replaces the deciduous teeth.
\medterm{Permease} A permease is a membrane transport protein that facilitates movement of a specific molecule or ion across a cell membrane.

\medterm{Permethrin} A skin-applied medicine that kills scabies mites and head lice. It is usually applied as directed and washed off later; itching may continue briefly after treatment even when parasites are gone.
\medterm{Permissive Hypercapnia} A ventilator approach that allows carbon dioxide in the blood to remain above normal so that gentler pressures and smaller breaths can protect injured lungs. It is used selectively and requires close monitoring because excessive acidity or carbon dioxide can be harmful.


\medterm{Pernicious Anaemia} A form of vitamin B12 deficiency caused by immune damage to stomach cells needed for B12 absorption. It can cause anemia, tiredness, sore tongue, numbness, balance problems, and nerve damage.
\medterm{Pernicious Anemia} Anemia caused by vitamin B12 deficiency when the immune system damages stomach cells that make a substance needed for B12 absorption. It may cause tiredness, sore tongue, numbness, balance problems, or permanent nerve damage if untreated.

\synonyms
Addisons Anemia (Pernicious Anemia)

\medterm{Pernio} Alternate index label; see \textit{Chilblains}.

\medterm{Peroneal} Relating to the fibula or the muscles and nerve on the outer side of the lower leg.
\medterm{Peroneal Muscular Atrophy} Alternate index label; see \textit{Charcot-Marie-Tooth disease}.

\medterm{Peroneal Nerve} The common fibular nerve, a branch of the sciatic nerve that winds around the fibular neck and supplies sensation and muscles responsible for ankle dorsiflexion and eversion.

\medterm{Peroxidase} An enzyme that uses hydrogen peroxide or another peroxide to oxidize a substrate; thyroid peroxidase is essential for thyroid-hormone synthesis.

\medterm{Peroxisomal Disorder} An inherited condition in which small cell structures that process certain fats do not form or work properly. It may affect the brain, nerves, liver, eyes, hearing, bones, or development.

\medterm{Peroxisome} A small structure inside a cell that breaks down certain fats and harmful substances. It also helps control hydrogen peroxide and makes substances needed for cell membranes and nerve function.

\medterm{Peroxisome Proliferation} An increase in the number or size of peroxisomes, small cell structures that process certain fats and harmful chemicals. It can be triggered by some medicines, chemicals, hormones, or metabolic conditions and may alter cell function.

\medterm{Perphenazine} A first-generation (typical) phenothiazine antipsychotic used for schizophrenia and, in some settings, severe nausea or vomiting. Adverse effects can include sedation, orthostatic hypotension, extrapyramidal symptoms, tardive dyskinesia, and neuroleptic malignant syndrome.
\medterm{Perseveration} Involuntary repetition of a word, idea, movement, or response after the appropriate stimulus has changed, associated with frontal or basal-ganglia disease, aphasia, delirium, or psychiatric illness.

\medterm{Persistent Asthma} Asthma symptoms or airway narrowing that occur regularly rather than only occasionally and usually require ongoing controller treatment. Severity is judged by symptoms, night waking, reliever use, lung function, and activity limits; treatment is adjusted over time.

\medterm{Persistent Cloaca} Persistent cloaca is a congenital malformation in which the rectum, vagina, and urinary tract join into one channel, requiring specialist evaluation and reconstructive care.

\medterm{Persistent Depressive Disorder} A long-lasting depressive disorder, formerly called dysthymia, with depressed mood or irritability for at least two years in adults or one year in children and adolescents. Other symptoms may include low energy, poor self-esteem, sleep or appetite changes, poor concentration, or hopelessness.

\medterm{Persistent Pulmonary Hypertension Of The Newborn} Failure of the normal fall in pulmonary vascular resistance after birth, causing right-to-left shunting and hypoxaemia. Management requires urgent neonatal respiratory and circulatory support.

\medterm{Persistent Pupillary Membranes} Strands left from a temporary membrane that covered the developing eye before birth. They may cross the pupil and are usually harmless, but dense strands can interfere with vision and occasionally need treatment.

\medterm{Person With Disability} A person with disability has a long-term physical, sensory, intellectual, or mental-health impairment that may interact with barriers to limit participation; care should be accessible and person-centered.

\medterm{Personal Emergency Response System} A wearable or nearby device that lets a person summon help during an emergency.
\synonyms
Medical alert system

\medterm{Personal Protective Equipment} Clothing or equipment that protects a person from infection, chemicals, blood, or other hazards. It may include gloves, gowns, masks, respirators, goggles, or face shields; the correct items depend on the hazard and must be put on, removed, and disposed of safely.

\medterm{Personality} Personality is the relatively enduring pattern of thoughts, emotions, motivations, and behaviors through which a person experiences and responds to the world.

\medterm{Personality Change} Personality change is a new or marked alteration in behavior, emotional responses, judgment, or interpersonal style, which may reflect neurologic disease, psychiatric illness, substances, or medical conditions.

\medterm{Personality Disorder} Alternate index label for Personality Disorders.

\medterm{Personality Disorder Antisocial (Antisocial Personality Disorder)} Alternate index label for Antisocial Personality Disorder.

\medterm{Personality Disorder, Antisocial} Alternate index label; see \textit{Antisocial personality disorder}.

\medterm{Personality Disorder, Borderline} Alternate index label; see \textit{Borderline personality disorder}.

\medterm{Personality Disorder, Narcissistic} Alternate index label; see \textit{Narcissistic personality disorder}.

\medterm{Personality Disorder, Schizoid} Alternate index label; see \textit{Schizoid personality disorder}.

\medterm{Personality Disorder, Schizotypal} Alternate index label; see \textit{Schizotypal personality disorder}.

\medterm{Personality Disorders} A group of mental disorders involving enduring, inflexible patterns of thought, emotion, behavior, and relationships that differ markedly from the person’s cultural context and cause distress or impaired functioning. Traditional groupings include Cluster A (paranoid, schizoid, and schizotypal), Cluster B (antisocial, borderline, histrionic, and narcissistic), and Cluster C (avoidant, dependent, and obsessive-compulsive personality disorders).

\medterm{Personality Test} A psychological assessment using structured questions or stimuli to describe enduring traits, patterns, or possible personality pathology; results require validated instruments and clinical interpretation.

\medterm{Perspiration} Sweat released by skin glands. Eccrine glands mainly help cool the body, while apocrine glands contribute to odor after their secretions are changed by skin bacteria; unusually heavy sweating is called hyperhidrosis.
\medterm{Persuasive Communication} Communication intended to influence beliefs, decisions, or behavior through information, reasoning, emotion, or framing; healthcare use should preserve informed choice and avoid coercion.

\medterm{Pertussis} Alternate index label for Whooping Cough.
\medterm{Pertussis Toxin} A toxin produced by Bordetella pertussis that disrupts cell signaling and contributes to the cough and systemic effects of whooping cough.

\medterm{Pertuzumab} A monoclonal antibody targeting the HER2 receptor and blocking its dimerization, used with other medicines for selected HER2-positive cancers.

\medterm{Pervasive Developmental Disorders} An obsolete group term for childhood conditions involving social communication difficulties and restricted or repetitive behavior. Modern classification uses autism spectrum disorder and other specific developmental diagnoses.
\medterm{Pes} Pes means foot; it is also used in terms describing foot shape or deformity, such as pes planus or pes cavus.

\medterm{Pes Anserinus} The combined tendon insertion of the sartorius, gracilis, and semitendinosus muscles on the anteromedial proximal tibia; inflammation there is pes-anserine bursitis.

\medterm{Pes Planus} Alternate index label; see \textit{Flatfeet}.

\medterm{Pes Planus (Flatfoot (Pes Planus))} Alternate index label for Flatfoot (Pes Planus).

\medterm{Pessaries} Devices placed inside the vagina to support organs that have descended, such as the bladder, uterus, or rectum. Ring, cube, and other designs suit different shapes and needs; regular removal and follow-up help prevent sores, discharge, or infection.

\medterm{Pessary} A device placed inside the vagina to support pelvic organs, treat prolapse, or sometimes prevent leakage of urine.
\medterm{Pessimism} A tendency to expect negative outcomes or interpret situations unfavorably.


\medterm{Pesticide} A substance used to control insects, weeds, fungi, or other pests; exposure can cause poisoning.
\medterm{Pesticides} Chemicals used to kill insects, weeds, fungi, rodents, or other pests. Exposure can irritate the skin, eyes, lungs, or digestive tract; some pesticides can cause sweating, vomiting, weakness, seizures, or breathing failure.

\medterm{Pesticide Poisoning} Harm caused by excessive exposure to an insecticide, herbicide, fungicide, or other pesticide. Effects depend on the chemical and may include skin or eye injury, vomiting, sweating, weakness, seizures, breathing failure, or abnormal heart rhythms.

\medterm{Pesticides On Fruits And Vegetables} Small pesticide residues on properly washed produce usually do not cause acute poisoning. Swallowing concentrated pesticide or developing vomiting, sweating, breathing difficulty, weakness, or confusion requires immediate poison-control or emergency advice.

\medterm{Pestilence} Pestilence is an old term for a severe, widespread, often fatal infectious disease outbreak; modern usage usually specifies the pathogen and epidemiologic event.

\medterm{Pet Allergy} An allergic reaction to proteins in an animal's skin flakes, saliva, or urine, commonly causing rhinitis, conjunctivitis, or asthma.

\synonyms
Allergy, Pet

\medterm{PET Scan} An imaging study using a positron-emitting tracer to measure metabolism, receptor activity, or tissue function.
\medterm{PET Scan (Positron Emission Tomography)} Nuclear medicine imaging using radioactive tracers (most commonly FDG, a glucose analog)
to detect metabolically active tissues.

\synonyms
Positron Emission Tomography

\medterm{PET Scan For Breast Cancer} Positron-emission tomography, usually combined with computed tomography, used selectively to assess metabolically active breast cancer, staging, recurrence, or treatment response.

\medterm{PET Scanner} A machine that detects radiation from a tracer inside the body and creates images showing how actively tissues use substances such as glucose. It is used to assess selected cancers and other disorders.
\medterm{PET Tracers} Radioactive substances used during a PET scan to highlight particular body processes. Different tracers show glucose use, bone activity, or specific cell receptors; the appropriate tracer depends on the clinical question.

\medterm{PET/CT} A scan that combines PET, which shows active body processes, with CT, which shows detailed anatomy. It helps locate disease and assess selected cancers, heart conditions, and brain disorders; it uses a radioactive tracer and X-rays.

\medterm{PET/MRI} A scan that combines information about tissue activity with detailed MRI images of soft tissues. It is useful for selected brain, cancer, and heart problems and does not use X-rays, but takes longer than many other imaging tests.

\medterm{Petechia} A petechia is a tiny nonblanching red, purple, or brown spot caused by bleeding from small vessels into the skin or mucosa; numerous or rapidly spreading petechiae need assessment.

\medterm{Petechiae} Tiny nonblanching red or purple spots caused by small areas of bleeding into the skin or mucosa.
\medterm{Petit Mal} An older term for absence seizures, brief episodes of impaired awareness with characteristic generalized EEG activity.
\medterm{Petit Mal Seizure} Alternate index label; see \textit{Absence seizure}.


\medterm{Petrissage} A massage technique involving kneading, squeezing, or lifting soft tissues. It should be adapted or avoided when there is acute injury, infection, clot risk, or fragile skin.
\medterm{Petrolatum} A semisolid purified mixture of hydrocarbons used on skin as a protective ointment and moisture barrier. It reduces water loss, shields minor irritation, and supports healing of dry, cracked, or chafed skin.
\medterm{Petroleum Jelly Overdose} Swallowing a small amount of petroleum jelly usually causes little harm, but larger amounts can cause diarrhea or vomiting. If it enters the lungs, it may cause inflammation and breathing problems; contact poison control after a significant exposure.

\medterm{Petrositis} Inflammation or infection of the petrous portion of the temporal bone, usually arising from complicated otitis media and potentially affecting nearby cranial nerves.

\medterm{Petrous} Describing dense, hard bone, especially the part of the temporal bone that houses inner-ear structures.
\medterm{Pets And The Immunocompromised Person} People with weakened immunity can often keep pets but should reduce exposure to animal waste, scratches, bites, and unsafe food and follow individualized medical advice.

\medterm{Peutz-Jeghers Syndrome} An inherited condition causing dark spots on the lips, mouth, fingers, or skin and multiple growths in the digestive tract. These growths can bleed or block the bowel and increase the risk of several cancers.

\synonyms
PJS

\medterm{Peutz-Jeghers's Syndrome} An inherited condition causing dark patches around the mouth and many growths in the stomach or intestines. The growths may bleed, cause blockage, or increase the risk of digestive and other cancers.

\medterm{PEWS} Alternate name; see \textit{Pediatric Early Warning Score}.

\medterm{Pexophagy} A cell-cleaning process that removes damaged or unneeded peroxisomes, the small structures that process certain fats and harmful chemicals. The cell breaks them down and reuses some of their materials.

\medterm{Peyer Patches} Clusters of immune tissue in the wall of the lower small intestine. They sample material from the gut and help the immune system recognize and respond to germs while tolerating normal food and bacteria.

\medterm{Peyer's Patches} Organized lymphoid tissue in the mucosa and submucosa of the ileum. These aggregates sample intestinal contents and contribute to immune surveillance of the gastrointestinal tract.

\synonyms
ileal lymphoid nodules

\medterm{Peyronie Disease} Scar tissue forming inside the penis, causing a new bend or indentation that may make erections painful or intercourse difficult. Symptoms may stabilize or worsen; treatment ranges from observation and injections to surgery in selected severe cases.

\medterm{Peyronie’S Disease} A condition in which scar tissue inside the penis causes curved, painful erections and sometimes erectile difficulty.

\synonyms
Peyronie disease

\medterm{Pfeiffer Syndrome} An autosomal dominant craniosynostosis syndrome caused by FGFR1 or FGFR2 mutations.
Features: craniosynostosis (bilateral coronal), midface hypoplasia, proptosis,
hypertelorism, broad thumbs and great toes (medial deviation), and variable syndactyly.

\synonyms
Syndrome Pfeiffer (Pfeiffer Syndrome)

\medterm{PH Meter} An instrument that measures acidity or alkalinity as pH in a solution or specimen; clinical applications include gastric, urinary, laboratory, and environmental measurements.

\medterm{PH Monitoring} Measurement of acidity, usually in the esophagus, over time with a small sensor. It can help detect and measure acid reflux when symptoms or endoscopy do not give a clear answer.


\medterm{Phacoemulsification} A cataract-surgery technique that uses ultrasound to break up and remove the cloudy lens through a small incision. An
intraocular lens is usually implanted to restore focusing ability.

\medterm{Phagocyte} An immune cell that surrounds, swallows, and breaks down germs, dead cells, or foreign particles. Phagocytes help prevent infection, remove damaged material, and alert other immune cells to danger.
\medterm{Phagocytes} Immune cells that surround, swallow, and break down germs, dead cells, and foreign particles. They help prevent infection, clear damaged material, and signal other immune cells during inflammation and tissue repair.

\medterm{Phagocytic Vesicle} A small membrane-bound compartment formed inside a cell after it surrounds and swallows a germ, dead cell, or particle. It later joins a digestive compartment so the material can be broken down.

\medterm{Phagocytosis} The process by which a cell surrounds, takes in, and breaks down a particle or microorganism. Immune cells use it to remove germs, dead cells, and harmful debris from tissues.
\medterm{Phagolysosome} A compartment formed when an immune cell’s swallowed-particle compartment joins a digestive compartment. Enzymes and other chemicals inside it destroy germs and break down cellular debris.

\medterm{Phagophore} A temporary membrane that grows around damaged or unnecessary material inside a cell during recycling. It closes around the material so the cell can break it down and reuse some of its parts.

\medterm{Phalanges} The bones of the fingers and toes. The thumb and big toe each have two, while the other digits usually have three. Fractures, arthritis, infection, or nerve problems in these bones can cause pain, swelling, stiffness, or difficulty using the hand or foot.
\medterm{Phalanx} One of the bones of a finger or toe. Each phalanx has a base, shaft, and end; the end bone supports the fingertip or toenail. The thumb and big toe have two phalanges, while other digits usually have three.
\medterm{Phalanx Unspecified} A phalanx is one of the bones of a finger or toe; “unspecified” indicates that the particular digit or proximal, middle, or distal segment is not identified.

\medterm{Phalen Test} A test in which the wrists are bent to see whether numbness or tingling develops. Reproduced symptoms may support carpal tunnel syndrome, but the result must be interpreted with examination and other tests.

\medterm{Phantom} A test object made to resemble human tissue or organs so medical imaging equipment can be checked and calibrated. Different phantoms test image quality, measurements, or specific body regions.

\medterm{Phantom Limb} The feeling that an amputated body part is still present. Sensations may include position, movement, itching, pressure, or pain and often change over time after amputation.

\synonyms
Phantom-limb sensation

\medterm{Phantom Limb Pain} Painful sensations perceived in a limb that has been amputated, caused by persistent or
reorganized activity in peripheral nerves and central sensory pathways. It may feel
burning, shooting, cramping, or electric.

\medterm{Phantom Pain} Pain perceived in a body part that has been amputated, caused by changes in peripheral nerves and central pain-processing pathways.

\medterm{Phantosmia} Perception of a smell when no corresponding odor is present. It may occur with sinonasal
disease, migraine, seizures involving olfactory networks, head injury, neurodegenerative
disease, infection, or other neurological conditions.

\medterm{Pharmaceutical Adjuvant} A pharmaceutical adjuvant is an inactive or auxiliary ingredient added to improve a medicine’s stability, absorption, preservation, delivery, or therapeutic effect.

\medterm{Pharmaceutical Excipient} A pharmaceutical excipient is a non-active ingredient in a dosage form, such as a filler, binder, coating, solvent, preservative, or disintegrant.

\medterm{Pharmaceutical Solutions} Liquid preparations in which one or more active pharmaceutical ingredients are dissolved in a suitable vehicle for administration, with concentration, sterility, stability, and route determining safe use.

\medterm{Pharmacist} A medicine expert who prepares and dispenses drugs, checks safety, and explains correct use and possible side effects.
\medterm{Pharmacodynamics} The study of the biochemical and physiological effects of drugs on the body and their mechanisms of action, describing what the drug does to the body.
\medterm{Pharmacogenetic Testing} Testing for genetic variants that influence a person’s response, metabolism, efficacy, or toxicity of particular medicines and may support individualized prescribing.

\medterm{Pharmacogenetics} Pharmacogenetics studies how inherited genetic variation alters a person's response to medicines, including efficacy, metabolism, and adverse-effect risk, to support individualized prescribing.

\medterm{Pharmacogenomic Testing} Testing that examines genetic variants affecting drug metabolism, transport, targets, or adverse reactions to help select or dose certain medicines.

\medterm{Pharmacogenomics} The study of how genetic variations influence individual responses to drugs, including
drug efficacy, toxicity, and metabolism. Genetic polymorphisms in drug-metabolizing
enzymes, transporters, and receptors can lead to significant inter-individual
variability in drug response.

\medterm{Pharmacokinetics} The study of what the body does to a medicine, including its absorption, distribution, metabolism, and elimination. It helps explain drug concentrations over time and supports dose selection.

\medterm{Peak Plasma Concentration} The highest measured concentration of a medicine in the blood after a dose. It is often called Cmax and is influenced by the dose, absorption speed, formulation, and route of administration.
\medterm{Pharmacologic Substance} A pharmacologic substance is a chemical that produces a biological effect by interacting with cells, receptors, enzymes, transporters, or other targets.

\medterm{Pharmacological Stress Test} A heart stress test using medicine to imitate exercise when a person cannot exercise adequately. Imaging or heart recordings show blood flow and function, while possible effects include chest discomfort, breathlessness, or abnormal rhythm.

\medterm{Pharmacologist} A pharmacologist studies how medicines and chemicals act in living systems, including mechanisms, metabolism, toxicity, and therapeutic effects.

\medterm{Pharmacology} The study of medicines and other drugs, including how they work, how the body handles them, their uses, side effects, and interactions. This knowledge helps guide safe choice, dosing, testing, and monitoring of treatment.
\medterm{Pharmacopoeia} A pharmacopoeia is an official reference specifying standards for identity, purity, strength, testing, and preparation of medicines.

\medterm{Pharmacotherapy} Pharmacotherapy is treatment with medicines chosen to prevent, relieve, or control a disease or symptom.
\medterm{Pharmacovigilance} The science and activities relating to the detection, assessment, understanding, and
prevention of adverse drug reactions and other drug-related problems.

\medterm{Pharmacy} The healthcare field concerned with preparing, checking, supplying, and explaining medicines. Pharmacists help choose safe doses, identify interactions and allergies, support treatment decisions, and teach people how to use medicines correctly.
\medterm{Pharmacy Facility} A pharmacy facility is an authorized location where medicines may be stored, prepared, dispensed, monitored, and supplied under applicable professional and safety standards.

\medterm{Pharyngeal} Relating to the pharynx, or throat. The pharynx connects the nose and mouth with the voice box, food pipe, and breathing passages. Infection, swelling, narrowing, or a tumor there may cause sore throat, swallowing difficulty, noisy breathing, or voice changes.
\medterm{Pharyngeal Carcinoma} Pharyngeal carcinoma is cancer arising in the throat and may cause persistent sore throat, swallowing difficulty, hoarseness, or a neck mass.

\medterm{Pharyngeal Disease} A disorder affecting the throat, including infection, inflammation, blockage, injury, muscle or nerve problems, or a tumor. Symptoms may include sore throat, swallowing difficulty, choking, voice changes, or breathing problems.

\medterm{Pharyngeal Hemorrhage} Bleeding from the pharynx caused by trauma, surgery, infection, vascular lesions, tumor, or severe mucosal disease and potentially threatening the airway.

\medterm{Pharyngeal Mucositis} Inflammation and ulceration of the throat lining, commonly caused by radiation, chemotherapy, infection, trauma, or medications and producing pain, difficulty swallowing, and reduced intake.

\medterm{Pharyngeal Stenosis} Narrowing of the pharyngeal passage caused by scarring, congenital abnormality, tumor, inflammation, or surgery, which may impair swallowing or breathing.

\medterm{Pharyngeal Tonsil} Adenoid tissue in the roof of the throat behind the nose that helps immune defense, especially in children.

\medterm{Pharyngectomy} Surgery to remove part or all of the throat, usually to treat cancer. Depending on the extent, reconstruction or a temporary feeding tube may be needed, and treatment can affect swallowing, speech, breathing, or appearance.
\medterm{Pharyngitis} Inflammation of the throat behind the mouth and nose. It commonly causes sore throat, painful swallowing, fever, cough, or swollen neck glands and is most often caused by a virus.

\synonyms
Sore Throat (Pharyngitis)

\medterm{Pharyngitis (Sore Throat (Pharyngitis))} Alternate index label for Sore Throat (Pharyngitis).

\medterm{Pharyngitis - Sore Throat} Pharyngitis is inflammation of the throat causing pain with swallowing, commonly from viral infection and sometimes from group A streptococcal infection or other causes.

\medterm{Pharyngitis - Viral} Viral pharyngitis is inflammation of the throat caused by a virus, commonly producing sore throat, cough, runny nose, or fever and usually resolving with supportive care.

\medterm{Pharyngolaryngeal Pain} Pain involving the pharynx and larynx, commonly associated with infection, inflammation, irritation, trauma, reflux, or malignancy.

\medterm{Pharyngomaxillary Space Abscess} A pharyngomaxillary space abscess is a deep infection containing pus near the pharynx and maxilla that can cause severe pain, swelling, difficulty swallowing, and airway risk.

\medterm{Pharynx} The muscular passage behind the nose and mouth that carries air toward the voice box and food toward the esophagus. It helps with breathing, swallowing, and speech.
\medterm{Pharynx Necrosis} Death of tissue in the throat caused by severe infection, ischemia, radiation, caustic injury, or tumor, with risk of deep infection, bleeding, airway compromise, and sepsis.

\medterm{Pharynx Neoplasm} A pharynx neoplasm is an abnormal growth in the throat that may be benign or malignant.

\medterm{Phase I Metabolism} An early step in the body’s breakdown of many medicines. Enzymes, mainly in the liver, chemically change a medicine by adding or exposing a reactive part; this may inactivate it, activate it, or prepare it for further breakdown.

\medterm{Phase II Metabolism} A stage of drug breakdown in which the body attaches another chemical group to a medicine or its earlier breakdown product. This usually makes the substance easier to remove in urine or bile.

\medterm{Phase-Contrast Microscopy} A microscope technique that makes transparent cells and other specimens easier to see by converting small differences in light passing through them into visible contrast. It is useful for examining living, unstained cells.

\medterm{Phe} Phe is the standard abbreviation for phenylalanine, an essential amino acid incorporated into proteins and metabolized through a pathway disrupted in phenylketonuria.

\medterm{Phelan-McDermid Syndrome} A genetic disorder that commonly causes low muscle tone, delayed development, intellectual disability, very limited speech, and autism-related behaviors. Some people have reduced sensitivity to pain, seizures, feeding problems, or kidney and heart abnormalities. It usually results from loss of genetic material near the end of chromosome 22 or a change in the \textit{SHANK3} gene.

\synonyms
22q13.3 Deletion Syndrome

\medterm{Phenacetin} Phenacetin is a former analgesic and fever-reducing drug withdrawn or restricted because long-term use can cause kidney damage and increase cancer risk.

\medterm{Phenazopyridine} Phenazopyridine is a urinary analgesic that temporarily relieves dysuria without treating its cause; it can turn urine orange and prolonged or excessive use may cause haemolysis or methemoglobinaemia.

\medterm{Phencyclidine} A dissociative drug that can cause altered perception, agitation, hallucinations, impaired coordination, or dangerous behavior.
\medterm{Phencyclidine Overdose} A dangerous amount of the dissociative drug phencyclidine can cause agitation, hallucinations, high body temperature, high blood pressure, seizures, loss of coordination, coma, or injury. Severe or unusual behavior requires emergency care.

\medterm{Phenelzine} An antidepressant that blocks monoamine oxidase, an enzyme breaking down several brain chemicals. It has important interactions with medicines and tyramine-rich foods and can cause dangerous high blood pressure if precautions are ignored.
\medterm{Phenethylamines} Phenethylamines are a chemical family including neurotransmitters, medicines, and psychoactive substances that can affect adrenergic, dopaminergic, or serotonergic signaling.

\medterm{Phenformin} Phenformin is a former biguanide diabetes medicine withdrawn because it can cause life-threatening lactic acidosis.
\medterm{Phenindione} An older oral anticoagulant that interferes with vitamin-K-dependent clotting factors.
\medterm{Pheniramine Overdose} Taking too much pheniramine, an antihistamine, can cause extreme sleepiness or agitation, dry mouth, large pupils, fast heartbeat, seizures, abnormal rhythms, or coma. Immediate poison-control or emergency advice is needed.

\medterm{Phenobarbital} Phenobarbital is a long-acting barbiturate antiseizure medicine that enhances GABA-A receptor inhibition; sedation, respiratory depression, cognitive effects, dependence, and enzyme-inducing drug interactions are important risks.

\medterm{Phenobarbital Overdose} Taking too much phenobarbital can cause severe drowsiness, poor coordination, low blood pressure, slowed breathing, coma, or death. Emergency treatment is required, especially if alcohol or another sedative was also taken.

\medterm{Phenobarbitone} The British name for phenobarbital, a barbiturate anticonvulsant and sedative.

\medterm{Phenol} A corrosive aromatic compound used in chemical and laboratory applications and in some medical preparations.

\medterm{Phenols} A group of aromatic chemicals that can burn skin and mucous membranes and cause systemic poisoning if absorbed or swallowed.

\medterm{Phenomenon} An observable event, sign, or pattern in the body. Examples include fingers changing color in Raynaud phenomenon, early-morning blood-sugar rises in the dawn phenomenon, and upward eye movement when the eyelids close in Bell phenomenon.
\medterm{Phenothiazine} Phenothiazine is a chemical class that includes several antipsychotic and antiemetic medicines, with possible sedation, movement, and cardiac adverse effects.

\medterm{Phenothiazine Overdose} Taking too much of a phenothiazine medicine can cause sleepiness, confusion, low blood pressure, seizures, abnormal heart rhythms, or severe muscle stiffness and fever. It requires immediate medical assessment.

\medterm{Phenothiazines} A class of medicines that includes antipsychotics and antiemetics. They can cause sedation, low blood pressure, movement disorders, heart-rhythm changes, and other adverse effects.
\medterm{Phenotype} The observable features of a person, such as appearance, growth, body functions, and behavior. These features result from the interaction between inherited genetic information and the environment.

\medterm{Phenoxybenzamine} A long-acting alpha-blocker used mainly before surgery for a catecholamine-secreting tumor such as pheochromocytoma. It can cause dizziness, low blood pressure, rapid heartbeat, and nasal congestion.
\medterm{Phenoxymethylpenicillin} Penicillin V, an oral penicillin used for selected susceptible infections.
\medterm{Phenprocoumon} Phenprocoumon is a long-acting vitamin-K antagonist anticoagulant that reduces clotting-factor production and requires coagulation monitoring.
\medterm{Phentermine} A stimulant medicine that reduces appetite and may be used short term with diet and activity changes for weight management. It can raise heart rate and blood pressure and may cause insomnia, anxiety, dependence, or dangerous interactions with some medicines.
\medterm{Phenthoate} Phenthoate is an organophosphate insecticide that can inhibit acetylcholinesterase and cause cholinergic poisoning after exposure.
\medterm{Phentolamine} Phentolamine is a short-acting alpha-adrenergic blocker used for catecholamine-related hypertension and to limit tissue injury after certain vasoconstrictor leaks.

\medterm{Phenylacetate} A salt or ester of phenylacetic acid used in some metabolic and experimental medical applications; it can bind nitrogen-containing metabolites.

\medterm{Phenylalanine} An essential amino acid incorporated into proteins and converted to tyrosine; impaired metabolism causes phenylketonuria and requires dietary or medical management.

\medterm{Phenylbutazone} An older anti-inflammatory pain medicine whose use is limited by serious toxicity. It can cause stomach bleeding, kidney or liver injury, fluid retention, blood-cell problems, allergic reactions, and other dangerous effects.
\medterm{Phenylbutyrate} A nitrogen-scavenging medicine converted to phenylacetate, which combines with glutamine and facilitates nitrogen excretion in selected urea-cycle disorders.

\medterm{Phenylephrine} A medicine that narrows blood vessels. It may reduce a blocked nose, raise dangerously low blood pressure in hospital care, or widen the pupil in an eye examination; it can worsen high blood pressure, heart problems, or urinary difficulty in some people.
\medterm{Phenylketonuria} An inherited metabolic disorder, usually caused by disease-causing variants in the \textit{PAH} gene that reduce phenylalanine hydroxylase activity. Phenylalanine accumulates in the blood and tissues, while tyrosine availability may decrease. Untreated or poorly controlled disease can affect brain development and may be associated with intellectual disability, seizures, eczema, behavioral or psychiatric symptoms, and a musty odor. Severity ranges from classic to milder forms, and the disorder is commonly identified through newborn screening.

\synonyms
Phenylketonuria (PKU)
PKU

\medterm{Phenylpropanolamine} A sympathomimetic decongestant formerly used in some products; it can raise blood pressure and has been restricted because of safety concerns.

\medterm{Phenytoin} A medicine used to prevent or control some seizures. It reduces repeated electrical activity in brain cells; levels in the blood may need monitoring because interactions and side effects can include dizziness, poor coordination, gum overgrowth, skin reactions, or serious toxicity.
\medterm{Phenytoin Overdose} Too much phenytoin can cause nausea, severe dizziness, double vision, involuntary eye movements, poor coordination, confusion, abnormal heart rhythms, or coma. Significant or intentional overdoses need urgent medical assessment.

\medterm{Phenytoin Sodium} The sodium salt of phenytoin, an anticonvulsant used for selected focal and generalized tonic-clonic seizures; toxicity can cause nystagmus, ataxia, arrhythmia, or severe skin reactions.

\medterm{Pheochromocytoma} A rare tumor, usually in an adrenal gland, that may release excess adrenaline-like hormones. It can cause episodes of high blood pressure, headache, sweating, a pounding heartbeat, shakiness, or paleness.

\synonyms
Tumor Adrenal Gland (Pheochromocytoma)

\medterm{Pheresis} A procedure that removes a selected part of blood, such as plasma, platelets, or white blood cells, and returns the remaining blood. It may treat disease or collect a needed blood component from a donor.

\medterm{Pheromone} A pheromone is a chemical signal released by an organism that alters behavior or physiology in another member of the same species.


\medterm{Phialophora} A genus of darkly pigmented fungi that can cause chromoblastomycosis or other subcutaneous and opportunistic infections.

\medterm{Phialophora Verrucosa} A dark-pigmented fungus that commonly causes chromoblastomycosis, a long-lasting infection of skin and tissue beneath it. Infection usually follows implantation through injured skin and produces slowly enlarging warty or raised lesions.

\medterm{Philadelphia Chromosome} A rearrangement of genetic material between chromosomes 9 and 22 that creates the BCR::ABL1 gene. It is found in most cases of chronic myeloid leukemia and some acute leukemias and helps guide treatment.

\medterm{Philodendron Poisoning} Chewing philodendron leaves releases irritating crystals that can cause immediate burning and swelling of the mouth, tongue, or throat. Breathing or swallowing difficulty requires emergency care; contact poison control after ingestion.

\medterm{Phimosis} Tightness of the foreskin that prevents it from being pulled back over the head of the penis. It can be normal in young children, but pain, repeated infection, scarring, or difficulty passing urine needs medical assessment.
\medterm{Phimosis And Paraphimosis (Penis Disorders)} Alternate index label for Penis Disorders.

\medterm{Phlebitis} Inflammation of a vein causing pain, warmth, redness, or tenderness along its course. It may follow irritation from an intravenous line or occur with a clot in a surface vein; sudden limb swelling or breathlessness needs urgent assessment.
\medterm{Phlebitis And Thrombophlebitis} Phlebitis is inflammation of a vein, while thrombophlebitis includes a clot within the inflamed vein. Superficial disease causes localized pain and redness; deep-vein thrombosis can embolize to the lungs and is suspected with limb swelling, pain, or breathlessness.

\medterm{Phlebography} An imaging test in which contrast dye is injected into a vein to show the veins on X-ray. It can detect a blood clot or an abnormal vein but is now often replaced by ultrasound.
\medterm{Phlebotomy} Taking blood from a vein, usually with a needle, for testing, donation, or treatment. The procedure may cause brief pain, bruising, bleeding, fainting, or rarely infection; trained staff use safe collection and labeling practices.

\synonyms
Venipuncture

\medterm{Phlebovirus} A genus of segmented RNA viruses, many transmitted by sandflies or mosquitoes; some cause fever, encephalitis, or hemorrhagic disease.

\medterm{Phlegm} Thick mucus brought up by coughing from the lower airways. Its amount and appearance can change with infection, asthma, smoking, chronic lung disease, or other conditions.
\medterm{Phlegmasia Cerulea Dolens} Phlegmasia cerulea dolens is a severe, extensive deep-vein thrombosis that obstructs venous outflow, causing painful swollen blue discoloration and risking limb ischemia or pulmonary embolism.

\medterm{Phobia} An excessive, persistent fear of a specific object or situation that leads to avoidance or distress.
\medterm{Phobia - Simple/Specific} A specific phobia is intense, persistent fear of a particular object or situation that is excessive, leads to avoidance, and interferes with life.

\medterm{Phobia, Social} Alternate index label; see \textit{Social anxiety disorder (social phobia)}.

\medterm{Phobias} Persistent, excessive fear of a specific object or situation that leads to avoidance or significant distress.

\medterm{Phobic Disorder} A disorder characterized by clinically significant fear and avoidance of a particular object or situation.
\medterm{Phocomelia} A rare congenital limb-development defect in which hands or feet are attached close to the trunk because of absent or severely shortened long bones.
\medterm{Phonation} Production of voiced sound when exhaled air makes the vocal folds vibrate in the larynx; disorders cause hoarseness, weak voice, or loss of voice.

\medterm{Phonocardiogram} A graphic recording of heart sounds and murmurs, traditionally made with a microphone over the chest. It can document timing and duration but is usually interpreted with examination and echocardiography.
\medterm{Phonocardiography} Recording heart sounds and murmurs with a microphone, producing a tracing that can help evaluate heart valve problems.

\medterm{Phonological Disorder} A communication disorder involving persistent difficulty producing or organizing speech sounds beyond what is expected for age and language development.

\medterm{Phonophobia} Unusual sensitivity to sound or distress caused by ordinary sounds. It often occurs with migraine but may also accompany concussion, infection, anxiety, or other conditions; new severe sound sensitivity with headache or neurological symptoms needs assessment.

\medterm{Phorate} Phorate is a highly toxic organophosphate pesticide whose exposure can produce excessive cholinergic activity and medical emergency.


\medterm{Phorbols} Phorbols are diterpene compounds related to phorbol esters that can activate cellular protein kinase C signaling.

\medterm{Phorometer} An instrument used to measure ocular alignment and the amount of prism correction needed for binocular vision disorders.

\medterm{Phosgene} A highly toxic industrial gas that reacts with moist respiratory tissues and can cause delayed pulmonary edema, cough, chest tightness, and respiratory failure after exposure.

\medterm{Phosgene Oxime} A highly corrosive chemical warfare agent that causes immediate severe pain, skin and eye injury, airway damage, and tissue necrosis on contact.

\medterm{Phosmet} Phosmet is an organophosphate insecticide that can inhibit acetylcholinesterase and cause cholinergic toxicity after substantial exposure.

\medterm{Phosphamidon} A poisonous organophosphate insecticide. Exposure can cause sweating, excess saliva, vomiting, diarrhea, small pupils, muscle twitching, breathing difficulty, seizures, or weakness because it disrupts nerve signals; suspected poisoning needs urgent medical help.

\medterm{Phosphate} An inorganic ion and component of nucleotides, phospholipids, proteins, and bone mineral. It participates in ATP energy transfer, phosphorylation-dependent signaling, intracellular buffering, and acid excretion by the kidney.
\medterm{Phosphate Binders} Medicines taken with meals that bind phosphate in food inside the intestine, reducing its absorption. They are used mainly in advanced kidney disease when blood phosphate is high. The choice depends on calcium levels, vascular calcification, other illnesses, and treatment response.

\medterm{Phosphate Homeostasis} The body keeps blood phosphate within a useful range through the coordinated actions of the parathyroid glands, vitamin D, kidneys, intestines, and bones. Kidney disease can disrupt this balance.

\medterm{Phosphatidic Acid} A membrane phospholipid and signaling intermediate that serves as a precursor for other glycerolipids and influences membrane curvature, trafficking, and cell growth.

\medterm{Phosphatidylcholine} A major membrane phospholipid and component of pulmonary surfactant, lipoproteins, and bile; it also participates in cell signaling.

\medterm{Phosphatidylethanolamines} Membrane phospholipids containing ethanolamine that help determine membrane structure, fusion, and cell signaling. A group of cell-membrane fats that help maintain membrane shape, join membranes, and support cell signaling.
\medterm{Phosphatidylinositol} A membrane phospholipid that serves as a precursor for signaling molecules and helps anchor proteins to cell membranes.

\medterm{Phosphatidylserine} A membrane phospholipid normally concentrated on the inner cell surface that becomes exposed during apoptosis and helps signal cell removal.

\medterm{Phosphatidylserines} Molecules of phosphatidylserine, a membrane phospholipid involved in apoptosis, coagulation, membrane structure, and cell signaling.
\medterm{Phosphaturic Mesenchymal Tumor} A rare tumor that releases a substance causing the kidneys to lose too much phosphate. Low phosphate can soften bones and cause bone pain, weakness, or fractures; removing the tumor often corrects the problem.

\medterm{Phosphine} A highly toxic gas used in some industrial or pest-control settings that can cause severe lung, heart, liver, kidney, and neurologic injury.

\medterm{Phosphodiesterase-5 Inhibitors} Medicines that block an enzyme called PDE-5, increasing a chemical signal that relaxes blood vessels. They are used for erectile dysfunction and some forms of high blood pressure in the lung; nitrates must not be combined with them.

\medterm{Phosphofructokinase-1} A key enzyme that controls how cells break down sugar to produce energy, especially when energy demand increases.

\medterm{Phospholipase C Pathway} A system cells use to turn some hormone or nerve signals at their surface into internal actions. It can change calcium levels and control contraction, secretion, growth, or other cell responses.

\medterm{Phospholipid} A type of lipid with a water-attracting phosphate-containing head and water-repelling fatty-acid tails. Phospholipids form the basic double-layered structure of cell membranes and also help with cell signaling and transport.
\medterm{Phosphorus} An essential element in bone, nucleic acids, phospholipids, and ATP.
\medterm{Phosphorus Blood Test} A blood test measuring inorganic phosphate, interpreted with calcium, parathyroid hormone, vitamin D, and kidney function to evaluate mineral and renal disorders.

\medterm{Phosphorus In Diet} Dietary phosphorus comes from protein-rich foods and many processed foods; people with kidney disease may need individualized limits.

\medterm{Phosphorylation} Addition of a phosphate group to a molecule, often regulating protein activity or energy transfer.
\medterm{Photoallergy} An immune-mediated skin reaction that occurs when ultraviolet light activates a sensitizing substance on or
in the skin, causing an eczematous rash.

\medterm{Photocoagulation} Treatment that uses intense light, usually a laser, to heat and seal targeted tissue. In the eye, it can close leaking or abnormal blood vessels and limit damage from conditions such as diabetic retinopathy.

\medterm{Photodermatosis} A skin disorder caused or aggravated by exposure to ultraviolet or visible light, including polymorphous light eruption, phototoxic reactions, and photoallergic reactions.

\medterm{Photodynamic Therapy} A treatment that combines a photosensitizing agent with light to generate reactive
oxygen species and destroy cancer cells, used for actinic keratosis, superficial skin
cancers, and selected visceral tumors.

\medterm{Photodynamic Therapy For Cancer} Treatment in which a light-sensitive medicine collects in selected abnormal cells and is activated by a specific light, producing substances that destroy nearby cancer cells. It is used for certain accessible cancers and may cause pain, swelling, or temporary light sensitivity.

\medterm{Photoelectric Effect} An interaction in which tissue absorbs an X-ray and releases an electron. It contributes to how much radiation is absorbed and helps create contrast between tissues on an X-ray image.

\medterm{Photographic Fixative Poisoning} Harm from swallowing or inhaling chemicals used to process photographs. Effects depend on the product but may include irritation, vomiting, breathing problems, or chemical burns; keep the container and contact poison control.

\medterm{Photomultiplier Tube} A sensitive device that converts tiny flashes of light into electrical signals. In some medical scanners, groups of these devices help detect and locate radiation emitted by a tracer.

\medterm{Photopeak} The part of a radiation-measurement graph showing photons that deposited their full energy in the detector. Clinicians use it to select the energy range for interpreting some nuclear medicine scans.

\medterm{Photopheresis} A treatment in which some blood cells are removed, mixed with a light-sensitive medicine, exposed to ultraviolet light, and returned to the body. It is used for selected immune, blood, and skin disorders under specialist care.

\medterm{Photophobia} Unusual discomfort or pain caused by light, sometimes leading a person to avoid bright surroundings. It may occur with migraine, corneal injury, eye inflammation, infection, concussion, or meningitis. Severe eye pain, vision loss, fever, stiff neck, or a sudden severe headache needs urgent care.
\medterm{Photoreceptors} Specialized cells in the retina that detect light. Rods support dim-light and peripheral vision, while cones provide color and sharp central vision; damage can reduce sight or cause blind areas.

\medterm{Photorefractive Keratectomy} Laser surgery that removes a very thin surface layer of the cornea and reshapes the layer beneath it. It can reduce short-sightedness, long-sightedness, or astigmatism, but recovery is slower than some procedures and pain, haze, infection, or imperfect vision can occur.

\medterm{Photosensitivity} An unusually strong reaction of the skin to sunlight, causing redness, burning, itching, blistering, or a rash on exposed areas. It may result from a disease, inherited tendency, medicine, or chemical.

\medterm{Photosensitization} Increased sensitivity to ultraviolet or visible light caused by a medication, chemical, or disease, producing an exaggerated skin reaction.

\medterm{Phototherapy} Treatment using carefully controlled light for conditions such as newborn jaundice, psoriasis, eczema, or seasonal depression. The wavelength, dose, duration, and eye or skin protection depend on the condition treated.
\medterm{Phototoxic Dermatitis} A sunburn-like skin injury caused when a medicine or chemical makes the skin react strongly to ultraviolet light. It affects exposed areas and may cause redness, pain, swelling, blisters, or later darkening.

\medterm{Phototoxicity} Skin injury occurring when a chemical makes skin unusually sensitive to ultraviolet light. It resembles severe sunburn and may cause redness, pain, blistering, or later darkening after medicine or chemical exposure.

\medterm{Phrenic} Relating to the diaphragm or the phrenic nerves, which arise from the neck and carry the main motor supply to the diaphragm for breathing.
\medterm{Phrenic Nerve} The nerve arising mainly from C3-C5 that supplies the diaphragm and carries sensory fibers from nearby structures.
\medterm{PHT} Alternate index label; see \textit{Portal hypertension}.

\medterm{Phthiriasis} Infestation with lice, especially pubic lice, causing itching and sometimes visible lice or eggs attached to hair. Treatment usually includes an approved lice-killing product, washing or sealing recently used clothing and bedding, and evaluation or treatment of close or sexual contacts.
\medterm{Phyllodes Tumor} An uncommon breast growth made of supporting tissue and glandular tissue. It may be benign, borderline, or cancerous and often appears as a rapidly enlarging lump; a tissue examination determines its type and guides surgery.

\medterm{Physiatrist} A doctor who specializes in rehabilitation and in improving movement, function, independence, and quality of life after illness, injury, or disability.

\medterm{Physical Activity} Any movement produced by skeletal muscles that increases energy expenditure, including occupational, household, transportation, and recreational movement. Regular activity supports cardiovascular, metabolic, musculoskeletal, and mental health.

\medterm{Physical Child Abuse} Physical child abuse is intentional physical injury or harmful force inflicted on a child, requiring immediate safety assessment, medical care, and safeguarding action.

\medterm{Physical Dependence} A state in which the body adapts to a substance or medication and develops withdrawal symptoms when it is reduced or stopped.

\medterm{Physical Exam Frequency} How often a person should receive a physical examination. The interval depends on age, symptoms, chronic conditions, medicines, pregnancy, and preventive-care recommendations rather than one schedule for everyone.

\medterm{Physical Examination} A systematic assessment of the body using inspection, palpation, percussion, auscultation, and other examination methods.
\medterm{Physical Half-Life} The time needed for half the radioactive atoms in a substance to change through radioactive decay. It describes the substance itself and does not include how quickly the body removes it.

\medterm{Physical Medicine} The medical specialty focused on function, rehabilitation, pain, and restoration of independence after illness or injury.
\medterm{Physical Therapy} Rehabilitation using exercise, movement training, manual techniques, and physical modalities to improve function, mobility, and reduce disability.

\medterm{Physical Therapy Assessment} A detailed evaluation of movement, strength, flexibility, sensation, balance, walking, pain, and daily activities. The findings guide treatment goals, exercise choices, equipment needs, safety advice, and progress monitoring.

\medterm{Physician} A doctor trained and licensed to diagnose illness, provide medical treatment, prevent disease, and coordinate care within the scope of the relevant specialty and jurisdiction.
\medterm{Physician Assistant} A nationally certified and state-licensed medical clinician who practices medicine collaboratively with or under the supervision of physicians. PAs obtain medical histories, perform physical examinations, order and interpret diagnostic studies, formulate medical diagnoses, prescribe pharmacotherapy, assist in surgery, and perform clinical procedures across primary care and surgical specialties.


\medterm{Physicians' Desk Reference} A reference source containing prescribing information, indications, contraindications, dosing, warnings, and adverse effects for medicines and selected products.

\medterm{Physiologic} Physiologic means relating to normal body function or a process occurring within the expected biological range.

\medterm{Physiologic Jaundice Of Newborn} A common temporary yellowing of a newborn’s skin and eyes that begins after the first day of life. It results from normal newborn red-cell breakdown and immature liver processing of bilirubin; jaundice that starts earlier or becomes severe needs urgent assessment.

\medterm{Physiologic Jaundice Of The Newborn} Physiologic newborn jaundice is a temporary bilirubin rise after birth caused by immature bilirubin processing, but early or severe jaundice needs urgent evaluation.

\medterm{Physiological Dead Space} The part of a breath that does not exchange oxygen and carbon dioxide with blood. It includes the air passages and lung areas that receive little or no blood flow.

\medterm{Physiology} The study of how healthy body parts, organs, and living organisms normally work together. It examines processes such as breathing, circulation, digestion, movement, hormone action, nerve signaling, and maintaining internal balance.
\medterm{Physiotherapy} Assessment and treatment using movement, exercise, manual techniques, education, and physical modalities to improve function.
\medterm{Physostigmine} A reversible acetylcholinesterase inhibitor that may reverse severe central anticholinergic delirium in carefully selected patients after assessment for contraindications and cardiac conduction abnormalities.


\medterm{Phytomenadione} A form of vitamin K needed by the liver to make several blood-clotting proteins. It is used to prevent or treat bleeding caused by vitamin K deficiency or by medicines that block vitamin K.

\medterm{PI-RADS} A standard scoring system for prostate MRI findings. It estimates how likely an area is to contain an important prostate cancer and helps guide decisions about follow-up or biopsy; it does not diagnose cancer by itself.
\medterm{Pia Mater} The thin, delicate inner covering that closely follows the surface of the brain and spinal cord. Together with the other coverings and fluid around the nervous system, it helps protect these organs.
\medterm{Pica} Repeated eating of things that have no useful nutritional value, such as soil, paper, ice, or hair, when this is not expected for a person's age or culture. It may cause poisoning, infection, blockage, or poor nutrition and can be linked with iron deficiency or mental-health conditions.
\medterm{Pickwickian Syndrome} An older name for obesity hypoventilation syndrome, involving obesity, daytime hypercapnia, and sleep-disordered breathing.
\medterm{PID} Alternate index label; see \textit{Pelvic inflammatory disease}.

\medterm{Piebaldism} A rare condition present from birth in which some areas of skin and hair lack pigment. It usually causes stable, sharply outlined white patches and often a white section of hair at the front; it is commonly inherited through a change in the KIT gene.

\medterm{Pierre Robin Sequence} A condition present at birth in which a small lower jaw allows the tongue to fall backward and may block breathing. Many affected infants also have a U-shaped opening in the roof of the mouth and feeding difficulty.

\medterm{Pierson Syndrome} A rare inherited condition causing severe kidney protein loss from birth and unusually small pupils. It can lead to swelling, kidney failure, poor vision, and developmental problems.

\medterm{Pigment} A substance that gives skin, hair, eyes, tissue, or other material its color. Pigments may be made by the body, received from outside, or altered by disease, injury, medicines, or aging.

\medterm{Pigment Cirrhosis} Liver cirrhosis associated with abnormal pigment or iron deposition, as may occur in hereditary hemochromatosis.
\medterm{Pigmentary Mosaicism} Uneven skin coloring caused by groups of skin cells carrying different genetic information. It may produce lighter or darker lines, swirls, or patches; most cases affect only the skin, but associated developmental problems should be assessed when present.

\medterm{Pigmentation} Coloration of skin, hair, eyes, or other tissue produced mainly by melanin, carotenoids, blood, or exogenous material; abnormal changes may reflect inflammation, endocrine disease, drugs, or tumors.

\medterm{Pigmentation Disorder} A condition causing abnormal increase, decrease, or uneven distribution of pigment in the skin, mucosa, hair, or eye.

\medterm{Pilar Cyst} A usually harmless, firm lump beneath the skin that forms from a hair follicle, most often on the scalp. It may be inherited, become inflamed or painful, and can usually be removed if it causes symptoms or repeated problems.

\medterm{Piles} Enlarged or swollen blood vessels and supporting tissue in the anus or lower rectum. Internal piles may cause painless bright-red bleeding or tissue protruding during bowel movements; external piles may cause pain, itching, or a tender lump.

\medterm{Piles (Hemorrhoids (Piles))} Alternate index label; see \textit{Hemorrhoids}.

\medterm{Piles (Hemorrhoids)} Alternate term; see \textit{Hemorrhoids}.

\medterm{Pill} A solid medicine designed to be taken by mouth. It may contain one or more active ingredients and may be made to release them immediately or slowly; pills should be taken as directed and not crushed unless the instructions allow it.

\synonyms
Tablet

\medterm{Pilocarpine} A medicine that activates muscarinic receptors, causing pupil constriction and increased secretions. It may lower eye pressure in glaucoma or relieve dry mouth, but can cause sweating, blurred vision, headache, or breathing problems.
\medterm{Pilocytic Astrocytoma} A usually slow-growing brain tumor arising from support cells, most often in children. It commonly affects the cerebellum or visual pathway and may cause headache, vomiting, seizures, balance problems, or vision changes.

\medterm{Piloerection} Elevation of hair shafts producing gooseflesh, caused by contraction of arrector pili muscles under sympathetic stimulation from cold, fear, pain, fever, or certain drugs.

\medterm{Pilomatrix Carcinoma} A rare skin cancer that develops from cells involved in hair formation. It usually appears as a firm lump on the head, neck, or arm and can grow into nearby tissue or spread elsewhere.

\medterm{Pilomatrixoma} Pilomatrixoma is a usually benign skin tumor derived from hair-matrix cells, presenting as a firm subcutaneous nodule, often on the head or neck.

\medterm{Pilonidal Cyst} A cyst or sinus in the cleft between the buttocks, usually containing hair and skin debris. It may become infected, painful, swollen, and drain pus.

\synonyms
Cyst, Pilonidal
Pilonidal Sinus (Pilonidal Cyst)

\medterm{Pilonidal Dimple} Alternate index label; see \textit{Sacral dimple}.

\medterm{Pilonidal Disease} A condition in which hair and skin debris collect in a small pit or tunnel in the cleft between the buttocks. It can cause repeated swelling, pain, drainage, or an abscess.

\medterm{Pilonidal Fistula} An abnormal tunnel from a pilonidal sinus to the skin, usually in the natal cleft, that can drain pus or recur after abscess formation.

\medterm{Pilonidal Sinus (Pilonidal Cyst)} Alternate index label for Pilonidal Cyst.

\medterm{Pilonidal Sinus Disease} A chronic sinus or cyst in the natal cleft, usually containing hair and debris. It may
cause recurrent pain, swelling, drainage, or abscess formation.

\medterm{Pilus} A tiny hair-like structure on some bacteria. It may help bacteria attach to human cells, move, or transfer genetic material to other bacteria, which can contribute to infection or antibiotic resistance.

\medterm{Pimozide} An antipsychotic medicine used for selected disorders involving severe tics or abnormal thoughts. It can cause movement problems, sleepiness, and dangerous heart-rhythm changes, so dosing and heart monitoring require careful medical supervision.
\medterm{Pimple} A common lay term for a small inflamed or blocked skin lesion, often an acne papule or pustule.

\medterm{Pimples (Acne)} Alternate index label for Acne.

\medterm{Pin Care} Pin care keeps the skin around external fixation or traction pins clean and monitored for redness, swelling, drainage, loosening, or infection.

\medterm{Pinch Strength} Pinch strength is the force produced when the thumb presses against one or more fingers and can help assess hand, nerve, muscle, and functional impairment.

\medterm{Pinched Nerve} Compression or irritation of a nerve that can cause pain, tingling, numbness, or weakness along the nerve's distribution.

\synonyms
Compressed Nerve
Nerve Compression

\medterm{Pindolol} Pindolol is a beta-blocker used for selected cardiovascular conditions; it can slow the heart rate and worsen asthma in susceptible people.

\medterm{Pine Oil Poisoning} Swallowing or inhaling concentrated pine oil can irritate the mouth and stomach and may cause drowsiness, seizures, or lung injury if aspirated. Do not induce vomiting; obtain immediate poison-control advice.

\medterm{Pineal Cyst} A fluid-filled sac in the small gland deep in the brain that helps regulate sleep-related signals. Most are found incidentally and cause no harm, but a large cyst can rarely cause headache, eye problems, or fluid buildup.

\medterm{Pineal Germinoma} A pineal germinoma is a malignant germ-cell tumor in the pineal region, often causing headache, eye-movement abnormalities, or hydrocephalus and usually responsive to radiotherapy and chemotherapy.

\medterm{Pineal Gland} A small hormone-producing gland deep in the brain that releases melatonin. Melatonin helps coordinate the body's day-night clock and influences sleep timing, seasonal rhythms, and related biological functions.
\medterm{Pineal Gland Tumor} A growth in or near the pineal gland. It may block the normal flow of brain fluid and cause headache, nausea, vomiting, sleepiness, or problems moving the eyes; symptoms and treatment depend on its type and size.

\medterm{Pineal Parenchymal Tumor} A tumor arising from cells of the pineal gland, ranging from slow-growing to aggressive. It may block the flow of brain fluid and cause headache, vomiting, sleep changes, or abnormal eye movements.

\medterm{Pineoblastoma} An aggressive cancer in the pineal region of the brain, occurring mainly in children. It may block brain-fluid flow and cause headache, vomiting, sleepiness, or eye-movement problems and requires specialist treatment.

\medterm{Pineocytoma} Pineocytoma is a usually slow-growing tumor of pineal cells that may obstruct cerebrospinal-fluid flow and cause raised intracranial pressure.

\medterm{Pinguecula} A pinguecula is a benign yellowish raised growth on the conjunctiva near the cornea, often associated with ultraviolet light, wind, and dust; it may irritate the eye.

\medterm{Pinhole Meatus} An unusually tiny external opening, most often the opening of the urethra. It can make urine flow weak or spraying, cause discomfort, or make catheter placement difficult.

\medterm{Pink Eye} A common name for conjunctivitis, inflammation of the eye covering causing redness, irritation, watering, or discharge.

\synonyms
Conjunctivitis (Pink Eye)

\medterm{Pink Puffer} An older, nonformal term for a person with emphysema who has severe breathlessness, a thin build, and pursed-lip breathing. The term is not a separate diagnosis and is less useful than describing the person's COPD and test results.


\medterm{Pinky Broken Finger (Broken Finger)} Alternate index label for Broken Finger.

\medterm{Pinna} The visible external part of the ear that collects sound.
\medterm{Pinocytosis} A process in which a cell takes in small drops of surrounding fluid and dissolved substances. The cell membrane folds inward to form tiny bubbles that carry the material inside for use, transport, or breakdown.
\medterm{Pint} A pint is a unit of volume; its exact amount varies by measurement system, so the local standard must be specified when measuring fluids or medicines.
\medterm{Pinta} A long-lasting skin infection caused by Treponema carateum bacteria. It spreads through close skin contact and produces changing stages of discolored, scaly, or thickened patches, usually without damage to internal organs.
\medterm{Pinworm} A small intestinal worm that commonly causes intense itching around the anus at night. Infection spreads when microscopic eggs are swallowed or transferred from contaminated hands or surfaces; treatment usually includes medicine, hygiene, and sometimes treating household contacts.
\medterm{Pinworm Infection} Pinworm infection is an intestinal infestation by Enterobius vermicularis, causing nighttime anal itching, disturbed sleep, and sometimes abdominal or genital irritation. Treatment of the person and close household contacts, hygiene, and repeat dosing reduce reinfection.

\synonyms
Infection Pinworms (Pinworm Infection)
Worms Pinworms (Pinworm Infection)

\medterm{Pinworm Test} A test for Enterobius vermicularis, usually the early-morning adhesive-tape test applied to the perianal skin before bathing or defecation.

\medterm{Pinworms} Pinworms are intestinal nematodes, Enterobius vermicularis, that cause nighttime anal itching; treatment, hand hygiene, and laundering help prevent household reinfection.

\medterm{Pioglitazone} A medicine for type 2 diabetes that helps muscle and fat cells respond better to insulin, lowering blood sugar. It can cause weight gain and swelling and may worsen heart failure, so it is unsuitable for some people.



\medterm{Piperacillin} Piperacillin is an extended-spectrum penicillin antibiotic often combined with tazobactam to treat serious infections including Pseudomonas; allergy, diarrhoea, cytopenias, electrolyte effects, and renal dosing are important.

\medterm{Piperazine} An older antiparasitic medicine that paralyzes certain intestinal worms, especially roundworms and pinworms. Its use depends on local guidance and the particular infection.
\medterm{Piperonyl Butoxide} A pesticide synergist that inhibits insect detoxification enzymes and increases the activity of certain insecticides; it can irritate skin, eyes, or airways.

\medterm{Piperonyl Butoxide With Pyrethrins Poisoning} Exposure to these insecticide ingredients can irritate the skin, eyes, or airways and may cause coughing, nausea, vomiting, dizziness, or wheezing. Severe breathing symptoms or swallowed product require urgent poison-control or emergency advice.



\medterm{Pipotiazine} A long-acting phenothiazine antipsychotic given by injection in selected settings. It can cause movement disorders, drowsiness, and other adverse effects and requires clinical monitoring.
\medterm{Pipradrol} An older central nervous system stimulant with limited current use.
\medterm{Piracetam} A nootropic compound used in some countries for selected neurologic indications; evidence and approval vary.
\medterm{Pirfenidone} Pirfenidone is an antifibrotic medicine used to slow decline in idiopathic pulmonary fibrosis; nausea, photosensitivity, fatigue, and liver-enzyme elevation may occur and require monitoring.

\medterm{Piriformis Syndrome} Pain or altered sensation in the buttock or down the leg when the piriformis muscle irritates or presses on the sciatic nerve. Other causes, especially a spinal nerve problem, must be considered before treatment.

\synonyms
PiriformSyndrome (Piriformis Syndrome)

\medterm{PiriformSyndrome (Piriformis Syndrome)} Alternate index label for Piriformis Syndrome.

\medterm{Pirimicarb} Pirimicarb is a selective carbamate insecticide targeting aphids; exposure can inhibit acetylcholinesterase and cause cholinergic symptoms.



\medterm{Piromidic Acid} An older quinolone-like antibacterial investigated or used for urinary and gastrointestinal infections; current use depends on local availability, susceptibility, and safety guidance.

\medterm{Piroxicam} A long-acting nonsteroidal anti-inflammatory drug used for pain and inflammation; gastrointestinal bleeding, kidney injury, cardiovascular events, and hypersensitivity are important risks.

\medterm{Piroxicam Overdose} Taking too much piroxicam, an anti-inflammatory pain medicine, can cause stomach pain, vomiting, bleeding, kidney injury, drowsiness, seizures, or breathing problems. Large or intentional overdoses need immediate poison-control or emergency advice.

\medterm{Pisa Syndrome} A sustained involuntary sideways flexion of the trunk, often associated with medicines, Parkinson disease, or other neurologic disorders.

\medterm{Pisiform} A small, pea-shaped wrist bone embedded in a tendon on the little-finger side of the palm. It helps form the passageway for the ulnar nerve and artery and may hurt after injury or with arthritis.

\medterm{Pisiform Bone} A small pea-shaped carpal bone embedded in the tendon of flexor carpi ulnaris on the ulnar side of the wrist; fracture or arthritis can cause localized wrist pain.

\medterm{Pits} Small depressions or holes in a surface. In medicine, pits may occur in nails, skin, teeth, or other tissues and can result from normal anatomy, injury, infection, inflammation, or inherited conditions.
\medterm{Pitting Edema} Swelling in which pressure applied to the skin leaves a temporary indentation, usually because of excess interstitial fluid.

\medterm{Pituicytoma} A rare, usually slow-growing tumor near the pituitary gland. It may cause headache, vision problems, or hormone disturbances by pressing on nearby structures and is usually assessed with brain imaging and tissue examination.

\medterm{Pituitary} The pituitary is an endocrine gland at the base of the brain that releases hormones controlling growth, thyroid and adrenal function, reproduction, lactation, and water balance.

\medterm{Pituitary Adenoma} A usually noncancerous growth arising from the front part of the pituitary gland. It may release too much of a hormone or press on nearby tissue, causing headache, vision changes, or hormone problems.

\medterm{Pituitary Apoplexy} Sudden bleeding or loss of blood supply inside a pituitary growth, causing severe headache, vision changes, vomiting, eye-movement problems, or confusion. It is an emergency because hormone failure can cause life-threatening low blood pressure.

\medterm{Pituitary Carcinoma} Pituitary carcinoma is a rare malignant pituitary tumor defined by spread to distant or separate sites, often producing excess hormones.

\medterm{Pituitary Disease} A disorder of the pituitary gland that alters hormone secretion or causes mass effects, including adenomas, inflammation, hemorrhage, genetic disease, or hypopituitarism.

\medterm{Pituitary Dwarfism} Pituitary dwarfism is severe short stature caused by deficient growth hormone or broader pituitary dysfunction, with evaluation focused on endocrine and genetic causes.

\medterm{Pituitary Gland} An endocrine gland at the base of the brain that regulates growth, thyroid, adrenal, reproductive, lactation, and water-balance hormones.
\medterm{Pituitary Hormone} A hormone secreted by the anterior or posterior pituitary that regulates growth, thyroid and adrenal function, reproduction, lactation, water balance, or other endocrine organs.

\medterm{Pituitary Insufficiency} Alternate index label; see \textit{Hypopituitarism}.

\medterm{Pituitary Neoplasm} A pituitary neoplasm is an abnormal growth in or near the pituitary gland and may affect hormones, vision, or nearby brain structures.

\medterm{Pituitary Neuroendocrine Tumor} Alternate index label; see \textit{Pituitary tumors and adenomas}.

\medterm{Pituitary Tumor} A pituitary tumor is a growth in the pituitary gland that may be noncancerous but can cause hormone excess, hormone deficiency, or visual problems.

\medterm{Pituitary Tumor (Prolactinoma)} Alternate index label for Prolactinoma.

\medterm{Pityriasis} A group of skin conditions marked by fine, flaky scaling. Examples include pityriasis rosea and pityriasis versicolor, which have different causes and may produce itching or lighter, darker, or pink patches.
\medterm{Pityriasis Alba} Pityriasis alba is a common benign skin condition causing faint, dry, lighter patches, especially on children’s faces, often more noticeable after tanning.

\medterm{Pityriasis Lichenoides} An uncommon inflammatory skin disorder producing recurrent small scaly papules, with forms ranging from mild chronic lesions to an acute ulceronecrotic eruption.

\medterm{Pityriasis Rosea} An acute, self-limited eruption often beginning with a herald patch followed by oval
scaly lesions distributed along skin-cleavage lines. It usually resolves within several
weeks; treatment is symptomatic, and atypical, severe, recurrent, or persistent disease
should be reassessed.

\synonyms
Pityriasis Rosea Gibert (Pityriasis Rosea)

\medterm{Pityriasis Rosea Gibert (Pityriasis Rosea)} Alternate index label for Pityriasis Rosea.

\medterm{Pityriasis Rubra Pilaris} A group of rare, usually long-lasting inflammatory skin disorders. Small rough bumps form around hair follicles and may join to make well-defined, reddish-orange, scaly patches. Areas of normal skin may remain between the patches. The palms and soles may become thick, dry, yellow-orange, and sometimes cracked. The cause is often unknown, and the condition can vary in extent and duration.

\medterm{Pityriasis Versicolor} Alternate index label; see \textit{Tinea versicolor}.

\medterm{Pivampicillin} An orally absorbed prodrug of ampicillin used for selected bacterial infections; its use is limited in some settings by adverse effects and availability.

\medterm{Pixantrone} Pixantrone is an anthracenedione antineoplastic drug developed for selected aggressive lymphomas, with potential marrow, infection, and cardiac risks.

\medterm{PJI} Alternate index label; see \textit{Periprosthetic joint infection}.

\medterm{PJS} Alternate index label; see \textit{Peutz-Jeghers syndrome}.

\medterm{PKD} Alternate index label; see \textit{Polycystic kidney disease}.

\medterm{Placebo} An inactive substance, treatment, or procedure designed to resemble an active treatment. In a clinical trial, it is given to a control group so researchers can compare outcomes and estimate the specific effect of the treatment being studied.

\synonyms
Dummy Treatment
Inactive Treatment
\medterm{Placebo Control} A placebo control is an inactive treatment designed to resemble the intervention in a clinical trial, helping separate treatment effects from expectation and natural change.

\medterm{Placebo Effect} A real change in symptoms or wellbeing that can occur after a person receives an inactive treatment, partly because of expectations, attention, and the treatment setting. It does not prove that the inactive treatment changed the disease itself.

\medterm{Placenta} The temporary organ connecting mother and fetus, allowing exchange of oxygen, nutrients, wastes, and hormones.
\medterm{Placenta Abruptio} Placental abruption is premature separation of the placenta from the uterine wall, causing bleeding, abdominal pain, uterine tenderness, fetal compromise, or maternal shock.

\medterm{Placenta Abruption - Definition} Placental abruption is premature separation of the placenta from the uterine wall before birth, causing bleeding, abdominal pain, uterine tenderness, and possible fetal or maternal compromise.

\medterm{Placenta Accreta} Placenta accreta spectrum occurs when placental tissue abnormally attaches to or invades the uterine wall, creating major risk of hemorrhage during attempted delivery or removal.

\medterm{Placenta Accreta Spectrum} Abnormal adherence or invasion of placental villi into the uterine wall, ranging from accreta to increta and percreta; it can cause severe postpartum hemorrhage and requires planned obstetric management.

\medterm{Placenta Disorder} A disorder of placental implantation, structure, perfusion, or separation, including previa, accreta spectrum, abruption, infarction, or placental insufficiency.

\medterm{Placenta Increta} Abnormal placental invasion into the myometrium without reaching the serosa, increasing
the risk of massive hemorrhage at delivery and often requiring cesarean hysterectomy.

\medterm{Placenta Percreta} Placenta percreta is the most invasive form of placenta accreta spectrum, with villi growing through the uterine wall and sometimes into nearby organs such as the bladder.

\medterm{Placenta Praevia} A condition in which the placenta lies over or near the opening of the cervix. It can cause painless bleeding later in pregnancy and may require monitoring, activity advice, or cesarean delivery.
\medterm{Placenta Previa} Implantation of the placenta over or near the internal cervical opening, which can cause painless bleeding during the second half of pregnancy.

\medterm{Placenta, Succenturiate} A succenturiate placenta has one or more accessory placental lobes separate from the main placenta, which can cause retained tissue or fetal-vessel complications after delivery.

\medterm{Placental Abnormalities} Disorders of placental location, attachment, structure, or separation that may affect pregnancy or cause bleeding.
\medterm{Placental Abruption} Premature separation of the placenta from the uterine wall before delivery. It may cause
vaginal bleeding, abdominal pain, uterine tenderness, fetal distress, and maternal or fetal compromise.

\medterm{Placental Accreta Spectrum} Abnormal placental attachment to the myometrium, ranging from accreta (adherent to
decidua basalis) to increta (invading myometrium) to percreta (invading through serosa).
Major risk factor is prior caesarean section with anterior placenta previa.

\medterm{Placental Hormone} A hormone produced by placental tissue, such as hCG, progesterone, estrogens, or human placental lactogen, that supports pregnancy, fetal growth, and maternal metabolic adaptation.

\medterm{Placental Infarction} A region of placental tissue damaged by inadequate maternal blood flow and coagulation, which may be incidental or, when extensive, associated with fetal growth restriction or pregnancy complications.

\medterm{Placental Insufficiency} Placental insufficiency is inadequate placental delivery of oxygen and nutrients, which can cause fetal growth restriction, abnormal Doppler findings, reduced fluid, or fetal distress.

\medterm{Placental Lactogen} Placental lactogen is a hormone made by the placenta that alters maternal metabolism and supports breast development and fetal nutrient availability during pregnancy.

\medterm{Placental Mesenchymal Dysplasia} A rare abnormal development of the placenta in which supporting tissue becomes enlarged and cyst-like. It can resemble a molar pregnancy on scans and may be associated with poor fetal growth, fetal abnormalities, or pregnancy complications.

\medterm{Placental Separation Signs} Clinical signs indicating the placenta has separated from the uterine wall and is ready
for delivery. Classic signs: lengthening of the umbilical cord (cord lengthens as the
placenta descends), a sudden gush of blood, and uterine elevation and firming (the
uterus rises and becomes globular as the placenta moves into the lower segment).

\medterm{Placental Site Nodule} A benign, usually incidental remnant of intermediate trophoblast at a prior implantation site, composed of small nests of chorionic-type or extravillous trophoblast.

\medterm{Placentation} The developmental process by which the placenta forms and attaches to the uterus, establishes maternal--fetal circulation, and develops the exchange and endocrine functions needed for pregnancy.

\medterm{Plagiocephaly} Plagiocephaly is asymmetric flattening or shaping of an infant’s skull, caused by position or premature suture fusion; its evaluation distinguishes benign molding from craniosynostosis.

\synonyms
Flat Head Syndrome; Deformational Plagiocephaly

\medterm{Plague} A serious infection caused by Yersinia pestis, transmitted mainly by infected fleas or respiratory droplets in pneumonic disease.


\medterm{Planar} Flat or arranged in one plane. In medical imaging, the term describes a two-dimensional view or projection, unlike three-dimensional imaging that shows depth and allows reconstruction from multiple directions.
\medterm{Plane Warts (Verruca Plana)} Small, smooth, relatively flat warts caused by HPV, often occurring in groups on the face, hands, or legs.

\synonyms
Verruca Plana





\medterm{Plant Fertilizer Poisoning} Fertilizer exposure may irritate the mouth, skin, eyes, stomach, or lungs. Some products can dangerously alter blood oxygen or body salts; keep the container and contact poison control after swallowing or significant exposure.



\medterm{Plant Poisoning} Illness caused by eating, touching, or inhaling a toxic plant or plant product; effects range from local irritation to organ failure.


\medterm{Plant Thorn Arthritis} Plant thorn arthritis is persistent joint inflammation after a thorn puncture introduces plant material, sometimes requiring removal of foreign material and treatment of infection or inflammation.

\medterm{Plant Thorn Synovitis} Inflammation of a joint after penetration by a plant thorn or splinter, sometimes involving an infection introduced by
the injury. Symptoms may develop gradually as persistent joint pain, swelling, and
limited movement; evaluation may require removal of retained material and targeted
treatment.

\synonyms
Arthritis Plant Thorn (Plant Thorn Synovitis)


\medterm{Plantar} Relating to the sole or underside of the foot. Plantar problems include heel pain from plantar-fascia injury, warts, pressure ulcers, and abnormal reflexes; diabetes and nerve damage can increase the risk of foot injury.
\medterm{Plantar Flexion} Downward movement of the foot at the ankle, pointing the toes away from the shin. It is produced mainly by the calf muscles and is used in standing on tiptoe, walking, running, and jumping.
\medterm{Plantar Fasciitis} Pain from irritation and small injuries in the strong band of tissue supporting the sole of the foot. Heel pain is often worst with the first steps after waking or resting and may worsen after prolonged standing.

\medterm{Plantar Nerve} A nerve supplying sensation to the sole of the foot and motor branches to intrinsic foot muscles; compression can cause burning pain, numbness, or altered toe movement.

\medterm{Plantar Neuroma} Alternate index label; see \textit{Morton neuroma}.

\medterm{Plantar Reflex} Movement of the toes after the sole of the foot is gently stroked. In most adults the toes bend downward; upward movement of the big toe may indicate a problem in brain or spinal nerve pathways controlling movement.
\medterm{Plantar Response} Movement of the toes and foot after the sole is stimulated. The pattern helps clinicians assess pathways between the brain, spinal cord, and leg; an unusual upward movement of the big toe in an adult may indicate neurologic disease.

\medterm{Plantar Wart (Verruca Plantaris)} A wart on the sole of the foot caused by human papillomavirus. Pressure from standing and walking can push it inward and make it painful; small black dots may represent tiny clotted blood vessels.

\synonyms
Verruca Plantaris

\medterm{Plantaris Muscle} A small posterior-leg muscle with a long tendon that assists weakly with knee flexion and plantar flexion; it can be strained during sudden athletic movement.

\medterm{Plaque} A localized buildup or raised area of material. In the skin, a plaque is a broad, raised lesion often formed by joined papules; in an artery, plaque is a fatty deposit in the vessel wall that can narrow the artery.
\medterm{Plaque And Tartar On Teeth} Dental plaque is a bacterial biofilm that forms on teeth, while tartar is hardened mineralized plaque; both contribute to cavities, gingivitis, and periodontitis.

\medterm{Plaque Assay} A laboratory test that estimates the number of infectious virus particles in a sample. The sample is placed on living cells, and each infectious particle may produce a visible area of damaged cells that can be counted.

\medterm{Plaque Psoriasis} The most common form of psoriasis, producing well-demarcated raised red plaques with silvery scale, commonly on elbows, knees, scalp, and lower back and sometimes associated with arthritis.

\medterm{Plasma} The liquid part of blood carrying cells, proteins, nutrients, hormones, wastes, and other dissolved substances through the body.

\medterm{Plasma Amino Acids} A blood test that measures amino acids in plasma. It can help evaluate inherited metabolic disorders, poor nutrition, liver disease, or other conditions affecting amino-acid processing.

\medterm{Plasma Cell} An immune cell that develops from a B cell and produces antibodies. Antibodies help recognize and control germs and other foreign substances.

\medterm{Plasma Cell Mastitis} A chronic benign inflammatory breast lesion with prominent plasma-cell infiltration, often involving ducts beneath the nipple and potentially causing a mass, nipple retraction, or discharge. It can clinically mimic carcinoma.

\medterm{Plasma Cell Myeloma} A cancer of antibody-producing plasma cells, usually in the bone marrow. It can produce an abnormal antibody and damage bones, kidneys, blood-cell production, or calcium balance, causing bone pain, fractures, anemia, infections, or kidney problems.

\medterm{Plasma Cell Neoplasm} A clonal disorder of antibody-producing plasma cells, ranging from localized plasmacytoma to multiple myeloma with marrow, bone, kidney, calcium, or blood abnormalities.

\medterm{Plasma Cells} Immune cells that make antibodies. They develop from B cells after the body recognizes an infection or other foreign substance. Abnormal plasma-cell growth can cause diseases such as multiple myeloma and may produce excessive abnormal antibodies.

\medterm{Plasma Concentration} The amount of a substance present in a defined volume of plasma at a given time. It is
used to interpret exposure, toxicity, therapeutic drug monitoring, and pharmacokinetic
relationships.

\medterm{Plasma Expanders} Fluids given through a vein to increase circulating blood volume. Crystalloid solutions are commonly used, while selected colloids have limited roles; treatment depends on the cause of fluid loss and patient condition.
\medterm{Plasma Kallikrein} A blood enzyme that helps produce bradykinin, a substance that widens blood vessels and increases leakage from them. It contributes to inflammation and swelling and is involved in some inherited swelling disorders.

\medterm{Plasma Membrane} The selectively permeable lipid bilayer surrounding a cell, containing proteins that control transport, signaling, adhesion, recognition, and communication with the extracellular environment.

\medterm{Plasma Osmolarity} The concentration of dissolved particles in the liquid part of blood. It mainly reflects water balance and substances such as sodium, glucose, and urea; abnormal values can indicate dehydration, excess water, or metabolic disease.

\medterm{Plasma Substitutes} Intravenous fluids intended to expand plasma volume, including crystalloids and colloid solutions; choice depends on the clinical indication and risks.

\medterm{Plasmablastic Lymphoma} Plasmablastic lymphoma is an aggressive B-cell lymphoma often associated with immune deficiency, HIV infection, or immunosuppression.

\medterm{Plasmacytoma} A solitary mass of neoplastic plasma cells, which may occur in bone (solitary plasmacytoma of bone) or soft tissue (extramedullary plasmacytoma). It can progress to multiple myeloma.
\medterm{Plasmalogen} A special type of fat found in cell membranes, especially in the brain, nerves, and heart. The body makes it partly in structures called peroxisomes; very low levels occur in some inherited metabolic disorders.

\medterm{Plasmapheresis} A specialized medical procedure that filters and separates the pale yellow liquid portion of blood (the plasma) from blood cells. In therapeutic plasma exchange, whole blood is drawn through an intravenous catheter into a machine that spins or filters it, separating out the patient's plasma---which often carries disease-causing autoantibodies, abnormal proteins, or inflammatory toxins---and discarding it. The patient's red blood cells, white blood cells, and platelets are then mixed with replacement fluids (such as donor plasma or sterile albumin solution) and safely returned to the bloodstream. Plasmapheresis is an established, life-saving treatment for severe autoimmune and neurological disorders where the immune system attacks the body's own tissues, including Guillain--Barr\'e syndrome, myasthenia gravis crisis, thrombotic thrombocytopenic purpura (TTP), chronic inflammatory demyelinating polyneuropathy (CIDP), and Goodpasture syndrome. It is also used in blood donation centers to harvest plasma from healthy donors.

\synonyms
Therapeutic Plasma Exchange; TPE; Plasma Exchange; Apheresis
\medterm{Plasmid} A plasmid is a small usually circular DNA molecule separate from a bacterial chromosome that can replicate independently and sometimes carries antibiotic-resistance genes.

\medterm{Plasmids} Small DNA rings found separately from the main chromosome in many bacteria. They can copy themselves and may carry genes that help bacteria survive, including genes for antibiotic resistance or disease-causing traits.

\medterm{Plasmin} An enzyme that breaks down fibrin, the protein mesh holding a blood clot together. It is formed from plasminogen and helps remove clots after a damaged blood vessel has been repaired.

\medterm{Plasminogen} An inactive blood protein that can be changed into plasmin, an enzyme that breaks down the protein framework of blood clots. The body activates it when a clot needs to be removed.

\medterm{Plasminogen Activator} An enzyme or medicine that changes plasminogen into plasmin, which breaks down the protein mesh of blood clots. These medicines can dissolve dangerous clots but may cause serious bleeding.

\medterm{Plasminogen Inactivators} Substances or medicines that reduce the conversion of plasminogen into plasmin, thereby slowing breakdown of blood clots. They may be used or studied when limiting bleeding is important.


\medterm{Plasmodium} A group of parasites spread through bites from infected Anopheles mosquitoes and responsible for malaria. The parasites first infect the liver, then red blood cells, causing fever and potentially severe organ complications.

\medterm{Plasmodium Vivax} A parasite spread by infected Anopheles mosquitoes that causes malaria. It infects red blood cells and can remain dormant in the liver, allowing illness to return months or years after the first infection.

\medterm{Plasmodium Vivax Malaria} Malaria caused by the parasite Plasmodium vivax and spread by infected mosquitoes. Fever attacks may recur about every two days, and dormant liver forms can cause later relapses; treatment must clear both blood and liver stages when safe.

\synonyms
Vivax malaria; Benign tertian malaria


\medterm{Plaster Of Paris} A powder that hardens when mixed with water. In healthcare it is used to make casts and splints that hold an injured bone or joint still while it heals; swelling, numbness, or increasing pain under a cast needs urgent review.
\medterm{Plastic Bronchitis} A rare condition producing large, branching mucous casts that form in the airways,
causing life-threatening airway obstruction, often associated with congenital heart
disease or post-fontan circulation.

\medterm{Plastic Casting Resin Poisoning} Exposure to uncured casting resin or its solvents can irritate or burn the skin, eyes, lungs, and digestive tract. Swallowing or breathing fumes may cause vomiting, dizziness, or lung injury and requires poison-control advice.

\medterm{Plastic Resin Hardener Poisoning} Resin hardeners may contain corrosive or irritating chemicals that burn the skin, eyes, mouth, throat, or stomach. Fumes can injure the lungs; rinse exposed areas and seek urgent poison-control advice after significant exposure.

\medterm{Plastic Surgery} Surgery that repairs, rebuilds, or changes body tissues to improve physical function, appearance, or both.
\medterm{Platelet} A small piece of a blood cell that helps stop bleeding. Platelets have no nucleus, gather at damaged blood vessels, and stick together to form an early blood clot. They are made in the bone marrow from larger cells called megakaryocytes.
\medterm{Platelet Aggregation} The clumping of platelets at an injured blood vessel to form an early plug that helps stop bleeding.

\medterm{Platelet Aggregation Test} A laboratory test measuring platelet clumping in response to agonists such as ADP, collagen, or epinephrine, used to evaluate platelet function disorders and medication effects.

\medterm{Platelet Antibodies Blood Test} A blood test detecting antibodies directed against platelet antigens, used selectively in suspected immune thrombocytopenia, neonatal alloimmune thrombocytopenia, or platelet-transfusion refractoriness.

\medterm{Platelet Count} The number of platelets in a given volume of blood, used to assess bleeding or clotting risk.
\medterm{Platelet Count Low (Thrombocytopenia (Low Platelet Count))} Alternate index label for Thrombocytopenia (Low Platelet Count).

\medterm{Platelet Disorder - Immune Thrombocytopenia} Alternate index label for ITP.

\medterm{Platelet Function Disorders} Conditions in which platelets do not adhere, activate, or aggregate normally, causing impaired primary
hemostasis. They may be inherited or acquired and can cause easy bruising, nosebleeds,
heavy menstrual bleeding, or bleeding after procedures.

\medterm{Platelet Plug Formation} The first step in stopping bleeding after a blood vessel is injured. Platelets attach to the damaged area, become active, and clump together to form a temporary plug.


\synonyms
Thrombocytes

\medterm{Platelets (Thrombocytes)} Small cell fragments made in the bone marrow that circulate in the blood and help stop bleeding. They attach to damaged blood vessels, become activated, and join together to form an initial plug.

\synonyms
Thrombocytes

\medterm{Platybasia} An unusually flattened base of the skull. It may cause no symptoms, but in some people it narrows the space around the lower brain or spinal cord and contributes to headache, imbalance, weakness, or other nerve problems.
\medterm{Platypnea} Dyspnea that occurs in the upright position and is relieved by lying down. It is the
opposite of orthopnea. It is rare and may be associated with hepatopulmonary syndrome,
intracardiac shunting, or pulmonary arteriovenous malformations.

\medterm{Platypnea-Orthodeoxia Syndrome} A rare condition in which breathing becomes difficult and blood oxygen falls when a person sits or stands, then improves when lying down. It is usually caused by abnormal blood flow through the heart or lungs. Testing may include oxygen measurements, heart imaging, and lung scans; treatment addresses the underlying cause.

\synonyms
Orthodeoxia-Platypnea Syndrome; Postural Dyspnea and Deoxygenation

\medterm{Platysma} A thin superficial muscle of the neck that tenses the skin and assists lower-jaw and lower-lip movement.
\medterm{Pleiotropic Cytokine} An immune-signaling protein that affects several types of cells or produces more than one effect, such as changing inflammation, blood-cell production, immune activity, or tissue repair.

\medterm{Pleiotropy} The ability of one gene change to influence several traits or body systems. For example, one inherited disorder may affect the brain, skin, blood, and other organs at the same time.

\medterm{Pleomorphic} Pleomorphic means having multiple forms or marked variation in cell or tissue appearance; in pathology, it may be a feature of malignancy but is not diagnostic alone.

\medterm{Pleomorphic Adenoma} A usually noncancerous tumor of a salivary gland, most often the parotid gland near the ear. It usually grows slowly as a painless lump, but long-standing tumors can rarely become cancerous.

\medterm{Pleomorphic Adenoma Of The Vulva} An exceptionally rare benign mixed epithelial and mesenchymal tumor of the vulva, resembling pleomorphic adenoma of salivary tissue and requiring histopathologic confirmation.

\medterm{Pleomorphic Dermal Sarcoma} An uncommon high-grade cutaneous sarcoma, usually on chronically sun-damaged skin of older adults, composed of markedly pleomorphic cells and capable of local recurrence or metastasis.

\medterm{Pleomorphic Xanthoastrocytoma} A circumscribed astrocytic tumor, often arising superficially in the cerebral hemispheres of younger patients and associated with seizures. Most are low grade, but anaplastic transformation can occur.

\medterm{Pleomorphism} Variation in the size and shape of cells or their nuclei within a tissue. Marked
pleomorphism is a feature of many high-grade malignancies but can also occur in some
reactive or degenerative conditions.

\medterm{Plethoric} Plethoric means unusually full, congested, or ruddy, often describing the appearance associated with increased blood volume or red-cell mass.

\medterm{Plethysmography} A test that measures changes in the size of a body part or the amount of blood or air it contains. It can assess circulation in the limbs, lung volume, or changes in blood flow during breathing or exercise.

\medterm{Pleura} A thin, two-layered lining surrounding each lung and covering the inside of the chest wall. A small amount of fluid between its layers allows the lungs to move smoothly during breathing.
\medterm{Pleural} Relating to the pleura, the membranes surrounding the lungs and lining the inside of the chest.

\medterm{Pleural Cavity} The narrow space between the two slippery layers surrounding each lung. A small amount of fluid normally allows smooth breathing, while extra air, blood, pus, or fluid can restrict lung expansion.

\medterm{Pleural Decortication} An operation that removes a thick scar-like layer from the lung’s covering so a trapped lung can expand again. It may be used for long-standing infected fluid, blood, or inflammation around the lung when other treatment is insufficient. Risks include bleeding, infection, air leakage, and breathing complications.

\synonyms
Decortication of the Lung; Pulmonary Decortication; Thoracoscopic Decortication

\medterm{Pleural Disorders} Conditions affecting the lung lining, including fluid buildup, trapped air, inflammation, infection, bleeding, and tumors. They may cause chest pain, cough, breathlessness, fever, or reduced lung expansion.

\medterm{Pleural Effusion} An abnormal buildup of fluid in the space between a lung and the chest wall. It may result from
heart failure, infection, inflammation, cancer, or injury, and can cause shortness of breath, cough, or chest pain.
\medterm{Pleural Effusion (Fluid In The Chest Or On Lung)} Alternate index label for Fluid In the Chest or On Lung.

\medterm{Pleural Fluid Analysis} Laboratory testing of fluid collected from around the lung. Cell counts, protein, acidity, glucose, cultures, and other tests help determine whether the fluid is caused by infection, cancer, heart failure, inflammation, or another problem.

\medterm{Pleural Fluid Culture} Culture of pleural fluid to detect bacteria, fungi, or mycobacteria in suspected pleural-space infection; specimens should be collected appropriately, often in blood-culture bottles.

\medterm{Pleural Fluid Glucose} Pleural-fluid glucose may be reduced in empyema, complicated parapneumonic effusion, rheumatoid
pleuritis, tuberculosis, and malignancy. It is interpreted with pH, cell count, protein, LDH,
microbiology, imaging, and the clinical setting.

\medterm{Pleural Fluid Gram Stain} A pleural-fluid Gram stain looks for bacteria in fluid around the lung and may support evaluation of empyema or bacterial pleural infection.

\medterm{Pleural Fluid PH} Pleural-fluid pH helps assess infection and other causes of effusion. A markedly low pH supports a
complicated parapneumonic effusion in the appropriate clinical setting, but drainage decisions
also depend on imaging, fluid appearance, microbiology, glucose, and the patient's condition.

\medterm{Pleural Fluid Smear} Microscopic examination of fluid removed from the space around the lung. It may show inflammatory cells, bacteria, fungi, cancer cells, or other material and is interpreted with fluid chemistry, cultures, imaging, and symptoms.

\medterm{Pleural Friction Rub} A scratchy or grating sound heard through a stethoscope when inflamed layers around the lung rub together during breathing. It may occur with infection, a blood clot in the lung, or other causes of pleural inflammation.

\medterm{Pleural Needle Biopsy} A pleural needle biopsy removes a small sample of the pleural lining through a needle for evaluation of cancer, infection, or inflammation.

\medterm{Pleural Neoplasm} A pleural neoplasm is an abnormal growth in the lining around the lung, such as mesothelioma or metastatic cancer.

\medterm{Pleural Mesothelioma} A cancer arising from the pleura, the thin lining around the lungs and inner chest wall. It is strongly associated with asbestos exposure and may cause chest pain, breathlessness, or fluid around the lung.

\medterm{Pleural Perforation} A breach in the pleural lining allowing air or fluid to enter the pleural space,
potentially causing pneumothorax or hydropneumothorax, often from trauma, procedures, or
underlying lung disease.

\medterm{Pleural Plaque} A usually harmless area of thickened or calcified tissue on the lining around the lung, most often caused by past asbestos exposure. It is not the same as asbestos-related lung scarring.

\medterm{Pleural Sac} The potential space and enclosing pleural membranes around a lung, normally containing a small amount of lubricating fluid.

\medterm{Pleural Space} The narrow potential space between the lung’s inner covering and the chest wall’s lining. It normally contains a small amount of lubricating fluid; air, blood, or excess fluid can impair breathing.

\medterm{Pleural Tuberculosis} Tuberculous infection or inflammation of the pleura, often producing a lymphocyte-rich pleural effusion, chest pain, fever, and breathlessness and requiring microbiologic or tissue evaluation.

\medterm{Pleurectomy} Surgical removal of part or all of the pleura, used in the treatment of malignant
pleural mesothelioma and recurrent malignant pleural effusion.

\medterm{Pleurisy} Inflammation of the thin layers around the lungs, often causing sharp chest pain that worsens with deep breathing or coughing. Causes include viral infection, pneumonia, a blood clot, autoimmune disease, or injury; associated features vary with the underlying cause.

\synonyms
Pleuritis (Pleurisy)


\medterm{Pleuritic Pain} Sharp chest pain that worsens with breathing or coughing, often caused by inflammation of the pleura or nearby tissues.

\medterm{Pleuritis} Alternate index label for Pleurisy.

\medterm{Pleuritis (Pleurisy)} Alternate index label for Pleurisy.

\medterm{Pleurodesis} A procedure that makes the two layers around a lung stick together so air or fluid cannot collect repeatedly. It may use a chemical or surgery and can cause pain, fever, infection, or breathing complications.

\medterm{Pleurodynia} Severe, repeated pain in the chest or upper abdomen caused by inflammation of breathing muscles, often during infection with coxsackievirus. Pain may worsen with breathing or movement and usually resolves as infection improves.
\medterm{Pleuropulmonary Blastoma} A rare aggressive childhood cancer arising in the lung or pleura, usually associated with DICER1-related tumor predisposition.

\medterm{Plexiform Fibrohistiocytic Tumor} A rare low-grade malignant fibroblastic and histiocytic tumor, usually arising in the dermis or subcutis of the upper extremities of young people; local recurrence is common and metastasis is uncommon.

\medterm{Plexiform Neurofibroma} A plexiform neurofibroma is a diffuse nerve-sheath tumor, strongly associated with neurofibromatosis type 1, that can grow along multiple nerve branches and sometimes become malignant.

\synonyms
PN
PNF

\medterm{Plexus} A network formed by intersecting nerves, blood vessels, lymphatic vessels, or other structures. Examples include the brachial nerve plexus and the celiac vascular plexus.
\medterm{Pliability} Pliability is the ability of a material or tissue to bend, stretch, or deform without breaking; it is relevant to skin, vessels, lungs, and medical devices.

\medterm{Plica} A fold of tissue, especially a fold in the lining of a joint. Plicae are often normal, but an irritated plica around the knee can cause pain, clicking, swelling, or a catching feeling with movement.
\medterm{Plitidepsin} Plitidepsin is an antineoplastic compound that targets eukaryotic elongation factor 1A and has been studied for selected cancers and viral infections.

\medterm{Ploidies} Ploidy describes the number of complete chromosome sets in a cell, such as haploid, diploid, or polyploid; abnormal ploidy can characterize tumors or developmental disorders.


\medterm{PLS} Alternate index label; see \textit{Primary lateral sclerosis}.

\medterm{Plug (Vascular)} A vascular occlusion device used to block a selected vessel under image guidance. A device placed in a blood vessel to block selected blood flow during image-guided treatment.
\synonyms
Vascular occlusion plug; embolization plug

\medterm{Plummer-Vinson Syndrome} A condition involving iron-deficiency anemia, difficulty swallowing, and thin webs in the upper swallowing tube. Treating the iron deficiency and stretching the web can improve swallowing and may reduce complications.
\medterm{Pluralibacter Gergoviae} A gram-negative environmental bacterium that can contaminate medicines or cosmetics and rarely cause opportunistic infection.

\medterm{Pluripotential} Able to develop into many, but not every, type of cell in the body. Pluripotent stem cells can produce cells from the three early tissue groups and are studied for development, disease, and regenerative treatment.
\medterm{Plutonium} A radioactive heavy metal whose inhalation or ingestion can cause internal radiation exposure and toxicity.
\medterm{Pm} An abbreviation that may mean afternoon, post-mortem, or picometer; meaning depends on clinical context.
\medterm{PMBCL} Alternate index label; see \textit{Primary mediastinal large B-cell lymphoma}.

\medterm{PMBL} Alternate index label; see \textit{Primary mediastinal large B-cell lymphoma}.

\medterm{PMDD} Premenstrual dysphoric disorder is a severe cyclic mood disorder occurring before menstruation with irritability, depression, anxiety, and functional impairment that improve after menses.

\medterm{PML} PML commonly means progressive multifocal leukoencephalopathy, a JC-virus demyelinating brain infection during severe immune suppression, but it may have other meanings.

\medterm{PN} Alternate index label; see \textit{Plexiform neurofibroma}.

\medterm{Pneumatic} Relating to air or gas under pressure. Medical examples include an air-pressure ear examination, inflatable compression sleeves used to reduce clot risk, and an inflatable tourniquet used during some operations.
\medterm{Pneumatic Otoscopy} Examination of the eardrum with a lighted instrument while gently changing the air pressure in the ear canal. Reduced eardrum movement may suggest fluid behind it or poor middle-ear ventilation.

\medterm{Pneumatic Retinopexy} A procedure for selected retinal detachments in which an intravitreal gas bubble is injected and laser or cryotherapy seals the retinal break while positioning helps the bubble support the repair.

\medterm{Pneumatic Tourniquet} An inflatable cuff used to temporarily reduce blood flow to a limb during surgery or selected procedures.
\medterm{Pneumatocele} A thin-walled, air-filled space in the lung that often develops after infection, inflammation, or injury. Many resolve naturally, but large or infected spaces can impair breathing or require drainage or surgery.
\medterm{Pneumatosis} Gas trapped within the wall of an organ or other tissue, most often the bowel. It may be harmless, but when caused by poor blood flow, infection, or a tear, it can signal a life-threatening abdominal emergency.
\medterm{Pneumatosis Cystoides Coli} Multiple gas-filled cysts within the wall of the large intestine. They may be harmless and found by chance, or may signal bowel blockage, poor blood flow, infection, or another serious illness when pain or other warning signs occur.

\medterm{Pneumatosis Cystoides Intestinalis} Gas-filled cysts within the bowel wall. They may be harmless and found incidentally, but can also result from poor blood flow, infection, blockage, lung disease, or medicines. Severe pain, fever, or shock requires urgent evaluation.

\medterm{Pneumatosis Intestinalis} Gas within the wall of the intestine, seen as small bubbles or lines on an imaging test. It may be harmless, but pain, fever, shock, or signs of poor blood flow require urgent evaluation for bowel injury.

\medterm{Pneumocephalus} Air or another gas inside the skull, usually after a head injury, brain surgery, or a skull fracture near the sinuses. A large or trapped collection can press on the brain and cause headache, vomiting, confusion, or reduced alertness.

\medterm{Pneumococcal Infection} Infection caused by Streptococcus pneumoniae, ranging from ear or sinus infection to pneumonia, meningitis, or bloodstream infection.

\medterm{Pneumococcal Meningitis} A dangerous infection of the brain and spinal-cord coverings caused by Streptococcus pneumoniae. It can cause fever, headache, stiff neck, confusion, seizures, hearing loss, or death and requires immediate treatment.

\medterm{Pneumococcal Pneumonia} An acute bacterial infection of the lung parenchyma caused by Streptococcus pneumoniae, the single most common cause of community-acquired pneumonia globally. It classically presents with sudden high fever, shaking chills (rigors), pleuritic chest pain, productive cough with rust-colored sputum, and lobar consolidation on chest imaging.

\synonyms
Streptococcus Pneumoniae Pneumonia; Pneumococcal Lung Infection

\medterm{Pneumococcal Sepsis} Sepsis caused by Streptococcus pneumoniae, often arising from pneumonia, meningitis, or another invasive pneumococcal infection.
\medterm{Pneumoconiosis} A group of work-related lung diseases caused by breathing mineral or industrial dust over time. Dust can scar lung tissue and cause cough, breathlessness, reduced oxygen, and progressive breathing disability.
\medterm{Pneumocystis} A genus of fungi that includes Pneumocystis jirovecii, which can cause pneumonia in people with weakened cellular immunity.

\medterm{Pneumocystis Carinii} An older name associated with Pneumocystis; in humans, the usual species causing pneumonia is Pneumocystis jirovecii.

\medterm{Pneumocystis Carinii Pneumonia} The historical name for Pneumocystis jirovecii pneumonia, an opportunistic lung infection in immunocompromised people.
\medterm{Pneumocystis Infection} Infection by Pneumocystis jirovecii, most often causing pneumonia with dry cough, fever, and progressive shortness of breath in immunocompromised people.

\medterm{Pneumocystis Jiroveci} A fungus that can cause a serious pneumonia when the immune system is weakened, especially by advanced HIV, cancer treatment, or transplantation. It often causes gradually worsening dry cough, fever, breathlessness, and low blood oxygen.

\medterm{Pneumocystis Jirovecii Pneumonia} Opportunistic pneumonia caused by Pneumocystis jirovecii, especially in people with
advanced HIV infection or other cellular immunodeficiency. It causes progressive
exertional dyspnea, dry cough, fever, and hypoxemia.

\medterm{Pneumocystis Pneumonia} Opportunistic pneumonia caused by Pneumocystis jirovecii, usually in people with impaired
cell-mediated immunity. It may cause progressive exertional dyspnea, dry cough, fever,
and hypoxemia with diffuse ground-glass changes. Diagnosis and treatment require clinical,
microbiologic, and radiologic assessment; trimethoprim-sulfamethoxazole is commonly used.

\medterm{Pneumocyte} A cell lining the air sacs of the lungs. Type I cells allow oxygen and carbon dioxide to pass between air and blood, while type II cells make surfactant, which helps keep the air sacs open.

\medterm{Pneumolysis} Surgical separation or freeing of lung tissue or pleural adhesions.
\medterm{Pneumomediastinum} Air trapped in the central space of the chest between the lungs. It may follow severe coughing, an injury, a medical procedure, or a tear in the airway or swallowing tube and can cause chest pain or breathlessness.
\medterm{Pneumonectomy} Surgery to remove an entire lung. It is usually performed for selected lung cancers or severe one-sided lung disease when removing only part of the lung is not suitable.
\medterm{Pneumonia} Infection of the lung tissue that fills air sacs with fluid or inflammatory material. It can cause cough, fever, chills, chest pain, breathlessness, low oxygen, or confusion and may be caused by bacteria, viruses, fungi, or aspiration.
\medterm{Pneumonia - Weakened Immune System} Lung infection in a person whose immune defenses are weakened by illness or treatment. Bacteria, viruses, fungi, and unusual organisms can cause it; symptoms may be subtle but disease can worsen quickly and needs prompt assessment.




\medterm{Pneumonia, Bilateral} Pneumonia affecting both lungs. Because more lung tissue may be involved, it can cause greater breathing difficulty and lower oxygen levels than disease limited to one lung, although severity depends on the cause and person’s health.

\medterm{Pneumonitis} Inflammation of lung tissue, often caused by a medicine, radiation, an immune reaction, inhaled irritants, or material entering the lungs rather than by infection. It may cause dry cough and breathlessness.
\medterm{Pneumopericardium} Air around the heart, usually after a chest injury, surgery, a medical procedure, or a tear in the esophagus. A large or trapped collection can prevent the heart from filling and pumping properly and requires urgent assessment.

\medterm{Pneumoperitoneum} Air or gas inside the abdominal cavity around the organs. It may be expected after abdominal surgery, but unexpected gas can indicate a torn intestine and requires urgent assessment, especially with pain or illness.
\medterm{Pneumoplasty} Surgery to repair or reshape part of the lung or an airway. It may be used to remove a damaged area, improve airflow, or restore the shape of tissue after injury or disease.

\medterm{Pneumotachometer} An instrument that measures the speed and amount of air moving during breathing. It can help assess airflow limitation and is used in lung-function testing, often together with other measurements.

\medterm{Pneumothorax} Air trapped between the lung and the chest wall, causing part or all of the lung to collapse. It may happen spontaneously, after injury, or after a procedure and can cause sudden chest pain and breathlessness; severe symptoms are an emergency.
\medterm{Pneumothorax - Infants} Air trapped between an infant’s lung and chest wall, which can compress the lung and cause rapid breathing, low oxygen, poor feeding, or circulatory instability. Severe cases require immediate drainage and breathing support.

\medterm{Pneumovirinae} Pneumovirinae is a subfamily of viruses that includes respiratory pathogens such as respiratory syncytial virus and human metapneumovirus.

\medterm{Pneumovirus} A group or genus of respiratory viruses; human pneumovirus terminology often refers to human metapneumovirus or related paramyxoviruses.

\medterm{PNF} Alternate index label; see \textit{Plexiform neurofibroma}.

\medterm{Pocket Erosion} Loss of tissue at the margin or base of a periodontal or surgical pocket, potentially exposing deeper structures and promoting infection, bleeding, or impaired healing.

\medterm{Podagra} A sudden gout attack in the big-toe joint, causing intense pain, redness, warmth, and swelling.
\medterm{Podiatrist} A healthcare professional specializing in diseases and injuries of the feet, ankles, and related lower limbs. Care may include examination, wound treatment, medicines, footwear advice, procedures, and prevention of complications.
\medterm{Podiatry} The healthcare field focused on the feet, ankles, and related lower-limb conditions. It includes prevention, diagnosis, treatment, wound care, mobility support, and management of problems linked with diabetes or circulation.
\medterm{Podocyte} A specialized cell covering the tiny blood vessels that filter blood in the kidney. Its delicate extensions help keep proteins in the bloodstream; damage can cause protein to leak into urine.

\medterm{Podofilox} Podofilox is a topical antimitotic medicine used to treat selected external anogenital warts caused by human papillomavirus; irritation and correct application are important.

\medterm{Podophyllin} Podophyllin is a plant-derived resin with antimitotic activity formerly used topically for warts; it can be severely toxic if absorbed and is generally replaced by safer products.

\medterm{Podophyllum} A plant genus containing podophyllotoxin, used in some topical wart treatments but toxic if misused.
\medterm{Podosome} A temporary attachment structure on the surface of some cells, including immune cells and bone-removing cells. It helps the cell grip tissue and break down nearby material while moving or carrying out its function.

\medterm{POEMS Syndrome} A rare disorder linked to abnormal plasma cells that affects several body systems. It may cause nerve damage, enlarged organs, hormone problems, skin changes, fluid buildup, eye swelling, and unusual bone changes.

\medterm{Poikilocytosis} The presence of red blood cells with unusually varied shapes on a blood smear. It may occur with anemia, inherited blood disorders, or damage to red blood cells.
\medterm{Poikiloderma} A skin pattern combining mottled pigmentation, visible small vessels, and epidermal thinning, seen in connective-tissue disease, chronic sun damage, inflammation, and some cutaneous lymphomas.

\medterm{Poikilotherm} An organism whose body temperature changes with surroundings, unlike humans, who normally maintain a relatively stable temperature.
\medterm{Poinsettia Plant Exposure} Chewing poinsettia usually causes mild mouth irritation, nausea, or stomach upset, not severe poisoning. Rinse the mouth and contact poison control if symptoms are severe, a large amount was swallowed, or another product was involved.

\medterm{Point} A specific location, spot, or value; in medicine it may refer to a tender point, pressure point, or measurement point.
\medterm{Point Mutation} A change affecting one DNA building block. Depending on its location, it may have no effect, change the protein made by a gene, stop protein production early, or contribute to an inherited condition or cancer.

\medterm{Point Tenderness - Abdomen} Abdominal point tenderness is sharply localized pain when a particular abdominal site is pressed, which may indicate inflammation, infection, obstruction, or another focal process.

\medterm{Point-Of-Care Testing} A test performed near the patient, such as in a clinic, ambulance, home, or hospital ward, instead of a central laboratory. Results are available quickly, but trained users, quality checks, and clinical interpretation remain essential.

\medterm{Point-Of-Care Ultrasound} Ultrasound performed and interpreted at the bedside by the treating clinician to answer a focused question or guide a procedure. It can assess the heart, lungs, abdominal fluid, blood vessels, and bladder, but does not replace a complete imaging study when one is needed.

\medterm{Poison} A substance that can cause injury, illness, or death when enough enters the body. Harm depends on the substance, amount, route, timing, and person's health; suspected exposure needs immediate poison-control advice.
\medterm{Poison Control} A specialist service that gives immediate advice after exposure to a medicine, chemical, plant, or other possible poison. They assess the substance and amount and advise on first aid, observation, emergency care, or further treatment.

\medterm{Poison Control Center - Emergency Number} A poison-control center provides immediate expert advice after a possible poisoning or toxic exposure; the correct emergency number depends on the country and life-threatening symptoms require emergency services.

\medterm{Poison Ivy - Oak - Sumac} Poison ivy, oak, and sumac cause allergic contact dermatitis through urushiol oil, producing intensely itchy red patches and blisters after skin exposure.

\medterm{Poison Ivy Rash} An itchy allergic contact dermatitis caused by exposure to urushiol, an oil found in poison ivy and related plants. The rash commonly appears as red streaks or patches with swelling, small bumps, or blisters; the pattern and severity vary with the amount and route of exposure.

\medterm{Poison Ivy - Oak - Sumac Rash} Alternate index label for Poison Ivy Rash.

\medterm{Poison Ivy, Oak, And Sumac} Alternate index label for Poison Ivy Rash.

\synonyms
Allergy Plant Contact (Poison Ivy, Oak, And Sumac)

\medterm{Poison-Center Consultation} Specialist toxicology advice used to identify hazards, interpret exposure and laboratory
data, select antidotes, guide observation or disposition, and coordinate public-health
response. Consultation is appropriate for significant, uncertain, mixed, pediatric, or
unusual exposures.

\medterm{Poisoning} Harmful effects caused by swallowing, inhaling, injecting, or absorbing a toxic substance. Suspected poisoning requires urgent advice from emergency services or a poison center and should not be treated by inducing vomiting.

\medterm{Polytrauma} Serious injuries affecting two or more body regions or organ systems, usually caused by a major accident, fall, blast, or violence. The injuries may include bleeding, fractures, brain injury, chest injury, or damage to internal organs.
\medterm{Poisoning - Fish And Shellfish} Toxin-related illness after eating contaminated fish or shellfish. Vomiting, diarrhea, flushing, tingling, weakness, or breathing difficulty may occur; sudden neurologic or breathing symptoms require emergency care.

\medterm{Poisoning Arsenic (Arsenic Poisoning)} Alternate index label for Arsenic Poisoning.

\medterm{Poisoning Ciguatera (Ciguatera Poisoning)} Alternate index label for Ciguatera Poisoning.

\medterm{Poisoning First Aid} Move away from fumes, rinse chemicals from skin or eyes, and call poison control with the product container available. Do not induce vomiting or give food unless advised; breathing problems, seizures, collapse, or inability to wake require emergency services.

\medterm{Poisoning Thallium (Thallium)} Alternate index label for Thallium.

\medterm{Pokeweed Poisoning} Eating pokeweed, especially its roots or unripe parts, can cause burning in the mouth, vomiting, diarrhea, weakness, low blood pressure, or seizures. Significant ingestion requires immediate poison-control or emergency advice.

\medterm{Pol Gene} The HIV gene encoding enzymes required for viral replication, including protease, reverse transcriptase, and integrase; antiretroviral drugs target these products.

\medterm{Poland Syndrome} A congenital condition characterized by absence or underdevelopment of part of the pectoral
muscles on one side, sometimes accompanied by abnormalities of the hand on the same side.

\medterm{Polar} Relating to a pole or direction, or describing a molecule with slightly different electrical charges at its ends. In biology and chemistry, polarity affects how substances dissolve, move across cell membranes, and interact with other molecules.
\medterm{Polarity} The arrangement or direction of positive and negative electrical, chemical, or biological properties. Cell polarity helps organize tissue structure and directs processes such as movement, division, absorption, and secretion.
\medterm{Polarography} Polarography is an electrochemical analytical method that measures current while voltage changes at an electrode to identify or quantify reducible substances.

\medterm{Polidocanol} Polidocanol is a detergent sclerosant injected to close selected varicose veins or vascular lesions by damaging the vessel lining and promoting fibrosis.

\medterm{Polio} An infection caused by poliovirus, also called poliomyelitis. It spreads mainly when virus from an infected person’s stool enters the mouth, and it can also spread through respiratory secretions. Most infected people have no symptoms or have a short, mild illness with fever, sore throat, tiredness, nausea, headache, or stomach pain. Fewer than 1\% of infections cause the virus to affect the brain or spinal cord, leading to meningitis, muscle weakness, or acute flaccid paralysis, which is often uneven between the limbs and may become permanent. Paralysis of the muscles used for breathing can be fatal. Vaccination prevents polio.



\medterm{Poliomyelitis} Alternate index label; see \textit{Polio}.
\medterm{Poliosis} A localized patch of white or depigmented hair caused by reduced melanin in follicles, occurring in vitiligo, alopecia areata, inflammatory disease, genetic syndromes, or after injury.

\medterm{Poliovirus} An enterovirus that causes poliomyelitis. It is transmitted mainly by the fecal-oral route; most infections are asymptomatic, but a small proportion cause meningitis or acute flaccid paralysis by damaging motor neurons in the spinal cord. Diagnosis uses molecular testing, usually reverse-transcription PCR of stool or other appropriate specimens.

\medterm{Polish (Dental)} Smoothing restorative surfaces to reduce plaque accumulation. Essential for longevity.

\synonyms
Dental polishing

\medterm{Pollakisuria} Pollakisuria is unusually frequent urination, often in small amounts, caused by irritation, infection, stones, bladder disorders, anxiety, or other urinary conditions.

\medterm{Pollakiuria} Frequent urination, often in small amounts. It may occur with urinary infection, irritation, diabetes, pregnancy, medicines, or anxiety; in some children it is a temporary benign pattern, but pain, fever, blood, thirst, or weight loss needs assessment.
\medterm{Pollen} Tiny plant grains released for reproduction that can trigger allergies, asthma symptoms, or eye inflammation.

\medterm{Pollen-Food Allergy Syndrome} An allergic reaction to certain raw fruits, vegetables, or nuts in people sensitized to related pollen proteins. It usually causes itching or tingling of the mouth and throat and is also called oral allergy syndrome.

\medterm{Pollex} The anatomical term for the thumb, the first and most mobile digit of the hand. It has two bones and can touch the fingertips, allowing precise grasping; injury or arthritis can cause pain and loss of hand function.
\medterm{Pollution} Contamination of air, water, soil, or surroundings by substances that may harm health or ecosystems.
\medterm{Poly A} A sequence of adenine nucleotides added to the end of many messenger RNA molecules, helping regulate stability, transport, and translation.

\medterm{Polyadenylation} Addition of a polyadenosine tail to the 3$^{\prime}$ end of many eukaryotic messenger RNAs, promoting stability, nuclear export, and efficient translation.

\medterm{Polyamine} A small positively charged molecule, such as putrescine, spermidine, or spermine, required for cell growth and nucleic-acid function.

\medterm{Polyamine Catabolism} The enzyme-controlled breakdown of polyamines, natural substances involved in cell growth and function. This process produces chemicals that can affect cell signals, oxidative stress, inflammation, development, and cancer biology.

\medterm{Polyarteritis} Inflammation of several arteries, which can reduce blood flow and damage organs. Causes and effects vary; polyarteritis nodosa is a rare form that mainly affects medium-sized arteries and may cause fever, muscle pain, skin changes, nerve problems, or abdominal pain.
\medterm{Polyarteritis Nodosa} A rare disease in which inflammation damages medium-sized arteries. It can reduce blood flow to the skin, nerves, kidneys, muscles, or intestines, causing fever, weight loss, nerve pain or weakness, skin changes, high blood pressure, or abdominal pain.

\synonyms
PAN (Polyarteritis Nodosa)

\medterm{Polyarthritis} Inflammation involving multiple joints, commonly defined as five or more. It is a clinical pattern rather than a single diagnosis and may result from autoimmune, infectious, crystal-related, or other inflammatory disorders.
\medterm{Polybrominated Biphenyl} A group of brominated industrial chemicals that can persist in the environment and cause toxic effects after exposure.

\medterm{Polychlorinated Biphenyl} A group of persistent industrial chemicals that can accumulate in body fat and are associated with cancer, endocrine, immune, and neurologic effects.

\medterm{Polychondritis} A relapsing immune-related disease that inflames and damages cartilage. It may affect the ears, nose, joints, voice box, windpipe, or inner ear, causing pain, deformity, breathing problems, hearing loss, or balance symptoms.

\medterm{Polychromasia} The presence of unusually bluish, immature red blood cells on a blood smear. It usually means bone marrow is releasing cells quickly, often because of blood loss, red-cell destruction, or treatment response.
\medterm{Polychromatophilia} Polychromatophilia is bluish-gray staining of red blood cells on a blood smear, reflecting increased circulating immature reticulocytes from stimulated red-cell production.

\medterm{Polycystic Kidney Disease} An inherited disorder in which numerous kidney cysts enlarge the kidneys and may cause hypertension and renal failure.

\synonyms
PKD

\medterm{Polycystic Liver Disease} An inherited condition in which many fluid-filled cysts develop in the liver. The cysts may enlarge the liver and cause abdominal fullness or pain; liver function is often preserved, but kidney cysts may also occur.

\medterm{Polycystic Ovarian Syndrome} A hormonal and metabolic disorder involving ovulatory dysfunction and androgen excess, often causing irregular periods, acne, excess hair, infertility, or metabolic risk.

\medterm{Polycystic Ovary} Polycystic ovaries contain many small follicles and may occur as a normal variant or with polycystic ovary syndrome. PCOS additionally involves ovulatory dysfunction, androgen excess, irregular periods, infertility, acne, or metabolic risk and requires clinical—not ultrasound alone—diagnosis.

\synonyms
PCOS (Polycystic Ovary)

\medterm{Polycystic Ovary Syndrome} A common hormonal and metabolic condition that affects ovulation and the levels or effects of certain hormones. It may cause irregular or absent periods, reduced or absent ovulation, acne, extra facial or body hair, scalp hair thinning, or difficulty becoming pregnant. It may also cause higher levels or effects of androgens, hormones found in everyone but usually present at higher levels in males. Some people have many small follicles in the ovaries, but this finding is not required for the condition. Insulin resistance, weight gain, darkened thickened skin, type 2 diabetes, and other metabolic problems may also occur. In adults, the diagnosis is based on at least two of three features---ovulation problems, androgen excess, or polycystic-appearing ovaries---after other causes have been considered.

\medterm{Polycythaemia} An increased red-cell mass or concentration, which may be primary or secondary to another condition.
\medterm{Polycythemia} An abnormally high red-cell concentration or total red-cell amount. It may result from a bone-marrow disorder, low oxygen, dehydration, smoking, medicines, or another illness and can increase the risk of blood clots.

\medterm{Polycythemia (High Red Blood Cell Count)} Alternate index label for High Red Blood Cell Count.

\medterm{Polycythemia - Newborn} Neonatal polycythemia is an abnormally high red-cell concentration in a newborn that can increase blood viscosity and cause poor feeding, lethargy, breathing problems, or low glucose.

\medterm{Polycythemia Vera} A long-term bone-marrow cancer that produces too many red blood cells and sometimes too many white cells and platelets. The thicker blood can cause headaches, itching, clots, bleeding, or an enlarged spleen.

\synonyms
Primary Polycythemia

\medterm{Polydactylism (Polydactyly)} A birth condition in which a person has one or more extra fingers or toes. The extra digit may be on the thumb or big-toe side, the little-finger or little-toe side, or between the usual digits. It may occur alone or as part of a genetic syndrome. Treatment depends on its size and function and may involve reconstructive surgery.

\medterm{Polydactyly} A birth condition in which a person has one or more extra fingers or toes. The extra digit may be small or fully formed and may occur alone or with a genetic syndrome; treatment depends on its structure, function, and effect on nearby digits.
\medterm{Polydipsia} Excessive thirst that leads to drinking unusually large amounts of fluid. Common causes include uncontrolled diabetes, water-balance disorders, medicines, and some mental health conditions.
\medterm{Polyene} A molecule containing multiple carbon-carbon double bonds; polyene medicines such as amphotericin B act against fungal cell membranes.

\medterm{Polyethylene} A plastic polymer used in packaging and medical devices; implanted polyethylene components may wear over time and produce foreign-body reactions.

\medterm{Polyethylene Glycol} A polymer used in medicines, especially laxatives that draw water into stool to ease constipation.
\medterm{Polygenic} Influenced by variants in many genes, often together with environmental factors. Height, blood pressure, and risk for many common diseases are polygenic traits.
\medterm{Polygenic Disease} A disease influenced by small effects from many genes together with environment and lifestyle. Examples include common diabetes, high blood pressure, heart disease, and some cancers; risk is inherited, not certainty.

\medterm{Polygenic Inheritance} Inheritance in which many genes, each with a small effect, together influence a trait or disease risk. Height, skin color, blood pressure, diabetes risk, and heart-disease risk are examples; environment also contributes.

\medterm{Polygenic Obesity} Obesity influenced by many small genetic differences together with eating patterns, activity, sleep, medicines, and the surrounding environment.

\medterm{Polygenic Risk Score} An estimate of a person’s inherited tendency toward a disease, calculated from many small genetic differences. It changes risk assessment but does not predict certainty and must be considered with family history, environment, and other health factors.

\medterm{Polyglandular Autoimmune Syndrome} An autoimmune disorder in which two or more endocrine glands are damaged, causing combinations of thyroid, adrenal, parathyroid, pancreatic, or gonadal dysfunction.

\medterm{Polygraphy} Polygraphy is the simultaneous recording of several physiologic signals, such as breathing, heart rate, oxygen level, and movement, often during sleep testing.

\medterm{Polyhydramnios} An unusually large amount of amniotic fluid during pregnancy. It may cause abdominal discomfort, breathlessness, or swelling and can increase the risks of early birth, abnormal fetal position, cord problems, or heavy bleeding after delivery.

\medterm{Polyketide Synthase} An enzyme system that builds complex natural products from repeated carbon units, producing many antibiotics, toxins, and other biologically active molecules.

\medterm{Polymagia (Polyphagia)} An unusually strong or excessive appetite. It may occur with diabetes, an overactive thyroid, certain medicines, genetic disorders, or some eating and mental health conditions.
\medterm{Polymastia} The presence of extra breast tissue, sometimes with an additional nipple, along the embryonic milk line. Accessory breast tissue can enlarge or develop the same benign or malignant conditions as normal breast tissue.
\medterm{Polymenorrhea} Menstrual periods that occur unusually often, commonly at intervals shorter than about 21 days. Causes include hormonal changes, thyroid disease, fibroids, bleeding disorders, medicines, and the menopausal transition; persistent frequent bleeding needs evaluation.

\medterm{Polymenorrhoea} British spelling of polymenorrhea: menstrual periods occurring unusually often, commonly at intervals shorter than about 21 days.
\medterm{Polymer} A large molecule made of repeating smaller units called monomers; polymers are widely used in biology, medicines, and medical materials.

\medterm{Polymerase Chain Reaction} A laboratory method that makes millions of copies of a selected piece of DNA so it can be detected and studied. It is used to identify infections, inherited conditions, cancer-related gene changes, and other genetic material. The test must be chosen and interpreted for the specific question.
\medterm{Polymerization} Chemical joining of many small monomers into a larger polymer, a process important in biomolecules, plastics, dental materials, drug delivery, and laboratory techniques.

\medterm{Polymicrobial Infections} Infections caused by two or more kinds of germs at the same time. They commonly affect contaminated wounds, the abdomen, or the lungs and may require treatment directed at several different organisms.

\medterm{Polymicrogyria} A brain-development abnormality in which the outer brain surface has too many small, unusually formed folds. It may cause seizures, developmental delay, movement problems, or learning difficulties, depending on the areas involved.

\medterm{Polymorphic} Having more than one form or appearance. The term may describe a rash, crystal, cell, gene variant, heart rhythm, or tumor. Its medical importance depends on the structure being described and the other clinical findings.

\medterm{Polymorphism} The presence of two or more inherited genetic forms or visible traits within a population.
\medterm{Polymorphonuclear Leukocyte} A white blood cell with a multilobed nucleus, especially a neutrophil, eosinophil, or basophil, involved in innate immune defense and inflammation.

\medterm{Polymorphous} Having lesions or clinical features that differ in shape, size, or stage and therefore show multiple forms.

\medterm{Polymorphous Light Eruption} Polymorphous light eruption is an itchy inflammatory rash triggered by sunlight, usually appearing hours to days after exposure and recurring seasonally in susceptible people.

\synonyms
Sun Poisoning

\medterm{Polymyalgia Rheumatica} An inflammatory disorder causing aching and morning stiffness around the shoulders and hips in older adults.

\synonyms
PMR (Polymyalgia Rheumatica)

\medterm{Polymyositis} An uncommon immune-related disease causing gradual, usually symmetrical weakness in muscles close to the body, especially the shoulders, hips, and thighs. It can make climbing stairs, lifting, or swallowing difficult and may raise muscle-enzyme levels.
\synonyms
PM (Polymyositis)

\medterm{Polymyxin} A group of antibiotics that disrupt bacterial cell membranes, used mainly for selected gram-negative infections.
\medterm{Polymyxin B} An antibiotic active against selected gram-negative bacteria; kidney and nerve toxicity limit its use to difficult infections.

\medterm{Polyneuritis} Inflammation or injury affecting multiple peripheral nerves, causing symmetric pain, numbness, weakness, or autonomic dysfunction; causes include immune disease, infection, toxins, and deficiency.

\medterm{Polyneuritis, Acute Idiopathic} Sudden inflammation or dysfunction of several peripheral nerves without an identified cause. It may cause tingling, weakness, pain, or loss of reflexes in the limbs and can sometimes affect breathing or swallowing. Rapidly worsening weakness requires urgent assessment because an immune nerve disorder such as Guillain--Barré syndrome may be responsible.

\medterm{Polyneuropathy} A generalized disorder affecting multiple peripheral nerves, often causing symmetric numbness, tingling, burning pain, weakness, or impaired reflexes in the limbs.

\medterm{Polynucleotide} A chain of linked nucleotides forming DNA, RNA, or a related laboratory molecule. Its sequence stores or carries genetic information and determines how cells make proteins or regulate biological processes.

\medterm{Polyoma Virus} A polyomavirus; some human polyomaviruses are widespread and cause serious disease mainly in immunocompromised people.
\medterm{Polyomavirus Hominis 2} A human polyomavirus designation associated with a virus found in people; clinically important disease is uncommon and mainly associated with immune suppression.

\medterm{Polyp} A growth projecting from a mucous membrane; it may be benign, precancerous, or rarely cancerous.
\medterm{Polyp Biopsy} A polyp biopsy removes part or all of a polyp for microscopic examination to determine whether it is benign, precancerous, or cancerous.

\medterm{Polyp Surveillance Colonoscopy} A follow-up colon examination after polyps have been found or removed. The timing depends on the number, size, type, and microscopic features of the polyps and on the person’s overall risk.

\medterm{Polypeptide} A chain of amino acids joined by peptide bonds. One or more folded polypeptide chains make up many proteins.
\medterm{Polyphagia} Unusually strong or persistent hunger that may lead to eating much more than usual. Causes include poorly controlled diabetes, an overactive thyroid, certain medicines, inherited conditions, and some mental-health disorders.

\medterm{Polypharmacy} Use of several medicines at the same time, often five or more. Some combinations are necessary, but taking many medicines increases the risk of interactions, side effects, confusion, missed doses, and prescribing errors, especially in older adults.



\medterm{Polyposis} A condition in which many polyps develop, especially in the bowel or other digestive organs. Polyposis may be inherited or acquired, and some types increase the risk of colorectal or other cancers.
\medterm{Polypropylene} A synthetic plastic polymer used in medical supplies, sutures, meshes, and containers; implanted materials may rarely cause infection or foreign-body reactions.

\medterm{Polyps Colon (Colon Polyps)} Alternate index label for Colon Polyps.

\medterm{Polyps Rectal (Colon Polyps)} Alternate index label for Colon Polyps.

\medterm{Polyps Uterus (Uterine Fibroids)} Alternate index label for Uterine Fibroids.

\medterm{Polyps, Colon} Alternate index label; see \textit{Colon polyps}.

\medterm{Polyps, Endometrial} Alternate index label; see \textit{Uterine polyps}.

\medterm{Polyps, Nasal} Alternate index label; see \textit{Nasal polyps}.

\medterm{Polyps, Stomach} Alternate index label; see \textit{Stomach polyps}.

\medterm{Polyps, Uterine} Alternate index label; see \textit{Uterine polyps}.

\medterm{Polyradiculitis} Inflammation or other disease affecting several spinal nerve roots. It may cause weakness, numbness, pain, and reduced reflexes; causes include infections, autoimmune disorders, and toxins.
\medterm{Polyradiculopathy} Disease affecting multiple spinal nerve roots, producing pain, weakness, sensory loss, or reflex changes in a distribution determined by the involved roots.

\medterm{Polysaccharide} A carbohydrate made of many linked simple sugar units. Examples include starch, glycogen, cellulose, and some bacterial capsules.
\medterm{Polysome} A group of ribosomes attached to one messenger RNA molecule at the same time. Together they make many copies of the protein instructions, helping cells produce proteins efficiently.

\medterm{Polysomnography} An overnight sleep study that records brain activity, eye movements, muscle activity, heart rhythm, airflow, breathing effort, and blood oxygen. It helps identify sleep-related breathing problems, unusual movements, seizures, and other sleep disorders.

\medterm{Polysomnography (Sleep Study)} An overnight recording of sleep, breathing, body movement, heart rhythm, and blood oxygen. It helps diagnose sleep apnea and other sleep disorders by showing how breathing and body functions change during sleep.

\synonyms
Sleep Study

\medterm{Polysplenia} A congenital condition with multiple small spleens and variable abnormalities of organ arrangement, heart, blood vessels, or intestinal rotation.

\medterm{Polystyrene} A synthetic plastic polymer used in containers and laboratory materials; burning it can release hazardous fumes.

\medterm{Polytene Chromosome} An enlarged chromosome formed by repeated DNA replication without cell division, displaying visible bands and commonly used to study gene activity in insect tissues.

\medterm{Polythelia} The presence of one or more additional nipples, usually along the embryonic milk line. It is often harmless, although accessory tissue can rarely have the same disorders as normal breast tissue.
\medterm{Polythiazide} A thiazide diuretic formerly used to treat high blood pressure and fluid retention by increasing renal salt and water excretion.

\medterm{Polyunsaturated Fat} A type of fat usually liquid at room temperature, found in fish, vegetable oils, nuts, and seeds. Some are essential because the body cannot make them.

\medterm{Polyurethane} A family of synthetic polymers used in medical tubing, implants, dressings, and other products; tissue response depends on formulation and site.

\medterm{Polyuria} Producing unusually large amounts of urine, commonly more than about three liters daily in an adult. Causes include diabetes, excess fluids, medicines, kidney problems, hormone disorders, and some inherited conditions.

\medterm{Pompe Disease} A rare inherited disorder in which a missing enzyme causes stored sugar to build up inside muscle cells. Babies may develop severe weakness, an enlarged heart, and breathing or feeding problems; later-onset disease mainly weakens the limbs and breathing muscles. Enzyme replacement and supportive care can slow or improve the disease.

\synonyms
Glycogen storage disease type II; Acid maltase deficiency; GSD II

\medterm{Polyvalent} Effective against or involving multiple types, strains, antigens, or targets.
\medterm{Polyvinyl} A family of polymers used in plastics and medical materials; properties vary according to attached chemical groups.

\medterm{Pomalidomide} An immunomodulatory medicine used with other treatments for selected multiple myeloma and related blood cancers; it can cause severe birth defects and blood abnormalities.

\medterm{Pompholyx} A form of dyshidrotic eczema causing intensely itchy small blisters on the hands or feet.
\medterm{Pompholyx Eczema} Pompholyx, or dyshidrotic eczema, causes intensely itchy small blisters on the palms, sides of fingers, or soles and may leave painful peeling or cracks.

\medterm{Pons} The middle part of the brainstem, containing pathways and nuclei involved in movement, sensation, breathing, and cranial-nerve function.
\medterm{Pontiac Fever} A mild, self-limited flu-like illness caused by Legionella bacteria without the pneumonia typical of Legionnaires’ disease.

\medterm{Pontine} Relating to the pons, a part of the brainstem between the midbrain and medulla. The pons helps carry movement and feeling signals and contains centers involved in eye movement, facial movement, hearing, balance, and breathing. A stroke or bleed there can cause weakness, loss of sensation, trouble speaking or swallowing, coma, or a locked-in state. Rapid correction of very low sodium can also damage it.
\medterm{Poor Circulation} Reduced or abnormal blood flow, usually referring to arterial insufficiency or venous disease, causing coolness, color change, pain, swelling, ulcers, or delayed wound healing.

\medterm{Poor Color Vision} Alternate index label; see \textit{Color blindness}.

\medterm{Popcorn Lung Bronchiolitis Obliterans} Bronchiolitis obliterans is irreversible scarring and narrowing of small airways that can follow toxic inhalation, infection, transplantation, or autoimmune disease. It causes persistent cough and breathlessness with obstructive lung testing; exposure avoidance and specialist care are important.

\medterm{Popliteal} Relating to the hollow behind the knee. This area contains the popliteal artery, vein, and nerves and can be affected by an artery aneurysm, a blood clot, or a Baker cyst.
\medterm{Popliteal Aneurysm} A weakened, bulging section of the artery behind the knee. It may form blood clots, block blood flow, or send clots into the lower leg, causing pain, coldness, numbness, or tissue damage. Larger aneurysms may need repair.

\medterm{Popliteal Artery} The main artery behind the knee. It continues from the thigh artery and divides into arteries supplying the lower leg and foot; blockage or injury can cause pain, coldness, numbness, or loss of circulation.

\medterm{Popliteal Vein} A deep vein behind the knee formed by the union of the anterior and posterior tibial veins. It continues as the femoral vein and can be affected by deep-vein thrombosis.

\medterm{Popliteal Artery Aneurysm} An abnormal enlargement of the artery behind the knee that may cause limb ischemia, clotting, or embolization.

\medterm{Popliteal Cyst} A fluid-filled swelling behind the knee, also called a Baker cyst, caused by communication between the joint and a bursa and often associated with arthritis, meniscal injury, or inflammation.

\medterm{Popliteal Cyst (Baker Cyst)} Alternate index label for Baker Cyst.

\medterm{Popliteal Fossa} Diamond-shaped space on the posterior knee, bounded by hamstrings superiorly and
gastrocnemius heads inferiorly. Contains popliteal artery/vein, tibial nerve, and common
peroneal nerve. The popliteal pulse is palpated at the inferior margin.

\medterm{Population Attributable Fraction} The estimated proportion of disease in a whole population that could be prevented if a specific harmful exposure were removed, assuming the exposure truly causes the disease and other conditions stayed the same.

\medterm{Population Attributable Risk} The estimated amount of disease in a population linked to a specific exposure. It helps show the potential public-health benefit of reducing that exposure, but it depends on accurate data and a causal relationship.

\medterm{Population Genetics} The study of how gene frequencies and genetic variation change within and between populations through inheritance, mutation, selection, migration, and genetic drift.

\medterm{Population Group} A defined set of people sharing a characteristic such as age, location, exposure, ancestry, occupation, or health condition for epidemiologic or clinical analysis.


\medterm{Porcelain Gallbladder} Hard calcium deposits in the wall of the gallbladder, usually after long-standing inflammation or gallstones. It may cause no symptoms but can be associated with gallbladder disease; imaging findings guide discussion of monitoring or removal.

\medterm{Porcine} Porcine means derived from or relating to pigs; porcine tissues, enzymes, and biologic products may be used in medicine after appropriate processing and screening.

\medterm{Pore} A pore is a small opening or channel in a surface, membrane, tissue, or material that permits passage of fluid, gas, particles, or cellular contents.
\medterm{Porencephaly} A fluid-filled cavity within the brain caused by abnormal development or damage before or after birth. Effects vary with its size and location and may include seizures, weakness, movement problems, or developmental delay.
\medterm{Porfimer Sodium} A photosensitizing medicine used with light therapy to destroy selected abnormal tissues, including some tumors; it causes prolonged light sensitivity.



\medterm{Porin} A porin is a channel-forming protein in an outer membrane, especially of Gram-negative bacteria, allowing passage of selected small molecules.

\medterm{Pork Tapeworm} A parasite acquired by eating undercooked pork; its larvae can invade tissues and cause cysticercosis.
\medterm{Porokeratosis} A group of skin conditions that cause one or more ring-shaped, scaly patches with a raised border. Some forms can rarely develop into skin cancer, so changing, bleeding, or persistent lesions need dermatology assessment.

\medterm{Poroma} A usually noncancerous tumor arising from a sweat gland. It often appears as a flesh-colored, pink, or red lump on the palms, soles, or other skin and may be removed if painful, bleeding, or uncertain in diagnosis.

\medterm{Porphyria} A group of mostly inherited disorders of heme biosynthesis in which reduced activity of a pathway enzyme causes porphyrins or their precursors to accumulate. Acute hepatic porphyrias can cause episodic abdominal, neurologic, or psychiatric symptoms, while cutaneous porphyrias can cause photosensitivity and skin injury; some forms have both patterns. Porphyria cutanea tarda is usually acquired rather than inherited.
\medterm{Porphyria Cutanea Tarda} The most common porphyria, caused by reduced activity of hepatic uroporphyrinogen decarboxylase (UROD), leading to accumulation of porphyrins in the liver and skin. It causes fragile, blistering skin on sun-exposed areas. Associated factors include hepatitis C, excess iron, alcohol use, estrogen exposure, and certain genetic UROD variants.

\medterm{Porphyrias} A group of metabolic disorders caused by reduced activity of enzymes in the heme-synthesis pathway. Acute hepatic forms can cause abdominal and neurologic symptoms; cutaneous forms can cause pain, burning, blistering, or other injury after light exposure. Examples include acute intermittent porphyria, variegate porphyria, hereditary coproporphyria, porphyria cutanea tarda, and erythropoietic protoporphyria.

\medterm{Porphyrin} A ring-shaped compound involved in heme and other biological molecules; abnormal accumulation occurs in porphyrias.
\medterm{Porphyrins Blood Test} A blood porphyrin test measures porphyrins or related compounds and helps evaluate suspected porphyria and disorders of haem production.

\medterm{Porphyrins Urine Test} A urine porphyrin test measures porphyrins or their precursors and helps investigate suspected porphyria or abnormal haem metabolism.

\medterm{Porphyromonas} A genus of anaerobic gram-negative bacteria found especially in the mouth and associated with periodontal and other infections.

\medterm{Porphyromonas Bennonis} An anaerobic oral bacterium rarely isolated from human specimens and occasionally associated with mixed infection.
\medterm{Porphyromonas Catoniae} An anaerobic gram-negative bacterium found in the mouth and occasionally associated with periodontal or mixed infection.

\medterm{Porphyromonas Endodontalis} A mouth bacterium that grows without oxygen and is associated with infected tooth pulp and root canals. It may contribute to dental infection, but diagnosis depends on symptoms, examination, imaging, and laboratory findings.

\medterm{Porphyromonas Levii} A Porphyromonas species found mainly in animals and rarely associated with human infection. An animal-associated bacterium rarely linked to human infection, usually in wounds or other opportunistic settings.
\medterm{Porphyromonas Somerae} A bacterium that grows without oxygen and is occasionally found in human wounds, abscesses, or other samples. When it causes infection, treatment depends on laboratory identification, infection location, and antibiotic susceptibility.

\medterm{Port-Wine Stain} A permanent birthmark caused by widened small blood vessels, appearing as flat pink, red, or purple skin.
\medterm{Portacaval Shunting} A surgical or image-guided connection that diverts blood around the liver to lower dangerously high pressure in the portal veins. It may reduce bleeding risk but can worsen confusion caused by liver disease.

\medterm{Portal Hypertension} Abnormally high pressure in the veins carrying blood from the digestive organs to the liver, most often because of cirrhosis. It can enlarge veins in the esophagus or stomach, cause abdominal fluid buildup, enlarge the spleen, or lead to bleeding.

\synonyms
PHT

\medterm{Portal Hypertensive Gastropathy} Changes in the stomach lining caused by high pressure in the veins around the liver. The lining may look mottled and bleed slowly or suddenly, causing anemia or vomiting of blood in severe cases.

\medterm{Portal System} A network of veins that carries blood from the stomach, intestines, spleen, pancreas, and gallbladder to the liver. The liver processes nutrients and many absorbed substances before the blood returns to the general circulation.

\medterm{Portal Vein} The large vein that carries blood from the digestive organs, spleen, pancreas, and gallbladder to the liver. A clot or increased pressure in this vein can cause enlarged abdominal veins, bleeding, or fluid buildup.

\medterm{Portal Vein Thrombosis} A blood clot that partly or completely blocks the portal vein or one of its branches. It may occur with cirrhosis, cancer, abdominal infection or inflammation, surgery, or a clotting disorder. It can cause portal hypertension, enlarged veins, spleen enlargement, intestinal injury, or gastrointestinal bleeding.

\medterm{Portion Of Cytosol} A region of the cell’s aqueous internal fluid outside membrane-bound organelles, containing ions, metabolites, proteins, and enzymes that support cellular reactions.

\medterm{Portion Size} The amount of food or drink served or eaten at one time, distinct from a standardized serving size and important for estimating energy, carbohydrate, sodium, and nutrient intake.

\medterm{Portopulmonary Syndrome} High blood pressure in the arteries of the lungs occurring with portal hypertension or serious liver disease. It can cause breathlessness, tiredness, chest discomfort, or fainting and affects decisions about liver treatment or transplantation.

\medterm{Position} The location or arrangement of a body part, organ, fetus, or patient. Position can affect examination findings, movement, comfort, circulation, breathing, childbirth, medical procedures, and interpretation of imaging.
\medterm{Positional Asphyxia} Fatal or potentially fatal impairment of breathing caused by body position that
restricts chest movement, airway patency, or venous return. Risk is increased by
restraint, intoxication, obesity, and inability to reposition; findings are nonspecific
and require scene and medical correlation.


\medterm{Positive} Positive means that a test, finding, result, or response is present or above a defined threshold; interpretation depends on the test and clinical setting.

\medterm{Positive Airway Pressure Treatment} Treatment that uses pressurized air delivered through a mask or similar device to keep the airway open during sleep or breathing difficulty. It is commonly used for obstructive sleep apnea; mask discomfort, dry airways, leaks, or skin irritation may occur.

\medterm{Positive End-Expiratory Pressure} Pressure kept in the airways at the end of breathing out, usually by a ventilator. It helps keep the lung’s air sacs open and improve oxygen levels, but excessive pressure can injure the lungs or lower blood pressure.

\medterm{Positive Feedback} A body-control process in which a response strengthens the original change, making the process grow until a specific event is completed. Examples include contractions during labor and blood clotting.
\medterm{Positive Finding} A positive finding is an observed symptom, examination sign, image, laboratory result, or test response that supports the presence of a specified condition or feature.

\medterm{Positive Predictive Value} The chance that a person with a positive test result truly has the condition being tested for. It depends on test accuracy and how common the condition is among the people tested.

\medterm{Positive Predictive Value In Surveillance} The proportion of reported signals or alerts that represent true events meeting the
surveillance case definition. It is influenced by the specificity of the signal and the
underlying prevalence of the condition.

\medterm{Positive Pressure Ventilation} Breathing support that pushes air or oxygen into the lungs under pressure. It may be delivered through a mask or a breathing tube and can support or replace a person’s own breathing during respiratory failure.

\medterm{Positive Psychology} The study of wellbeing, strengths, resilience, relationships, and factors that help people function and thrive. Its practices may support coping and quality of life but do not replace treatment for illness.

\medterm{Positive Symptoms} Symptoms reflecting an excess or distortion of normal mental functions, such as
hallucinations, delusions, disorganized speech, or grossly disorganized behavior. They
are contrasted with negative symptoms, which reflect loss or reduction of function.

\medterm{Positive-Strand RNA Virus} A virus whose RNA can be read directly by an infected cell as instructions for making viral proteins. This allows the virus to begin producing new parts soon after entering the cell.

\medterm{Positron} A particle with the same mass as an electron but a positive charge. Some medical tracers release positrons; when one meets an electron, energy is produced and detected during a PET scan.
\medterm{Positron Camera} A positron camera detects paired gamma rays from positron annihilation and is used in positron-emission tomography to map tracer distribution.

\medterm{Positron Emission Mammography} PET imaging of breast. Higher sensitivity than mammography for dense breasts. Limited availability. A breast imaging test combining positron emission scanning with mammography to show tissue activity and structure.
\medterm{Positron Emission Tomography} An imaging test that uses a small amount of radioactive tracer to show tissue activity, such as glucose use, blood flow, or specific proteins. It is often combined with CT or MRI to locate abnormal areas more precisely.

\medterm{Post Cibum} A medication instruction meaning that a medicine should be taken after meals. Clear instructions should spell this out to prevent confusion with other uses of the abbreviation P.C.

\synonyms
P.C.; After Meals

\medterm{Post-Acute Rehabilitation} Time-limited interdisciplinary care after an acute illness, injury, or hospitalization
to restore mobility, self-care, cognition, communication, and safe community function.
Intensity and setting depend on medical stability and rehabilitation potential.

\medterm{Post-Cardiac Arrest Management} Care after return of spontaneous circulation that supports breathing and circulation, identifies and treats the cause of
the arrest, manages temperature and seizures when indicated, and assesses neurologic
recovery over time.

\medterm{Post-Intensive Care Syndrome} New or worsened physical, thinking, emotional, or mental-health problems that persist after critical illness or an intensive-care stay. It may include muscle weakness, poor memory or concentration, anxiety, depression, sleep problems, or post-traumatic stress.

\medterm{Post-Cardiac Injury Syndrome} Alternate index label; see \textit{Dressler syndrome}.

\medterm{Post-Chemotherapy Cognitive Impairment} Alternate index label; see \textit{Chemo brain}.

\medterm{Post-Concussion Syndrome} Ongoing symptoms after a concussion, such as headache, dizziness, tiredness, poor concentration, sleep problems, or mood changes. Symptoms usually improve gradually; evaluation and a stepwise return to activity help guide recovery.
\medterm{Post-Concussional Syndrome} Persistent physical, cognitive, emotional, or sleep symptoms after a concussion or mild traumatic brain injury.
\medterm{Post-Dural Puncture Headache} A headache that begins after a spinal tap or accidental puncture during an epidural. It is usually worse when sitting or standing and improves when lying down because spinal fluid has leaked; neck stiffness, nausea, ringing ears, or light sensitivity may occur.

\medterm{Post-Exposure Prophylaxis} Medicine or vaccination given after a possible exposure to prevent infection from developing. It is used for selected exposures to HIV, hepatitis B, rabies, and other infections; the correct treatment depends on the pathogen and must begin promptly.

\medterm{Post-Hepatic Jaundice (Obstructive Jaundice)} Jaundice caused by obstruction of bile flow after bilirubin has been processed by the liver, commonly from gallstones, strictures, or tumours.

\synonyms
Obstructive Jaundice

\medterm{Post-Hoc} Performed or analyzed after an event; in statistics, a post-hoc analysis follows the planned primary analysis.
\medterm{Post-Infectious Cough} A chronic cough persisting 3--8 weeks after a URI, caused by transient bronchial
hyperresponsiveness and airway inflammation. It is self-limited and usually resolves
without treatment. Inhaled corticosteroids, antitussives, and time are the mainstays of
management.

\medterm{Post-Kala-Azar Dermal Leishmaniasis} A dermatological sequela of visceral leishmaniasis (kala-azar) caused by Leishmania donovani that develops months to years after successful or incomplete antimonial or amphotericin treatment, predominantly observed in South Asia and East Africa. It manifests as an extensive cutaneous eruption progressing from asymptomatic hypopigmented macules to erythematous papules, plaques, and parasite-rich nodular lesions, representing a vital epidemiologic reservoir for anthroponotic sandfly transmission.

\synonyms
PKDL; Post-Kala-Azar Dermatosis

\medterm{Post-Mortem} Occurring after death or relating to examination of a body afterward. A post-mortem examination can help determine the cause of death, identify disease, and provide information for families or legal investigations.
\medterm{Post-Nasal Drip (Chronic Rhinitis)} Alternate index label for Chronic Rhinitis.

\medterm{Post-Natal Depression} Depression beginning during pregnancy or after childbirth, with persistent low mood or loss of interest and possible effects on sleep, bonding, and daily functioning.

\medterm{Post-Obstructive Pneumonia} Lung infection developing beyond a blocked airway. A tumor, inhaled object, thick mucus, or narrowing can prevent drainage and allow infection; antibiotics may not cure it until the blockage is found and relieved.

\medterm{Post-Operative Atelectasis} Partial collapse of air sacs in the lung after surgery. Anesthesia, pain, shallow breathing, mucus, and lying still can contribute, causing shortness of breath or low oxygen. Pain control, movement, and breathing exercises may help prevent or treat it.

\medterm{Post-Orgasmic Illness Syndrome} A rare condition in which flu-like symptoms, extreme tiredness, headache, nasal or eye symptoms, and difficulty thinking begin soon after orgasm and last for days. The cause is uncertain; repeated episodes need medical assessment to exclude other conditions.

\medterm{Postcoital Bleeding} Bleeding from the vagina after sexual intercourse. Causes include infection, dryness, injury, polyps, pregnancy-related conditions, or cancer.

\synonyms
POIS

\medterm{Post-Partum Haemorrhage} Excessive bleeding after childbirth, most commonly caused by failure of the uterus to contract, retained placental tissue, genital-tract trauma, or clotting problems.

\medterm{Post-Pneumonectomy Syndrome} A rare problem after removal of one lung in which the remaining chest structures shift and press on the airway or major blood vessels. It may cause breathlessness, noisy breathing, repeated infections, or difficulty swallowing.

\medterm{Post-Prandial} Occurring after a meal. Post-prandial blood glucose helps assess diabetes, while post-prandial low blood pressure or chest pain may indicate a circulation or heart problem and needs assessment.
\medterm{Post-Splenectomy Complications} Post-splenectomy complications include infection by encapsulated bacteria, blood clots, bleeding, and changes in blood-cell counts; vaccination and follow-up reduce risk.

\medterm{Post-Streptococcal Glomerulonephritis} Kidney inflammation that can develop after a group A streptococcal throat or skin infection. It may cause dark urine, swelling, reduced urine, high blood pressure, tiredness, or kidney failure and needs medical evaluation.

\medterm{Post-Surgical Pain} Pain after an operation caused by tissue injury and inflammation. Its severity depends on the procedure and the person; uncontrolled pain can limit breathing, movement, sleep, and recovery and should be discussed with the care team.

\medterm{Post-Surgical Rehabilitation} A planned program to restore movement, strength, independence, and safe return to daily activities after surgery. It may include wound care, pain control, exercises, mobility training, and advice tailored to the operation and the person’s recovery.

\medterm{Post-Term Pregnancy} Pregnancy lasting 42 completed weeks or longer. The placenta may provide less support, the baby may become unusually large, and amniotic fluid may decrease. Clinicians monitor parent and baby and discuss delivery timing to reduce complications.

\medterm{Post-Test Probability} The estimated chance that a person has a disease after a test result is known. It depends on the original likelihood of disease and how accurately the test distinguishes people with and without the condition.

\medterm{Post-Thrombotic Syndrome} Chronic pain, heaviness, edema, skin change, and sometimes ulceration developing after
deep-vein thrombosis because of persistent venous obstruction and valvular damage.
Compression, exercise, skin care, and selected interventions may reduce disability.

\medterm{Post-Transfusion Purpura} A rare immune reaction that causes a severe fall in platelets about 5–10 days after a blood transfusion. It can produce pinpoint skin bleeding, bruising, nosebleeds, or dangerous internal bleeding and requires urgent treatment.

\medterm{Post-Traumatic Headache} A headache beginning or continuing after a head or neck injury. It may resemble migraine or tension headache and can occur with dizziness, poor concentration, sleep problems, nausea, or light sensitivity.

\medterm{Post-Traumatic Stress Disorder} A mental-health disorder that may follow actual or threatened death, serious injury, sexual violence, or another traumatic event. It causes ongoing unwanted memories, nightmares, avoidance, emotional changes, and heightened alertness that interfere with daily life.

\medterm{Post-Vasectomy Pain Syndrome} Persistent or recurring testicular or scrotal pain lasting at least three months after vasectomy. The pain may result from nerve irritation, inflammation, or pressure in the reproductive tract.
\synonyms
PVPS; Chronic Post-Vasectomy Pain

\medterm{Post-Void Residual} The amount of urine remaining in the bladder after urination, measured by ultrasound or
catheterization. A persistently high residual may indicate incomplete emptying, but its
significance depends on the volume, timing, symptoms, and clinical context.

\medterm{Post-Void Residual Volume} The amount of urine remaining in the bladder immediately after urination, measured by
ultrasound or catheterization. An elevated residual suggests incomplete emptying but
must be interpreted with timing, bladder volume, symptoms, and repeated measurements.


\medterm{Postanesthesia Care Unit} The postanesthesia care unit is a monitored recovery area where patients are
observed during the immediate postoperative period. Admission assessment includes vital
signs, level of consciousness, pain assessment, and surgical site evaluation.

\medterm{Postcholecystectomy Syndrome} Persistence or recurrence of symptoms after cholecystectomy. Causes include retained
common bile duct stones, sphincter of Oddi dysfunction, bile salt malabsorption,
irritable bowel syndrome, and other gastrointestinal conditions. Evaluation includes
liver function tests, MRCP, and endoscopic ultrasound.

\medterm{Postconcussion Syndrome} Persistent symptoms after a concussion, such as headache, dizziness, fatigue, sleep disturbance, slowed thinking, or mood change. Recovery varies, and worsening neurologic symptoms need urgent assessment.

\medterm{Posterior} Located behind or toward the back of the body or a body structure. The term is used to describe anatomical direction, position, movement, imaging findings, and areas affected by disease or injury.
\medterm{Posterior Chamber} The eye space between the iris and lens, filled with clear fluid that nourishes tissues and helps maintain pressure.


\medterm{Posterior Cruciate Ligament Injury} A stretch or tear of the ligament deep inside the knee that prevents the shinbone from moving too far backward. It may follow a blow to the shin, a fall on a bent knee, or severe twisting and can cause pain, swelling, and instability.

\medterm{Posterior Cruciate Ligament Tear} A partial or complete tear of the strong ligament deep inside the knee that limits backward movement of the shinbone. It often follows a forceful blow to the shin or a fall on a bent knee and may cause pain, swelling, and knee instability. Treatment depends on severity and other injuries.

\medterm{Posterior Descending Artery} A heart artery that runs along the back groove between the lower chambers and supplies part of the back and lower heart muscle. It usually branches from the right coronary artery but may arise from another heart artery.

\medterm{Posterior Drawer Test} A knee examination in which the lower leg is gently pushed backward while the knee is bent. Excessive backward movement may indicate injury to the ligament that stabilizes the back of the knee.

\medterm{Posterior Fossa Tumor} A growth in the lower back part of the skull, where the cerebellum and brainstem are located. It may cause headache, vomiting, unsteady movement, eye problems, or changes in breathing by pressing on nearby tissue or blocking brain fluid.

\medterm{Posterior Heel Bursitis} Inflammation of a bursa near the back of the heel, causing pain and swelling.
\synonyms
Retrocalcaneal bursitis

\medterm{Posterior Keyhole Foraminotomy} An operation through a small opening in the back that enlarges the passage where a spinal nerve leaves the spine. It relieves pressure from causes such as a disk bulge or arthritis and may reduce radiating pain or weakness.
\medterm{Posterior Myocardial Infarction} A heart attack affecting the back wall of the heart, usually because a coronary artery is blocked. It may be missed on a standard ECG, so additional ECG leads or imaging may be needed when symptoms suggest it.

\medterm{Posterior Pituitary} The back part of the pituitary gland, which stores and releases hormones made in the brain. Vasopressin helps control body water and blood pressure, while oxytocin supports uterine contractions and release of breast milk.

\medterm{Posterior Prolapse} Alternate index label; see \textit{Posterior vaginal prolapse (rectocele)}.

\medterm{Posterior Reversible Encephalopathy Syndrome} Brain swelling caused by disturbed blood-flow regulation, often linked to very high blood pressure, pregnancy-related hypertension, kidney disease, severe infection, or immune-suppressing medicines. Headache, seizures, confusion, or vision loss require emergency treatment because prompt care can prevent lasting injury.

\medterm{Posterior Rhinorrhea} Posterior rhinorrhea is mucus draining from the nose into the throat, commonly causing throat clearing or cough from allergy, infection, irritation, or sinus disease.

\medterm{Posterior Subcapsular Cataract} A cloudy area at the back of the eye’s lens, just beneath its outer capsule. Because it often lies in the path of light, it can cause glare, blurred vision, and difficulty reading, especially in bright light.

\medterm{Posterior Tibial Artery} A major artery running down the back of the lower leg and behind the inner ankle. It supplies the lower leg and sole of the foot; its pulse is checked when assessing blood flow to the foot.

\medterm{Posterior Urethral Valves} Abnormal tissue folds present in the urethra of some male babies that block urine leaving the bladder. They can cause a weak urine stream, swollen kidneys, urinary infection, or kidney damage and require specialist treatment.

\medterm{Posterior Uveitis} Inflammation at the back of the eye involving the retina, the layer beneath it, or both. It may cause floaters, blurred or missing vision, or eye pain and can result from infection, immune disease, or another inflammatory condition.

\medterm{Posterior Vitreous Detachment} Separation of the jelly inside the eye from the retina, usually as an age-related change. It may cause new flashes or floaters; urgent eye examination is needed because a retinal tear or detachment can occur at the same time.

\medterm{Posteroanterior} Posteroanterior means directed from back to front; a PA chest radiograph uses this projection to reduce magnification of the heart.

\medterm{Postganglionic Fiber} An autonomic nerve fiber extending from a ganglion to its target organ after the synapse, carrying sympathetic or parasympathetic signals to smooth muscle, glands, or the heart.

\medterm{Postherpetic Neuralgia} Persistent nerve pain in an area affected by shingles after the rash has healed.
\synonyms
Neuralgia, Postherpetic


\medterm{Posthitis} Inflammation of the foreskin, causing redness, swelling, itching, pain, discharge, or difficulty retracting it; infection, irritation, diabetes, and poor hygiene are possible contributors.

\medterm{Posthumous} Occurring, discovered, or published after a person's death. A posthumous diagnosis may be based on records, stored samples, or autopsy findings.
\medterm{Postinflammatory Hyperpigmentation} Darkening of skin after inflammation or injury, caused by increased melanin production or its deposition in the epidermis or dermis.

\medterm{Postmenopausal} Postmenopausal describes the period after menopause, when ovarian estrogen production and menstrual periods have permanently ceased.

\medterm{Postmenopausal Bleeding} Vaginal bleeding after menopause. Even light or one-time bleeding needs medical evaluation because causes include tissue thinning, polyps, hormone effects, and cancer.

\medterm{Postmenopausal Osteoporosis} Bone loss caused by lower estrogen levels associated with menopause. Sometimes called type I osteoporosis.
\medterm{Postmenopause} The stage of life after menopause, when menstrual periods have permanently stopped because the ovaries no longer release eggs regularly and produce much less estrogen.

\medterm{Postmortem Artifact} A change that happens to a body after death and can look like disease or injury. Causes include decomposition, medical resuscitation, embalming, insects, animals, or movement of blood after death; examiners must distinguish it from changes before death.

\medterm{Postmortem Interval} The estimated time between death and the examination or discovery of a body. Investigators consider body temperature, stiffness, pooling of blood, chemical changes, environment, and insect activity, but the estimate is not exact.

\medterm{Postmortem Toxicology} Testing blood, tissues, or other samples after death for medicines, alcohol, poisons, or their breakdown products. Results must be interpreted with the medical history, dose, sample site, decomposition, and possible movement of substances after death.

\medterm{Postmyocardial Infarction Syndrome} Alternate index label; see \textit{Dressler syndrome}.

\medterm{Postnatal} Relating to the period after birth, especially the newborn and infant period. Postnatal care includes checking breathing, feeding, temperature, weight, jaundice, development, and signs of infection or other illness.
\medterm{Postnatal Care} Health care for the parent and baby after birth. It includes checking bleeding, blood pressure, wounds, pain, feeding, mood, newborn growth, and warning signs, while providing guidance about recovery, contraception, vaccination, and follow-up.

\medterm{Postnatal Depression} A depressive episode occurring after childbirth, with persistent low mood, anhedonia,
sleep or appetite disturbance, guilt, impaired bonding, or suicidal thoughts.

\medterm{Postop} Postop is an abbreviation for postoperative, meaning relating to the period after a surgical procedure.
\medterm{Postoperative} Postoperative means occurring or provided after surgery, including recovery, wound care, pain control, rehabilitation, and complication monitoring.

\medterm{Postoperative Care} Care after surgery that monitors breathing, circulation, pain, fluids, the wound, movement, eating, and possible complications. It may include early walking, exercises, medicines, and prevention of blood clots.

\medterm{Postoperative Cognitive Dysfunction} Problems with memory, attention, thinking, or planning after surgery. They may be temporary or persistent, especially after major operations; assessment compares the person with their usual abilities and checks for delirium, medicines, infection, or other causes.

\medterm{Postoperative Complications} Problems that develop after an operation, including infection, bleeding, blood clots, breathing problems, organ dysfunction, wound separation, or medication reactions. Risk depends on the procedure and the person's health; severe pain, fever, breathlessness, or heavy bleeding needs urgent care.

\medterm{Postoperative Delirium} A sudden change in attention and thinking after surgery. A person may become confused, restless, unusually sleepy, or see things that are not there; symptoms may fluctuate and require prompt assessment for medical causes.

\medterm{Postoperative Hemorrhage} Bleeding after surgery that may occur at the operative site or internally and can cause shock, anemia, or the need for urgent intervention.

\medterm{Postoperative Ileus} Temporary slowing or stopping of bowel movement after surgery, especially abdominal surgery. It can cause bloating, nausea, vomiting, pain, and inability to pass gas or stool; medicines, electrolyte problems, and the operation may contribute.

\medterm{Postoperative Infection} An infection developing after surgery, involving the incision, an organ space, implanted material, or the bloodstream.

\medterm{Postoperative Nausea And Vomiting} Nausea, retching, or vomiting occurring after surgery. Risk is influenced by factors
such as previous motion sickness or postoperative nausea, anesthetic technique, opioid
use, and patient characteristics. Prevention and treatment are tailored to the individual
risk and procedure.

\medterm{Postoperative Pain Management} Management of pain after surgery using a combination of medicines, regional
techniques, physical measures, and rehabilitation. The plan is tailored to the procedure, pain severity, medical
conditions, and risks of adverse effects.

\medterm{Postoperative Period} The postoperative period is the time after an operation when recovery, pain control, wound care, mobility, nutrition, complications, and follow-up are monitored.

\medterm{Postoperative Pulmonary Complication} A breathing problem developing after surgery, such as partial lung collapse, pneumonia, aspiration, worsening asthma or COPD, or respiratory failure. Risk is reduced by suitable pain control, movement, breathing exercises, and careful monitoring.

\medterm{Postpartum} Relating to the period after childbirth. Recovery includes changes in bleeding, pain, breasts, hormones, sleep, mood, and body function; important complications include heavy bleeding, infection, high blood pressure, blood clots, depression, and psychosis.

\synonyms
Postnatal

\medterm{Postpartum Anxiety} Excessive, persistent, or distressing anxiety during pregnancy or after childbirth that interferes with sleep, daily
function, bonding, or self-care. It may occur with panic, intrusive worries, or depression and deserves clinical
assessment; thoughts of self-harm or harm to the baby require emergency help.

\medterm{Postpartum Blues} A common, short-lived period of tearfulness, mood changes, irritability, or anxiety beginning
soon after childbirth and usually improving within two weeks. Support, rest, and practical
help are important; persistent, severe, psychotic, or suicidal symptoms require prompt
assessment for a postpartum mental-health disorder.

\medterm{Postpartum Cardiomyopathy (Peripartum Cardiomyopathy)} A dilated cardiomyopathy causing heart failure toward
the end of pregnancy or during the months after delivery. Symptoms may include breathlessness, difficulty lying flat,
swelling, fatigue, or chest discomfort and require prompt assessment.

\synonyms
Peripartum Cardiomyopathy

\medterm{Postpartum Care} Clinical care during the weeks after childbirth that assesses bleeding, uterine involution, blood pressure, infection, wound healing, breastfeeding, mental health, contraception, and recovery.

\medterm{Postpartum Contraception} Contraceptive counseling and care after childbirth, taking account of breastfeeding, venous-thromboembolism risk, medical conditions, fertility intentions, and the timing of method initiation.

\medterm{Postpartum Depression} Depression that begins during pregnancy or after childbirth and lasts beyond ordinary short-term mood changes. It may cause persistent sadness, loss of interest, guilt, anxiety, sleep or appetite changes, difficulty bonding, or thoughts of self-harm and needs treatment.

\synonyms
Depression, Postpartum
PPD (Postpartum Depression)

\medterm{Postpartum Endometritis} Infection of the inner lining of the uterus after childbirth, more common after cesarean birth. It may cause fever, lower-abdominal pain, uterine tenderness, fast heartbeat, or foul-smelling discharge and needs prompt treatment.

\medterm{Postpartum Haemorrhage} Excessive bleeding after childbirth, commonly defined by blood loss of at least 1,000 mL or bleeding with hypovolemia after birth.
\medterm{Postpartum Hemorrhage} Excessive bleeding after childbirth, usually caused by poor uterine contraction, retained placental tissue, genital-tract
trauma, or a clotting disorder. It can cause shock and requires urgent obstetric
assessment and treatment.

\medterm{Postpartum Hemorrhage (Secondary)} Heavy or abnormal bleeding occurring more than 24 hours and up to 12 weeks after childbirth. Causes include retained placental tissue, infection, poor healing where the placenta was attached, or a bleeding disorder. Heavy bleeding, dizziness, fainting, fever, or severe pain requires urgent care.

\synonyms
Secondary

\medterm{Postpartum Infection} An infection after childbirth, commonly involving the uterus, surgical wound, bladder, breasts, or bloodstream. Fever, worsening abdominal or pelvic pain, foul-smelling discharge, wound redness, breast redness, or feeling very unwell requires prompt medical assessment.

\medterm{Postpartum Preeclampsia} New-onset hypertension with pre-eclampsia features after childbirth, sometimes occurring several days or weeks postpartum. Severe headache, visual symptoms, chest pain, shortness of breath, or seizures require urgent assessment.

\medterm{Postpartum Psychosis} A rare but severe psychiatric emergency that can begin suddenly during the days or weeks after childbirth. Symptoms may include hallucinations, delusions, paranoia, severe confusion, mania, depression, disorganized behavior, or rapidly changing moods. It can worsen quickly and may threaten the safety of the mother or baby, so immediate medical assessment is required.

\synonyms
Postnatal Psychosis
Puerperal Psychosis

\medterm{Postpartum Thyroiditis} Inflammation of the thyroid during the first year after childbirth, caused by an immune reaction. It may first cause anxiety, sweating, and a fast heartbeat, followed by tiredness, cold intolerance, or low thyroid function; recovery varies.

\medterm{Postpartum Urinary Retention} Inability to empty the bladder adequately after childbirth, sometimes associated with epidural anesthesia, perineal trauma, swelling, or prolonged labor. Prompt assessment prevents bladder overdistension and injury.

\medterm{Postprandial} Postprandial means occurring after a meal, especially the changes in glucose, insulin, digestion, and circulation after eating.

\medterm{Postprandial Period} The postprandial period is the interval after eating when digestion and nutrient absorption occur and glucose, insulin, and other metabolic responses change.

\medterm{Postrenal AKI} Sudden loss of kidney function caused by blocked urine flow. Prostate enlargement, stones, tumors, scarring, or bladder-nerve problems may be responsible; kidney damage becomes more likely when the blockage is severe or lasts a long time.

\medterm{Poststreptococcal Glomerulonephritis} Poststreptococcal glomerulonephritis is kidney inflammation occurring after certain streptococcal infections, causing blood or protein in urine, swelling, high blood pressure, and reduced kidney function.

\medterm{Poststreptococcal Pustulosis} A rare immune-related outbreak of pus-filled spots after a streptococcal throat infection. The spots usually appear on reddened skin of the trunk or limbs and should be assessed to confirm the cause and exclude infection or a medicine reaction.

\medterm{Postsynaptic Membrane} The receiving surface of a nerve or muscle cell where a neighboring nerve releases a chemical signal. Receptors in this surface detect the signal and help start or change an electrical response.

\medterm{Postterm Pregnancy} Pregnancy continuing beyond 42 completed weeks. The placenta may work less effectively and the baby may become unusually large or pass stool before birth, so fetal wellbeing and delivery timing require medical monitoring.

\medterm{Posttraumatic Stress Disorder} Alternate index label for Post-Traumatic Stress Disorder.

\synonyms
PTSD (Posttraumatic Stress Disorder)

\medterm{Postural} Postural means relating to body position or posture, including symptoms or vital signs that change when a person sits, stands, or lies down.

\medterm{Postural Drainage} Using specific body positions so gravity helps mucus move from smaller airways into larger ones, where it can be coughed out. A respiratory clinician selects positions based on the affected lung area and the person’s condition.
\medterm{Postural Hypotension} A drop in blood pressure after standing that may cause dizziness, blurred vision, weakness, or fainting. It can result from dehydration, medicines, nerve disorders, blood loss, or heart problems and increases fall risk.
\medterm{Postural Orthostatic Tachycardia Syndrome} A disorder of the autonomic nervous system, which helps control heart rate and blood flow. When a person stands, the heart rate rises unusually quickly while blood pressure does not fall substantially, causing ongoing symptoms of poor tolerance of upright posture. In adults, the heart rate typically rises by at least 30 beats per minute within 10 minutes of standing; the increase is at least 40 beats per minute in adolescents. Symptoms may include dizziness, fainting or near-fainting, racing heartbeat, tiredness, weakness, headache, blurred vision, nausea, trouble thinking, or poor ability to exercise.

\medterm{Posture} The alignment and position of the body while resting or moving. Posture depends on bones, joints, muscles, nerves, vision, and balance and can affect comfort, breathing, movement, and injury risk.
\medterm{Potassium} An essential mineral and body salt needed for nerve signals, muscle movement, and normal heart rhythm.
\medterm{Potassium (K)} The chemical symbol for potassium, a major intracellular electrolyte needed for nerve signaling, muscle contraction, and cardiac electrical activity.
\medterm{Potassium Acetate} A potassium salt used to replace potassium or provide acetate in selected medical and laboratory applications; dosing requires monitoring of kidney function and serum potassium.

\medterm{Potassium Binders} Medicines taken by mouth that trap potassium in the intestine so it leaves the body in stool. They help lower persistently high blood potassium in selected people, especially with kidney disease, but require monitoring for constipation, swelling, and other electrolyte changes.

\medterm{Potassium Carbonate Poisoning} Swallowing potassium carbonate can irritate or burn the mouth, throat, and stomach and may raise blood potassium dangerously. Vomiting, weakness, abnormal heartbeat, or breathing difficulty requires emergency care.

\medterm{Potassium Channel} A membrane protein that allows potassium ions to cross cells and helps control electrical excitability, heart rhythm, and hormone release.

\medterm{Potassium Chloride} A salt used to prevent or treat low blood potassium. It may be taken by mouth or given into a vein; too much can cause dangerous heart-rhythm problems, especially when kidney function is reduced.

\medterm{Potassium Citrate} A medicine that makes urine less acidic and increases urinary citrate, helping prevent some kidney stones. It can dangerously raise blood potassium in people with kidney disease or those taking certain medicines.

\medterm{Potassium Cyanide} A highly poisonous chemical that prevents cells from using oxygen. Swallowing or inhaling it can rapidly cause headache, confusion, seizures, collapse, heart failure, and death and requires emergency treatment.

\medterm{Potassium Deficiency} An abnormally low blood potassium level, which can cause weakness, cramps, constipation, abnormal heart rhythms, or paralysis.

\medterm{Potassium Dichromate} A toxic, corrosive chemical that can burn the skin, eyes, mouth, and digestive tract. Swallowing or inhaling it may damage the kidneys, liver, blood, or lungs and requires urgent poison-control or emergency care.

\medterm{Potassium Hydroxide Poisoning} A severe chemical burn caused by swallowing or contacting potassium hydroxide, a strong alkali. It can damage the mouth, throat, stomach, skin, eyes, and lungs; rinse exposed areas and seek emergency help immediately.

\medterm{Potassium In Diet} Dietary potassium comes from fruits, vegetables, beans, dairy, and other foods; intake may need adjustment in kidney disease or with medicines that change potassium levels.

\medterm{Potassium Iodide} Potassium iodide supplies iodide and is used in selected thyroid conditions, radiation emergencies, and some expectorant preparations; excess can disturb thyroid function.

\medterm{Potassium Ion} The positively charged form of potassium, the principal intracellular cation that helps establish resting membrane potential, nerve transmission, muscle contraction, and cardiac rhythm.

\medterm{Potassium Permanganate} A chemical sometimes used as a very dilute skin antiseptic or drying solution. Concentrated material can burn tissue, stain skin, and poison if swallowed, so it must be prepared and used carefully.
\medterm{Potassium Test} A potassium test measures potassium in blood, usually to assess heart, muscle, kidney, or medication-related electrolyte problems.

\medterm{Potassium Urine Test} A urine potassium test measures potassium excretion and helps determine whether abnormal blood potassium is related to kidney losses or other causes.

\medterm{Potassium-Sparing Diuretic} A diuretic that promotes sodium excretion while reducing potassium loss, such as spironolactone or amiloride.
\medterm{Potassium-Sparing Diuretics} Water tablets that help remove salt and water while limiting potassium loss. They are used for selected blood-pressure, heart, liver, or hormone conditions, but can raise potassium dangerously, especially in kidney disease or with certain medicines.

\medterm{Potato Plant Poisoning - Green Tubers And Sprouts} Green potatoes, sprouts, and leaves can contain high levels of natural toxins that cause nausea, vomiting, diarrhea, abdominal pain, confusion, weakness, or rarely seizures. Contact poison control after significant ingestion.

\medterm{Potbellies And Toddlers} A rounded abdomen in a toddler is often normal body shape or posture, but persistent distension with pain, vomiting, poor growth, constipation, or a mass needs medical assessment.

\medterm{Potency} The amount of a drug required to produce a specified effect. A more potent drug produces that effect at a
lower dose, but potency does not indicate the maximum effect.


\medterm{Potential} Potential describes a capacity, possibility, or stored energy that may become clinically relevant depending on circumstances and additional findings.

\medterm{Potential Acuity Meter} An eye-testing device that estimates the vision a person might achieve after treatment of an eye problem such as a cataract. Results can be limited by retinal, nerve, or other disease.


\medterm{Potocki-Lupski Syndrome} A genetic condition caused by an extra copy of a small section of chromosome 17. It may cause low muscle tone, feeding and growth problems, developmental or learning difficulties, autism-related features, behavior differences, and congenital heart defects.

\synonyms
PTLS; Duplication 17p11.2 Syndrome

\medterm{Pott Disease} Tuberculosis infection of the bones of the spine. It can cause persistent back pain, collapse or bending of the spine, an abscess, and pressure on the spinal cord leading to weakness or paralysis.

\medterm{Pott's Disease} Tuberculosis of the spine, causing vertebral destruction, back pain, deformity, or neurologic compromise.
\medterm{Pott's Fracture} An ankle fracture-dislocation involving the malleoli, the bony prominences on each side of the ankle. The historical term is used variably, so the exact fracture pattern should be described on imaging.
\medterm{Potter Syndrome} A pattern of birth abnormalities caused by very low amniotic fluid, often because both fetal kidneys are absent or cannot make urine. It may cause underdeveloped lungs, distinctive facial features, and limb changes and is usually life-threatening.

\medterm{Pouch} A sac-like body structure or a surgically created reservoir. Examples include a cheek pouch, a pouch in the digestive tract, and an ileal pouch made from small intestine.
\medterm{Pouch Of Douglas} The pouch of Douglas is the peritoneal recess between the uterus and rectum, also called the rectouterine pouch, where fluid may collect.

\medterm{Pound} A unit measuring weight or force, with meaning determined by context. In healthcare, pounds may describe body weight, pressure, equipment loads, or medicine instructions and should be distinguished from metric units.
\medterm{Povidone} Povidone is a water-soluble polymer used in medicines, eye products, and other preparations; it is also part of povidone-iodine antiseptic.

\medterm{Povidone-Iodine} Povidone-iodine is an iodine-containing antiseptic used on skin or tissue before selected procedures; it can irritate skin and affect the thyroid with extensive use.

\medterm{Powassan Encephalitis} Inflammation of the brain caused by Powassan virus, a tick-borne flavivirus. It can cause fever, headache, vomiting, confusion, seizures, or weakness and may require hospitalization.

\medterm{Powassan Virus Disease} Powassan virus disease is a rare tick-borne infection that can cause fever, headache, vomiting, confusion, seizures, meningitis, or encephalitis. Treatment is supportive, and prevention relies on avoiding ticks and removing them promptly.

\medterm{Powder Diffraction} An analytical method that identifies crystalline substances by measuring how an X-ray or other beam is diffracted by powdered material.

\medterm{Power} The rate at which work is performed; in exercise, it combines force and movement speed.
\medterm{Power Training} Exercise designed to improve the ability to produce force rapidly.

\synonyms
Explosive training

\medterm{Pox} A disease producing vesicles, pustules, or pitted scars; the term is used in names such as smallpox and cowpox.
\medterm{Poxviridae} A family of large DNA viruses that includes orthopoxviruses, molluscipoxvirus, and parapoxviruses; some cause significant human disease.

\medterm{Poxviridae Infection} An infection caused by a poxvirus, ranging from localized skin lesions to severe systemic disease depending on the virus and host.

\medterm{Poxvirus} A family of large DNA viruses that includes variola, vaccinia, mpox, and molluscum contagiosum viruses.
\medterm{Ppar Alpha} A protein inside cells that helps control how the body uses fats. Fibrate medicines activate it to lower some blood fats, especially triglycerides.

\medterm{Ppar Delta} A protein inside cells that helps regulate fat use, energy balance, inflammation, and the development of some tissues. It is mainly studied in laboratory and drug research.
\medterm{PPD} PPD may mean purified protein derivative used in tuberculosis skin testing or post-partum depression; context determines the intended meaning.

\medterm{PPD Skin Test} A PPD skin test places a small amount of tuberculin under the skin to assess immune evidence of tuberculosis exposure; interpretation depends on risk and the size of the reaction.

\medterm{PR Interval} The ECG interval from the start of atrial depolarization to the start of ventricular depolarization, reflecting atrioventricular conduction.

\medterm{Practitioner} A person qualified and authorized to practice a healthcare profession.
\medterm{Prader-Labhart-Willi Syndrome} Alternate index label; see \textit{Prader-Willi syndrome}.

\medterm{Prader-Willi Syndrome} A genetic disorder that often causes weak muscle tone and feeding difficulty in infancy, followed later by unusually strong hunger, obesity, short stature, learning difficulties, hormone problems, and behavioral changes.

\synonyms
Prader-Labhart-Willi Syndrome

\medterm{Pragmatic} Focused on what works in ordinary real-world conditions. A pragmatic clinical trial studies how an intervention performs in usual healthcare rather than under highly controlled conditions.
\medterm{Pragmatic Trial} A pragmatic trial evaluates an intervention under usual real-world care conditions to estimate how well it works in routine practice.

\medterm{Pralatrexate} A folate-blocking chemotherapy medicine used for some relapsed or treatment-resistant peripheral T-cell lymphomas. It can cause mouth sores, low blood counts, infection, fatigue, nausea, and other serious treatment-related effects.

\medterm{Pralidoxime} An oxime that can reactivate inhibited acetylcholinesterase when administered early
after organophosphate poisoning, particularly for nicotinic weakness and
respiratory-muscle dysfunction. It is used with atropine and supportive care according
to exposure and local guidance.

\medterm{Pramipexole} A medicine that stimulates dopamine receptors and is used mainly for Parkinson disease or restless legs syndrome. It can cause nausea, sleepiness, dizziness, sudden sleep, hallucinations, or impulse-control problems.
\medterm{Prandial} Prandial means relating to a meal, as in preprandial before eating or postprandial after eating.
\medterm{Prasugrel} An antiplatelet medicine that prevents blood cells from clumping. It is used in selected people with acute coronary syndrome receiving coronary stents, but it increases bleeding risk and is unsuitable for some patients.
\medterm{Pravastatin} A statin that lowers LDL cholesterol by inhibiting HMG-CoA reductase.
\medterm{Pravastatin Sodium} A statin medicine that lowers harmful LDL cholesterol by reducing cholesterol production in the liver. It helps reduce heart attack and stroke risk alongside diet, exercise, and other risk-factor treatment.

\medterm{Praziquantel} A medicine used to treat several infections caused by parasitic worms, including schistosomiasis, tapeworms, and flukes. It damages the worms so the body can remove them; dizziness, sleepiness, nausea, or abdominal discomfort may occur.
\medterm{Pre-Cancerous} Describing abnormal cells or tissue that are not cancer but have a higher-than-usual chance of becoming cancer. The risk varies by the tissue and the severity of change; monitoring or removal may reduce the risk.
\medterm{Pre-Eclampsia} A disorder that can develop during pregnancy, usually after 20 weeks, or during the weeks after childbirth. It involves new high blood pressure together with protein in the urine or signs of organ injury. Possible findings include a low platelet count, abnormal kidney or liver function, severe headache, vision changes, upper-abdominal pain, or fluid in the lungs. During pregnancy, it can affect the placenta and fetal growth. It may progress to eclampsia, which is pre-eclampsia with seizures.
\medterm{Pre-Existing Diabetes And Pregnancy} Pre-existing type 1 or type 2 diabetes during pregnancy increases risks to parent and fetus, so glucose control, medicines, fetal growth, blood pressure, and eye and kidney health need coordinated care.

\medterm{Pre-Hepatic Jaundice (Haemolytic Jaundice)} Jaundice caused by excessive breakdown of red blood cells before bilirubin reaches the liver, producing increased unconjugated bilirubin.

\synonyms
Haemolytic Jaundice

\medterm{Pre-MiRNA} A pre-miRNA is an early hairpin-shaped microRNA transcript processed in the nucleus before becoming a mature regulator of gene expression.

\medterm{Pre-Synaptic} Located before a synapse, especially on the nerve ending that releases a chemical messenger. Presynaptic activity controls how much signal reaches the next nerve cell, muscle, or gland.
\medterm{Prealbumin} Another name for transthyretin, a protein made mainly by the liver that carries thyroid hormone and vitamin A in the blood. Its blood level changes with inflammation, illness, kidney disease, and nutrition, so it is not a nutrition measure by itself.

\medterm{Prebiotics} Food components, usually nondigestible fibers, that feed selected helpful bacteria in the large intestine. Sources include garlic, onions, leeks, asparagus, bananas, and whole grains; increasing them gradually may reduce gas.

\medterm{Prebiotics And Probiotics} Prebiotics are food components that feed certain helpful gut bacteria. Probiotics are live microorganisms found in some foods or supplements. Their effects depend on the specific product and health condition; they are not interchangeable treatments.

\medterm{Precancerous Condition} A precancerous condition contains abnormal changes that may progress to cancer but are not cancer themselves.

\medterm{Precancerous Polyp} A precancerous polyp is a growth with cellular changes that can become cancerous over time if not removed or monitored appropriately.

\medterm{Precaution} An action taken in advance to lower the chance of harm, infection, or error. Examples include handwashing, protective equipment, checking a medicine, using a seat belt, and avoiding a known trigger.
\medterm{Precipitate Labour} Unusually rapid labour, often defined as delivery within three hours of the onset of regular contractions.

\medterm{Precipitating Factor} A precipitating factor is an event or exposure that triggers the onset or worsening of symptoms or disease in a susceptible person.

\medterm{Preclinical} Before human testing, or describing disease changes that occur before symptoms become noticeable or cause recognized illness.

\medterm{Precocious Puberty} Development of puberty-related body changes earlier than expected. It may result from early brain signals, hormone-producing disorders, medicines, or other conditions and requires evaluation to protect growth and emotional wellbeing.

\synonyms
Early Puberty
Premature Puberty

\medterm{Preconception Care} Health care before pregnancy to improve the chance of a healthy pregnancy. It reviews medical conditions, medicines, vaccinations, genetic risks, nutrition, folic acid, substance exposure, and lifestyle, and helps plan changes before conception.

\medterm{Preconception Counseling} Healthcare advice before pregnancy to improve the health of the pregnant person and future baby. It reviews medicines, vaccines, chronic diseases, nutrition, genetic risks, harmful exposures, and recommended supplements.

\medterm{Precordial Catch Syndrome} Brief, sharp, localized chest pain near the heart that worsens with breathing and usually occurs in otherwise healthy children or adolescents.

\medterm{Precordium} The area of the front chest wall over the heart and nearby large blood vessels. A clinician examines this area by looking, feeling for the heartbeat or unusual movement, and listening for heart sounds and murmurs.

\medterm{Precursor} A substance that the body can change into another substance with a useful biological effect. For example, some vitamins, hormones, and medicines begin as precursors and are changed by enzymes in particular tissues.

\medterm{Prediabetes} Blood glucose higher than normal but below the diagnostic threshold for diabetes, usually identified by A1C, fasting glucose, or an oral glucose-tolerance test; it increases future diabetes and cardiovascular risk.

\synonyms
Diabetes Pre (Prediabetes)

\medterm{Prediabetes Syndrome} Blood glucose levels above normal but below the diabetes threshold, associated with increased future risk of type 2 diabetes and cardiovascular disease.

\medterm{Predictive Test} A test that estimates whether a person or disease is likely to respond to a particular treatment.

\medterm{Predispose} To predispose is to increase a person’s susceptibility or likelihood of developing a disease or response when additional triggers occur.

\medterm{Prednimustine} Prednimustine is a former antineoplastic drug combining a corticosteroid-related group with an alkylating component and causing immunosuppression and marrow toxicity.

\medterm{Prednisolone} Prednisolone is a corticosteroid used to reduce inflammation and immune activity; prolonged or high-dose use can cause infection, bone loss, and other effects.

\medterm{Prednisone} Prednisone is a corticosteroid used for inflammatory, allergic, and autoimmune conditions; it should often be tapered after prolonged use.

\medterm{Preeclampsia} A disorder that can develop during pregnancy, usually after 20 weeks, or during the weeks after childbirth. It involves new high blood pressure together with protein in the urine or signs of organ injury. Possible findings include a low platelet count, abnormal kidney or liver function, severe headache, vision changes, upper-abdominal pain, or fluid in the lungs. During pregnancy, it can affect the placenta and fetal growth. It may progress to eclampsia, which is preeclampsia with seizures.

\synonyms
Pregnancy-Related Hypertension


\medterm{Preeclampsia With Severe Features} A severe form of preeclampsia involving very high blood pressure or damage to organs such as the brain, liver, kidneys, lungs, or blood-clotting system. Severe headache, vision changes, breathing difficulty, or upper-abdominal pain requires emergency care.

\medterm{Preexcitation Syndrome} A heart-rhythm condition in which an extra electrical pathway activates the lower heart chambers too early. It can cause episodes of a very fast heartbeat, dizziness, or fainting and may require specialist assessment.

\medterm{Prefrontal} Relating to the anterior part of the frontal lobe involved in planning, judgment, working memory, and behavior.
\medterm{Prefrontal Cortex} The prefrontal cortex is the front part of the frontal lobe involved in planning, judgment, working memory, impulse control, and social behavior.

\medterm{Pregabalin} Pregabalin is used for certain nerve pain and seizures and may cause dizziness, sleepiness, swelling, or dependence.

\medterm{Preganglionic Fiber} A nerve fiber that carries automatic body signals from the brain or spinal cord to a relay station before they reach an organ. These signals help control the heart, blood vessels, glands, digestion, and bladder.

\medterm{Pregnancy} The period during which an embryo or fetus develops in the uterus, usually lasting about 40 weeks from the first day of the last menstrual period until birth.

\medterm{Pregnancy - Health Risks} Health risks in pregnancy vary with medical conditions, medicines, infections, age, exposures, and obstetric history; prenatal assessment identifies preventable complications and needed monitoring.

\medterm{Pregnancy - Identifying Fertile Days} Fertile days are the interval around ovulation when intercourse can lead to pregnancy; cycle timing, cervical mucus, basal temperature, and ovulation tests can provide estimates but are imperfect.

\medterm{Pregnancy And Herpes} Genital herpes during pregnancy can affect delivery planning and rarely infect the newborn; antiviral treatment and examination of lesions near labor help reduce transmission risk.

\medterm{Pregnancy And The Flu} Influenza during pregnancy can be more severe and affect the fetus; vaccination is recommended in pregnancy, and suspected illness should receive prompt clinical advice about antiviral treatment.

\medterm{Pregnancy And Travel} Travel during pregnancy depends on gestational age, health, destination infections, transport, and access to care; prevention includes vaccines as appropriate, food and water safety, movement, and medical planning.

\medterm{Pregnancy And Work} Work during pregnancy is usually possible but may require adjustments for heavy lifting, prolonged standing, heat, chemicals, radiation, infection, night work, or other occupational hazards.

\medterm{Pregnancy at Age 35 or Older} Pregnancy in a person who will be 35 years or older at the estimated time of delivery. It is associated with increased risks of some chromosomal conditions and pregnancy complications; individual risk depends on overall health and pregnancy history.

\synonyms
Advanced Maternal Age


\medterm{Pregnancy Complication} A maternal, fetal, placental, or labor-related problem that disrupts an otherwise expected pregnancy, such as preeclampsia, hemorrhage, diabetes, growth restriction, or preterm labor.


\medterm{Pregnancy Ectopic (Ectopic Pregnancy)} Alternate index label for Ectopic Pregnancy.

\medterm{Pregnancy High Blood Pressure (Pregnancy: Preeclampsia And Eclampsia)} Alternate index label for Pregnancy: Preeclampsia and Eclampsia.

\medterm{Pregnancy Nutrition} Food and supplements during pregnancy that support the pregnant person and fetal growth. Needs vary, but a varied diet usually includes enough protein, folate, iron, iodine, calcium, vitamin D, and other nutrients; individual advice is important.

\medterm{Pregnancy Of Unknown Location} A positive pregnancy test without a confirmed intrauterine or ectopic pregnancy on ultrasound. Serial human chorionic gonadotropin measurements and repeat imaging help determine whether the pregnancy is early, failing, or ectopic.

\medterm{Pregnancy Outcome} Pregnancy outcome describes what happens to a pregnancy, such as live birth, miscarriage, stillbirth, preterm birth, or other maternal or fetal result.

\medterm{Pregnancy Status} Pregnancy status indicates whether a person is pregnant, not pregnant, or of unknown status and can affect medication, imaging, and treatment decisions.

\medterm{Pregnancy Test} A urine or blood test that detects human chorionic gonadotropin, a hormone made soon after an embryo attaches to the uterus. Testing too early can be falsely negative; an unexpected result may need repeat testing or medical assessment.
\medterm{Pregnancy Trimester} One of three approximate stages of pregnancy: first through 13 weeks, second from 14 through 27 weeks, and third from 28 weeks until birth. The stages help organize fetal development, screening, and care.

\medterm{Pregnancy Tubal (Ectopic Pregnancy)} Alternate index label for Ectopic Pregnancy.


\medterm{Pregnancy-Induced Hypertension} New-onset hypertension during pregnancy without the proteinuria or organ
dysfunction that defines preeclampsia. Diagnosis and treatment depend on repeated blood-pressure measurements,
gestational age, symptoms, and evidence of maternal or fetal complications.

\medterm{Pregnancy-Related Hypertension} Alternate index label; see \textit{Preeclampsia}.

\medterm{Pregnancy: Bleeding During The First Trimester} First-trimester bleeding may occur with implantation, cervical changes, miscarriage, ectopic pregnancy, or other causes. Assessment uses symptoms, examination, serial pregnancy hormone testing, and ultrasound; heavy bleeding, severe pain, faintness, or shoulder pain is urgent.

\medterm{Pregnancy: Pain Relief Options For Birth} Labor pain relief may include breathing and movement, water, massage, nitrous oxide, intravenous medicines, epidural analgesia, or spinal anesthesia. Choice depends on labor stage, maternal and fetal status, contraindications, availability, and informed preference.

\medterm{Pregnancy: Placenta Previa} Placenta previa occurs when the placenta lies over or near the cervical opening and can cause painless bright-red bleeding, especially later in pregnancy. Digital vaginal examination is avoided until location is known; recurrent bleeding requires obstetric planning and sometimes cesarean delivery.

\medterm{Pregnancy: Preeclampsia And Eclampsia} Preeclampsia is new hypertension during pregnancy with proteinuria or organ dysfunction; eclampsia is the occurrence of seizures. Severe headache, visual symptoms, upper-abdominal pain, shortness of breath, or seizure requires emergency obstetric treatment.

\synonyms
Eclampsia Pregnancy (Pregnancy: Preeclampsia And Eclampsia)
Pregnancy High Blood Pressure (Pregnancy: Preeclampsia And Eclampsia)

\medterm{Pregnane} Pregnane is a 21-carbon steroid hydrocarbon that forms the structural backbone of progesterone and several corticosteroid and progestin compounds.

\medterm{Pregnanolone} Pregnanolone is a neuroactive steroid metabolite that can modulate GABA-A receptors and influence neuronal excitability, mood, and sedation.

\medterm{Pregnant} Pregnant means carrying a developing embryo or fetus in the uterus after conception; pregnancy status affects medicines, imaging, and clinical decisions.

\medterm{Pregnenolone} Pregnenolone is a steroid hormone precursor made from cholesterol that can be converted into progesterone, corticosteroids, and sex hormones.

\medterm{Prehn Sign} A test in which lifting the painful testicle is said to relieve pain in epididymitis but not testicular torsion. This finding is unreliable and cannot safely rule out torsion. Sudden severe testicle pain, swelling, nausea, or a high-riding testicle requires urgent assessment, usually with Doppler ultrasound and immediate surgery when torsion cannot be excluded.

\synonyms
Prehn's Sign

\medterm{Prehypertension} An older term for blood pressure that is higher than ideal but below the level once used to diagnose high blood pressure. Current guidelines use different categories; repeated readings, overall heart risk, and lifestyle are considered when deciding what help is needed.

\synonyms
Elevated blood pressure

\medterm{Preimplantation Genetic Testing} Genetic testing of cells from an embryo created through in-vitro fertilization before possible transfer to the uterus, used in selected situations to identify specific inherited disorders or chromosome-number abnormalities.

\medterm{Preload} The amount of blood filling the heart’s lower chambers and stretching their muscle just before each heartbeat. It mainly depends on blood returning to the heart and influences how strongly the heart can pump.

\medterm{Preload Reduction} A decrease in the amount of blood returning to the heart, which lowers the heart’s filling pressure and workload. It may occur with fluid loss or medicines and can help some forms of heart failure.
\medterm{Premalignant} Describing cells or tissue changes that have a higher chance of becoming cancer but have not grown into nearby tissues. Not every premalignant change becomes cancer; follow-up or removal depends on its type, location, size, and risk.
\medterm{Premature} Occurring earlier than expected or before the usual developmental or pregnancy-related time. A premature birth occurs before thirty-seven completed weeks and may increase risks involving breathing, feeding, temperature, and development.
\medterm{Premature (Early) Ejaculation} A persistent or recurrent pattern of ejaculation that occurs sooner than desired,
with limited control and associated personal distress. Assessment considers the person's typical experience and sexual
context.

\synonyms
Early

\medterm{Premature Adrenarche} Early appearance of pubic or axillary hair, body odor, or mild acne from adrenal androgen production without true central puberty; evaluation excludes virilizing disorders.

\medterm{Premature Atrial Contraction} An early heartbeat that starts in the upper chambers of the heart rather than at the usual pacemaker. It may feel like a skipped or extra beat and is often harmless, but frequent episodes or fainting need assessment.

\medterm{Premature Baby} An infant born before 37 completed weeks of pregnancy, with risks related to immature lungs, brain, feeding, temperature control, immunity, and other organ systems.

\medterm{Premature Birth} Birth before 37 completed weeks of pregnancy; the earlier the birth, the greater the risks of respiratory, neurological, and developmental complications.

\medterm{Prematurity} Birth before 37 completed weeks of pregnancy. Prematurity can affect breathing, feeding, temperature control, immunity, brain development, and growth, especially when birth occurs very early.

\medterm{Premature Closure} Closing of a body opening, growth plate, or developmental connection earlier than expected. The result depends on the site and may include blockage, abnormal body shape, or restricted growth; treatment depends on the cause.

\medterm{Premature Ejaculation} Ejaculation that occurs sooner than desired with poor control and causes persistent personal or relationship distress.
\medterm{Premature Infant} An infant born before 37 completed weeks of pregnancy. Earlier birth increases the risk of breathing, feeding, temperature-control, infection, and developmental problems.
\medterm{Premature Labor} Regular uterine contractions with cervical change before 37 completed weeks of pregnancy, which can lead to preterm birth and requires assessment of membranes, infection, and fetal well-being.

\medterm{Premature Labour (Preterm Labour)} Labour that begins before 37 completed weeks of pregnancy and may result in preterm birth.
\synonyms
Preterm Labour

\medterm{Premature Menopause} Loss of ovarian function and menstruation before age 40, now often termed primary ovarian insufficiency.
\medterm{Premature Menopause (Medical Procedural Causes)} Alternate index label for Medical Procedural Causes.

\medterm{Premature Ovarian Failure} Premature ovarian failure is loss of normal ovarian function before age 40, causing irregular or absent periods, infertility, and low-estrogen symptoms; causes may be genetic, autoimmune, toxic, or unknown.

\medterm{Premature Ovarian Insufficiency} Reduced or intermittent ovarian function before age 40, causing irregular or absent periods, reduced fertility, and sometimes hot flushes, vaginal dryness, or low bone density. Causes include genetic conditions, autoimmune disease, chemotherapy, radiation, surgery, or no identifiable cause.

\medterm{Premature Puberty} Alternate index label; see \textit{Precocious puberty}.

\medterm{Premature Rupture Of Membranes} Breaking of the amniotic sac before labor begins. If it happens before 37 weeks, it is called preterm rupture; a gush or continuing leak of fluid needs prompt obstetric assessment because infection and early birth may follow.

\medterm{Premature Thelarche} Isolated breast development in a young child without other signs of puberty, often transient but assessed for growth acceleration, estrogen exposure, or true precocious puberty.

\medterm{Premature Ventricular Contraction} An early heartbeat that starts in a ventricle instead of following the heart's usual rhythm. It may feel like a skipped or extra beat and is often harmless, but frequent PVCs may occur with fainting, chest pain, or breathlessness.

\synonyms
PVC

\medterm{Premature Ventricular Contractions} Ectopic beats originating from the ventricles, appearing as wide, bizarre QRS complexes
that occur earlier than expected. PVCs are common in the general population and are
usually benign. Frequent PVCs (>10\% of total beats) may cause left ventricular
dysfunction.

\synonyms
Extrasystole (Premature Ventricular Contractions)
PVC (Premature Ventricular Contractions)

\medterm{Premature Ventricular Contractions (PVCs)} Early extra heartbeats originating in the heart's ventricles, often felt as fluttering, skipped beats, or pounding and frequently benign.

\synonyms
Extrasystole
PVCs
Ventricular Premature Contraction
VPCs

\medterm{Premenstrual} Occurring during the days before menstruation; the term commonly refers to cyclical symptoms such as mood change, breast tenderness, bloating, headache, or fatigue.

\medterm{Premenstrual Breast Changes} Premenstrual breast changes include cyclic fullness, tenderness, swelling, or lumpiness caused by hormonal variation; a persistent one-sided mass, skin change, or bloody discharge needs evaluation.

\medterm{Premenstrual Dysphoric Disorder} A severe form of premenstrual symptoms that repeatedly causes depression, irritability, anxiety, mood swings, or loss of control along with physical symptoms and major disruption of work, relationships, or daily life before menstruation.

\medterm{Premenstrual Syndrome} A recurring group of physical and emotional symptoms in the days or weeks before menstruation that improves when bleeding begins. Symptoms may include irritability, mood changes, bloating, breast tenderness, headache, tiredness, or appetite changes.

\medterm{Premolar (Bicuspid)} One of the permanent teeth between the canine teeth and molars. Premolars usually have two raised points and help tear and crush food; decay, crowding, or infection may require dental treatment.
\medterm{Prenatal} Prenatal means occurring during pregnancy before birth, including care, testing, development, counseling, and treatment of the pregnant person or fetus.

\medterm{Prenatal Care} Regular health care during pregnancy to monitor the pregnant person and developing baby. Visits may include blood-pressure and urine checks, blood tests, ultrasound, screening, vaccination, nutrition and medicine review, and advice about warning signs.




\medterm{Prenatal Development} Prenatal development is growth and differentiation of the embryo and fetus from conception to birth, influenced by genes, nutrition, maternal health, and exposures.

\medterm{Prenatal Diagnosis} Testing during pregnancy to identify a possible genetic, chromosome, or structural condition in the fetus. It may use ultrasound, blood-based screening, or diagnostic tests such as chorionic-villus sampling and amniocentesis, depending on the clinical question and stage of pregnancy.

\medterm{Prenatal Exposure} Contact of a fetus with a medicine, substance, infection, or environmental factor during pregnancy.
Effects depend on the agent, dose, timing, and fetal susceptibility.

\medterm{Prenatal Genetic Counseling} A consultation that reviews the risk of inherited or chromosomal conditions, available screening and diagnostic tests, possible results, and reproductive options before or during pregnancy.

\medterm{Prenatal Testing} Screening and diagnostic tests performed during pregnancy to assess fetal health and development or identify genetic, chromosomal, and structural conditions. Tests may include ultrasound, maternal blood tests, cell-free DNA screening, chorionic-villus sampling, and amniocentesis.

\medterm{Prenatal Vitamins} Supplements designed to provide nutrients needed during pregnancy, commonly including folic acid, iron, iodine, and vitamin D. The appropriate product and dose depend on diet, medical conditions, laboratory results, and pregnancy needs; excessive doses can be harmful.

\medterm{Prenylation} A cell process that attaches a fat-like group to a protein, helping direct where the protein works and how it interacts with other proteins. Abnormal prenylation can contribute to some inherited disorders and cancers.

\medterm{Preop} Preop is an abbreviation for preoperative, meaning relating to preparation and assessment before surgery.
\medterm{Preoperative Assessment} A medical review before surgery that checks illnesses, medicines, allergies, previous anesthesia problems, heart and lung function, and the airway. Tests are ordered when indicated, and correctable risks are addressed before anesthesia and the operation.

\medterm{Preoperative Care} Assessment and preparation before an operation. It may include reviewing illnesses and medicines, checking allergies, performing needed tests, explaining anesthesia and recovery, planning fasting and medicines, and reducing risks such as infection or blood clots.

\medterm{Preoperative Evaluation} Medical assessment before surgery to identify and reduce risks. It reviews health history, medicines, allergies, physical findings, needed tests, anesthesia concerns, and preparations that improve safety during and after surgery.

\medterm{Preoptic Area} The preoptic area is a hypothalamic region involved in temperature regulation, sleep, reproductive behavior, and hormone control.

\medterm{Prepubertal} Prepubertal means before the physical and hormonal changes of puberty have begun or progressed to reproductive maturity.

\medterm{Prepuce} A fold of skin covering and protecting the end of an external sex organ. The penile prepuce is the foreskin over the head of the penis; the clitoral prepuce is the hood over the clitoris.
\medterm{Prerenal Azotemia} A rise in kidney waste products because the kidneys receive too little effective blood flow, without initial damage to the kidney tissue itself. Dehydration, blood loss, severe infection, heart failure, or liver disease may cause it and can lead to kidney injury if untreated.

\medterm{Presbyacusis} Gradual age-related sensorineural hearing loss, often affecting high-pitched sounds first and making speech hard to understand in noise. Hearing aids and communication strategies may help.
\medterm{Presbyacusis (Presbycusis)} Age-related, usually gradual sensorineural hearing loss, particularly affecting the ability to hear high-frequency sounds and understand speech in noise.

\synonyms
Presbycusis

\medterm{Presbycusis} Bilateral, usually symmetric age-related sensorineural hearing loss, initially affecting
high frequencies and speech discrimination in noise. Management includes hearing
rehabilitation, hearing aids, communication strategies, and cochlear implantation in
selected severe cases.

\medterm{Presbyphonia} Age-related changes in the voice caused by thinning or weakening of the vocal folds and changes in the voice box. It may produce a weak, breathy, rough, or quieter voice.

\medterm{Presbyopia} The gradual loss of the eye’s ability to focus on nearby objects with age. It commonly causes difficulty reading small print or doing close work and is corrected with reading glasses, contact lenses, or other prescribed lenses.
\medterm{Preschooler Development} Preschooler development includes rapid gains in language, motor skills, thinking, self-care, play, and social-emotional regulation; milestones vary, but loss of skills warrants assessment.

\medterm{Preschooler Test Or Procedure Preparation} Preparing a preschooler for a test or procedure uses simple explanations, play, comfort objects, and caregiver support while following any fasting or medicine instructions.

\medterm{Prescription} A written or electronic instruction from an authorized healthcare prescriber for a medicine, device, or treatment. It states what is needed and how to use it, allowing a pharmacy or service to provide care safely.
\medterm{Prescription Medications} Pharmaceutical drugs that require authorization from a licensed healthcare provider
before they can be dispensed to a patient. These medications are regulated because they
carry risks that necessitate professional oversight, including potential side effects,
drug interactions, and the need for precise dosing.

\medterm{Presenile Dementia} Presenile dementia is an older term for dementia beginning before the usual older-adult age range; modern diagnosis specifies the underlying disease and age at onset.
\medterm{Presenilin} Presenilin is a component of the gamma-secretase complex that processes membrane proteins; inherited variants can cause familial early-onset Alzheimer disease.

\medterm{Presentation} The part of a fetus that is directed toward the birth canal, or the symptoms and signs with which a patient seeks care.
\medterm{Preseptal Cellulitis} An infection of the eyelid and tissues around the eye, without infection inside the eye socket. It causes redness, swelling, pain, and fever; pain with eye movement, reduced vision, or a bulging eye may indicate deeper orbital infection and needs emergency assessment.

\synonyms
Periorbital Cellulitis

\medterm{Preservation} Maintenance of tissue, cells, organs, or function by preventing deterioration.
\medterm{Preservative} A substance added to food, medicine, or cosmetics to prevent microbial growth or chemical breakdown.
\medterm{Pressor} Raising blood pressure, usually by tightening blood vessels or increasing heart output. Pressor medicines may treat dangerously low blood pressure but require monitoring because they can reduce circulation to some tissues.
\medterm{Pressoreceptor} A pressure-sensitive receptor, especially a baroreceptor in the carotid sinus or aortic arch, that detects arterial stretch and helps regulate heart rate and blood pressure.

\medterm{Pressure} Force applied over a particular area. Medical pressure measurements include blood pressure, pressure inside the skull, eye pressure, and airway pressure, and abnormal values can damage organs or impair function.
\medterm{Pressure Atrophy} Thinning or wasting of tissue caused by prolonged pressure. Pressure can compress small blood vessels, reduce oxygen delivery, and gradually damage the affected tissue; the result depends on its location and duration.


\medterm{Pressure Injury} Localized damage to the skin and/or underlying soft tissue caused by prolonged pressure,
shear, or friction, most commonly over bony prominences.

\medterm{Pressure Sore} Alternate index label; see Bedsores (pressure ulcers). A localized injury to skin and underlying tissue caused by prolonged pressure or shear, also called a pressure ulcer.
\medterm{Pressure Support Ventilation} A ventilator mode in which the machine gives extra pressure with each breath the person starts, making breathing easier and reducing the work of the breathing muscles. It is often used while reducing ventilator support.

\medterm{Pressure Ulcer} An area of skin and deeper tissue damaged by prolonged pressure, sliding, or rubbing, usually over a bony point such as the heel, hip, or tailbone. It may begin as persistent discoloration and progress to an open wound or deep tissue damage.

\medterm{Pressure Ulcer Prevention} Measures that reduce prolonged pressure and tissue injury, including repositioning, support
surfaces, skin inspection, moisture control, nutrition, and mobility.

\medterm{Pressure Urticaria} A form of chronic inducible urticaria in which sustained pressure causes localized
swelling, erythema, and pain, usually after a delay of several hours. Lesions may last 8
to 72 hours and commonly affect the hands, feet, buttocks, or areas compressed by
clothing or sitting.

\medterm{Pressure-Flow Study} A bladder test that measures pressure while the bladder fills and while a person urinates, together with urine flow. It helps identify blockage or weak bladder muscle when difficulty emptying cannot be explained by symptoms alone.

\medterm{Presynaptic Terminals} Presynaptic terminals are the ends of neurons that release neurotransmitters into synapses in response to electrical activity.

\medterm{Presyncope} A feeling that fainting is about to occur without actually losing consciousness. It may cause light-headedness, weakness, sweating, nausea, or blurred vision and can result from low blood pressure, dehydration, abnormal heart rhythm, medicines, or other illness.

\medterm{Preterm Birth} Birth before 37 completed weeks of gestation. It may follow spontaneous labor or rupture of the
membranes, or may be medically indicated because of maternal or fetal complications.

\medterm{Preterm Infant} A baby born before 37 completed weeks of pregnancy. Earlier birth increases the chance of breathing, feeding, temperature-control, infection, and developmental problems; care depends on the baby’s maturity, weight, and health.

\medterm{Preterm Labor} Regular uterine contractions causing cervical change before 37 weeks gestation. Risk
factors include prior preterm delivery, multiple gestation, infection, and cervical
insufficiency. Diagnosis requires regular contractions with documented cervical change.

\medterm{Preterm Premature Rupture Of Membranes} Breaking of the amniotic sac before labor begins and before 37 completed weeks of pregnancy. It may cause a gush or continuing leak of fluid and increases the risks of infection, early birth, and problems from too little amniotic fluid.

\medterm{Pretest Probability} The estimated probability that a patient has a disease before the result of a diagnostic
test is known. It is based on prevalence, symptoms, examination findings, risk factors,
and prior test results and should guide test selection and interpretation.

\medterm{Pretibial Myxedema} Thickened, firm skin over the shins caused by an immune reaction associated with Graves disease. The skin may look swollen, waxy, or dimpled and usually does not leave a lasting indentation when pressed.

\medterm{Pretibial Myxoedema} Thickened, waxy swelling of the skin over the shins, most often associated with Graves disease.
\medterm{Prevalence} The proportion of people in a defined population who have a disease or characteristic at a particular time or during a stated period. It includes both newly diagnosed and previously existing cases.
\medterm{Preventing Falls} Fall prevention combines strength and balance activity, vision and medication review, safe footwear, adequate lighting, handrails, removal of hazards, and appropriate assistive devices.


\medterm{Preventing Food Poisoning} Reduce foodborne illness by washing hands, separating raw and cooked foods, cooking foods thoroughly, refrigerating them promptly, using safe water, and avoiding unpasteurized or contaminated products.

\medterm{Preventing Head Injuries In Children} Children’s head injuries can be reduced with age-appropriate car seats, helmets, supervision, safe play areas, window and stair protection, and prevention of abusive shaking or impact.

\medterm{Preventing Hepatitis A} Hepatitis A prevention relies on vaccination, handwashing, safe food and water, sanitation, and post-exposure advice when close contact with an infected person occurs.

\medterm{Preventing Hepatitis B Or C} Hepatitis B prevention includes vaccination and avoiding blood or sexual exposure; hepatitis C prevention centers on sterile injection equipment, screened blood, and testing and treatment of infection.


\medterm{Preventing Pressure Injuries} Pressure-injury prevention uses regular repositioning, skin inspection, moisture control, nutrition, pressure-redistributing surfaces, and minimizing friction and shear in people with limited mobility.

\medterm{Preventing Stroke} Stroke prevention includes controlling blood pressure, diabetes and cholesterol, avoiding tobacco, treating atrial fibrillation when indicated, exercising, and addressing other vascular risks.

\medterm{Prevention} Actions that lower the chance of disease, detect problems early, prevent complications, or reduce disability. Examples include vaccination, healthy habits, screening, safety measures, controlling chronic conditions, and avoiding harmful exposures.

\medterm{Preventive Dentistry} Dental care that lowers the risk of tooth decay, gum disease, and tooth loss. It includes brushing with fluoride toothpaste, cleaning between teeth, protective sealants when appropriate, healthy eating advice, and regular dental examinations.
\medterm{Preventive Health Care} Preventive health care aims to prevent disease or detect it early through vaccination, screening, healthy behaviors, and routine medical care.

\medterm{Preventive Medicine} The medical field focused on preventing illness and promoting health through individual care, screening, vaccination, and community programs.
\medterm{Prevertebral Air} Air seen in the soft tissues in front of the spine on a neck X-ray. It can indicate injury or a hole in the swallowing or breathing passages, a deep infection, or recent surgery and may require urgent evaluation.

\medterm{Priapism} Prolonged, painful erection lasting longer than 4 hours, occurring independently of sexual desire or stimulation. Medical emergency categorized into two types: ischemic (low-flow, veno-occlusive) and non-ischemic (high-flow, arterial).
\synonyms
Prolonged erection; persistent erection
Low-Flow Priapism
Nonischemic Priapism

\medterm{Priapism (Penis Disorder)} Alternate index label for Penis Disorder.

\medterm{Prick Test} A skin test in which small amounts of suspected allergens are introduced into the skin.
\synonyms
Skin-prick test

\medterm{Prickly Heat} Miliaria, an itchy rash caused by blockage of sweat ducts.
\medterm{Primary} First in time, origin, importance, or site; in disease classification it often means arising independently rather than secondary to another cause.
\medterm{Primary Adrenal Insufficiency} Alternate index label; see \textit{Addison's disease}.

\medterm{Primary Adrenal Insufficiency (Addison's Disease)} Destruction or dysfunction of the adrenal cortex causing deficiency of cortisol and
aldosterone, presenting with fatigue, hyperpigmentation, weight loss, and salt craving.
Lifelong hormone replacement is required.

\synonyms
Addison's disease

\medterm{Primary Adrenal Lymphoma} A rare lymphoma arising primarily in the adrenal gland, most commonly diffuse large
B-cell lymphoma. It often presents bilaterally and may cause adrenal insufficiency.

\medterm{Primary Aldosteronism} Excess production of aldosterone by the adrenal glands, usually from an adrenal nodule or enlargement of both glands.
It increases sodium retention and may cause hypertension, low potassium, or suppressed
renin.

\medterm{Primary Alveolar Hypoventilation} A rare disorder of inadequate automatic breathing, often related to PHOX2B or other autonomic-control abnormalities, causing elevated carbon dioxide especially during sleep.

\medterm{Primary Amebic Meningoencephalitis} A fulminant, rapidly fatal central nervous system infection caused by the free-living thermophilic amoeba Naegleria fowleri, contracted when warm contaminated freshwater is inhaled into the nasal cavity during swimming, diving, or nasal irrigation. The amoebae penetrate the olfactory neuroepithelium and cribriform plate to reach the brain, causing acute hemorrhagic necrotizing meningoencephalitis characterized by severe frontal headache, high fever, meningismus, altered olfaction, seizures, and rapid coma, carrying a mortality rate exceeding 97\%.

\synonyms
PAM; Naegleriasis; Primary Amoebic Meningoencephalitis; Fatal Amebic Meningoencephalitis

\medterm{Primary Amyloidosis} Primary amyloidosis, usually AL amyloidosis, results from abnormal plasma-cell proteins depositing as amyloid in organs such as the heart, kidneys, nerves, or gastrointestinal tract.

\medterm{Primary And Secondary Hyperaldosteronism} Excess aldosterone production, from an adrenal source in primary disease or from increased renin stimulation in secondary disease, causing hypertension, sodium retention, and sometimes low potassium.

\medterm{Primary Angioplasty} Angioplasty used as the initial treatment to reopen an artery during an acute heart attack.
\medterm{Primary Biliary Cholangitis} A chronic autoimmune liver disease in which small intrahepatic bile ducts are
progressively damaged. It can cause cholestasis, itching, fatigue, and eventually cirrhosis; antimitochondrial antibodies
often support the diagnosis.

\medterm{Primary Biliary Cirrhosis} Alternate index label; see \textit{Primary biliary cholangitis}.

\medterm{Primary Biliary Cirrhosis Treatment} Alternate index label for PBC.

\medterm{Primary Care} The usual first place people go for healthcare. It includes prevention, check-ups, treatment of common illnesses, care of long-term conditions, basic mental-health support, and referral or coordination with specialists when needed.

\medterm{Primary Ciliary Dyskinesia} An inherited disorder in which tiny moving hairs lining the airways do not clear mucus properly. It can cause repeated sinus and lung infections, permanent airway widening, hearing problems, infertility, and organs arranged differently from usual.

\medterm{Primary Clear Cell Adenocarcinoma Of The Vulva} A rare cancer of the vulva made up of clear-looking cells. Diagnosis requires examination of a tissue sample and tests that distinguish it from cancer that has spread from the kidney or another organ.

\medterm{Primary Cutaneous B-Cell Lymphoma} A type of lymphoma that starts in the skin and has not spread elsewhere when diagnosed. It may cause one or more painless lumps or plaques; behavior and treatment depend on the subtype.

\medterm{Primary Endpoint} The main outcome chosen in advance to judge whether a treatment or other intervention works in a clinical trial. It may be a symptom, laboratory result, complication, quality-of-life measure, or survival.

\medterm{Primary Hyperaldosteronism} Excess production of the hormone aldosterone by the adrenal glands, usually from a small gland growth or enlargement of both glands. It can cause difficult-to-control high blood pressure, low potassium, muscle weakness, or abnormal heart rhythms.

\medterm{Primary Hyperoxaluria} A group of inherited disorders in which the liver makes too much oxalate. This can cause repeated kidney stones, calcium deposits, kidney damage, and eventually failure of the kidneys.

\medterm{Primary Hyperparathyroidism} Overproduction of parathyroid hormone by one or more abnormal parathyroid glands, usually from a noncancerous gland growth. It raises blood calcium and may cause kidney stones, bone loss, abdominal symptoms, thirst, tiredness, or no symptoms at all.

\medterm{Primary Hypersomnia} Excessive daytime sleepiness that is not explained by too little sleep, another sleep disorder, a medicine, substance use, or another illness. It can cause unintended sleep, long sleep, and difficulty becoming fully alert after waking.

\medterm{Primary Hypertension} Persistently high blood pressure without one identifiable underlying disease. It results from a combination of genetic, environmental, vascular, and lifestyle factors and can often be controlled with lifestyle changes and medicines.

\medterm{Primary Hypertrophic Osteoarthropathy} A rare inherited condition causing enlarged, rounded fingertips, thickened facial or scalp skin, and new bone growth along the limbs. It may cause joint pain, stiffness, swelling, and excessive sweating.

\medterm{Primary Hypoparathyroidism} Inadequate parathyroid hormone production causing low calcium and high phosphate, with possible tingling, cramps, tetany, seizures, cataracts, or chronic bone and dental effects.

\medterm{Primary Hypothyroidism} Reduced thyroid-gland production of thyroxine and triiodothyronine, usually causing high TSH and symptoms such as fatigue, cold intolerance, constipation, dry skin, and slowed heart rate.

\medterm{Primary Immune Response} The body’s first targeted immune response to a new germ or other foreign substance. It develops over days, produces antibodies and immune cells, and creates memory that allows a faster response next time.

\medterm{Primary Immunodeficiency} A group of usually inherited disorders in which part of the immune system is missing or does not work properly. They can cause repeated or unusually severe infections and may also cause autoimmune disease or inflammation.

\medterm{Primary Immunodeficiency Disease} An inherited or genetic disorder of immune development or function that causes unusual, recurrent, severe, or persistent infections and sometimes autoimmunity.

\medterm{Primary Immunodeficiency Disease PIDD} Primary immunodeficiency diseases are inherited disorders of immune development or function that cause unusual, severe, recurrent, or persistent infections and sometimes autoimmunity or inflammation. Evaluation may include blood counts, immunoglobulins, functional tests, and genetic testing.

\synonyms
Deficiency Primary Immunodeficiency (Primary Immunodeficiency Disease PIDD)

\medterm{Primary Lactase Deficiency} Reduced lactase enzyme activity developing with age, causing lactose malabsorption and possible bloating, abdominal pain, gas, or diarrhea after dairy consumption.

\medterm{Primary Lateral Sclerosis} A rare progressive nerve disorder causing increasing muscle stiffness and weakness, usually beginning in the legs. Speech, swallowing, walking, and emotional control may become affected over time.

\medterm{Primary Lymphoma Of The Brain} Primary central-nervous-system lymphoma is a lymphoma arising in the brain, spinal cord, eyes, or surrounding fluid without systemic disease at diagnosis.

\medterm{Primary Metabolism} The essential chemical activity inside cells that produces energy, builds needed substances, breaks down waste, and keeps cells alive and functioning.
\medterm{Primary Microrna} The first RNA copy made from a microRNA gene. It is processed inside and outside the cell nucleus into a shorter molecule that helps control which proteins cells make.

\medterm{Primary Myelofibrosis} A long-term bone-marrow disorder in which scar tissue replaces healthy marrow. Blood-cell production may move to the spleen or liver, causing anemia, tiredness, easy bleeding or infection, bone pain, and an enlarged spleen.

\synonyms
Agnogenic Myeloid Metaplasia

\medterm{Primary Open-Angle Glaucoma} A slowly progressive disease that damages the optic nerve, usually without pain or early symptoms. Vision is gradually lost from the sides inward; lowering pressure inside the eye can slow or prevent further damage.

\medterm{Primary Osteoporosis} Weakening and thinning of bone caused mainly by aging, loss of estrogen after menopause, or age-related changes, rather than by another disease or medicine. It increases the risk of fractures, especially in the hip, spine, and wrist.
\medterm{Primary Ovarian Insufficiency} Loss of normal ovarian activity before age 40, characterized by irregular or absent
menses, reduced estrogen, and intermittently or persistently elevated gonadotropins.
Causes include genetic, autoimmune, iatrogenic, toxic, infectious, and unknown factors.

\medterm{Primary Polycythemia} Alternate index label; see \textit{Polycythemia vera}.

\medterm{Primary Prevention} Actions taken before a disease begins to prevent it, such as vaccination, avoiding tobacco, reducing harmful exposures, using sun protection, and controlling health risks. It differs from screening, which looks for disease that may already be present.
\medterm{Primary Progressive Multiple Sclerosis} A form of multiple sclerosis that gradually worsens from the beginning, without clear attacks and recoveries. It commonly affects walking, strength, balance, bladder control, and thinking; medicines and rehabilitation may slow disability.

\synonyms
PPMS; Primary-progressive MS

\medterm{Primary Progressive Aphasia} A neurodegenerative disorder causing gradual worsening of language abilities, including speaking, naming, comprehension, reading, or writing.

\medterm{Primary Sclerosing Cholangitis} A long-term disease in which inflammation and scarring narrow the bile tubes inside and outside the liver. It can cause itching, tiredness, jaundice, repeated infections, and liver damage and is often associated with inflammatory bowel disease.

\medterm{Primary Tuberculosis} The first infection with the bacterium that causes tuberculosis. It often begins in the lungs and nearby lymph nodes; the immune system may contain it without symptoms, or it may progress to active disease.

\medterm{Primary Tumor} The original tumor, or site where a cancer began, from which cancer cells may spread to form secondary or metastatic tumors. A cancer may be described as having an unknown primary when the original site cannot be identified.


\medterm{Primidone} Primidone is an antiseizure medicine used for epilepsy and sometimes essential tremor; it can cause sleepiness, dizziness, and drug interactions.

\medterm{Primigravida} A person who is pregnant for the first time. First pregnancies may have a longer labor and require monitoring for conditions such as high blood pressure, but the term itself does not indicate a problem.

\medterm{Primipara} A person who has given birth once after the fetus could survive outside the womb, regardless of outcome.
\medterm{Primitive} Relating to an early developmental stage or a basic, less specialized form. In medicine, primitive cells may be immature and capable of developing into several more specialized cell types.
\medterm{Primitive Node} The primitive node is an embryonic organizer at the cranial end of the primitive streak that helps establish body axes and germ-layer formation.

\medterm{Primitive Streak} The primitive streak is an embryonic structure through which cells migrate during gastrulation to form the three germ layers and establish body organization.

\medterm{Primordial} Present at the earliest stage of development or existing in an original, basic form. Examples include primordial germ cells and primordial follicles.
\medterm{Primordial Prevention} Preventing the emergence of disease risk factors by addressing social, economic,
behavioural, and environmental conditions early. Examples: promoting healthy diets and
physical activity in children, reducing tobacco exposure, and creating safe environments
for exercise.


\medterm{Prion} An infectious protein folded into an abnormal shape that can make other proteins fold incorrectly. Prions cause rare, serious brain diseases that worsen over time; they are not ordinary bacteria or viruses and currently have no cure.
\medterm{Prion Disease} A group of rare, rapidly progressive, and fatal brain disorders caused by prions---infectious agents made entirely of misfolded protein that contain no genetic material (no DNA or RNA). Normally, healthy prion proteins exist in brain cells, but when a protein misfolds into an abnormal shape, it triggers a chain reaction that forces neighboring healthy proteins to misfold into the same defective shape. These clumped proteins destroy brain tissue and leave microscopic holes, giving the brain a sponge-like appearance (transmissible spongiform encephalopathies). Prion diseases develop in three ways: sporadic (occurring spontaneously in late adulthood without a known trigger, accounting for approximately 85\% of cases such as sporadic Creutzfeldt--Jakob disease), genetic or inherited (caused by mutations in the \textit{PRNP} gene passed through families, such as familial Creutzfeldt--Jakob disease, fatal familial insomnia, and Gerstmann--Sträussler--Scheinker disease), or acquired (contracted through exposure to infected tissue, such as eating beef from cattle with bovine spongiform encephalopathy or ``mad cow disease'', contaminated surgical instruments, or ritual cannibalism as in kuru). Although the incubation period can last years or decades without symptoms, once illness begins it causes rapid memory loss, progressive dementia, severe unsteadiness and balance loss (ataxia), sudden muscle jerks (myoclonus), and loss of speech and movement. Prions are extraordinarily resistant to standard heat and chemical sterilization. There is currently no cure, vaccine, or disease-halting treatment, and the conditions are universally fatal, with medical care centered on comfort and palliative support.

\synonyms
Prion Diseases; Transmissible Spongiform Encephalopathies; TSE; Prion Disorders

\medterm{Prism} A clear optical device that bends light. In eye care, prisms can measure or correct eye misalignment, reduce double vision, and be built into glasses or used during examinations.
\medterm{Privacy} Privacy is a person’s right and expectation that personal information, communications, and bodily matters are protected from unauthorized access or disclosure.

\medterm{Privacy Breach} A privacy breach is unauthorized access, use, disclosure, loss, or exposure of personal or health information.

\medterm{Private Mutation} A genetic variant confined to a family, individual, tumor, or small population rather than commonly seen in the general population; its significance requires genetic and clinical interpretation.

\medterm{PRK} Photorefractive keratectomy is laser eye surgery that reshapes the cornea to reduce short-sightedness, long-sightedness, or astigmatism. The surface layer is removed and grows back during healing; pain, haze, infection, or imperfect vision can occur.

\medterm{Pro Re Nata} A medication or treatment instruction meaning that it should be used only when a specified symptom or clinical circumstance occurs, within clearly prescribed dose and frequency limits.

\synonyms
P.R.N.; As Needed


\medterm{Probe} A blunt instrument used to examine an opening, cavity, wound, tooth, or narrow passage. Different probes can measure depth, check whether a passage is open, guide a procedure, or trace a sinus or fistula.

\medterm{Probenecid} A uricosuric medicine used to prevent gout attacks in selected patients and to alter renal excretion of some drugs.
\medterm{Probiotic} A live microorganism that may provide a health benefit when taken in sufficient amounts. Benefits depend on the specific strain and condition; a probiotic is not automatically effective or safe for every person.

\medterm{Probiotics} Probiotics are live microorganisms that may provide a health benefit when consumed in adequate amounts, although effects are strain-specific and vary by condition.

\medterm{Problem} A symptom, sign, diagnosis, or issue requiring assessment or management.
\medterm{Problems Sleeping During Pregnancy} Sleep problems during pregnancy may result from discomfort, reflux, urination, anxiety, hormonal change, or sleep apnea; safe sleep habits and clinician-guided treatment are important.

\medterm{Probucol} Probucol is a lipid-lowering drug that reduces cholesterol but may prolong the QT interval and is not widely used in current practice.

\medterm{PROC} PROC is the gene symbol for protein C, an anticoagulant protein that limits clot formation; inherited deficiency increases thrombosis risk.

\medterm{Procainamide} An antiarrhythmic medicine that can slow abnormal electrical conduction in the heart. It is used in selected rhythm disorders and requires monitoring because it can affect blood pressure, the heart rhythm, and immune or blood systems.
\medterm{Procaine} An ester local anesthetic that temporarily blocks nerve signals and reduces pain. It is now used mainly in limited or historical settings because newer anesthetics are commonly preferred and allergic reactions can occur.
\medterm{Procalcitonin} A biomarker that may rise during some bacterial infections and systemic inflammation. It can support
assessment of infection and antibiotic decisions, but it cannot by itself diagnose or
exclude bacterial infection or determine sepsis severity.

\medterm{Procarbazine} A chemotherapy medicine used in some lymphomas and brain tumors. It damages the DNA of rapidly growing cells and can cause nausea, low blood counts, infertility, or other side effects; alcohol and many medicines may interact with it.
\medterm{Procarbazine Hydrochloride} A salt form of procarbazine, a chemotherapy medicine that damages cancer cells. It is used for selected lymphomas and brain tumors and can cause low blood counts, nausea, infertility, or serious medicine and food interactions.

\medterm{Procedural Memory} Memory for how to perform learned skills, such as riding a bicycle, typing, or playing an instrument. It usually works without consciously recalling each step and differs from memory for facts or events.

\medterm{Procedural Sedation} Use of medicines to reduce anxiety, discomfort, and awareness during a medical procedure while maintaining breathing and protective reflexes. Continuous monitoring is required because sedation can unexpectedly become deeper and slow breathing.

\medterm{Procedure} A planned medical or surgical action performed to diagnose, treat, prevent, monitor, or relieve a health problem. It may involve an examination, sample collection, medicine delivery, device placement, or operation, with risks depending on the method and person.

\synonyms
Medical Procedure; Clinical Procedure; Intervention

\medterm{Procentriole} An immature structure that forms next to an existing centriole as a cell prepares to divide. It later matures into a centriole, which helps organize the fibers that separate chromosomes.

\medterm{Procerus Muscle} A small facial muscle over the bridge of the nose that draws the medial eyebrows downward and produces horizontal wrinkles between the eyebrows.

\medterm{Process} A process is an ordered series of biological, physical, chemical, or operational changes that produces a result.

\medterm{Processing Bodies} Processing bodies are cytoplasmic RNA-protein granules where messenger RNAs may be stored, decapped, degraded, or regulated.

\medterm{Prochlorperazine} A phenothiazine medicine used mainly for severe nausea or vomiting and, in selected cases, psychotic disorders. It can cause drowsiness, low blood pressure, and movement disorders.
\medterm{Prochlorperazine Overdose} Taking too much prochlorperazine can cause severe sleepiness, confusion, low blood pressure, abnormal heart rhythms, seizures, or muscle stiffness with fever. It requires immediate emergency or poison-control advice.

\medterm{Procidentia} The most severe form of uterine prolapse, in which the uterus and cervix descend through and outside the vaginal opening. It may cause pressure, bleeding, sores, difficulty urinating, or bowel problems and requires gynecologic assessment.

\synonyms
Complete Uterovaginal Prolapse; Third-Degree Uterine Prolapse


\medterm{Procollagen} An unfinished form of collagen made inside cells. After chemical changes and removal of its end sections outside the cell, it joins with other molecules to form strong collagen fibers in skin, bone, tendons, and other tissues.

\medterm{Proctalgia} Pain in or around the rectum or anal canal. Causes include fissures, hemorrhoids, abscesses, inflammation, injury, and functional muscle pain; sudden severe pain or pain with fever, bleeding, or a lump needs medical assessment.
\medterm{Proctitis} Inflammation of the rectum. It can cause a constant urge to pass stool, rectal pain, bleeding, mucus or pus, and pelvic discomfort. Causes include inflammatory bowel disease, sexually transmitted infections, radiation treatment, and diversion of stool through an opening in the abdomen. Diagnosis may require an examination, stool or swab tests, and a small tissue sample. Treatment depends on the cause.

\synonyms
Rectal Inflammation

\medterm{Proctocolitis} Inflammation of the rectum and nearby lower colon. It may cause bleeding, mucus, diarrhea, urgency, pain, or difficulty passing stool and can result from infection, inflammatory bowel disease, reduced blood flow, radiation, or other causes.

\medterm{Proctology} The part of medicine concerned with diseases of the anus and rectum. Care may include assessment and treatment of hemorrhoids, fissures, fistulas, abscesses, bowel-control problems, infections, inflammation, and cancers.
\medterm{Proctoscope} An instrument used to inspect the anal canal and rectum.
\medterm{Proctoscopy} Examination of the anal canal and rectum with a short viewing instrument. It can help investigate bleeding, pain, inflammation, ulcers, polyps, or a lump and may allow a small tissue sample to be taken.
\medterm{Proctostomy} An operation that creates an opening from the rectum through the abdominal wall or nearby skin so stool can leave the body. It may be temporary or permanent and requires care of the opening and surrounding skin.

\medterm{Prodromal} Describing early, nonspecific signs or symptoms that occur before the characteristic clinical features of a disease become apparent.

\medterm{Prodrome} An early symptom or group of symptoms that appears before a disease episode or major event. Examples include tiredness before migraine, tingling before a seizure, or fever before an infection becomes obvious.
\medterm{Prodrug} A medicine given in a form that the body changes into its active form. This design may improve absorption, delivery to a target tissue, duration of action, or tolerability.

\medterm{Productive Cough} A cough that brings up mucus or phlegm from the airways. It can occur with infections, asthma, chronic lung disease, or other conditions and does not by itself identify the cause.
\medterm{Profunda} A deep artery or vessel, commonly referring to the profunda femoris artery in the thigh. This artery supplies muscles and nearby tissues and can be important in bleeding, circulation, or vascular procedures.
\medterm{Profunda Miliaria (Heat Rash)} Alternate index label for Heat Rash.

\medterm{Progenitor} A partly developed cell that can multiply and mature into one or a limited range of specialized cell types. Progenitor cells help replace or repair tissues but have less developmental potential than stem cells.
\medterm{Progeny} Biological offspring or descendants. In genetics, the term usually means children or later descendants whose traits may be studied to understand how characteristics or disorders are inherited. In cell biology, it can mean daughter cells produced by one parent cell.
\medterm{Progeria} A rare genetic disorder in which features resembling accelerated aging begin during childhood. The usual form is Hutchinson--Gilford progeria syndrome, caused by a disease-causing variant in the \textit{LMNA} gene that produces abnormal lamin A. Most cases result from a new variant rather than inheritance from an affected parent. Children are often healthy-looking at birth and later develop slowed growth, loss of body fat, hair loss, thin aged-looking skin, joint stiffness, and characteristic facial features. Accelerated atherosclerosis can cause early cardiovascular disease, while intellectual development is usually preserved.

\synonyms
Hutchinson--Gilford Progeria Syndrome
Hutchinson--Gilford Syndrome
HGPS
\medterm{Progestational} Relating to progesterone-like effects or preparation of the uterus for pregnancy. Progestational medicines can alter menstrual bleeding, prevent ovulation, support the uterine lining, or protect it from excessive estrogen stimulation.
\medterm{Progesterone} A steroid hormone produced mainly by the corpus luteum and placenta that supports the endometrium and pregnancy.
\medterm{Progesterone Test} A blood test that measures the level of progesterone, a hormone involved in ovulation and pregnancy. It may help assess ovulation, infertility, abnormal uterine bleeding, fertility treatment, or certain pregnancy problems. Results depend on the timing of the menstrual cycle and whether the person is pregnant.

\synonyms
Serum Progesterone
Progesterone Blood Test

\medterm{Progesterone Antagonist} A progesterone antagonist blocks progesterone receptors and may be used for selected reproductive, gynecologic, or other medical purposes.

\medterm{Progesterone Supplementation} Use of progesterone during pregnancy in selected people at risk of preterm birth, especially when ultrasound shows a short cervix. Vaginal progesterone may be considered in that situation; routine use after a previous preterm birth without a short cervix is not recommended, and 17-hydroxyprogesterone caproate is no longer FDA-approved for this purpose.

\medterm{Progestin} A progestin is a synthetic progesterone-like hormone used in contraception, hormone therapy, and treatment of selected gynecologic conditions.

\medterm{Progestogen} A natural or synthetic hormone with effects like progesterone. Progestogens help regulate the menstrual cycle and pregnancy and are used in some contraceptives and hormone treatments.
\medterm{Proglottid} One segment of a tapeworm. Each mature segment contains reproductive organs and may release eggs in stool.
\medterm{Prognathism} Forward projection of the upper or lower jaw relative to the facial skeleton, affecting bite, appearance, chewing, speech, and sometimes airway function.

\medterm{Prognosis} The expected course and outcome of a disease or condition.
\medterm{Prognostic Factor} A prognostic factor is a clinical, pathologic, or biologic feature associated with the likely course, outcome, or survival of a disease.

\medterm{Prognostic Test} A test that provides information about a disease's likely course or outcome. Results may estimate risks, guide treatment choices, or monitor response, but they cannot predict an individual future with certainty.

\medterm{Progression} Worsening or advancement of a disease, such as increasing tumor size, new lesions, or spread to another site.

\medterm{Progression-Free Survival} The length of time after a defined starting point, such as diagnosis or treatment, during which a person remains alive without the cancer worsening. It is a group measure used in research and cannot predict exactly what will happen to one person.

\medterm{Progressive} Increasing, worsening, or spreading over time, or moving forward in a particular direction. A progressive disease may cause gradually increasing symptoms, but the word alone does not explain the cause, speed, or expected outcome.
\medterm{Progressive Disease} A disease that is worsening, spreading, or becoming more active over time. In cancer, progression may mean that tumors have grown, new tumors have appeared, or existing disease has clearly worsened on examination or imaging.

\medterm{Progressive Lenticular Degeneration} An older name for Wilson disease, an inherited disorder causing copper buildup. Excess copper can damage the liver, brain, eyes, kidneys, and other organs unless treatment lowers copper levels.
\medterm{Progressive Multifocal Leukoencephalopathy} A serious brain infection caused by reactivation of JC virus in people with severely weakened immune systems. It damages the brain’s nerve insulation and may cause worsening weakness, clumsiness, vision or speech problems, confusion, or seizures.

\medterm{Progressive Muscular Atrophy} A motor-neuron disorder predominantly affecting lower motor neurons and causing progressive muscle wasting and weakness.
\medterm{Progressive Supranuclear Palsy} A progressive brain disorder that causes balance problems, frequent falls, stiffness, slowed movement, speech or swallowing difficulty, and trouble moving the eyes, especially looking down. It resembles Parkinson disease but has different features and treatment needs.

\synonyms
Steele-Richardson-Olszewski Syndrome
Supranuclear Palsy

\medterm{Proinflammatory} Describing a substance, cell, or immune response that starts or increases inflammation. Such activity can help fight infection or repair injury, but excessive or prolonged inflammation can damage healthy tissue.

\medterm{Proinsulin} An unfinished hormone made by pancreatic beta cells before it is split into insulin and C-peptide. Measuring proinsulin can provide information about insulin production and certain disorders of glucose regulation.

\medterm{Projectile Vomiting} Forceful vomiting expelled with unusual intensity, which may occur with gastric outlet obstruction, raised intracranial pressure, infection, migraine, or other causes; severity and associated signs guide evaluation.

\medterm{Projection} A view or image made by directing X-rays or other imaging energy through the body toward a detector. Different projections show body structures from different angles and can help locate an abnormality.
\medterm{Prokaryote} A prokaryote is a cell or organism without a membrane-bound nucleus, including bacteria and archaea.
\medterm{Prokinetic} A medicine that promotes gastrointestinal motility and movement of contents through the digestive tract; use depends on the underlying disorder and safety profile.
\medterm{Prolactin} A hormone released mainly by the pituitary gland that supports breast development and milk production. Abnormal levels can affect periods, fertility, sexual function, pregnancy, breastfeeding, and sometimes vision because of pituitary growths.
\medterm{Prolactin Blood Test} A prolactin blood test measures the hormone prolactin and helps evaluate abnormal milk production, menstrual or fertility problems, and some pituitary disorders.

\medterm{Prolactin Physiology} Prolactin is a hormone made by the pituitary gland that supports breast development and milk production. Levels normally rise during pregnancy, breastfeeding, sleep, stress, and breast stimulation.

\medterm{Prolactin Receptor} A cell-surface receptor activated by prolactin to regulate mammary development, milk production, reproduction, immune effects, and other tissue responses.

\medterm{Prolactinoma} A usually benign pituitary tumor that produces too much prolactin. It may cause milk production, menstrual or fertility problems, low sex hormones, headache, or vision loss and is often treated with dopamine-agonist medicines.

\synonyms
Pituitary Tumor (Prolactinoma)

\medterm{Prolactinoma Treatment} Treatment may include a dopamine agonist, surgery, radiation, or observation,
depending on tumor size, symptoms, hormone levels, treatment response, and patient preferences. Dopamine agonists often
lower prolactin levels and reduce tumor size.

\medterm{Prolapse} Downward displacement or protrusion of an organ or tissue from its normal position, often through a supporting structure or body opening.
\medterm{Prolapse Of Orbital Fat} Bulging of the eye’s normal cushioning fat through a weak spot in the thin tissue around the eye. It usually appears as a soft, painless lump in an upper or lower eyelid, often with aging, but other lumps may look similar.

\medterm{Prolapse Of The Uterus (Uterine Prolapse)} Descent of the uterus into the vaginal canal due to weakening of pelvic floor supports,
most commonly the uterosacral and cardinal ligaments.

\synonyms
Uterine Prolapse

\medterm{Prolapse, Mitral Valve} Alternate index label; see \textit{Mitral valve prolapse}.

\medterm{Prolapsed Bladder} Alternate index label; see \textit{Anterior vaginal prolapse (cystocele)}.

\medterm{Prolapsed Cord} A prolapsed umbilical cord descends beside or ahead of the fetus after membrane rupture, risking cord compression and reduced fetal oxygen; it is an obstetric emergency.

\medterm{Prolapsed Disk} A spinal disk that bulges or ruptures and may compress a nearby nerve; also called a herniated disk.

\synonyms
a herniated disk.

\medterm{Prolapsed Uterus} Alternate index label; see \textit{Uterine prolapse}.

\medterm{Proliferating Trichilemmal Tumor} A usually low-grade follicular tumor arising from a trichilemmal cyst, most often on the scalp; large or recurrent lesions may show locally aggressive or malignant behavior.

\medterm{Proliferation} Increase in the number of cells or organisms through growth and division. It may be normal, part of healing, or abnormal as in uncontrolled tumor growth.
\medterm{Proliferative} Characterized by increased cell multiplication or new tissue growth. The meaning depends on context: some proliferative changes are normal or healing, while others are precancerous or cancerous.
\medterm{Proliferative Diabetic Retinopathy} An advanced form of diabetic eye disease in which poor blood supply causes fragile new blood vessels to grow on the retina or nearby structures. They can bleed, scar, detach the retina, and threaten vision; specialist treatment can reduce these risks.

\medterm{Proliferative Fasciitis} A benign, reactive fibroblastic proliferation resembling nodular fasciitis but containing distinctive ganglion-like cells; it usually grows quickly and resolves or stabilizes after local treatment.

\medterm{Proliferative Myositis} A benign reactive proliferation within skeletal muscle that may present as a rapidly enlarging mass and mimic sarcoma. Imaging and biopsy are used to establish the diagnosis.

\medterm{Proliferative Retinopathy} Advanced retinal disease in which abnormal new blood vessels grow and may bleed or impair vision.

\medterm{Proline} An amino acid used to build proteins. Its ring-shaped side chain helps give collagen and other proteins their structure.
\medterm{Prolonged Deceleration} A temporary slowing of the fetal heart rate lasting at least two minutes but less than ten minutes. It may result from reduced blood flow, excessive contractions, cord problems, or placental separation and requires prompt assessment during labor.

\medterm{Prolonged Second Stage} The pushing stage of labor lasting longer than expected. The safe duration depends on whether it is a first birth, use of epidural anesthesia, progress, and the health of the parent and fetus; the team reassesses before deciding on further care.

\medterm{Prolymphocyte Count} The number or proportion of prolymphocytes, immature-appearing lymphoid cells, in blood or marrow; an increased count may occur in selected lymphoid leukemias and requires morphology and immunophenotyping.

\medterm{Prolymphocytic Leukemia} Prolymphocytic leukemia is a rare aggressive leukemia with large abnormal lymphocytes in the blood, bone marrow, spleen, or lymph nodes.

\medterm{Promazine} A first-generation phenothiazine antipsychotic with sedative properties; it is now used infrequently in many settings and availability varies by country. It can cause drowsiness, orthostatic hypotension, anticholinergic effects, extrapyramidal symptoms, and other serious antipsychotic adverse effects.
\medterm{Promecarb} Promecarb is a carbamate insecticide that inhibits acetylcholinesterase and can cause cholinergic poisoning after substantial exposure.


\medterm{Prometaphase} Prometaphase is the stage of mitosis after nuclear-envelope breakdown when spindle microtubules attach to chromosome kinetochores.

\medterm{Promethazine} Promethazine is an antihistamine used for allergy symptoms, nausea, motion sickness, and sometimes sedation; it can cause marked sleepiness and breathing problems in young children.

\medterm{Promethazine Overdose} Too much promethazine can cause extreme drowsiness, confusion, agitation, fast heartbeat, seizures, breathing problems, or coma. Children are particularly vulnerable; significant or intentional overdoses need immediate emergency advice.

\medterm{Prominent Iris Collarette} A visible thickening or ridge of the iris at the junction of the pupillary and ciliary
zones, representing the site where the embryonic pupillary membrane was attached. It is
typically a normal anatomical variant but can be a clinical sign in conditions such as
Axenfeld-Rieger syndrome or iridocorneal endothelial syndrome.

\medterm{Promonocyte Count} The number or proportion of promonocytes, immature monocyte precursors, in blood or marrow; increased cells may support evaluation of acute or chronic myeloid leukemia.

\medterm{Promontory} A part of the body that projects outward like a small ridge or bump. Examples include the sacral promontory in the pelvis and the cochlear promontory in the inner ear; the exact meaning depends on the structure named.
\medterm{Promoter} A region of DNA near a gene where proteins attach to start copying the gene into RNA. Promoters help determine when a gene is active and how much of its protein is made.

\medterm{Promoters} DNA control regions that help start gene activity when regulatory proteins attach. They influence how much messenger RNA and protein a gene makes and are important in development, normal function, and disease.
\medterm{Promyelocytes} Immature cells in the bone-marrow pathway that normally develops into several types of white blood cells. Large numbers in blood or marrow may occur in leukemia and need examination with other blood and cell tests.

\medterm{Pronation} Turning the forearm so the palm faces down when the elbow is bent, or backward when the arm hangs at the side. During this movement, the radius rotates across the ulna.
\medterm{Pronator} A muscle that turns a limb or part of it into a pronated position. In the forearm, pronator muscles rotate the palm downward and help stabilize the elbow during movement.
\medterm{Prone} Prone means lying face downward with the front of the body against a surface; positioning affects breathing, pressure injury, and procedures.

\medterm{Prone Position} A position in which the patient lies face down, commonly providing access to the back, spine, buttocks, or posterior legs.

\medterm{Prone Positioning} Placing a person face down to improve breathing, especially during severe acute respiratory distress syndrome. It can improve oxygen levels by helping more of the lung take part in breathing, but requires careful monitoring of the airway, skin, eyes, and pressure points.

\medterm{Pronephric} Relating to the earliest temporary kidney system formed in an embryo. In humans it does not provide effective kidney function and is replaced during development by later kidney structures.

\medterm{Pronephros} The first temporary kidney structure formed in a human embryo. It is not functional as a kidney and largely disappears as the permanent kidneys develop.
\medterm{Propafenone} A heart-rhythm medicine that slows electrical signals through the heart. It is used for selected irregular rhythms, including episodes of atrial fibrillation, but can worsen or cause dangerous rhythms and is unsuitable for some people with structural heart disease.

\medterm{Propane} Propane is a flammable hydrocarbon gas used as fuel; inhalation or displacement of oxygen can cause suffocation, and liquid contact can cause cold burns.

\medterm{Propane Poisoning} Breathing concentrated propane can displace oxygen and cause dizziness, headache, confusion, unconsciousness, abnormal heart rhythms, or suffocation. Leave the area if safe, avoid flames, and call emergency services for symptoms or significant exposure.

\medterm{Propanol} Propanol is an alcohol solvent with two structural isomers; ingestion or substantial exposure can depress the nervous system and irritate tissues.

\medterm{Propanolol} A nonselective beta-adrenergic blocker used for hypertension, arrhythmias, angina, tremor, migraine prevention, and other indications.
\medterm{Propantheline} An older medicine that reduces certain nerve signals to relax the gut and decrease sweating. It can cause dry mouth, blurred vision, constipation, difficulty urinating, fast heartbeat, or confusion and is rarely used today.
\medterm{Propellant} A propellant is a gas or volatile substance that expels a product from a container or delivers medicine in an inhaler.

\medterm{Properdin} A blood protein that strengthens one part of the body’s early immune defense. It helps mark and damage certain germs; too little properdin can increase susceptibility to serious bacterial infections.

\medterm{Properdin Deficiency} A rare inherited immune-system disorder in which a complement protein is missing or low. It increases the risk of severe infections, especially meningococcal disease, and is evaluated with immune tests and a history of recurrent infections.

\medterm{Propetamphos} Propetamphos is an organophosphate insecticide that can inhibit acetylcholinesterase and cause cholinergic toxicity after exposure.
\medterm{Prophage} A prophage is bacteriophage genetic material integrated into a bacterial chromosome or maintained in a dormant state inside a bacterium.

\medterm{Prophase} Prophase is the early stage of cell division when chromosomes condense, the spindle begins forming, and the nucleolus disappears.

\medterm{Prophenofos} Prophenofos is an organophosphate pesticide whose toxicity results mainly from acetylcholinesterase inhibition and excess cholinergic activity.

\medterm{Prophylactic} Prophylactic means intended to prevent disease, infection, injury, or another unwanted outcome before it occurs.
\medterm{Prophylactic Mastectomy} Prophylactic mastectomy removes one or both breasts to greatly reduce cancer risk in selected people at very high genetic or clinical risk.

\medterm{Prophylaxis} Medicine or other preventive care used to lower risk of disease, infection, or recurrence before it occurs.

\medterm{Propionates} Salts or esters of propionic acid used in food preservation, metabolism, and some medicines; clinical effects depend on the specific compound and exposure.


\medterm{Propionibacterium} A former bacterial genus name now largely replaced by Cutibacterium; its members are common skin organisms and can cause acne or device-related infection.


\medterm{Propionic Acidemia} A rare inherited disorder in which the body cannot properly break down certain proteins and fats. Harmful acids can build up, especially during illness or fasting, causing vomiting, extreme sleepiness, poor feeding, seizures, or coma. Newborn screening can detect it; treatment uses a controlled diet, avoids fasting, and manages attacks urgently.

\synonyms
Propionic Aciduria (PA); Propionyl-CoA Carboxylase Deficiency

\medterm{Propofol} A short-acting medicine given through a vein to start or maintain anesthesia and provide procedural sedation. It acts quickly but can lower blood pressure, slow breathing, or cause dangerous reactions without monitoring.
\medterm{Propofol Infusion Syndrome} A rare, potentially fatal illness after prolonged or high-dose propofol infusion. It may cause severe blood acidity, muscle breakdown, high potassium, kidney failure, heart failure, or dangerous heart rhythms. Stopping propofol and urgent intensive care are required when it is suspected.

\medterm{Propoxur} Propoxur is a carbamate insecticide that inhibits acetylcholinesterase and may cause sweating, salivation, vomiting, breathing difficulty, or weakness after exposure.

\medterm{Propoxycaine} A local anesthetic formerly used in some topical or dental preparations; toxicity can affect the nervous system and heart.

\medterm{Propoxyphene} Propoxyphene is a former opioid analgesic withdrawn in many countries because of overdose toxicity, respiratory depression, and dangerous cardiac arrhythmias.

\medterm{Propranolol} A beta-blocking medicine that slows heart rate and reduces the body's response to adrenaline. It treats selected heart conditions, tremor, migraine prevention, anxiety symptoms, and other problems, but can worsen asthma or low blood sugar.
\medterm{Proprioception} Awareness of body position and movement without looking, based on signals from muscles, tendons, joints, skin, and balance organs. Reduced proprioception can cause clumsiness, poor coordination, falls, or difficulty walking.
\medterm{Proprioceptor} A sensory receptor that detects position, movement, or muscle stretch.
\medterm{Proptosis (Exophthalmos)} Forward bulging of one or both eyes. Causes include thyroid eye disease, infection, inflammation, bleeding, or a growth behind the eye; sudden onset requires urgent medical assessment.
\medterm{Propulsion} Movement that drives the body or another object forward. In some neurologic disorders, impaired forward movement can cause short, shuffling steps or a tendency to lean and fall forward.
\medterm{Propyl Alcohol} Propyl alcohols are flammable solvents; ingestion or substantial exposure can irritate tissues and depress the central nervous system, with toxicity depending on the isomer and dose.

\medterm{Propylene Glycol} Propylene glycol is a solvent and humectant used in medicines, foods, and cosmetics; large exposures can cause metabolic or kidney problems.

\medterm{Propylthiouracil} Propylthiouracil is an antithyroid drug that reduces thyroid-hormone synthesis and conversion of T4 to T3; liver injury and fetal risks require careful use.

\medterm{Prosencephalon} The front part of the early developing brain. It later forms the cerebral hemispheres, thalamus, hypothalamus, and related structures that control thought, sensation, movement, hormones, and other functions.

\medterm{Prosopagnosia} Difficulty or inability to recognize familiar faces despite adequate vision and general intelligence. It may be developmental or result from brain injury or disease.
\medterm{Prosoplasia} Development of a more specialized cell function or structure within a tissue. The change may occur during normal development or as an adaptation, and its medical meaning depends on the tissue and surrounding disease process.
\medterm{Prosoposchisis} A rare congenital facial cleft or fissure caused by abnormal fusion of facial processes during embryonic development, potentially affecting the nose, mouth, eyes, and feeding or breathing.

\medterm{Prostacyclin Analogs} Medications that mimic prostacyclin (PGI2), causing vasodilation and inhibiting platelet
aggregation in the pulmonary vasculature. They include epoprostenol (continuous IV),
treprostinil (subcutaneous, IV, inhaled, oral), and iloprost (inhaled). They are used
for severe PAH.

\medterm{Prostaglandin} A locally acting lipid mediator derived from arachidonic acid that influences pain, inflammation, fever, vascular tone, platelet function, gastric protection, kidney function, and uterine contraction.

\medterm{Prostaglandin Analogs} Eye-drop medicines that lower pressure inside the eye by helping fluid drain. They treat glaucoma or ocular hypertension and may cause eye redness, eyelash growth, or darkening of the iris or skin around the eye.

\medterm{Prostaglandin Degradation} Enzymatic inactivation and removal of prostaglandins, which limits the duration and intensity of their local effects on vessels, inflammation, smooth muscle, kidneys, stomach, and reproductive tissues.

\medterm{Prostaglandin E2} A prostaglandin that mediates inflammation, pain, fever, vasodilation, gastric protection, cervical ripening, and uterine contraction and is used in selected obstetric and cardiovascular treatments.

\medterm{Prostaglandin F2-Alpha} A prostaglandin that promotes smooth-muscle contraction and luteolysis; analogues are used in glaucoma, postpartum hemorrhage, labor-related care, and veterinary medicine.

\medterm{Prostaglandin Receptor} A protein on or inside a cell that recognizes a prostaglandin and starts signals affecting blood-vessel width, inflammation, pain, smooth-muscle activity, platelets, or reproduction.

\medterm{Prostaglandins} Locally acting lipid mediators involved in pain, inflammation, vascular tone, gastric protection, and reproductive physiology.
\medterm{Prostaglandins E} The E series of prostaglandins, lipid mediators that influence inflammation, vascular tone, gastric protection, kidney blood flow, cervical ripening, and uterine contraction.

\medterm{Prostaglandins F} The F series of prostaglandins, lipid mediators involved in smooth-muscle contraction, luteal function, vascular responses, inflammation, and labor-related uterine activity.

\medterm{Prostate} A gland below the bladder surrounding part of the urine tube. It adds fluid to semen and may enlarge, become inflamed, or develop cancer, causing urinary, sexual, or pelvic symptoms.
\medterm{Prostate Adenocarcinoma} The most common type of prostate cancer, arising from gland cells in the prostate. It often grows slowly but can spread, and its growth is usually stimulated by male sex hormones; grade and stage help determine treatment.

\medterm{Prostate Biopsy} Removal of small samples of prostate tissue with a needle to look for cancer under a microscope. The needle may pass through the skin between the scrotum and anus or through the rectum; bleeding, infection, and temporary difficulty urinating are possible risks.

\medterm{Prostate Brachytherapy} Radiation treatment in which small radioactive sources are placed inside or next to the prostate to destroy cancer cells. It is used for selected cancers and may cause temporary urinary or bowel symptoms and later erectile problems.


\medterm{Prostate Cancer} Cancer that forms in the prostate gland. It may grow slowly or behave aggressively and may cause no symptoms early on. Risk increases with age, family history, inherited gene changes, and certain racial or ethnic backgrounds. Prostate-specific antigen testing and sometimes a digital rectal examination may help evaluate risk, but diagnosis requires a prostate biopsy. Tumor grade and stage help guide treatment.

\medterm{Advanced Prostate Cancer} Prostate cancer that has grown beyond the prostate or spread to distant parts of the body. It may involve nearby tissues, lymph nodes, bones, or other organs.

\medterm{Early-Stage Prostate Cancer} Prostate cancer that remains confined to the prostate or nearby tissues and has not spread to distant organs. The term includes cancers with different grades and risks of progression.

\synonyms
Cancer, Prostate
Prostatic Carcinoma (Prostate Cancer)


\medterm{Prostate Cancer Screening} Prostate cancer screening may use a PSA blood test with or without a digital rectal examination after discussing possible benefits, false positives, and harms.

\medterm{Prostate Cancer Staging} Determining how far prostate cancer has grown or spread. Clinicians combine examination, PSA level, tumor grade, biopsy findings, and imaging to decide whether it is confined to the prostate, nearby, or in distant organs.

\medterm{Prostate Cancer Treatment} Care for cancer of the prostate, chosen according to cancer grade, stage, growth rate, age, and health. Options may include active surveillance, surgery, radiation, hormone-blocking treatment, or medicines for advanced disease.

\medterm{Prostate Carcinoma} Prostate carcinoma is cancer arising from the prostate gland and may grow slowly or behave aggressively.

\medterm{Prostate Disease} Any disorder affecting the prostate, including benign enlargement, inflammation, infection, or cancer. It may cause difficulty urinating, weak flow, urgency, pelvic pain, blood in urine, or no symptoms initially.

\medterm{Prostate Enlargement} Alternate index label; see \textit{Benign prostatic hyperplasia}.

\medterm{Prostate Gland} The prostate, a gland below the bladder that surrounds the urethra and contributes secretions to semen.
\medterm{Prostate Gland Enlargement} Alternate index label; see \textit{Benign prostatic hyperplasia}.

\medterm{Prostate Infection} Infection or inflammation of the prostate, usually causing pelvic or urinary symptoms and sometimes fever or bloodstream infection.


\medterm{Prostate Neoplasm} An abnormal mass of cells arising in the prostate. It may be noncancerous, such as a glandular growth, or cancerous; symptoms and treatment depend on its type, size, and spread.


\medterm{Prostate Resection - Minimally Invasive} A procedure that removes part of an enlarged prostate through the urine channel or through small openings. It aims to improve urine flow while avoiding a large incision; bleeding, infection, scarring, or leakage can occur.


\medterm{Prostate-Specific Antigen} A protein made mainly by the prostate and measured in blood. A high or rising level can occur with prostate cancer, but also with prostate enlargement, infection, urine retention, ejaculation, or recent procedures. The result is considered with age, examination, trends, and other tests rather than used alone.

\medterm{Prostate-Specific Antigen Blood Test} A PSA blood test measures prostate-specific antigen; an elevated result may reflect enlargement, inflammation, infection, or cancer and is not diagnostic by itself.

\medterm{Prostatectomy} Surgery to remove part or all of the prostate, commonly for prostate cancer or severe urinary obstruction. Possible effects include urinary leakage, erectile dysfunction, bleeding, and infertility.
\medterm{Prostatic Carcinoma (Prostate Cancer)} Alternate index label for Prostate Cancer.

\medterm{Prostatic Fluid} Prostatic fluid is a mildly alkaline secretion that forms part of semen and contains substances supporting sperm function and survival.

\medterm{Prostatic Obstruction} Blockage of urine flow caused by enlargement, inflammation, scarring, stones, or a tumor of the prostate or nearby structures.

\medterm{Prostatic Urethral Stent} A small device placed in the prostatic urethra to help maintain urine flow when the prostate obstructs it.

\medterm{Prostatism} Urinary symptoms caused by blockage or poor function at the bladder outlet, often related to an enlarged prostate. Symptoms include weak flow, straining, urgency, frequent urination, nighttime urination, or incomplete emptying.
\medterm{Prostatitis} Inflammation of the prostate gland. It may be caused by a bacterial infection or by inflammation involving the prostate, pelvic muscles, or nerves without a proven infection. Symptoms may include pain in the pelvis, area between the scrotum and anus, penis, scrotum, lower abdomen, or lower back; painful or frequent urination; difficulty emptying the bladder; painful ejaculation; and, in acute bacterial prostatitis, fever and chills. The main types are acute bacterial prostatitis, chronic bacterial prostatitis, and chronic prostatitis/chronic pelvic pain syndrome.

\medterm{Prostatitis - Bacterial} Bacterial prostatitis is infection and inflammation of the prostate, causing pelvic or urinary pain, fever, urinary symptoms, or recurrent infection.


\medterm{Prostatitis - Nonbacterial} Persistent pain or urinary symptoms in the prostate and pelvis without evidence of a bacterial infection. Symptoms may fluctuate and can involve pelvic muscles, nerves, inflammation, or bladder function; care is tailored to the person’s findings.


\medterm{Prostatodynia} Prostatodynia is an older term for pelvic or prostate-region pain without proven bacterial infection, now often considered within chronic pelvic pain syndrome.

\medterm{Prosthesis} An artificial device replacing or supporting a missing body part or function. Prostheses may replace limbs, joints, teeth, eyes, or heart valves and require fitting, care, adjustment, and monitoring.
\medterm{Prosthetic} Relating to an artificial device that replaces or supports a missing or damaged body part. Examples include artificial limbs, joints, heart valves, teeth, and blood-vessel grafts; complications can include wear, loosening, skin problems, or infection.
\medterm{Prosthetic Graft} A prosthetic graft is an artificial tube or patch used to replace or bypass a damaged blood vessel or other tubular structure; infection and thrombosis are important risks.

\medterm{Prosthetic Heart Valve} An artificial valve that replaces a diseased heart valve and keeps blood moving in the correct direction. Mechanical valves usually require lifelong blood-thinning medicine; tissue valves may wear out sooner and have different medication needs.

\medterm{Prosthetic Limb} An artificial limb used to replace some or all of a missing arm or leg. It may improve movement, balance, or daily activities and is fitted to the person’s remaining limb, strength, goals, and comfort. Training helps with safe use and skin care.

\medterm{Prosthetic Training} Rehabilitation teaching a person to use an artificial limb safely and effectively. Training includes skin and residual-limb care, putting it on, walking, balance, stairs, transfers, daily activities, and community mobility.

\medterm{Prosthetics} The design, fitting, and use of artificial replacements for missing body parts, such as arms, legs, teeth, or eyes. Successful use depends on comfort, strength, balance, skin health, goals, and rehabilitation.

\medterm{Prosthodontics} The dental specialty that replaces missing teeth and restores chewing, speech, and appearance. Treatment may use crowns, bridges, implants, dentures, or facial prostheses and is planned around oral health and remaining tissues.

\medterm{Protamine} A medicine that binds to and reverses the blood-thinning effect of heparin. It is given carefully because it can cause low blood pressure, allergic reactions, slow heart rate, or dangerous clotting.
\medterm{Protanopia} Protanopia is complete absence of red-cone function, causing a form of red-green color-vision deficiency and altered perception of red and related colors.

\medterm{Protean} Highly variable in form, appearance, or behavior; in medicine it may describe a disease with changing or diverse manifestations.

\medterm{Proteaphagy} The process by which a cell uses its recycling system to remove whole proteasomes, the structures that normally break down unwanted proteins. It helps adjust protein disposal during stress or changes in nutrition.
\medterm{Protease} An enzyme that breaks proteins into smaller pieces by cutting the chemical links between their amino acids. Proteases help digest food and control many processes inside and outside cells.
\medterm{Protease Inhibitors} Antiviral medicines that block an enzyme viruses need to cut newly made proteins into working parts. They are used for selected viral infections and may interact with many other medicines.

\medterm{Proteasome} A barrel-shaped protein machine inside cells that breaks down unwanted, damaged, or short-lived proteins. This controls cell growth and signaling and prevents harmful protein buildup.

\medterm{Proteasome Binding} The attachment of a medicine or other molecule to a cell’s protein-destroying machine. This can increase, reduce, or otherwise change the breakdown of proteins and may alter cell survival or function.

\medterm{Protecting Yourself From Cancer Scams} Protection from cancer scams includes checking claims with qualified clinicians, avoiding guaranteed cures, reviewing evidence, and being cautious about costly unproven treatments.

\medterm{Protective} Serving to prevent or reduce injury, infection, or other harm.
\medterm{Protective Agent} A substance, device, or treatment used to reduce harm to the body. Examples include sunscreen protecting skin from ultraviolet light, barrier creams protecting skin, and medicines that protect the stomach or bones during selected treatments. The appropriate agent depends on the hazard and the person’s health.


\medterm{Protegrin} A small antimicrobial peptide, usually derived from pig immune cells, that can disrupt microbial membranes.
\medterm{Protein} A molecule made of amino acids that performs structural, enzymatic, transport, signaling, or immune functions.
\medterm{Protein Acetylation} Addition of an acetyl group to a protein, often changing its activity, localization, stability, or interactions; histone acetylation is a major regulator of gene expression.

\medterm{Protein Binding} The temporary attachment of a medicine to proteins in the blood. The unattached portion can usually leave the bloodstream, reach tissues, and produce effects; changes in protein levels or other medicines can alter the amount available.

\medterm{Protein Biosynthesis} The way cells make proteins: they copy information from DNA into messenger RNA, then use ribosomes to join amino acids in the instructed order. The new chain folds and may be modified before it works.

\medterm{Protein C} A natural blood-clotting regulator that helps prevent excessive clotting. Low levels may be inherited or acquired and increase clot risk; starting warfarin in severe deficiency can rarely cause skin injury.

\medterm{Protein C Blood Test} A protein C blood test measures a natural anticoagulant and may help evaluate unexplained or recurrent blood clots or suspected inherited deficiency.

\medterm{Protein C Deficiency} A condition in which too little or poorly functioning protein C reduces the body’s natural control of blood clotting. It may be inherited or acquired and increases the risk of deep-vein thrombosis and pulmonary embolism.

\medterm{Protein Deficiency} Inadequate protein intake, digestion, absorption, or use, causing loss of muscle, impaired immunity and wound healing, edema, or growth failure when severe or prolonged.

\medterm{Protein Denaturation} Loss of a protein’s normal three-dimensional structure from heat, chemicals, pH, or mechanical stress, often impairing its biological function without breaking its peptide backbone.

\medterm{Protein Dephosphorylation} Removal of a phosphate group from a protein by a phosphatase, often reversing or modifying signaling, enzyme activity, localization, or cell-cycle control.

\medterm{Protein Electrophoresis} A laboratory test that separates proteins in serum or urine according to their electrical properties, helping detect abnormal protein patterns.

\medterm{Protein Electrophoresis - Serum} Serum protein electrophoresis separates blood proteins into patterns and helps evaluate abnormal antibodies, myeloma, inflammation, and some liver or kidney disorders.

\medterm{Protein Engineering} Design or modification of proteins through genetic, biochemical, or computational methods to alter activity, stability, specificity, expression, or therapeutic usefulness.

\medterm{Protein Expression} The process by which a cell uses genetic instructions to make a protein, then folds, changes, and transports that protein so it can perform its function.

\medterm{Protein Folding} The process by which a newly made protein bends into the three-dimensional shape needed for its job. Incorrect folding can stop the protein from working or cause harmful protein deposits in cells.

\medterm{Protein Glycosylation} The attachment of sugar chains to proteins. This helps proteins fold, remain stable, reach the correct part of a cell, and work properly; inherited problems with this process can affect many organs.

\medterm{Protein Import} Movement of a newly made protein into a cell structure or across a cell membrane. Address labels on the protein and transport systems guide it to the correct location, where it can fold and work properly.

\medterm{Protein In Diet} Dietary protein supplies amino acids needed for muscle, enzymes, hormones, immunity, and tissue repair; requirements vary with age, health, and kidney function.

\medterm{Protein In Urine (Proteinuria)} An abnormal amount of protein in the urine, which may be temporary or indicate kidney disease, diabetes, high blood pressure, infection, or another condition.

\medterm{Protein Isoform} One of several related protein forms produced from different genes, alternative splicing, or post-translational processing, with potentially different locations or functions.

\medterm{Protein Kinase} An enzyme that adds a small chemical group to a protein and changes how it works. Protein kinases help control cell signals, metabolism, growth, division, and other functions.

\medterm{Protein Malnutrition} Poor health caused by inadequate protein or energy intake, leading to weight loss, muscle wasting, weak immunity, or impaired healing.

\medterm{Protein Metabolism} The synthesis, folding, modification, breakdown, and recycling of proteins and amino acids in cells and tissues.

\medterm{Protein Methylation} Addition of small chemical groups called methyl groups to selected parts of a protein. This can change the protein’s activity or location and can alter how cells switch genes on or off.

\medterm{Protein Overexpression} Production of an abnormally large amount of a protein, caused by gene amplification, transcriptional activation, altered degradation, or other mechanisms and potentially driving disease or serving as a biomarker.

\medterm{Protein Phosphorylation} Attachment of a small phosphate group to a protein by a cell enzyme. This acts like a reversible switch that can change the protein’s activity, location, or interactions and helps control cell signals and division.

\medterm{Protein Prenylation} Attachment of a fat-like chemical group to a protein, helping position it in a cell membrane. This placement allows some proteins to send signals or help move materials within the cell.

\medterm{Protein Purification} Separation of a desired protein from cells or a mixture using properties such as charge, size, solubility, affinity, or binding, while preserving its structure and activity.

\medterm{Protein Refolding} Return of a denatured or newly synthesized protein to its functional three-dimensional structure, sometimes requiring chaperones or controlled laboratory conditions.

\medterm{Protein S} A natural blood protein that helps activated protein C limit clotting. Too little or abnormal protein S can increase the risk of clots in deep veins or the lungs.

\medterm{Protein S Blood Test} A protein S blood test measures a natural anticoagulant and may help evaluate suspected inherited or acquired tendency to form blood clots.

\medterm{Protein S Deficiency} A condition in which the blood has too little or poorly functioning protein S, reducing natural control of clotting. It may be inherited or acquired and can increase the risk of deep-vein thrombosis or pulmonary embolism.

\medterm{Protein Secretion} Release of a protein from a cell or into a particular cell compartment. Cells use this process to send hormones, digestive enzymes, antibodies, and other proteins to where they are needed.

\medterm{Protein Sequence} The precise order of amino acids joined together to make a protein. This order helps determine the protein’s three-dimensional shape, location, interactions, and biological job.

\medterm{Protein Splicing} A process in which a newly made protein removes an internal section and joins the remaining sections together. The finished protein may then fold into the form needed for its normal function.

\medterm{Protein Subunit} One protein component of a larger protein assembly. Several subunits may join to create the final structure and allow the assembly to perform its biological function.

\medterm{Protein Synthesis} The process by which a cell uses messenger RNA instructions to join amino acids into a protein. The resulting protein folds into a working shape and may be changed or moved before performing its function.

\medterm{Protein Translocation} Movement of a newly synthesized or mature protein across a membrane or to a different cellular compartment, directed by targeting signals and transport machinery.

\medterm{Protein Transport} Cellular movement of proteins between compartments, across membranes, or outside the cell, ensuring that each protein reaches its functional location.

\medterm{Protein Unfolding} Loss of a protein’s native three-dimensional structure, caused by heat, chemicals, mutations, or cellular stress, which can impair function or promote aggregation.

\medterm{Protein-Energy Malnutrition} Severe lack of calories, protein, or both. It can cause weight loss, muscle wasting, weakness, swelling, poor growth, weak immunity, and slow healing; nutrition must often be restored gradually and carefully.

\medterm{Protein-Losing Enteropathy} Excessive loss of plasma proteins through the gastrointestinal tract, causing low albumin, swelling, diarrhea, or nutritional deficiency.

\medterm{Proteinaceous Plug} A blockage made primarily of protein-rich material, mucus, cellular debris, or secretions that can obstruct an airway, duct, catheter, or other passage.

\medterm{Proteinuria} An unusually high amount of protein in urine. It can result from kidney disease, diabetes, high blood pressure, infection, exercise, fever, or pregnancy and may require repeat testing to find the cause.
\medterm{Proteobacteria} A large group of bacteria containing many gram-negative species, including Escherichia, Salmonella, Pseudomonas, and Vibrio. Some are harmless or helpful, while others cause intestinal, urinary, lung, wound, or bloodstream infections.

\medterm{Proteogenomics} Integration of genomic, transcriptomic, and proteomic data to identify proteins, discover coding variants, characterize tumors or microbes, and connect DNA changes with protein function.

\medterm{Proteoglycan} A large molecule made of a protein with many attached sugar chains. It helps connective tissues and cartilage hold water, resist pressure, maintain structure, and send signals between cells.

\medterm{Proteolipid Protein} A protein containing a substantial fat component, especially important in the myelin sheath around nerve fibers.

\medterm{Proteolipids} Proteins tightly associated with lipid molecules or lipid membranes, important in myelin, mitochondrial membranes, ion transport, and other membrane structures.

\medterm{Proteolysis} Enzymatic breakdown of proteins into smaller peptides or amino acids. It is needed for digestion, cell regulation, tissue remodeling, and removal of damaged proteins.
\medterm{Proteolytic} Relating to the breakdown of proteins by cutting the links between their amino acids. Proteolytic enzymes help digest food, activate clotting and immune proteins, remove damaged proteins, and control cell death; medicines can block some of these enzymes.
\medterm{Proteome} The complete set of proteins produced by a cell, tissue, organism, or biological condition at a particular time.

\medterm{Proteomic Profile} A proteomic profile is a pattern of proteins measured in a biological sample to study disease, treatment response, or biological processes.

\medterm{Proteomics} The large-scale study of proteins, including their abundance, structure, modifications, interactions, and functions. Large-scale study of proteins, including their amounts, structures, changes, interactions, and roles in cells or disease.
\medterm{Proteus} A group of Gram-negative bacteria that can cause urinary, wound, bloodstream, or other infections. Proteus urinary infections may produce urease, raising urine pH and encouraging certain kidney stones.
\medterm{Proteus Hauseri} A gram-negative bacterium in the Proteus group that is rarely isolated from human clinical specimens.
\medterm{Proteus Mirabilis} A gram-negative bacterium that commonly causes urinary-tract infection and can produce urease-related kidney or bladder stones.

\medterm{Proteus Penneri} A gram-negative Proteus species that can rarely cause urinary, wound, bloodstream, or other opportunistic infections.
\medterm{Proteus Syndrome} A rare mosaic overgrowth disorder, usually caused by AKT1 variants, producing asymmetric enlargement of bones, skin, or other tissues and an increased risk of tumors.

\medterm{Proteus Vulgaris} A gram-negative bacterium that can cause urinary, wound, bloodstream, and other infections and often produces urease.

\medterm{Prothionamide} A thioamide antimycobacterial drug used as part of selected multidrug regimens for drug-resistant tuberculosis; gastrointestinal, hepatic, neurologic, and thyroid adverse effects require monitoring.

\medterm{Prothrombin} A blood protein made with the help of vitamin K. When a blood vessel is injured, it is changed into thrombin, which helps form a stable blood clot.
\medterm{Prothrombin Deficiency} A rare bleeding disorder caused by too little or poorly functioning prothrombin, a protein needed for blood clotting. It can cause easy bruising, nosebleeds, heavy menstrual bleeding, or prolonged bleeding after injury or surgery.

\medterm{Prothrombin Time} A blood test measuring how quickly a sample forms a clot. It is often reported as the INR and helps assess clotting disorders, liver disease, vitamin K deficiency, and the effect of warfarin.

\synonyms
PT

\medterm{Proto-Oncogene} A normal gene that helps control cell growth, division, or survival. If it becomes abnormally active because of a genetic change or extra copies, it can act as an oncogene and contribute to cancer.

\medterm{Protocol} A written plan describing the steps, timing, responsibilities, and safety rules for a study, examination, or treatment. Protocols improve consistency, protect participants, guide documentation, and define how results or complications are handled.

\medterm{Proton Pump Inhibitors} Proton pump inhibitors suppress gastric acid secretion and treat reflux, ulcers, and other acid-related disorders; prolonged or inappropriate use should be periodically reviewed.

\medterm{Proton Therapy} Radiation treatment using high-energy protons to damage cancer cells. The beam can deliver much of its dose at a selected depth, potentially reducing radiation to nearby healthy tissue; suitability depends on the tumor and treatment plan.

\medterm{Protozoa} A traditional group of single-celled organisms with nuclei. Some cause human disease, including malaria, giardiasis, and amoebiasis; modern biological classification places them in several different groups.

\medterm{Protozoal Infection} Infection caused by a disease-producing protozoan parasite. Protozoal infections may affect the intestine, blood, brain, eyes, skin, or genital tract and include malaria, amoebiasis, giardiasis, leishmaniasis, toxoplasmosis, trichomoniasis, and trypanosomiasis; transmission and treatment depend on the organism and site of disease.

\synonyms
Protozoal Infections

\medterm{Protractile} Able to move forward or outward from its usual position. The tongue and lower jaw are examples of body structures that can be actively protruded.
\medterm{Protractor} An instrument or muscle that moves a body part forward or outward. In anatomy, a protractor muscle may move the jaw, shoulder blade, or another structure away from its usual position.
\medterm{Protrusion} A part that projects outward or a disc or tissue displaced beyond its normal boundary.
\medterm{Protuberance} A protuberance is a structure or area that projects outward from a surface, such as a normal bone prominence, swelling, or mass.

\medterm{Provisional Diagnosis} A temporary working diagnosis made from available evidence before all test results are known, and subject to change.
\medterm{Provisional Tic Disorder} A disorder with motor or vocal tics lasting less than one year, beginning in childhood and not better explained by another condition.

\medterm{Provocation} An action or stimulus that brings on a symptom, response, or test finding. Clinicians may use controlled provocation during examination, but it should be stopped if it causes dangerous pain or other harm.
\medterm{Proximal} Proximal means nearer to the point of attachment or origin; in a limb, the shoulder or hip end is proximal to the hand or foot.
\medterm{Proximal Humerus Fractures} Breaks near the upper end of the upper-arm bone, close to the shoulder. They may cause pain, swelling, bruising, and difficulty moving the arm; treatment depends on displacement, joint involvement, bone quality, nerve or blood-vessel injury, and the person’s needs.

\medterm{Proximal Neuropathy Diabetic (Diabetic Neuropathy)} Alternate index label for Diabetic Neuropathy.

\medterm{Proximal Radioulnar Joint} The joint near the elbow where the two forearm bones rotate around each other. It allows the palm to turn up or down; injury or dislocation can limit rotation and cause pain.

\medterm{Proximal Renal Tubular Acidosis} Proximal renal tubular acidosis occurs when the kidney cannot reclaim enough bicarbonate, causing excess acid in the blood and often low potassium.

\medterm{Proximal Tubule} The first part of the kidney’s microscopic filtering tube after the filtering unit. It returns much of the filtered water, salt, bicarbonate, sugar, amino acids, and phosphate to the blood.

\medterm{Proximal Tubule Reabsorption} The kidney’s return of much of the filtered water and useful substances, including salt, bicarbonate, glucose, amino acids, and phosphate, from the first tubule back into the blood.

\medterm{Prune Belly Syndrome} A condition present at birth, usually in boys, involving weak or absent abdominal muscles, undescended testicles, and abnormalities of the urinary tract. Kidney problems, repeated infections, breathing difficulty, and impaired growth may occur.

\medterm{Prurigo} A group of intensely itchy skin disorders with papules or nodules caused by inflammation and scratching.
\medterm{Prurigo Nodularis} A long-lasting skin disorder with intensely itchy, firm bumps called nodules. Repeated scratching or rubbing causes and maintains an itch--scratch cycle; treatment aims to reduce itching, heal the skin, and prevent new nodules.
\medterm{Pruritic} Itchy or causing an urge to scratch; pruritus can result from skin inflammation, allergy, infection, systemic disease, medications, or neuropathic mechanisms.

\medterm{Pruritic Eruption} An itchy skin eruption that may result from allergic, inflammatory, infectious, drug-related, or autoimmune disease; diagnosis depends on its morphology, distribution, timing, and associated findings.

\medterm{Pruritis Ani} Alternate index label; see \textit{Anal itching}.

\medterm{Pruritus} An unpleasant itchy feeling that creates the desire to scratch. It may result from dry skin, allergy, infection, medicines, liver or kidney disease, nerve problems, or other conditions; severe or widespread itching needs assessment.
\medterm{Pruritus (Itch)} Alternate index label for Itch.

\medterm{Pruritus Ani} Persistent itching around the anus with multiple potential causes including hemorrhoids,
fissures, infections such as pinworms and candidiasis, dermatologic conditions including
psoriasis and lichen sclerosus, and idiopathic causes.

\medterm{Pruritus In CKD} Persistent itching associated with advanced chronic kidney disease, often worse in people receiving dialysis. It may occur without a visible rash and can disturb sleep and quality of life; treatment is individualized after other causes are assessed.

\medterm{Prussian Blue} A blue pigment used in laboratory staining and as an antidote for selected radioactive cesium or thallium poisoning.
\medterm{PSA} Prostate-specific antigen, a protein made mainly by prostate cells and measured in blood. A raised level may occur with enlargement, inflammation, infection, or cancer and does not diagnose cancer alone.

\synonyms
Prostate-specific antigen

\medterm{PSA Velocity} The rate at which prostate-specific antigen concentration changes over time. It may provide additional information in prostate evaluation but should not be used alone to diagnose cancer.

\medterm{Psammoma} A tumor or lesion containing rounded, layered calcium deposits called psammoma bodies. These deposits can occur in several tumors and are identified by microscopic examination rather than symptoms alone.
\medterm{Pseudarthrosis} Failure of a broken bone to heal, leaving persistent movement or a false joint at the fracture site. It can cause pain, weakness, deformity, and loss of function and may require surgery.
\medterm{Pseudoaneurysm} A weak area or leak in an artery in which blood collects outside the normal artery but remains connected to it. It may follow injury, surgery, infection, or a catheter procedure and can enlarge, burst, or send clots downstream.

\medterm{Pseudoangiomatous Stromal Hyperplasia} A benign breast stromal proliferation forming slit-like spaces within dense collagen that may mimic vascular channels. It can be incidental or present as a palpable or imaging-detected mass.

\medterm{Pseudobulbar Affect} Sudden, difficult-to-control episodes of laughing or crying that do not match how a person feels or the situation. It can occur after stroke or with other brain and nerve diseases and is different from depression.
\synonyms
Pathological Laughter And Crying
PBA

\medterm{Pseudoclubbing} Nail-tip broadening or flattening that resembles clubbing but does not have the characteristic change in the angle between the nail and surrounding skin. It may occur with some bone, joint, skin, or inherited conditions and needs assessment of the underlying cause.

\medterm{Pseudocyesis} A condition in which a person has pregnancy-like symptoms, such as missed periods, abdominal enlargement, nausea, or breast changes, without an actual pregnancy. A pregnancy test and examination distinguish it from pregnancy and identify other possible causes.
\medterm{Pseudodementia} Memory and thinking complaints caused by depression or another potentially reversible problem rather than progressive brain-cell loss. Careful assessment is needed because depression and true dementia can occur together or resemble each other.
\medterm{Pseudoephedrine} An oral decongestant that shrinks swollen blood vessels in the nose and improves blockage from colds or sinus problems. It may raise blood pressure, cause sleeplessness or palpitations, and interact with some medicines.

\medterm{Pseudogout} A joint disease caused by calcium pyrophosphate crystals collecting in cartilage. It can cause sudden pain, swelling, warmth, and stiffness, often in a knee or wrist, and may also produce longer-term joint symptoms.

\synonyms
Arthritis Pseudogout (Pseudogout)

\medterm{Pseudogout Flare} A sudden attack of joint pain and swelling caused by calcium pyrophosphate crystals. The knee or wrist is often affected; the joint may be red and warm, and fluid testing can confirm the diagnosis.

\medterm{Pseudohypertrophic Muscular Dystrophy} An older name for Duchenne or related muscular dystrophy with enlarged-appearing muscles caused by fat and connective tissue.
\medterm{Pseudohypoparathyroidism} A group of inherited disorders in which the body does not respond properly to parathyroid hormone. This can cause low calcium, high phosphate, muscle cramps, seizures, and sometimes short fingers or other physical features.

\medterm{Pseudomelanosis Duodeni} A usually harmless dark discoloration of the first part of the small intestine seen during endoscopy. It results from pigment-containing material in the lining and may be associated with chronic illness, medicines, or mineral deposits.

\medterm{Pseudomembranous Colitis} Severe inflammation of the colon, usually caused by overgrowth of Clostridioides difficile after antibiotics. It can cause frequent watery diarrhea, fever, abdominal pain, dehydration, and serious complications.

\synonyms
C. Difficile Colitis
Colitis, Pseudomembranous

\medterm{Pseudomonas} A group of Gram-negative bacteria that can cause opportunistic infections, especially in damaged tissues or people with weakened immunity.
\medterm{Pseudomonas Aeruginosa} A gram-negative bacterium that can cause pneumonia, bloodstream, urinary, wound, and ear infections, especially in hospitalized or immunocompromised people; it often has antibiotic resistance.

\medterm{Pseudomonas Aeruginosa Infection} Infection caused by Pseudomonas aeruginosa, an opportunistic gram-negative bacterium
that can affect the lungs, wounds, urinary tract, blood, or other sites.

\medterm{Pseudomonas Aeruginosa Resistance} The ability of Pseudomonas aeruginosa to survive medicines that would normally kill it. Resistance may be natural or develop through genetic changes, so laboratory testing is needed to choose an effective antibiotic.

\medterm{Pseudomyxoma Peritonei} A rare condition in which jelly-like tumor material gradually fills the abdomen, usually after a mucus-producing appendix tumor ruptures. It can cause increasing abdominal size, pain, bowel blockage, or poor nutrition and requires specialist treatment.

\medterm{Pseudophakia} Having an artificial lens inside the eye after the natural lens has been removed, usually during cataract surgery. The replacement lens focuses light, but later clouding of its supporting capsule or other eye problems can still affect vision.

\medterm{Pseudorabies} Aujeszky disease, a herpesvirus infection mainly affecting pigs and other animals. People are rarely infected, but exposed humans may develop serious nervous-system disease; it is different from rabies.
\medterm{Pseudothrombocytosis Due To Red Cell Fragments} A falsely high platelet count caused when automated equipment mistakes very small broken red-cell fragments for platelets. A blood-smear examination can distinguish this error from a genuinely high platelet count.

\medterm{Pseudotumor Cerebri} An older name for idiopathic intracranial hypertension, a condition in which pressure around the brain rises without a tumor. It can cause daily headache, brief dimming of vision, whooshing in the ears, and swelling of the optic nerves that may threaten sight.

\medterm{Pseudoxanthoma Elasticum} A rare inherited disorder in which elastic tissue becomes hardened and damaged. It can cause yellowish or loose skin, bleeding or vision loss from eye blood vessels, and narrowing or bleeding in arteries.

\synonyms
PXE (Pseudoxanthoma Elasticum)

\medterm{Psittacosis} An infection caused by Chlamydia psittaci, usually acquired by breathing dust from infected birds or their droppings. It commonly causes fever, headache, cough, and pneumonia and requires appropriate antibiotic treatment.
\medterm{PSMA} A protein found in large amounts on many prostate cancer cells and in smaller amounts in some normal tissues. Special tracers can use it to show prostate cancer on scans or deliver radiation to cancer cells.

\medterm{Psoas} A large muscle running from the lower spine to the thigh bone. It helps bend the hip, lift the leg, maintain upright posture, and stabilize the spine during movement.
\medterm{Psoas Abscess} A collection of pus in or around the psoas muscle, caused by infection spreading through the blood or from nearby tissues.
It may cause back or hip pain, fever, difficulty walking, or few symptoms and usually
requires imaging and targeted treatment.

\medterm{Psoriasis} A long-term immune-mediated skin disease in which skin cells are produced too quickly, causing inflammation and a buildup of skin. It commonly causes raised, thick, dry, scaly patches that may be red, purple, brown, darker, or lighter than nearby skin. The patches often affect the scalp, elbows, knees, trunk, palms, or soles, and the nails may become pitted, ridged, thick, or separated from the nail bed. Symptoms may flare and then improve. Genetic and environmental factors can contribute. Psoriasis is not contagious and may occur with psoriatic arthritis, an inflammatory disease of the joints and tendon attachments.


\medterm{Psoriatic Arthritis} An inflammatory arthritis associated with psoriasis of the skin or nails. It may affect
the spine, peripheral joints, digits, or tendon insertions, and diagnosis uses the clinical pattern together with
appropriate laboratory and imaging findings.

\synonyms
Arthritis Psoriatic (Psoriatic Arthritis)

\medterm{PSP} PSP commonly means progressive supranuclear palsy, a neurodegenerative disorder causing early postural instability, vertical gaze impairment, rigidity, and cognitive change.

\medterm{Psychiatric} Relating to the assessment, diagnosis, treatment, or prevention of disorders of mood, thought, perception, behavior, or emotional functioning.

\medterm{Psychiatric Emergency} A mental-health condition involving an immediate and serious risk to life, safety, basic functioning, or the safety of others. Examples include imminent suicide risk, severe psychosis, dangerous agitation, catatonia, or severe substance intoxication or withdrawal. It is a subset of mental-health crises distinguished by the immediacy and severity of risk.

\medterm{Psychiatrist} A medical doctor who assesses, diagnoses, and treats mental-health conditions. A psychiatrist may provide talking therapy, prescribe medicines, coordinate other treatments, and assess how mental or physical illness affects safety and daily life.

\medterm{Psychiatry} The medical specialty concerned with mental health and the diagnosis and treatment of mental disorders.
\medterm{Psycho-Oncology} The clinical field addressing psychological, social, behavioral, and communication needs of people with cancer and their families during diagnosis, treatment, survivorship, and end-of-life care.

\medterm{Psychoacoustics} The study of how people perceive and interpret sound, including loudness, pitch, masking, localization, speech, and hearing thresholds.

\medterm{Psychoactive} Affecting mood, thought, perception, behavior, alertness, or other mental processes. Psychoactive substances include prescribed medicines, alcohol, and recreational drugs.
\medterm{Psychoactive Substance} A substance that alters mood, perception, cognition, behavior, or consciousness.

\medterm{Psychoanalysis} A psychological theory and talking treatment emphasizing unconscious thoughts, feelings, memories, and conflicts. Traditional treatment explores patterns in relationships and behavior, while modern approaches vary in methods, goals, and evidence.
\medterm{Psychodermatology} The area of care that considers links between skin conditions and emotional health. Stress can worsen problems such as eczema, psoriasis, acne, or hives, while visible skin disease can cause anxiety or low mood; both aspects may need treatment.

\medterm{Psychodidae} A family of small hairy flies that includes sand flies; some species transmit pathogens such as Leishmania and therefore have public-health importance.

\medterm{Psychodrama} A therapy using guided acting to explore emotions, relationships, and behavior in a supportive setting.
\medterm{Psychodynamic Therapy} Talking therapy that explores how past experiences, relationships, feelings, and patterns outside a person’s immediate awareness may influence current distress and behavior. The therapist and patient use the therapeutic relationship to develop insight and healthier ways of coping.

\medterm{Psychogenic} Caused or strongly influenced by psychological factors such as stress, trauma, or emotional conflict. The symptoms are real and not deliberately produced, and physical causes should still be considered.

\medterm{Psychogenic Erectile Dysfunction} Difficulty getting or keeping an erection mainly related to stress, anxiety, depression, relationship concerns, or other emotional factors. Physical causes can coexist, so medical assessment should consider health, medicines, and psychological wellbeing.

\medterm{Psychogenic Non-Epileptic Seizures} Episodes that resemble epileptic seizures but are not caused by the abnormal electrical discharges of epilepsy. They are real and involuntary and may involve shaking, stiffening, unresponsiveness, or altered movement; video-EEG can help establish the diagnosis.

\medterm{Psycholinguistics} The study of how the brain and mind produce, understand, acquire, and remember language, including speech, reading, grammar, and meaning.

\medterm{Psychological} Relating to thoughts, feelings, behavior, or the way these affect health. Psychological factors can influence pain, sleep, coping, treatment choices, and recovery, but describing a factor as psychological does not mean symptoms are imagined.

\medterm{Psychological Trauma} Emotional or psychological effects that may follow an event or ongoing situation involving actual or threatened death, serious injury, violence, abuse, or another overwhelming experience. Responses vary and may include fear, intrusive memories, avoidance, emotional numbness, sleep problems, or changes in mood and behavior.
\medterm{Psychological Adjustment} The process of adapting thoughts, emotions, and behavior to stress, illness, disability, loss, developmental change, or other life circumstances.

\medterm{Psychological Dependence} A pattern in which a person feels a strong emotional or mental need for a substance or behavior, often with craving, preoccupation, or distress when it is unavailable.

\medterm{Psychological Support} Emotional and practical help for people coping with illness, disability, pain, or difficult treatment. It may include counseling, coping skills, support groups, and assessment or treatment of anxiety, depression, and adjustment problems.

\medterm{Psychological Transfer} The process in which feelings, expectations, or relationship patterns from earlier experiences become directed toward someone in the present, especially a therapist. Recognizing this pattern can help understand distress and relationships during therapy.

\medterm{Psychologist, Clinical} A licensed mental-health professional who assesses and treats psychological, emotional, behavioral, and cognitive problems using psychotherapy and other evidence-based interventions within local scope of practice.

\medterm{Psychology} The scientific study of behavior, emotion, learning, and mental processes. Psychological care may include assessment, counseling, and behavioral therapies.

\medterm{Psychometrics} The design and use of tests that measure mental abilities, feelings, behavior, or personality. It also examines whether a test gives consistent results and truly measures the quality it claims to measure.
\medterm{Psychomotor} Relating to the connection between mental activity and physical movement, including the speed, coordination, and purposeful control of actions.
\medterm{Psychomotor Agitation} Restless, excessive, or purposeless movement associated with inner tension, anxiety, mania, delirium, depression, substance effects, or medication reactions.

\medterm{Psychomotor Disorder} A condition affecting the speed, starting, coordination, or expression of movement and mental activity. A person may move unusually slowly or quickly, seem restless, or have difficulty organizing purposeful actions.

\medterm{Psychomotor Performance} The ability to coordinate mental processing with physical movement, assessed through tasks involving speed, precision, reaction time, balance, or sequencing.

\medterm{Psychoneuroimmunology} The study of how thoughts, emotions, the brain, hormones, and the immune system influence one another in health and illness. It examines links involving stress, inflammation, infection, and recovery.

\medterm{Psychopathic} Describing severe, persistent antisocial traits such as limited empathy, deceitfulness, and disregard for others; it is not a precise diagnosis.
\medterm{Psychopathology} The study and description of abnormal mental processes, symptoms, behaviors, and mechanisms underlying psychiatric disorders. The study or description of abnormal mental states and behavior.
\medterm{Psychopharmacology} The study of how medicines affect mood, behavior, thinking, sleep, and other mental functions in people.
\medterm{Psychophysics} The study of how physical stimuli, such as light, sound, pressure, or heat, are experienced by a person. It measures detection thresholds, ability to distinguish stimuli, and perceived strength.

\medterm{Psychophysiology} The study of links between psychological processes and measurable bodily functions such as heart rate, skin conductance, muscle activity, and brain signals.

\medterm{Psychosexual} Relating to how psychological factors interact with sexual development, feelings, behavior, relationships, or function. Mental health, past experiences, body image, culture, medicines, and physical conditions can all influence sexuality.
\medterm{Psychosexual Development} Development of sexual identity, interest, relationships, and behavior across the lifespan as influenced by biology, cognition, attachment, culture, and experience.

\medterm{Psychosis} A state in which a person loses contact with reality. It may involve strongly held false beliefs, hearing or seeing things others do not, confused speech, or markedly disorganized behavior and can result from mental, medical, neurologic, or substance-related causes.
\medterm{Psychosis ICU (ICU Psychosis)} Alternate index label for ICU Psychosis.

\medterm{Psychosocial Rehabilitation} Support that helps a person with mental illness build skills, relationships, confidence, and independence. It may include daily-living training, education, work support, peer support, therapy, and help taking part in family and community life.

\medterm{Psychosocial Support} Practical, emotional, and social help for people facing illness or difficult treatment. It may include counseling, support groups, help with work or money worries, family support, and care for anxiety, depression, fear, loneliness, or grief.

\medterm{Psychosomatic} Relating to physical symptoms influenced by psychological processes or stress.

\medterm{Psychosurgery} Surgery intended to alter brain circuits for severe, treatment-resistant psychiatric illness; it is rarely used and highly specialized.
\medterm{Psychotherapy} Structured treatment using therapeutic communication and psychological techniques to reduce distress, change unhelpful patterns, and improve coping or relationships. Common approaches include cognitive behavioral and interpersonal therapies.
\medterm{Psychotic} Describing psychosis, a syndrome involving impaired reality testing. It may include delusions, hallucinations, disorganized thinking or speech, or markedly disorganized behavior. Psychosis can occur with schizophrenia-spectrum disorders, severe mood disorders, substances or medicines, delirium, neurologic disease, or other medical conditions.
\medterm{Psychotic Disorder} Alternate index label for Psychotic Disorders.

\medterm{Psychotic Disorder Brief (Brief Psychotic Disorder)} Alternate index label for Brief Psychotic Disorder.

\medterm{Psychotic Disorders} A group of mental disorders characterized by impaired reality testing, such as delusions, hallucinations, disorganized thought or speech, or markedly disorganized behavior. The group includes schizophrenia-spectrum disorders and other conditions in which psychosis is a defining feature; psychotic symptoms can also occur with mood disorders, substances, medicines, neurologic disease, or other medical conditions.


\medterm{PTCA} Percutaneous transluminal coronary angioplasty, a procedure using a small balloon to open a narrowed heart artery. A stent is often placed afterward to help keep the artery open and restore blood flow.

\synonyms
Percutaneous coronary intervention

\medterm{PTCL} Alternate index label; see \textit{Peripheral T-cell lymphoma}.

\medterm{PTE} Alternate name; see \textit{Pulmonary Thromboendarterectomy}.

\medterm{Pteridine} A nitrogen-containing chemical structure found in folate vitamins, natural pigments, and other active body compounds. Pteridine-containing molecules participate in cell growth, chemical reactions, color production, and inherited metabolic disorders.

\medterm{Pterygium} A raised, often triangular growth of the clear covering of the eye that extends onto the cornea. Long-term sun, wind, dust, or dry-eye exposure may contribute; it can cause irritation, redness, blurred vision, or astigmatism.
\medterm{PTH-Related Peptide} A hormone-like substance sometimes made in excess by cancer cells. It acts partly like parathyroid hormone and can raise blood calcium, causing thirst, constipation, weakness, confusion, dehydration, or kidney problems.

\medterm{Pthirus Pubis} The insect commonly called the pubic louse. It lives mainly in coarse genital hair and causes intense itching and visible lice or eggs; it spreads mainly through close body or sexual contact.

\medterm{Ptilosis} An abnormal hair pattern or hair disorder, including excess growth, unusual direction, or loss. The term is nonspecific, so the body site, appearance, timing, and underlying cause must be assessed.
\medterm{PTLD} Post-transplant lymphoproliferative disorder, an abnormal overgrowth of immune cells after organ or stem-cell transplantation. It ranges from a reversible response to lymphoma and may cause fever, swollen lymph nodes, organ problems, or unexplained weight loss.

\medterm{Ptosis} Drooping of an eyelid or another body structure. Eyelid ptosis may result from muscle, nerve, tendon, aging, injury, or a neurologic disorder.
\medterm{Ptosis - Infants And Children} Ptosis in an infant or child is drooping of the upper eyelid, which may obstruct vision or cause amblyopia and can reflect congenital, muscular, nerve, or other disease.

\medterm{PTSD} Alternate index label for Post-Traumatic Stress Disorder.

\medterm{Ptyalin} An older name for salivary amylase, the enzyme that begins starch digestion in the mouth.
\medterm{Ptyalism (Sialorrhoea)} Excessive saliva or drooling. It may result from increased saliva production or difficulty swallowing, and can occur with pregnancy, mouth irritation, reflux, medicines, or disorders affecting the brain, nerves, or muscles.
\medterm{Pubertal Failure} Puberty that does not begin or stops progressing as expected. Causes may involve the brain’s hormone centers, pituitary gland, testes or ovaries, chromosomes, chronic illness, inadequate nutrition, excessive exercise, or other factors.

\medterm{Puberty} The stage when hormone changes cause the body to mature sexually and become capable of reproduction. It includes growth, development of breasts or testes and genitals, body hair, voice changes, and menstrual periods or sperm production.
\medterm{Puberty In Boys} Puberty in boys is the hormone-driven development of testes, penis, body hair, voice, muscle, and fertility, usually progressing through predictable physical stages.

\medterm{Puberty In Girls} Puberty in girls is the hormone-driven development of breasts, body hair, growth, menstrual cycles, and reproductive maturity, usually progressing through predictable stages.

\medterm{Pubes} The area in front of the pelvis over the pubic bone, or the coarse hair that grows there during puberty. Changes in pubic hair can help assess sexual development, but must be interpreted with other physical findings.
\medterm{Pubic} Relating to the pubis, the front part of the pelvis, or the area covered by hair above the genitals.
\medterm{Pubic Lice} Pubic lice are parasitic insects infesting coarse body hair, causing intense itching and visible nits; treatment and assessment of close or sexual contacts are needed.

\medterm{Pubic Lice (Crabs)} Infestation of coarse body hair by lice, usually causing intense itching and requiring treatment and contact precautions.

\synonyms
Crabs

\medterm{Pubis} The front part of the pelvic bone, meeting its opposite side at the pubic symphysis.
\medterm{Public Assistance} Government or community support providing eligible people with food, housing, income, healthcare, disability, or other essential services.

\medterm{Public Exposure} Radiation received by other people from a patient who has undergone nuclear medicine treatment. The amount depends on the radioactive substance and treatment; the medical team gives specific instructions to limit exposure when needed.

\medterm{Public Health} The science and practice of protecting and improving the health of communities through
education, policy-making, and research for disease and injury prevention.

\medterm{Public Health Emergency} A serious event threatening the health of many people, such as a pandemic, disaster, or deliberate biological attack.



\medterm{Puborectalis Muscle} A loop-shaped pelvic muscle that surrounds the lower rectum and helps control bowel movements. At rest it pulls the rectum forward to make a bend that helps hold stool; relaxing it helps stool pass.

\medterm{Pudendal} Relating to the external genitalia, anus, or perineum. The pudendal nerve carries sensation from these areas and helps control several muscles involved in continence and sexual function.
\medterm{Pudendal Arteries} Arteries supplying blood to the external genital and perineal region.

\medterm{Pudendal Block} Injection of local anesthetic near the pudendal nerve to numb the lower vagina, vulva, anus, and perineum. It may reduce pain during late labor, assisted delivery, or repair of a birth-related tear.

\medterm{Pudendal Nerve} The nerve carrying sensation and motor signals to much of the perineum and pelvic floor.
\medterm{Puerperal} Relating to the period after childbirth, usually the first six weeks. During this time the body recovers from pregnancy and birth, the womb becomes smaller, bleeding gradually settles, and problems such as infection, heavy bleeding, or blood clots may occur.
\medterm{Puerperal Disorder} A physical or mental health disorder occurring during the postpartum period, including infection, bleeding, hypertension, depression, or psychosis.

\medterm{Puerperal Fever} Fever occurring after childbirth, often caused by infection of the uterus, birth canal, urinary tract, breast, wound, or bloodstream. It requires prompt assessment because untreated infection can spread and become life-threatening.
\medterm{Puerperal Infection} An infection after childbirth, commonly involving the uterus, surgical wound, urinary tract, breast, or bloodstream, often causing fever, pain, or foul discharge.

\medterm{Puerperal Psychosis} A rare, severe mental-health emergency that begins soon after childbirth. Rapid mood changes, confusion, false beliefs, hallucinations, or disorganized behavior can develop quickly and may threaten the safety of the parent or baby, requiring immediate assessment.

\medterm{Puerperal Sepsis} A serious infection of the reproductive tract during labor or the six weeks after childbirth. Fever, pelvic pain, foul-smelling discharge, weakness, or fast breathing require urgent medical assessment and antibiotic treatment.

\synonyms
Postpartum Sepsis; Puerperal Septicemia

\medterm{Puerperium} The period after childbirth, usually about six weeks, when the body recovers from pregnancy and birth. The uterus shrinks, bleeding settles, milk production develops, and health problems may still arise.
\medterm{Puffy Eyes} Alternate index label; see \textit{Bags under eyes}.

\medterm{Pugilistic Stance} A flexed posture of the arms, hands, and legs sometimes seen in bodies exposed to intense heat. It results from heat-related muscle and tendon contraction, not from purposeful movement, and does not by itself show when or how death occurred.

\medterm{Pugilistica Dementia (Dementia)} Alternate index label for Dementia.

\medterm{Pulled Hamstring (Hamstring Injury)} Alternate index label for Hamstring Injury.

\medterm{Pulling A Patient Up In Bed} A safe repositioning technique using adequate staff, friction-reducing equipment, and coordinated movement to move a patient toward the head of the bed while preventing falls, shear, and staff injury.

\medterm{Pulmonary} Relating to the lungs or the blood vessels and airways that support breathing. Pulmonary disease may involve infection, asthma, chronic obstructive lung disease, high pressure in lung vessels, blood clots, fluid, scarring, or cancer.

\medterm{Pulmonary Actinomycosis} Pulmonary actinomycosis is a chronic bacterial infection of the lungs caused by Actinomyces and may form mass-like lesions or spread across tissue planes.

\medterm{Pulmonary Alveolar Hemorrhage} Bleeding into the lung’s air sacs. It may cause coughing up blood, breathlessness, low oxygen, anemia, and widespread lung shadows and can result from immune disease, infection, medicines, injury, or a clotting problem.

\medterm{Pulmonary Alveolar Microlithiasis} A rare inherited lung disease in which tiny mineral deposits build up inside the air sacs. It may cause no symptoms at first, but later can lead to cough, breathlessness, low oxygen, and progressive lung scarring.

\medterm{Pulmonary Alveolar Proteinosis} A rare lung disease in which a fatty material builds up inside the air sacs and interferes with oxygen exchange. It can cause gradually worsening breathlessness, cough, tiredness, and low oxygen and may be immune-related or secondary to another illness.

\medterm{Pulmonary Angiography} Contrast imaging of the pulmonary arteries, usually by catheter. It can diagnose pulmonary vascular abnormalities but is now used less often for pulmonary embolism because CT pulmonary angiography is less invasive.

\medterm{Pulmonary Arterial Hypertension} High blood pressure in the arteries carrying blood from the heart to the lungs because these vessels become narrowed or damaged. It can cause breathlessness, tiredness, chest discomfort, dizziness, or fainting and may strain the right side of the heart.

\medterm{Pulmonary Arteries} The two large blood vessels carrying oxygen-poor blood from the right side of the heart to the lungs. In the lungs, the blood releases carbon dioxide and picks up oxygen.

\synonyms
pulmonary artery

\medterm{Pulmonary Arteriovenous Fistula} An abnormal direct connection between an artery and vein in the lung that lets blood bypass normal oxygen exchange. It can cause low oxygen, breathlessness, blue skin, stroke-causing clots, or rarely a brain abscess.

\medterm{Pulmonary Arteriovenous Malformation} An abnormal direct connection between an artery and vein in the lung. Blood bypasses normal oxygen exchange, which can cause low oxygen, breathlessness, coughing blood, stroke, or a brain abscess; some cases are linked to an inherited blood-vessel disorder.

\medterm{Pulmonary Artery} The large blood vessel carrying oxygen-poor blood from the right side of the heart to the lungs. There, blood releases carbon dioxide, receives oxygen, and returns through pulmonary veins.
\medterm{Pulmonary Artery Catheter} A thin tube passed through the right side of the heart into a lung artery to measure pressures, blood flow, and heart performance. It is used selectively in severe heart or lung disease and critical illness because insertion has risks.

\medterm{Pulmonary Artery Pulsatility Index} A measurement calculated during right-heart catheterization from pulmonary-artery pressure and right-atrial pressure. A low result may indicate weak right-ventricle pumping and a higher risk of right-sided heart failure in seriously ill patients.

\synonyms
PAPi; PAP Index

\medterm{Pulmonary Aspergilloma} A pulmonary aspergilloma is a ball of Aspergillus fungus and debris growing in a pre-existing lung cavity, sometimes causing coughing of blood.

\medterm{Pulmonary Atresia} A heart defect present from birth in which the valve or passage from the right ventricle to the pulmonary artery does not open normally. Blood cannot reach the lungs through the usual route, so newborns may become blue or severely unwell and need urgent specialist care.

\synonyms
Tetralogy of Fallot With Pulmonary Atresia

\medterm{Pulmonary Blastoma} A rare malignant lung tumor composed of primitive epithelial and mesenchymal tissue, usually treated with specialist oncology care.

\medterm{Pulmonary Capillary Wedge Pressure} A pressure measured with a special catheter in a small lung artery to estimate pressure on the left side of the heart. It can help assess heart failure or fluid overload in critically ill patients, but breathing support and other conditions affect its meaning.

\medterm{Pulmonary Circuit} The circulation in which blood travels from the right ventricle to the lungs for gas exchange and returns to the left atrium.

\medterm{Pulmonary Circulation} The blood-flow circuit from the right ventricle through the pulmonary arteries and lung capillaries to the left atrium, where blood releases carbon dioxide and gains oxygen.

\medterm{Pulmonary Compliance} How easily the lungs expand when pressure changes during breathing. Low compliance means the lungs are stiff and harder to inflate; high compliance can mean they expand easily but may not recoil effectively.

\medterm{Pulmonary Contusion} Bruising of the lung after a blow or crush injury to the chest. Bleeding and swelling inside the lung can reduce oxygen exchange and worsen during the first day or two; treatment supports breathing and addresses associated injuries.

\medterm{Pulmonary Edema} A buildup of fluid in the lungs that can make breathing difficult. It is often caused by
heart failure but may also result from kidney problems, lung injury, medicines, or high altitude.
Sudden severe breathing difficulty is an emergency.

\synonyms
Edema, Pulmonary
Edema Pulmonary (Pulmonary Edema)

\medterm{Pulmonary Embolism} A blockage of a pulmonary artery, usually caused by a blood clot that traveled from a deep vein. It may cause sudden breathlessness, chest pain that worsens with breathing, rapid heartbeat, coughing blood, dizziness, or fainting. A large embolism can severely strain the heart or cause cardiac arrest and requires emergency assessment.

\synonyms
Blood Clot In The Lung (Pulmonary Embolism)
Embolism Pulmonary (Pulmonary Embolism)

\medterm{Pulmonary Embolism Treatment} Treatment for a blood clot in the lung usually uses anticoagulant medicine to stop growth and prevent new clots. Severe or unstable cases may need clot-dissolving medicine or a procedure; choice and duration depend on the person’s risks and cause.

\medterm{Pulmonary Embolus} A blockage of a pulmonary artery, usually by a blood clot that traveled from a deep vein, causing sudden breathlessness, chest pain, low oxygen, or collapse.

\medterm{Pulmonary Emphysema} Alternate index label; see \textit{Emphysema}.

\medterm{Pulmonary Fibrosis} A group of lung diseases in which scar tissue forms in or around the air sacs and other lung tissue. The scarring makes the lungs thick and stiff, so they cannot expand normally and oxygen has more difficulty moving into the blood. It may cause gradually worsening shortness of breath, a dry cough, tiredness, and reduced ability to exercise. Causes include autoimmune disease, some medicines, workplace dust or chemicals, radiation, infection, or an unknown cause; the unknown-cause form is called idiopathic pulmonary fibrosis.

\medterm{Pulmonary Fistula} An abnormal passage between the lung or an airway and another structure, such as the space around the lung, skin, or swallowing tube. It can allow air or infected fluid to travel where it should not.

\medterm{Pulmonary Function Tests} Breathing tests that measure how much air the lungs hold, how quickly air can be moved, and how well oxygen passes into the blood. They help assess asthma, chronic obstructive lung disease, scarring, and other breathing problems. Results depend on effort, technique, age, and body size.

\medterm{Pulmonary Hemorrhage} Bleeding into the airways or lung tissue, which may cause coughing blood, low oxygen, anemia, or respiratory failure.

\medterm{Pulmonary Hemosiderosis} A condition in which repeated bleeding into the lung’s air sacs leaves iron deposits behind. It may cause cough, breathlessness, coughing blood, anemia, and low oxygen and can result from immune disease or other causes.

\medterm{Pulmonary Hypertension} Abnormally high blood pressure in the blood vessels that carry blood from the right side of the heart to the lungs. Narrowed, blocked, or damaged vessels make the right side of the heart work harder to pump blood and may eventually weaken it. Pulmonary hypertension may occur on its own or result from left-sided heart disease, lung disease or low oxygen, long-term blood clots, or other conditions affecting the blood vessels. Symptoms may include breathlessness, tiredness, chest discomfort, dizziness, fainting, or swelling from heart strain.

\synonyms
PH (Pulmonary Hypertension)

\medterm{Pulmonary Hypertension - At Home} Home care for pulmonary hypertension includes taking prescribed medicines, monitoring symptoms, conserving energy, following oxygen advice, and seeking urgent help for severe breathlessness or chest pain.

\medterm{Pulmonary Hypertension Classification} A clinical grouping based on cause: disease of the pulmonary arteries, left-sided heart disease, lung disease or low oxygen, chronic blood clots, or several other mechanisms. Identifying the group guides testing and treatment because medicines useful for one group may be harmful in another.

\medterm{Pulmonary Infarction} Death of a small area of lung tissue after its blood supply is blocked, usually by a pulmonary embolism. It may cause sharp pain with breathing, coughing up blood, breathlessness, or a shadow on imaging.

\medterm{Pulmonary Infection} An infection of the lungs or lower respiratory tract caused by bacteria, viruses, fungi, or other organisms.

\medterm{Pulmonary Langerhans Cell Histiocytosis} A rare lung disease strongly linked to cigarette smoking, in which unusual immune cells damage small airways and form lung cysts. It can cause cough, breathlessness, chest pain, or a collapsed lung. Stopping smoking is essential and may stabilize the disease; specialist treatment is needed if it progresses.

\synonyms
PLCH; Histiocytosis X of the Lung; Pulmonary Eosinophilic Granuloma

\medterm{Pulmonary Lobes} The anatomical divisions of the lungs separated by interlobar fissures. The right lung
has superior, middle, and inferior lobes; the left lung has superior and inferior lobes
and a lingula.

\medterm{Pulmonary Mass} A space-occupying lesion in or near the lung, which may represent infection, inflammation, benign growth, or cancer and requires appropriate imaging assessment.

\medterm{Pulmonary Nocardiosis} A lung infection caused by Nocardia bacteria, usually affecting people with weakened immunity and sometimes spreading to the brain or skin.

\medterm{Pulmonary Nodule} A rounded opacity in the lung, usually smaller than 3 cm, caused by scar, infection, inflammation, or tumor and assessed according to size and risk factors.

\medterm{Pulmonary-Renal Syndrome} A group of diseases that cause inflammation or bleeding in the lungs together with inflammation or injury in the kidneys. It is often associated with small-vessel vasculitis or another autoimmune disease.

\medterm{Pulmonary Oedema} Fluid collecting in lung tissue or air sacs, making breathing difficult and reducing oxygen transfer. It often results from heart failure but can also follow infection, toxins, kidney disease, or severe injury.
\medterm{Pulmonary Reactance} A measure of how the stretch and movement of the lungs and chest wall affect airflow during breathing. It is obtained during specialized tests and can help describe stiffness or altered movement of the respiratory system.
\medterm{Pulmonary Rehabilitation} A supervised, multidisciplinary program combining exercise training, education,
behavioral support, and self-management for chronic respiratory disease. It improves
exercise capacity, dyspnea, functional status, and quality of life.

\medterm{Pulmonary Sequestration} A birth-related area of lung tissue that does not connect normally to the airways and receives blood from an unusual artery. It may cause repeated lung infections, cough, or breathing problems, although some cases cause no symptoms.

\medterm{Pulmonary Stenosis} Narrowing that obstructs blood leaving the right side of the heart on its way to the lungs, usually at the pulmonary valve. It is often present from birth and may cause a heart murmur, tiredness, breathlessness, or fainting.
\medterm{Pulmonary Surfactant} A slippery mixture made by cells lining the lung’s air sacs that reduces surface tension and helps keep the sacs open. Too little in premature babies can cause serious breathing difficulty.

\medterm{Pulmonary Thromboendarterectomy} Major surgery to remove old clots and scar tissue from the arteries of the lungs. It may improve or cure long-term clot-related high blood pressure in the lungs in selected people, but it carries serious surgical risks.

\synonyms
PTE; Pulmonary Endarterectomy; PEA (Pulmonary Endarterectomy)

\medterm{Pulmonary Trunk} The large vessel leaving the right ventricle that divides into the right and left pulmonary arteries.

\medterm{Pulmonary Tuberculosis} Tuberculosis infection involving the lungs, caused by Mycobacterium tuberculosis. It may cause prolonged cough, fever, night sweats, weight loss, or coughing blood and can spread through the air.
\medterm{Pulmonary Valve} A one-way valve between the right lower heart chamber and the artery leading to the lungs. It opens to let blood leave the heart and closes to prevent backward flow; narrowing or leakage can strain the right side of the heart.

\medterm{Pulmonary Valve Regurgitation} Leakage through the pulmonary valve that allows blood to flow back into the right lower heart chamber as the heart relaxes. Mild leakage may cause no symptoms; severe leakage can enlarge and weaken the right side of the heart.

\synonyms
Pulmonic regurgitation; Pulmonary valve insufficiency


\medterm{Pulmonary Valve Insufficiency} Leakage of blood backward through the pulmonary valve into the right lower heart chamber while the heart relaxes. Mild leakage may cause no symptoms, but severe disease can enlarge the heart and reduce exercise ability.

\medterm{Pulmonary Vascular Resistance} The opposition to blood flow through the lung blood vessels. High resistance makes the right side of the heart pump harder and can lead to high pressure in the lungs, right-heart enlargement, and heart failure.

\medterm{Pulmonary Vasodilator Therapy} Treatment with medicines that widen blood vessels in the lungs and lower pressure on the right side of the heart. It is used for selected types of pulmonary hypertension after specialist testing because it may be harmful in other types.

\medterm{Pulmonary Vein} The four veins (two from each lung) returning oxygenated blood from the lungs to the
left atrium, a unique exception in the venous system (carrying oxygenated blood).
Superior and inferior pulmonary veins enter the left atrium posteriorly.

\medterm{Pulmonary Vein Isolation} A catheter procedure for selected people with atrial fibrillation. The clinician creates small areas of scar around the veins entering the upper heart chamber to block abnormal electrical signals. It can reduce episodes but may not cure them; bleeding, stroke, narrowing of a vein, or other complications are possible.

\synonyms
PVI; Antral Pulmonary Vein Isolation; AFib Catheter Ablation

\medterm{Pulmonary Veins} Usually four veins that carry oxygen-rich blood from the lungs to the heart's left upper chamber. They are unusual veins because they carry oxygen-rich rather than oxygen-poor blood.

\medterm{Pulmonary Veno-Occlusive Disease} A rare progressive disease in which small veins in the lungs become narrowed or blocked. Pressure then rises in the lung circulation, causing breathlessness and low oxygen; some medicines for pulmonary hypertension can trigger life-threatening fluid in the lungs.

\medterm{Pulmonary Ventilation} The movement of air into and out of the lungs. It brings oxygen into the air sacs and removes carbon dioxide; inadequate ventilation can cause low oxygen or a buildup of carbon dioxide.

\medterm{Pulmonary Ventilation/Perfusion Scan} An imaging test comparing air reaching different lung regions with blood flow through those regions. A mismatch may suggest a blood clot in the lung, while other patterns can indicate lung or blood-vessel disease.

\medterm{Pulmonary Wedge Pressure} A pressure measured through a special catheter in a small lung artery to estimate the pressure in the left side of the heart. It helps assess fluid overload or heart failure in seriously ill patients.

\medterm{Pulmonary-Artery Dissection} A rare, life-threatening tear in the wall of an artery carrying blood from the heart to the lungs. It may cause sudden chest pain, breathlessness, low blood pressure, or collapse and requires emergency specialist care.

\medterm{Pulmonic Valve Stenosis} Pulmonic valve stenosis is narrowing of the valve between the right ventricle and pulmonary artery, creating outflow obstruction and sometimes right-heart strain.

\medterm{Pulmonologist} A pulmonologist is a physician specializing in diseases of the lungs and breathing system, including asthma, COPD, infections, and sleep-related breathing disorders.

\medterm{Pulp} The soft living tissue in the center of a tooth, containing nerves and blood vessels. It can become inflamed or infected after decay, cracks, or injury and may then require root-canal treatment or removal.
\medterm{Pulp Canal} The pulp canal is the passage within a tooth root containing dental pulp, nerves, and blood vessels; infection may require root-canal treatment.

\medterm{Pulp Disease} A problem affecting the living tissue inside a tooth, including inflammation, infection, loss of blood supply, or tissue death. It may cause toothache or sensitivity and can follow decay, a crack, a filling, or injury.

\medterm{Pulp Stone} A small area of hardened mineral inside the soft center of a tooth. It often causes no symptoms and is found on dental imaging, but it may complicate root-canal treatment if it blocks the canal.

\medterm{Pulpectomy} Complete removal of infected or damaged pulp from inside a tooth, usually a primary tooth. The space is cleaned and filled with a material that can safely resorb as the adult tooth develops.

\medterm{Pulpitis} Inflammation of the dental pulp, usually caused by caries, trauma, cracks, or restorative treatment. Reversible pulpitis causes brief pain provoked by cold or sweet stimuli; irreversible pulpitis causes spontaneous or lingering pain and may progress to pulp necrosis, requiring endodontic treatment or extraction.
\medterm{Pulposus} Relating to the nucleus pulposus, the gelatinous central portion of an intervertebral disk. Relating to the soft, jelly-like center of an intervertebral disk, which helps cushion spinal movement.
\medterm{Pulpotomy} Removal of inflamed soft tissue from the crown of a baby tooth while keeping the root tissue alive. It treats suitable decay or injury and preserves the tooth until it naturally falls out.

\medterm{Pulse} A pulse is the rhythmic expansion of an artery produced by each heartbeat; its rate, rhythm, strength, and symmetry provide clues about circulation and cardiac function.
\medterm{Pulse - Bounding} A bounding pulse feels unusually strong or full and can occur with fever, anemia, hyperthyroidism, exercise, anxiety, aortic regurgitation, or high-output states.

\medterm{Pulse Examination} Checking arterial pulses for speed, regularity, strength, and equality on both sides. It helps assess blood flow, heart rhythm, circulation to limbs, and possible blockage or reduced blood supply.

\medterm{Pulse Irregular} An irregular pulse has uneven timing or varying beats and may reflect premature beats, atrial fibrillation, other arrhythmias, or measurement variation.

\medterm{Pulse Oximeter} A small device, often placed on a fingertip, that estimates blood oxygen and pulse rate using light. Movement, cold fingers, poor circulation, nail products, and device limitations can make the reading inaccurate.

\medterm{Pulse Oximetry} A painless method of estimating blood oxygen with a light-based sensor, usually on a finger, toe, or ear. The result must be interpreted with symptoms and circulation because it can be inaccurate in some situations.

\medterm{Pulse Pressure} The systolic blood pressure minus the diastolic blood pressure. A widened or narrowed value may provide clues about circulation but must be interpreted with the full blood-pressure reading and clinical context.


\medterm{Pulse Rate} Pulse rate is the number of palpable arterial beats per minute and usually approximates heart rate, although some arrhythmias produce a pulse deficit.

\medterm{Pulseless Disease} Takayasu arteritis, an inflammation of large arteries that can narrow or block blood flow. It may cause weak pulses, unequal blood pressure, limb pain, dizziness, vision problems, or organ damage.
\medterm{Pulseless Electrical Activity} A cardiac-arrest rhythm in which the heart’s electrical activity continues but the heart does not pump enough blood to produce a pulse. It requires immediate resuscitation.
\medterm{Pulsus} A term used in medicine to describe the pulse and named patterns of arterial beats. The pattern may provide clues about heart rhythm, pumping strength, blood-vessel resistance, breathing, or circulation.
\medterm{Pulsus Paradoxus} An unusually large drop in blood pressure during breathing in, usually more than 10 millimeters of mercury. It can occur when the heart is compressed by fluid, or during severe asthma, chronic lung disease, or other heart conditions.

\medterm{Pulvinar} A region at the back of the thalamus, a deep brain structure. It helps direct visual attention, combine sensory information, and communicate with areas of the cerebral cortex.

\medterm{Pumice} A porous volcanic stone used as an abrasive; excessive dust inhalation can irritate the lungs.
\medterm{Pump Failure} Severe weakening of the heart’s pumping ability, causing inadequate blood flow to organs.
\synonyms
Cardiogenic pump failure

\medterm{Pump-Oxygenator} The part of a heart-lung machine that moves blood and adds oxygen while removing carbon dioxide during some heart operations. It temporarily replaces the pumping and gas-exchange work of the heart and lungs. Use can be associated with bleeding, inflammation, clots, or organ complications.

\medterm{Punch Biopsy} A circular, usually full-thickness skin specimen obtained with a punch instrument. It is useful for inflammatory, infectious, melanocytic, and other cutaneous lesions; small wounds may be closed with a suture.

\medterm{Punctal Occlusion} Partial or complete blocking of the small tear-drain openings near the inner eyelids, often with plugs or a procedure, to keep tears on the eye surface. It may help selected people with dry eye; watering, irritation, or infection can occur.

\medterm{Punctate} Made up of, or marked by, small spots, dots, or points. The term describes appearance and does not by itself identify a diagnosis.
\medterm{Punctate Hemorrhages (Tardieu Spots)} Tiny pinpoint areas of bleeding under the skin or on internal surfaces. They can form when small blood vessels rupture during severe pressure or oxygen deprivation, but their presence alone cannot establish the cause or timing of death.

\synonyms
Tardieu Spots

\medterm{Punctate Palmoplantar Keratoderma} An inherited condition causing many small, hard areas of thickened skin on the palms and soles. It often begins in adolescence or adulthood and can make walking or using the hands painful; rarely, a skin cancer develops.

\medterm{Puncture} A small hole or wound made by a pointed object, needle, bite, or sharp instrument. A puncture may be intentional for an injection or sample, or accidental and may carry infection, bleeding, or deeper injury risks.
\medterm{Puncture Biopsy} Alternate index label for needle biopsy: removal of cells or tissue through a percutaneous puncture for microscopic, molecular, or microbiologic examination; imaging may guide the target and route.

\medterm{Puncture Wound} A wound produced when a pointed object penetrates the skin and underlying tissue, creating a small surface opening with potentially deep tissue damage and infection risk.

\medterm{Pungent} Describes a strong, sharp smell or taste that may irritate the nose, mouth, eyes, or other moist linings.
\medterm{Punicic Acid} A conjugated polyunsaturated fatty acid found mainly in pomegranate seed oil and studied for lipid and inflammatory effects; clinical supplement benefits remain uncertain.

\medterm{Punishment} Punishment is a consequence intended to reduce a behavior; in healthcare, coercive or harmful punishment is inappropriate and behavior support should prioritize safety and dignity.

\medterm{Pupil} The central opening in the colored iris through which light enters the eye. Muscles change its size to control incoming light, and unequal or poorly reactive pupils may indicate eye or nerve disease.
\medterm{Pupil - White} A white pupil, or leukocoria, may result from reflection, cataract, retinal disease, infection, or retinoblastoma and requires prompt eye assessment, especially in a child.

\medterm{Pupillary} Relating to the pupil, the opening in the iris that controls how much light enters the eye. The pupil normally becomes smaller in bright light and larger in dim light; unequal size or poor reaction can indicate an eye, nerve, or brain problem.
\medterm{Pure Red Cell Aplasia} A rare disorder in which the bone marrow stops making red blood cells while other blood cells may remain normal. It causes severe anemia, tiredness, breathlessness, and pale skin and may be inherited or acquired from infection, immune disease, cancer, or medicines.

\synonyms
PRCA; Pure red-cell aplasia

\medterm{Pure Tone Audiometry} A hearing test in which a person signals whenever they hear tones of different pitches and loudness. It measures the quietest sounds detectable through air and bone and helps identify the type and degree of hearing loss.

\medterm{Pure-Tone Audiometry} A hearing test that measures the softest tones a person can hear at selected frequencies through air and bone conduction.
The results are plotted on an audiogram and help distinguish conductive from sensorineural hearing loss.

\synonyms
Pure-tone threshold audiometry; pure-tone hearing test.

\medterm{Purgative} A strong laxative that causes rapid emptying of the bowel. It may be used for bowel preparation before some procedures, but excessive use can cause diarrhea, dehydration, and dangerous changes in body salts.
\medterm{Purge} To empty the bowel, stomach, or body of a substance; the term may also mean forceful vomiting or diarrhea.
\medterm{Purging} Deliberate behaviors intended to compensate for eating, such as self-induced vomiting or misuse of laxatives,
diuretics, or excessive exercise.

\medterm{Purging Disorder} Repeated self-induced vomiting or misuse of laxatives, diuretics, or other methods to influence weight or body shape without recurrent binge eating. It can cause dangerous salt disturbances, tooth damage, dehydration, heart-rhythm problems, and tears in the food pipe.

\medterm{Purine} A nitrogen-containing building block of DNA, RNA, and energy-carrying molecules. The body breaks down purines into uric acid, which can accumulate and contribute to gout or kidney stones.
\medterm{Purine Nucleosides} Molecules made from a purine base, such as adenine or guanine, joined to a sugar. Examples include adenosine and guanosine, which help form genetic material and participate in cell energy and signaling.

\medterm{Purine Nucleotides} Molecules containing adenine or guanine joined to a sugar and phosphate groups. They build DNA and RNA and include ATP and GTP, which store energy and help transmit signals inside cells.
\medterm{Purines} Nitrogen-containing building blocks, mainly adenine and guanine, used to make DNA, RNA, and energy-carrying molecules. The body breaks them down into uric acid, which can form crystals and contribute to gout or kidney stones.

\medterm{Purinoceptor} A protein on a cell’s surface that detects substances such as ATP, ADP, or adenosine outside the cell. It helps regulate cell communication, inflammation, blood-vessel activity, and other body responses.

\medterm{Purkinje Cell} A major nerve cell in the cerebellum, the brain region that coordinates movement and balance. It sends signals that help fine-tune posture, coordination, and learned movements.

\medterm{Purkinje Fibers} Special electrical fibers spread through the lower chambers of the heart. They rapidly carry each heartbeat’s signal so the ventricles contract together and pump blood efficiently.

\medterm{Purkinje Fibres} Specialized heart-muscle fibers that rapidly spread electrical signals through the ventricles. Their coordinated activity helps the lower chambers contract together and pump blood efficiently to the lungs and body.

\medterm{Purpura} Purple or red spots caused by blood leaking under skin or mucous membranes; they do not fade when pressed.
\medterm{Purpura Fulminans} A rapidly worsening, life-threatening condition in which widespread blood clots and clotting-system failure cause purple skin patches, bleeding, and tissue death. It can follow severe infection, loss of natural anticoagulants, or another serious illness.

\medterm{Purulent} Containing or producing pus, a thick fluid made of immune cells, dead tissue, and often germs. Purulent drainage can occur in an abscess or infected wound, but its appearance alone does not identify the organism or the treatment needed.


\medterm{Purulent Pericarditis} A serious bacterial infection of the sac around the heart, causing pus and fluid to collect there. Fever, chest pain, breathlessness, and pressure on the heart may occur; urgent drainage and antibiotic treatment are needed.

\medterm{Pus} A thick fluid made of dead cells, germs, immune cells, and damaged tissue. It commonly collects during infection and may appear in wounds, abscesses, infected organs, or other inflamed body spaces.
\medterm{Push Enteroscopy} An examination of the upper small intestine using a long flexible camera passed through the mouth. It allows clinicians to see the lining, take samples, stop bleeding, or treat selected lesions.

\medterm{Pustule} A small raised skin lesion filled with pus. Pustules can occur with acne, follicle infection, impetigo, psoriasis, or a medicine reaction; their cause is determined from their location, appearance, symptoms, and timing.
\medterm{Pustules} Pustules are small raised skin lesions filled with pus, occurring in acne, folliculitis, infections, inflammatory disorders, or drug reactions.

\medterm{Pustulosis} A skin disorder characterized by sterile or infectious pustules, which may be localized or widespread and associated with inflammation, fever, psoriasis, drug reactions, or infection.

\medterm{Putamen} A deep brain structure that helps select and smooth purposeful movements and supports learning of repeated movement habits. Damage to it can cause abnormal movements, weakness, or changes in behavior.
\medterm{Putrefaction} Decomposition of body tissues after death caused mainly by bacteria and natural enzymes. It can produce color changes, swollen tissues from gas, visible veins, a strong odor, and gradual breakdown into liquid material; the rate depends on the environment.
\medterm{Putrescine} A chemical with a strong unpleasant smell made when proteins break down. It is found in decaying material and is also produced in small amounts during normal cell processes; its presence alone does not diagnose disease.
\medterm{PUVA} A light treatment that combines a light-sensitizing medicine called psoralen with ultraviolet-A exposure. It can treat selected skin diseases but may cause burns, eye injury, premature skin aging, and increased skin-cancer risk.

\medterm{PVC} A PVC, or premature ventricular contraction, is an early heartbeat arising in a ventricle; occasional PVCs are common, while frequent or symptomatic beats need evaluation.


\medterm{PVD} Alternate index label; see \textit{Peripheral artery disease}.

\medterm{PVI} Alternate name; see \textit{Pulmonary Vein Isolation}.

\medterm{PVPS} Alternate index label; see \textit{Post-vasectomy pain syndrome}.

\medterm{Pyromania} A mental disorder involving recurrent deliberate fire-setting, mounting tension or arousal before the act, and pleasure, gratification, or relief afterward, without a primary external motive such as financial gain, revenge, concealment of a crime, or political purpose. The behavior is not better explained by conduct disorder, mania, psychosis, or substance use.

\medterm{Pyarthrosis} Infection of a joint with pus in the joint space, usually bacterial septic arthritis, causing acute pain, swelling, restricted movement, fever, and risk of cartilage destruction.

\medterm{Pycnodysostosis} A rare inherited bone disorder in which bones are unusually dense but fragile. It can cause short stature, frequent fractures, short fingers, delayed closure of skull bones, crowded teeth, and problems with the jaw.

\medterm{Pyelitis} Inflammation of the renal pelvis, the funnel-shaped part of the kidney that drains urine. It is often associated with urinary infection and may cause flank pain, fever, or painful urination.
\medterm{Pyelo} Pyelo is an informal shortening usually referring to pyelonephritis, an infection of the kidney and renal pelvis that can cause fever and flank pain.

\medterm{Pyelogram} An image of the kidney drainage area and ureters after contrast material is given. It can show blockage, stones, narrowing, leaks, or structural changes in the urinary collecting system.
\medterm{Pyelography} Imaging of the renal pelvis and ureters using contrast, by X-ray or another imaging method. It can show obstruction, stones, leaks, strictures, or structural abnormalities.
\medterm{Pyelolithotomy} Surgical removal of a kidney stone through an incision into the renal pelvis. Less invasive stone treatments are often preferred when suitable.
\medterm{Pyelonephritis} Infection and inflammation of the kidney and its urine-collecting area, usually caused by bacteria traveling upward from the bladder. Symptoms may include fever, chills, side pain, nausea, and painful urination.
\medterm{Pyelotomy} An incision into the renal pelvis, usually made to remove an obstruction or kidney stone.
\medterm{Pylephlebitis} An infected blood clot in the large vein carrying blood from the abdominal organs to the liver. It may follow appendicitis, diverticulitis, or another abdominal infection and can cause fever, abdominal pain, jaundice, and life-threatening sepsis.

\medterm{Pyloric Antrum} The pyloric antrum is the lower stomach region that grinds food and regulates its passage toward the pyloric canal and duodenum.

\medterm{Pyloric Sphincter} A ring of muscle at the outlet of the stomach. It opens in a controlled way to let partly digested food enter the first part of the small intestine and helps prevent the intestine's contents from flowing backward.

\synonyms
Pylorus

\medterm{Pyloric Stenosis} Narrowing of the pyloric outlet, classically causing projectile vomiting in infants or gastric-outlet symptoms in adults.

\synonyms
Stenosis, Pyloric

\medterm{Pyloric Stenosis In Infants} Infantile pyloric stenosis is narrowing of the stomach outlet that causes forceful, non-bilious vomiting in young infants and usually requires surgical treatment.

\medterm{Pyloroplasty} Surgery to widen the pylorus, the outlet of the stomach, to improve passage of stomach contents. It may be used with treatment for obstruction or selected gastric-motility disorders.
\medterm{Pylorus} The distal part of the stomach leading into the duodenum.
\medterm{Pyoderma} A bacterial infection of the skin that can produce pus-filled spots, crusts, open sores, or deeper inflammation. Staphylococcus and Streptococcus are common causes; the infection may spread and require medical treatment, especially with fever, pain, or rapidly worsening redness.
\medterm{Pyoderma Gangrenosum} A rare inflammatory skin disease that causes very painful ulcers, often with a dark purple edge. It may occur with bowel inflammation, arthritis, or blood disorders; injury can make it worse, so treatment differs from ordinary wound infection.

\medterm{Pyogenesis} The formation of pus, a thick fluid made mainly of immune cells, dead tissue, and germs. It commonly occurs during bacterial infection or severe inflammation.
\medterm{Pyogenic} Producing pus or associated with a pus-forming infection. Pyogenic infections commonly cause redness, warmth, swelling, pain, and a collection of pus, although the cause and treatment vary.
\medterm{Pyogenic Granuloma} A noncancerous, rapidly growing lump made of extra small blood vessels. It is usually red, soft, and prone to bleeding and may occur after minor injury or during pregnancy; examination or biopsy confirms the diagnosis and guides removal.

\medterm{Pyogenic Liver Abscess} A pyogenic liver abscess is a pus-filled infection in the liver, usually from biliary, abdominal, or bloodstream spread, causing fever, pain, and possible sepsis.

\medterm{Pyometra} A collection of pus inside the uterus, usually caused by infection combined with blockage of drainage through the cervix. It may cause lower abdominal pain, fever, tenderness, or foul-smelling discharge; in older adults it can be linked to cervical narrowing or a tumor and needs prompt evaluation.
\medterm{Pyonephrosis} Pus trapped in the kidney's urine-collecting system, usually because infection occurs behind a blockage. It is a medical emergency requiring urgent drainage, antibiotics, and treatment of the obstruction.

\medterm{Pyosalpinx} A fallopian tube filled with pus, usually because of pelvic infection. It may cause lower abdominal pain, fever, unusual discharge, infertility, or an abscess and generally requires prompt treatment.
\medterm{Pyostomatitis Vegetans} A rare inflammatory mouth condition causing many small yellowish pus-filled spots on red, swollen tissue, especially the inner cheeks and gums. It is strongly associated with inflammatory bowel disease and may appear before bowel symptoms.

\medterm{Pyramid} A structure shaped like a pyramid. In the brain, the medullary pyramids contain nerve fibers that carry movement commands from the brain to the spinal cord; other pyramids include pyramid-shaped areas in the kidney.
\medterm{Pyramidal Cell} A triangular-shaped excitatory neuron with a prominent apical dendrite, found especially in the cerebral cortex and hippocampus and involved in long-range brain signaling.

\medterm{Pyramidal Tract} The pyramidal tract is a major motor pathway from the cerebral cortex to the brainstem and spinal cord that controls voluntary movement.

\medterm{Pyramidal Tracts} Nerve pathways from brain movement areas down the brainstem and spinal cord, carrying signals controlling voluntary movement.
\medterm{Pyrazinamide} An antibiotic used with other medicines during the initial treatment of tuberculosis. It helps shorten treatment but can injure the liver or raise uric acid, so monitoring and prescribed dosing are important.
\medterm{Pyrazines} A family of chemical compounds found naturally in some foods and used in flavors, medicines, and industry. Their effects on health depend on the specific compound, amount, and route of exposure.

\medterm{Pyrazole} A small five-membered chemical structure used to build medicines, pesticides, and other compounds. The word names a chemical family and does not identify one particular treatment or effect.

\medterm{Pyrazoles} Pyrazoles are a family of chemical compounds containing a pyrazole ring, including some medicines and pesticides.


\medterm{Pyrenochaeta} A group of fungi found mainly in soil and plants that rarely infect people. When infection occurs, it usually affects tissue beneath injured skin or people with weakened immunity.

\medterm{Pyrethrins} Natural insect-killing chemicals derived from chrysanthemum flowers and used in some pest-control products. Exposure can irritate skin and airways and, in large amounts, may cause vomiting, tremors, seizures, or breathing problems.
\medterm{Pyrethrum} The dried flower heads of certain Chrysanthemum species, used as a natural source of pyrethrins, insecticidal compounds that affect insect sodium channels. Pyrethrum and pyrethrins differ from synthetic pyrethroids; exposure may irritate the skin or airways and can cause toxicity with substantial exposure.
\medterm{Pyrexia (Fever)} An abnormally high body temperature, usually $38^\circ\text{C}$ ($100.4^\circ\text{F}$) or higher. Fever is a controlled response in which the brain raises the body’s temperature setting, often because of infection or inflammation. It differs from hyperthermia, in which temperature rises without that change in setting, such as during heat stroke. Treatment depends on the cause and the person’s symptoms.
\medterm{Pyridine} A nitrogen-containing ring-shaped chemical structure found in many medicines, vitamins, and biologically active substances. Pyridine groups can change how compounds dissolve, react, travel through the body, and produce medical effects.

\medterm{Pyridostigmine} A medicine that increases acetylcholine, a chemical helping nerves activate muscles. It mainly treats myasthenia gravis and may cause stomach cramps, diarrhea, sweating, increased saliva, muscle twitching, or slow heartbeat.
\medterm{Pyridostigmine Bromide} Pyridostigmine bromide increases communication between nerves and muscles and is used mainly for myasthenia gravis; it may cause abdominal cramps, diarrhea, or excess secretions.

\medterm{Pyridoxal Phosphate} The active form of vitamin B6 used by many enzymes. It helps the body process amino acids, make certain brain chemicals and red blood cells, and maintain normal nerve function.

\medterm{Pyridoxine} One form of vitamin B6, converted in the body to pyridoxal phosphate, a coenzyme for amino-acid metabolism. It is needed for processing food, making red blood cells, and supporting nerves and the immune system; deficiency can cause anemia or nerve problems, while excessive supplements can damage nerves.

\synonyms
Vitamin B6

\medterm{Pyridoxine (Vitamin B6)} Alternate index label; see \textit{Pyridoxine}.

\medterm{Pyrimethamine} A parasite-fighting medicine that blocks folate use in certain protozoa. It is used for selected infections, often with another medicine, and can suppress bone marrow or cause serious skin reactions.
\medterm{Pyrimidine} A six-membered nitrogen-containing ring that forms the bases cytosine, thymine, and uracil in nucleic acids.
\medterm{Pyrimidine Analog} A medicine or chemical resembling a pyrimidine nucleotide that can interfere with DNA or RNA synthesis, often for cancer or antiviral treatment.

\medterm{Pyrimidine Dimer} A DNA lesion formed when adjacent pyrimidine bases bond abnormally after ultraviolet exposure, potentially causing mutations if unrepaired.

\medterm{Pyrimidine Nucleosides} Nucleosides containing cytosine, thymine, or uracil bases linked to a sugar, serving as components or precursors of DNA and RNA.

\medterm{Pyrimidine Nucleotides} Genetic and energy-related building blocks containing pyrimidine bases, including cytidine, uridine, and thymidine. Cells use them to make DNA, RNA, cell membranes, and other substances needed for growth and repair.

\medterm{Pyrimidines} Nitrogenous bases containing a single six-membered heterocyclic ring, serving as
structural building blocks for nucleic acids. The primary pyrimidines are cytosine (C),
thymine (T, present in DNA), and uracil (U, present in RNA).

\medterm{Pyrogen} A pyrogen is a substance that induces fever, commonly through immune signaling; bacterial endotoxin is an important example in contaminated medicines or fluids.

\medterm{Pyroglyphidae} Pyroglyphidae is a mite family that includes house-dust mites whose allergens can trigger allergic rhinitis, asthma, or atopic dermatitis.


\medterm{Pyrophosphatase} An enzyme that hydrolyzes inorganic pyrophosphate and helps drive many biosynthetic reactions forward. An enzyme that breaks down pyrophosphate, helping many chemical reactions needed to build cellular materials.
\medterm{Pyropoikilocytosis} A rare inherited disorder that makes red blood cells unusually small, irregular, and fragile. The cells break down easily, causing hemolytic anemia with tiredness, jaundice, pale skin, or an enlarged spleen, often beginning in infancy or childhood.

\medterm{Pyroptosis} A form of controlled cell death that causes inflammation. The cell swells, breaks open, and releases signals that alert the immune system, often in response to infection or severe cell stress.

\medterm{Pyrosequencing} A laboratory method for reading selected DNA sequences by detecting light produced as DNA building blocks are added. It can identify some genetic changes or microbes but is not suitable for every sequencing need.

\medterm{Pyrosis} Heartburn, a burning feeling behind the breastbone usually caused by stomach contents flowing backward into the swallowing tube. Frequent symptoms may reflect reflux disease and can damage the esophagus without treatment.

\medterm{Pyrroles} A group of five-membered nitrogen-containing ring compounds found in heme, bile pigments, medicines, and other biologically active molecules.

\medterm{Pyrrolidonecarboxylic Acid} A cyclic amino-acid derivative, also called pyroglutamic acid, found in proteins and metabolic pathways; abnormal accumulation may occur in selected metabolic or medication-related conditions.

\medterm{Pyrrolnitrin} Pyrrolnitrin is an antimicrobial compound produced by some bacteria and studied for antifungal activity and biological control.

\medterm{Pyruvate} A small molecule produced when cells break down glucose. Cells can use it to make more energy, convert it into other substances, or produce lactate when oxygen or energy-processing pathways are limited.
\medterm{Pyruvate Decarboxylase} Pyruvate decarboxylase is an enzyme that converts pyruvate to acetaldehyde and carbon dioxide during fermentation in some organisms.

\medterm{Pyruvate Dehydrogenase} An enzyme complex that changes pyruvate, made when cells break down sugar, into acetyl-CoA for energy production. Inherited deficiency can cause excess lactic acid, poor growth, developmental problems, or neurologic disease.

\medterm{Pyruvate Kinase} An enzyme that helps red blood cells make energy from glucose. If it is inheritedly deficient, red cells may break down too early, causing chronic hemolytic anemia, jaundice, gallstones, tiredness, or an enlarged spleen.

\medterm{Pyruvate Kinase Blood Test} A pyruvate kinase blood test measures the enzyme in red blood cells and may help evaluate pyruvate kinase deficiency, an inherited cause of hemolytic anemia.

\medterm{Pyruvate Kinase Deficiency} A rare inherited condition in which red blood cells cannot make enough energy and break down early. It causes anemia of varying severity, jaundice, tiredness, gallstones, an enlarged spleen, and sometimes excess iron in the body.

\medterm{Pyuria} Presence of white blood cells or pus in urine, often indicating urinary-tract inflammation or infection.

\medterm{Paget's Disease} An ambiguous term that may refer to Paget disease of bone or Paget disease of the breast. Paget disease of bone is a chronic disorder of abnormal bone breakdown and regrowth that can enlarge, weaken, or deform bone; mammary Paget disease is a malignancy involving the nipple and areola, often associated with underlying breast cancer.

\medterm{Paget's Disease of the Nipple} A rare breast cancer involving the skin of the nipple and areola. It may cause redness, scaling, itching, discharge, or a lump and requires breast evaluation.

\medterm{Paget's Disease of the Vulva} A rare cancer or precancerous condition affecting the skin of the vulva. It may cause persistent itching, burning, redness, or a scaly patch and needs biopsy assessment.


\medterm{Paleo Diet} An eating pattern based on foods thought to resemble those available to early humans, such as meat, fish, vegetables, fruits, nuts, and seeds. It excludes or limits grains, legumes, and dairy and may not provide balanced nutrition for everyone.

\medterm{Papillomatosis} A condition in which multiple papillomas, or wart-like overgrowths of epithelial tissue, develop. Causes and risks vary by body site, virus, inherited condition, or other disease.

\medterm{Paracetamol} Also called acetaminophen, paracetamol reduces pain and fever. Excessive doses can cause serious liver injury, especially with alcohol or other products containing paracetamol.

\medterm{Paraneoplastic Neurologic Syndromes} Rare disorders in which an immune response to cancer also attacks the nervous system. Symptoms may include weakness, imbalance, seizures, memory change, or sensory problems and can precede cancer diagnosis.

\medterm{Paraneoplastic Pemphigus} A rare, often severe autoimmune blistering disease associated with an underlying tumor, especially a lymphoid cancer. It causes painful mouth sores and skin lesions and requires specialist treatment.

\medterm{Parechovirus} A group of viruses that commonly cause mild respiratory or gastrointestinal illness but can cause severe sepsis-like illness, meningitis, or encephalitis in newborns and young infants.

\medterm{Parenting} The care, protection, guidance, and support provided to a child. Healthy parenting includes meeting basic needs, setting safe limits, nurturing development, and seeking help when concerns arise.

\medterm{Paresthesia (Pins and Needles)} An abnormal sensation such as tingling, prickling, burning, or numbness. Temporary pressure on a nerve is common, but persistent or spreading symptoms may indicate nerve, metabolic, vascular, or spinal disease.

\medterm{Paruresis} Difficulty urinating in the presence of other people or in public places, often called shy bladder syndrome. It is related to anxiety and can improve with behavioral therapy and gradual exposure.

\medterm{Passive Heat Therapy} Use of external heat, such as warm packs, baths, or saunas, to relax muscles and temporarily reduce pain or stiffness. Excessive heat can cause burns, dehydration, or fainting.

\medterm{Passive Smoking} Breathing smoke from another person's tobacco or nicotine product. Secondhand smoke increases the risk of respiratory disease, heart disease, stroke, and lung cancer and is especially harmful to children and during pregnancy.

\medterm{Patau's Syndrome} A congenital disorder caused by an extra copy of chromosome 13. It can cause severe developmental problems, brain and heart defects, cleft lip or palate, and other organ abnormalities.

\medterm{Paternal Depression} Depression occurring in a father during pregnancy or after a child is born. Symptoms may include low mood, irritability, withdrawal, sleep or appetite changes, and loss of interest; treatment can help both parent and family.

\medterm{Paternal Postnatal Depression} Depression in a father during the first year after a baby's birth. It may be linked to stress, sleep loss, previous depression, or a partner's depression and deserves professional assessment.

\medterm{Pelvic Inflammatory Disease (PID)} Infection and inflammation of the upper female reproductive organs, often caused by sexually transmitted bacteria. Untreated PID can cause chronic pelvic pain, infertility, or ectopic pregnancy.

\medterm{Pemphigus Foliaceus} An autoimmune blistering skin disease caused by antibodies against proteins that hold skin cells together. It causes superficial fragile blisters and crusted sores but usually does not affect mucous membranes.

\medterm{Penile Intraepithelial Neoplasia (PeIN or PIN)} Abnormal precancerous cells in the surface lining of the penis. It may appear as a persistent red, white, or thickened patch and can progress to invasive cancer if untreated.

\medterm{Percutaneous Closure} Closing a heart or blood-vessel opening using a device delivered through a catheter inserted through the skin. It can avoid open surgery in selected patients but has procedure-related risks.

\medterm{Personalized Medicine} Healthcare that uses a person's genes, biological measurements, environment, and lifestyle to guide prevention or treatment. It can improve treatment selection but does not guarantee a response.

\medterm{Peyronie's Disease} A condition in which scar tissue forms in the penis and causes a curved or painful erection, shortening, or difficulty with sexual activity. Treatment depends on severity and symptoms.

\medterm{pH} A measure of how acidic or alkaline a solution is, usually on a scale from 0 to 14. Body fluids have tightly controlled pH ranges, and major changes can be dangerous.

\medterm{Phantom Pregnancy} A false belief or physical state in which a person has pregnancy-like symptoms without an ongoing pregnancy. Symptoms can include abdominal enlargement, breast changes, nausea, or missed periods and require medical evaluation.

\medterm{Pharmacoepidemiology} The study of how medicines are used and affect health in large populations. It helps assess effectiveness, safety, adverse effects, and patterns of prescribing.


\medterm{Pharmacoproteomics} The study of how a person's protein patterns influence medicine response, drug effects, or toxicity. It is an emerging part of precision medicine.

\medterm{Photoplethysmography} A noninvasive optical method that detects blood-volume changes in tissue. It is used in devices such as pulse oximeters and some heart-rate monitors.

\medterm{Phyllodes Tumors} Uncommon fibroepithelial tumors of the breast that can be benign, borderline, or malignant. They often appear as a rapidly growing lump and are treated mainly with complete surgical removal.


\medterm{Pick's Disease} An older term often referring to frontotemporal dementia, a group of disorders causing progressive changes in behavior, personality, language, or executive function. The term is now used less consistently.

\medterm{Piles (Haemorrhoids)} Alternate spelling; see \textit{Hemorrhoids}.

\medterm{Pins and Needles} A common name for paresthesia, an abnormal tingling, prickling, or numb sensation. Brief episodes often follow pressure on a nerve; persistent symptoms need evaluation.

\medterm{Polymorphic Light Eruption} An itchy rash triggered by sunlight, often appearing in spring or after the first strong sun exposure. It may improve as the skin adapts, and sun protection reduces recurrences.

\medterm{Polypharmacology} The study or use of a medicine that affects several targets in the body. This can be helpful when one disease involves multiple pathways, but it may also increase side effects and medicine interactions.

\medterm{Portal Hypertensive Colopathy} Changes in the colon caused by high pressure in veins around the liver, usually from advanced liver disease. Enlarged fragile blood vessels may cause visible or hidden bleeding and anemia.

\medterm{Post-Adoptive Depression} Depression that begins after adopting a child. Stress, sleep loss, adjustment demands, past trauma, or unmet expectations may contribute; persistent sadness, loss of interest, or thoughts of self-harm need prompt mental-health assessment.

\medterm{Posterior Cruciate} Usually referring to the posterior cruciate ligament, a strong knee ligament that limits backward movement of the shinbone. Injury can cause pain, swelling, and knee instability.

\medterm{Postnatal Anxiety} Excessive anxiety occurring after childbirth. It may involve persistent worry, panic, intrusive thoughts, restlessness, or physical symptoms and can be treated with therapy, support, and sometimes medicine.

\medterm{Post-Traumatic Stress Disorder (PTSD)} Alternate index label; see \textit{Post-Traumatic Stress Disorder}.

\medterm{Powassan Virus} A tick-borne virus that can cause encephalitis or meningitis. Severe infection may cause headache, fever, confusion, seizures, or weakness; there is no specific antiviral treatment.

\medterm{Precision Medicine} An approach that matches prevention or treatment to biological differences between people, often using genetic, molecular, or clinical information. It is related to but not identical to personalized medicine.

\medterm{Precision Oncology} An approach to cancer care that uses information about a tumor’s genes, proteins, or other features to help select prevention, testing, or treatment. It aims to match treatment to the biology of the cancer rather than to its location alone.

\medterm{Prenatal Depression} Depression during pregnancy. Symptoms may include persistent sadness, loss of interest, guilt, hopelessness, sleep or appetite changes, and thoughts of self-harm; prompt professional support is important.

\medterm{Programmed Cell Death} A controlled process by which cells die as part of normal development, tissue maintenance, or disease. Apoptosis is the best-known form; other regulated forms include necroptosis and pyroptosis.

\medterm{Proteasome Inhibitor} A medicine that blocks the proteasome, a cell structure that breaks down unwanted proteins. These drugs are used especially in multiple myeloma and some blood cancers.

\medterm{Proteostasis} The balance of protein production, folding, movement, and breakdown that keeps cells functioning. Failure of proteostasis can contribute to aging, neurodegeneration, and other diseases.

\medterm{Pruritis} The medical term for itching, an unpleasant urge to scratch. Causes include dry skin, allergy, inflammation, infection, medicines, nerve disorders, and internal disease.

\medterm{Psychosomatic Disorder} An older broad term for physical symptoms influenced by psychological and social factors. The symptoms are real; modern assessment considers both body and mind rather than assuming symptoms are imagined.

\medterm{Psychosomatic Medicine} Medical care that considers how body processes, thoughts, emotions, relationships, and social conditions interact in illness. Physical symptoms are real, and assessment looks for both medical and psychological factors without assuming symptoms are imagined.

\medterm{Pulmonary Hypoplasia} Underdevelopment of one or both lungs, resulting in too little lung tissue for normal breathing. It may occur with congenital abnormalities, low amniotic fluid, or compression before birth.

\medterm{Pulmonology} The medical specialty concerned with the lungs and breathing system. Pulmonologists diagnose and treat asthma, chronic lung disease, infections, sleep-related breathing disorders, and other respiratory conditions.

\medterm{PYY} Peptide YY is a hormone released by the intestine after eating. It helps signal fullness and slows movement of food through the digestive tract; abnormal signaling is being studied in obesity and other metabolic disorders.
